\documentclass[11pt,a4paper]{article}
\usepackage[margin=27mm]{geometry}
\usepackage{fontspec}
\setmainfont{Latin Modern Roman}
\usepackage{amsmath,amssymb,amsthm,mathtools}
\usepackage{unicode-math}
\setmathfont{Latin Modern Math}
\setcounter{MaxMatrixCols}{32}
\emergencystretch=4em
\usepackage{booktabs,longtable,array,microtype}
\usepackage[round,authoryear]{natbib}
\usepackage[unicode,colorlinks=true,linkcolor=blue,citecolor=blue,urlcolor=blue]{hyperref}
\usepackage{xurl}
\newcommand{\R}{\mathbb R}
\newcommand{\Z}{\mathbb Z}
\newcommand{\Q}{\mathbb Q}
\newcommand{\C}{\mathbb C}
\newcommand{\F}{\mathbb F}
\newcommand{\N}{\mathbb N}
\newcommand{\OO}{\mathcal O}
\newcommand{\dd}{\,\mathrm d}
\DeclareMathOperator{\Imn}{Im}
\DeclareMathOperator{\Ren}{Re}
\DeclareMathOperator{\Aut}{Aut}
\DeclareMathOperator{\End}{End}
\DeclareMathOperator{\Tr}{Tr}
\DeclareMathOperator{\Nm}{Nm}
\DeclareMathOperator{\Sha}{Sha}
\newtheorem{statement}{Statement}[section]
\newtheorem{theorem}[statement]{Theorem}
\newtheorem{proposition}[statement]{Proposition}
\newtheorem{lemma}[statement]{Lemma}
\newtheorem{corollary}[statement]{Corollary}
\theoremstyle{definition}
\newtheorem{definition}[statement]{Definition}
\newtheorem{remark}[statement]{Remark}
\newtheorem{question}[statement]{Question}
\title{S6\_05\_ARITHMETIC\\Arithmetic counts, lattice neighbours, and elliptic period norms}
\author{AI-produced by GPT-6 Astra\\OpenAI, ChatGPT Pro mode}
\date{15 November 2026}
\begin{document}
\maketitle
\begin{center}\textbf{INCOMPLETE research manuscript}\end{center}
\begin{abstract}
For the specified sectorial count we prove the error bound
$O_\epsilon(X^{2/3+\epsilon})$, the dyadic mean-square bound
$O_\epsilon(X^{2+\epsilon})$, and both signs of square-root
oscillation. The latter follows from an exact sharp-arc resolvent
kernel and a localized relative-wave calculation with summable
endpoint terms. An exact census gives $R(10^9)=35\,363\,970$.
We classify all $129$ Conway-group orbits of isotropic planes in
$\Lambda/3\Lambda$, including their eight unimodular extensions
and integral stabilizer generators. The determinant-$81$ lattice
$J$ has automorphism group of order $12\,597\,120$.
For the elliptic period norm, we prove the torsion-value and
global-height identities, determine all nine rank-one rational
leading factors, and establish the maximum and global level-set
geometry without an interval-box certificate. We prove the
transcendence assertion under the stated period conjecture.
Two requested conclusions remain unresolved here: a verified
implication from a numbered known conjecture to the sharp-arc
square-root upper bound, and the unconditional arithmetic nature
of the norm at nontorsion algebraic points. Exact route tests and
the smallest undecided instance are retained. This AI-produced
manuscript makes no claim of priority.
\end{abstract}
\tableofcontents
\IfFileExists{sections/A.tex}{\section{The arithmetic count}\label{sec:A}
We write
\[
 F(s,t,z)=s^2-3t^2-18z^2,\quad
 \varepsilon_0=2+\sqrt3,\quad L_0=\log\varepsilon_0,
 \quad b_0=\log\frac3{2\sqrt2}.
\]
The letter $t$ in this section is an integer coordinate, unrelated to the
elliptic-surface parameter in Part C. All areas on the hyperbolic plane use
curvature $-1$ and the measure $dx\,dy/y^2$.

\subsection{The coordinates, ideals, and all boundary terms}
\paragraph{Statement A.1.1 (Proved).}
For every real $X\geq1$, the count in the prompt equals
\[
 \#\{(s,t,z)\in\mathbb Z^3:F(s,t,z)=-2,\ 1\leq z\leq X,
 \ t\geq0,\ t+4z<s\leq17z/4\}.
\]
The comparison map from the original variables is the integral isometry
\[
 \begin{pmatrix}s\\t\\z\end{pmatrix}=
 \begin{pmatrix}17&0&72\\0&1&0\\4&0&17\end{pmatrix}
 \begin{pmatrix}a\\q\\n\end{pmatrix},\qquad
 (a,q,n)=(17s-72z,t,17z-4s).
\]
It is a bijection of the indicated integer sets, with no parity restriction
lost. In the original $(x,y,z)$ coordinates,
\[
 s=x+y+4z,\quad t=y-x,\quad
 x=(s-4z-t)/2,\quad y=(s-4z+t)/2,
\]
and $2\{4(xy+xz+yz)-x^2-y^2-z^2\}=F(s,t,z)$.
\begin{proof}
The matrix has determinant $1$ and preserves $\operatorname{diag}(1,-3,-18)$.
The upper inequality is precisely $n\geq0$, and the strict lower inequality
is precisely $a+4n>q$. Conversely these inequalities force $a>0$:
if $a\leq0$ and $n=0$ the strict inequality is impossible; if $n\geq1$, it
implies $-4n<a\leq0$, whereas the equation gives
$a^2=3q^2+18n^2-2\geq16n^2$, a contradiction. The transformed height is
exactly $z=4a+17n$. Reduction of $F=-2$ modulo $2$ gives $s\equiv t\pmod2$,
so the displayed inverse to $(x,y,z)$ is integral. Direct expansion proves
its polynomial identity, and proves every asserted equivalence.
\end{proof}

\paragraph{Statement A.1.2 (Proved).}
Every nonzero integral ideal of $\mathbb Z[\sqrt3]$, of absolute ideal
norm $m$, has a unique generator $\alpha=s+t\sqrt3$ satisfying
\[
 \sqrt m\leq\alpha<\varepsilon_0\sqrt m.
\]
Its field norm may be $+m$ or $-m$. For example $(1+\sqrt3)$ has ideal
norm $2$ and field norm $-2$; a generator of field norm $+2$ is impossible
modulo $3$. Thus the interval assertion in Lead 2 is correct as written,
but it must not be silently strengthened to positive field norm for every
$m$. For the norms actually used here,
\[
 m_z=18z^2-2\quad(z\geq1),
\]
every ideal of norm $m_z$ has positive field-norm generator, so the unique
interval generator satisfies
\[
 N_{\mathbb Q(\sqrt3)/\mathbb Q}(\alpha)=m_z,\qquad s>0,\quad t\geq0.
\]
The number of these ideals is
\[
 \rho(m_z)=\sum_{d\mid m_z}\chi_{12}(d)
          =\rho(3z-1)\rho(3z+1).
\]
\begin{proof}
The full ring of integers is $\mathbb Z[\sqrt3]$. It is norm Euclidean:
rounding the two coefficients of an element of $\mathbb Q(\sqrt3)$ gives
an error $u+v\sqrt3$ with $|u|,|v|\leq1/2$ and
$|u^2-3v^2|\leq3/4<1$. Thus all ideals are principal. The units are
$\{\pm\varepsilon_0^k:k\in\mathbb Z\}$: a unit of norm $-1$ is impossible
modulo $3$, and the least unit greater than $1$ with positive conjugate is
$2+\sqrt3$. The latter assertion also follows directly since
$1<\alpha<\varepsilon_0$ would give an integral coefficient
$0< (\alpha-\alpha^{-1})/(2\sqrt3)<1$.

For any $m$, choose a generator, change its real sign to positive, and
multiply by a unique power of $\varepsilon_0$ to put it in the stated
half-open interval. This proves the general assertion without any
assumption on the sign of its field norm.
A generator of an ideal of norm $m_z$ has field norm $\pm m_z$.
Since $m_z\equiv1\pmod3$, the negative sign is impossible modulo $3$.
Choose its positive sign in the real embedding and multiply by a unique
power of $\varepsilon_0$ to put it in the stated half-open interval.
Then $s>0$ and $t\geq0$. The standard ideal factorization is particularly
transparent here: a split prime contributes $e+1$ ideals of norm $p^e$,
an inert prime contributes one if $e$ is even and none otherwise, and a
ramified prime contributes one. These are exactly the Euler factors of
$1*\chi_{12}$. The two linear factors $3z-1$ and $3z+1$ have greatest common
divisor dividing $2$; the extra factor $2$ and their common factors at $2$
make no difference because $2$ is ramified and $\rho(2^e)=1$.
\end{proof}

\paragraph{Statement A.1.3 (Proved).}
Put $\theta(\mathfrak a)=L_0^{-1}\log(\alpha/\sqrt{m_z})$ for the generator
in Statement A.1.2 and $\lambda_0=b_0/L_0$. For $X\geq1$,
\begin{equation}\label{A:arc-exact}
 R(X)=\sum_{1\leq z\leq X}
 \#\{\mathfrak a:N\mathfrak a=m_z,\ 0\leq\theta(\mathfrak a)\leq\lambda_0\}
 +R_0(X)-1,
\end{equation}
where
\[
 R_0(X)=\#\{k\geq1:4a_k\leq X\},\qquad
 a_k+q_k\sqrt3=(1+\sqrt3)\varepsilon_0^k.
\]
In particular $R_0(X)=O(\log(2X))$. The only points with $z\geq1$,
$t\geq0$, $s>0$, $s\leq17z/4$, and $F=-2$ which fail $s-4z>t$ are
$(4,0,1)$ and $(17,1,4)$.
\begin{proof}
Since
\[
 \cosh b_0=\frac{17}{12\sqrt2},
\]
the fixed-arc condition is exactly
\begin{equation}\label{A:arc-integer}
 288s^2\leq289(18z^2-2).
\end{equation}
It implies $4s<17z$. Conversely let $r=17z-4s\geq1$. Then
$r<17z$ and
\[
 18r(34z-r)\geq18(34z-1)\geq594>578,
\]
which is precisely the strict version of \eqref{A:arc-integer} after
substitution. Consequently the only points between the fixed arc and
the moving upper wall are the points $4s=17z$.

If $s-4z\leq t$, the equation implies
$(6z-s)^2\leq3z^2+1$. The upper wall gives
$(6z-s)^2\geq49z^2/16$, so $z\leq4$. For $z=1,2,3,4$ the interval
$\lceil\sqrt{18z^2-2}\rceil\leq s\leq\lfloor17z/4\rfloor$
gives respectively $(s,t)=(4,0)$, no point, no point, and $(17,1)$.
The first lies in the fixed arc; the second is on the moving wall.

On that wall write $z=4a$, $s=17a$. Then $a^2-3t^2=-2$.
Every solution with $a>0,t\geq0$ is
$(a,t)=(a_k,q_k)$ for a unique $k\geq0$. Indeed, reduction by the unit to
$\sqrt2\leq a+t\sqrt3<\varepsilon_0\sqrt2$ gives
$a\geq0$, $2a<3t$, and hence $3t^2<8$; its only possibility is $(1,1)$.
The strict wall excludes $k=0$ and includes all $k\geq1$.
Finally $a_k$ grows geometrically, proving the logarithmic estimate and
\eqref{A:arc-exact}, including both exceptional points.
\end{proof}

\paragraph{Statement A.1.4 (Proved by exact exhaustive computation).}
The independently computed counts are
\[
 \begin{array}{c|r}
 X&R(X)\\\hline
 10^6&35\,381\\
 10^7&353\,598\\
 10^8&3\,536\,002\\
 10^9&35\,363\,970
 \end{array}
\]
The first three confirm Lead 5. The fourth extends the census by a factor
of ten. The exact algorithm, its exhaustiveness proof, and its numerical
discrepancy analysis are given below in Statement A.5.1.

\subsection{The full integral group and the two closed geodesics}
Let $B_F(x,y)=x^{\mathsf T}\operatorname{diag}(1,-3,-18)y$ and
\[
 \mathbb H_F=\{x:F(x)=1,\ x_1>0\},\qquad
 O^+(F,\mathbb Z)=\{g\in O(F,\mathbb Z):g\mathbb H_F=\mathbb H_F\}.
\]
The orientation-preserving Fuchsian group is
$\Gamma^+=O^+(F,\mathbb Z)\cap SO(F,\mathbb R)$.
Thus the literal $SO(F,\mathbb Z)$ in Lead 3 must be restricted to its
index-two subgroup preserving the chosen component before it acts on
$\mathbb H_F$.

\paragraph{Statement A.2.1 (Proved).}
For the five primitive vectors
\[
 r_0=(0,-1,0),\ r_1=(0,0,-1),\ r_2=(1,1,0),\
 r_3=(9,3,2),\ r_4=(4,0,1),
\]
let $S_jx=x-2B_F(x,r_j)r_j/F(r_j)$. Then
$O^+(F,\mathbb Z)=\langle S_0,S_1,S_2,S_3,S_4\rangle$.
A fundamental domain is the compact right-angled pentagon
\[
 \mathcal P=\{(s,t,z)\in\mathbb H_F:
 t,z\geq0,\ s\geq3t,\ s\geq t+4z,\ 2s\geq9z\}.
\]
Its area is $\pi/2$. The group $\Gamma^+$ has area $\pi$, is generated by
$S_0S_j$ ($1\leq j\leq4$), and has signature $(0;2,2,2,2,2)$.
\begin{proof}
The five reflection matrices, in their displayed order, are
\[
 \begin{gathered}
 S_0=\operatorname{diag}(1,-1,1),\quad
 S_1=\operatorname{diag}(1,1,-1),\quad
 S_2=\begin{pmatrix}2&-3&0\\1&-2&0\\0&0&1\end{pmatrix},\\
 S_3=\begin{pmatrix}10&-9&-36\\3&-2&-12\\2&-2&-7\end{pmatrix},\qquad
 S_4=\begin{pmatrix}17&0&-72\\0&1&0\\4&0&-17\end{pmatrix}.
 \end{gathered}
\]
They are integral isometries, have determinant $-1$, and preserve
$\mathbb H_F$. In projective coordinates $(u,v)=(t/s,z/s)$ the pentagon
has vertices
\[
 (0,0),\ (1/3,0),\ (1/3,1/6),\ (1/9,2/9),\ (0,2/9).
\]
Each lies strictly inside $3u^2+18v^2<1$. Consecutive normals are
orthogonal; hence all five angles are $\pi/2$. Reflection across the sides
produces a tiling: glue copies along the paired sides, with four copies
around every vertex. The relations $S_j^2=1$ and
$(S_jS_{j+1})^2=1$ (indices modulo $5$) remove the edge and vertex loops,
so the resulting simply connected development is a hyperbolic surface.
It is complete: a bounded path cannot cross infinitely many nonadjacent
sides, since the pentagon is compact and these sides are separated by a
positive distance; only four chambers meet at a vertex. Its developing
map to the complete simply connected plane is consequently an isometry.
Thus the interiors of the pentagons are disjoint. The area formula gives
$3\pi-5\pi/2=\pi/2$.

We also prove that no integral isometries have been omitted. Let
$x=(a,q,n)\in\mathbb Z^3$ have $F(x)=1$, $a>0$. Reflections $S_0,S_1$
make $q,n\geq0$. If one of the other three defining inequalities fails,
reflection in that side strictly decreases the positive integer $a$.
It remains positive because the reflection preserves $\mathbb H_F$.
The process terminates in $\mathcal P$. There
\[
 a^2=(1-3u^2-18v^2)^{-1}\leq27/2;
\]
the maximum follows by convexity of $3u^2+18v^2$ on the projective polygon
and evaluation at its vertices. For integral $a\leq3$, the only norm-one
point in the pentagon is $(1,0,0)$: $a=2$ would give $(q,n)=(1,0)$,
which violates $a\geq3q$, and $a=3$ is impossible modulo $3$.

Apply this descent to $g(1,0,0)$ for any $g\in O^+(F,\mathbb Z)$.
After multiplication by a word in the $S_j$, it fixes $(1,0,0)$.
An integral isometry of the remaining positive binary form
$3q^2+18n^2$ only changes the two signs: its first column of norm $3$
is $(\pm1,0)$, and its second column is then $(0,\pm1)$.
These sign changes are $S_0,S_1$. This proves the asserted equality.
The orientation subgroup has index two, and doubling the pentagon gives
its stated orbifold signature and area.
\end{proof}

\paragraph{Statement A.2.2 (Proved).}
The integral vectors of norm $-2$ form exactly two $\Gamma^+$-orbits,
represented by $v_1=r_2=(1,1,0)$ and $v_2=r_4=(4,0,1)$.
They are distinguished by
\[
 \operatorname{div}(s,t,z)=\gcd(s,3t,18z),
\]
which equals $1$ and $2$, respectively. The stabilizers of the oriented
normals are primitive infinite cyclic groups. The associated primitive
closed-geodesic lengths are $2L_0$ and $L_0$. The oriented stabilizer of
$e_z=(0,0,1)$ is generated by
\[
 T=\begin{pmatrix}2&3&0\\1&2&0\\0&0&1\end{pmatrix}
\]
and its primitive length is $L_0$.
\begin{proof}
Move a point of the geodesic $v^\perp\cap\mathbb H_F$ into $\mathcal P$.
If $p$ is the resulting point and $B_F(p,v)=0$, the negative-definite
Cauchy--Schwarz inequality on $p^\perp$ gives
\[
 |v_1|^2=|B_F(e_s,v)|^2\leq2(p_1^2-1)\leq25.
\]
For $|s|\leq5$, the norm equation has only the possibilities
\[(s,t,z)=(\pm1,\pm1,0),(\pm4,0,\pm1),(\pm5,\pm3,0).\]
Checking the linear function $s-3tu-18zv$ at the five projective vertices
shows that only $\pm r_2,\pm r_4$ have geodesics meeting the pentagon.
These are exactly its two norm-$-2$ sides. A side reflection negates its
normal; an adjacent side reflection fixes it. Their product is in
$\Gamma^+$, so signs do not split the orbits. Moreover a determinant-$-1$
word mapping a representative to a vector can be multiplied by the
adjacent reflection fixing the representative. Thus passage from the
full reflection group to $\Gamma^+$ does not split either orbit.
The invariant divisibility separates them.

On any side, its two adjacent reflections act on its line by reflections
in the endpoints of a fundamental segment. The line stabilizer is the
resulting infinite dihedral group; its translation subgroup is generated
by their product. This follows directly by moving a point into that
segment in the reflection tiling; an additional stabilizer would identify
two interiors of fundamental segments. For $r_1,r_2,r_4$ the products are,
respectively, $S_0S_2$, $S_1S_3$, $S_3S_0$. Their traces are $5,15,5$.
For a hyperbolic isometry in the three-dimensional Lorentz representation,
$\operatorname{tr}g=1+2\cosh\ell(g)$. Since
$\cosh L_0=2$ and $\cosh(2L_0)=7$, the lengths follow. The source
stabilizer also follows immediately from the Pell equation
$u^2-3v^2=1$ for its upper left $2\times2$ block.
\end{proof}

\paragraph{Statement A.2.3 (Proved).}
For $v=(s,t,z)$ with $F(v)=-2$, $s>0$, $z\geq1$, let
$g_v=v^\perp\cap\mathbb H_F$ and $g_0=e_z^\perp\cap\mathbb H_F$.
They are disjoint and
\[
 \cosh d(g_v,g_0)=3z.
\]
Writing $m_z=s^2-3t^2$, the foot of their common perpendicular on $g_0$
has arclength coordinate
\[
 u=\log\frac{s+\sqrt3t}{\sqrt{m_z}},\qquad
 p(u)=(\cosh u,\sinh u/\sqrt3,0).
\]
In particular $\theta=u/L_0$ modulo $1$, and the arc in
\eqref{A:arc-exact} has length $b_0$ in the source circle of length $L_0$.
If $H_0=\langle T\rangle$ and $H_j=\operatorname{Stab}_{\Gamma^+}(v_j)$,
then, with cyclic stabilizers and both normal orientations retained,
\begin{equation}\label{A:double-coset-exact}
 \sum_{j=1}^2
 \#\{H_0\gamma H_j:\ 3|z(\gamma v_j)|\leq3X\}
 =4\sum_{1\leq z\leq X}\rho(m_z)+2.
\end{equation}
\begin{proof}
The unit spacelike normals have absolute inner product
$|B_F(v,e_z)|/\sqrt{(-2)(-18)}=3z>1$; the hyperboloid formula for
ultraparallel geodesics therefore gives the distance. To verify this
formula directly, restrict to their common two-dimensional Lorentz
plane, in which unit normals have coordinates $(0,1)$ and
$(\sinh d,\cosh d)$.

The minimum along $g_0$ occurs where
$s\sinh u-\sqrt3t\cosh u=0$. Solving gives the stated logarithm;
$T$ translates $u$ by $L_0$. The map
$\Gamma^+/H_j\to\Gamma^+v_j$, $\gamma H_j\mapsto\gamma v_j$,
is bijective by its stabilizer definition. Quotienting on the left by
$H_0$ gives precisely unit orbits of the norm equation in $(s,t)$.
For each $z\geq1$, there are $\rho(m_z)$ with $s>0$; the independent signs
of $z$ and $s$ give four such sets. At $z=0$, the norm-$-2$ equation has
two $H_0$-orbits, according to the sign of $s+\sqrt3t$. These account for
the last $2$ in \eqref{A:double-coset-exact}. This also tracks the
orbifold half-turns: they are not silently included in the cyclic
stabilizers or divided out a second time.
\end{proof}

\subsection{Comparison with the quaternion orders}
\paragraph{Statement A.2.4 (Proved).}
Let $B=(3,-1)_{\mathbb Q}$ with $i^2=3$, $j^2=-1$, $ij=k=-ji$.
Then $B$ has discriminant $6$. The orders
\[
 \mathcal O_{12}=\mathbb Z[1,i,j,k],\qquad
 \mathcal O_6=\mathbb Z[1,i,j,(1+i+j+k)/2],\qquad
 \mathcal O_{108}=\mathbb Z[1,i,3j,3k]
\]
have the indicated reduced discriminants, and $\mathcal O_6$ is maximal.
Under the pure-coordinate map
\[
 V(a,q,n)=\frac{3q+a}{2}i+\frac{3(q+a)}2j+3nk,
 \qquad \operatorname{nrd}V=\frac32F(a,q,n),
\]
the spinor kernel of $\Gamma^+$ is
$\operatorname{Ad}(\mathcal O_{12}^{1})$, of index $2$ in $\Gamma^+$.
The image of the workbench's norm-one order is
$\operatorname{Ad}(\mathcal O_{108}^{1})$, of index $6$ in this spinor
kernel and index $12$ in $\Gamma^+$.
\begin{proof}
The Hilbert symbol $(3,-1)_3$ is $(-1/3)=-1$, and the odd-unit formula
at $2$ gives $(3,-1)_2=(-1)^{((3-1)/2)((-1-1)/2)}=-1$.
At all other finite primes the symbol is $+1$, as at infinity.
Thus the algebra is an indefinite division algebra of discriminant $6$.
Direct multiplication checks that $\mathcal O_6$ is an order and that
its trace-pairing determinant is $-36$; it is maximal since its reduced
discriminant equals the algebra discriminant. The other discriminants
follow either from the trace pairings or from lattice indices $2$ and
$18$ in $\mathcal O_6$.

For comparison with the original presentation set
$I=3i+3j$, $J=i+3j$, $K=6k$. Then $I^2=18$, $J^2=-6$, $IJ=K=-JI$,
and the original basis $(1,(I+J)/2,(I-J)/2,K/2)$ gives exactly
$\mathcal O_{108}$. The map $V$ sends the parity lattice $a\equiv q\pmod2$
onto its trace-zero lattice, but the adjoint norm-one action preserves the
whole diagonal lattice as well. This last assertion, absent from the
workbench bridge, follows by direct multiplication. More generally, for
$g=r+si+cj+dk\in\mathcal O_{12}^1$, the only possible denominator in its
adjoint matrix in the $(a,q,n)$ coordinates is $3$, in a term $rc/3$.
The norm equation implies $r^2+c^2\equiv1\pmod3$, hence $3\mid rc$.
Thus $\operatorname{Ad}(\mathcal O_{12}^1)\subseteq\Gamma^+$.

At $2$, the quotient of the maximal order by its radical is $\mathbb F_4$,
and $\mathcal O_{12}$ is the preimage of $\mathbb F_2$. The norm-one
unit index is $3$: the unit $(1+i+3j+k)/2$ has norm $1$ and its reduction
generates $\mathbb F_4^\times$. At other primes the orders coincide.
Therefore $[\mathcal O_6^1:\mathcal O_{12}^1]=3$.
Voight's Theorem 39.1.2, in the opened version 0.9.14 of
\emph{Quaternion algebras}, states that a maximal order in an indefinite
rational quaternion algebra of discriminant $D$ has projective norm-one
covolume $\frac\pi3\prod_{p\mid D}(p-1)$ \cite{Voight2018}.
Here the hypotheses have just been verified and $D=6$, giving $2\pi/3$.
Thus the adjoint image of $\mathcal O_{12}^1$ has area $2\pi$ and index
$2$ in the already identified group of area $\pi$.

The quaternion $3+k$ has norm $6$ and adjoint matrix $T^{-1}$.
Its spinor class is the non-square $6$, whereas norm-one elements have
trivial class. Hence the index-two subgroup just computed is exactly
the spinor kernel.
Finally $\mathcal O_{108}^1$ is the subgroup of $\mathcal O_{12}^1$
with $c\equiv d\equiv0\pmod3$. The norm-one reductions modulo $3$
have $r^2+c^2=1$ and arbitrary $s,d$, hence size $36$; the indicated
subgroup has size $6$. Surjectivity follows from the reductions of the
actual norm-one units $j$, $2+i$, and $2+k$: the latter two generate the
two additive directions and $j$ generates the four residue units.
This proves index $6$, including the local conditions at both ramified
primes. The exact matrix calculations are reproduced in
\path{code/A_quaternion.py}.
\end{proof}

\subsection{Absence of exceptional spectrum}
\paragraph{Statement A.2.5 (Proved).}
Every compact hyperbolic orbifold of signature $(0;2,2,2,2,2)$,
independently of its moduli, has Cheeger constant at least $1$.
On $M=\Gamma^+\backslash\mathbb H_F$, every nonconstant Laplace
eigenfunction has eigenvalue strictly greater than $1/4$. In particular no
exceptional spectral term obstructs the exponent $2/3$ in either
the unrestricted or the arc-restricted problem.
\begin{proof}
We prove the stronger geometric fact that every compact hyperbolic
orbifold of signature $(0;2,2,2,2,2)$ has Cheeger constant at least $1$.
Here the Cheeger constant is the infimum of
$\operatorname{length}(\partial U)/\operatorname{area}(U)$ over
smooth domains of area at most half the total area. Boundaries can be
taken disjoint from the finite set of cone points, by approximation.

First let $D$ be a disk with a smooth boundary containing at most two
cone points. Its orbifold fundamental group maps into the ambient
Fuchsian group with image either trivial, cyclic of order two, or
generated by two half-turns about distinct points. In the last case
the image is infinite dihedral: the product of the half-turns is a
translation by twice the distance between their centres. The disk
lifts injectively and isometrically to the corresponding complete
orbifold cover. The injectivity of the lift follows immediately from
the identity obtained after composition with the covering map; no
claim about an embedded developed image in $\mathbb H_F$ is needed.

Each of these three complete covers has a smooth orbifold vector
field $V$ with $|V|\leq1$ and $\operatorname{div}V=1$.
For the trivial and finite elliptic cases use polar coordinates about
a fixed point and
$V=\tanh(r/2)\partial_r$ in the metric
$dr^2+\sinh^2r\,d\phi^2$; its divergence is
$1/(\cosh r+1)+\cosh r/(\cosh r+1)=1$.
For the infinite dihedral case use Fermi coordinates about the common
axis, with metric $dr^2+\cosh^2r\,ds^2$, and
$V=\tanh r\,\partial_r$. The field is invariant under translations
and under $(r,s)\mapsto(-r,-s)$, and its divergence is
$\operatorname{sech}^2r+\tanh^2r=1$.
The divergence theorem, or its application after deleting disks of
radius tending to zero about the cone points, now gives
\begin{equation}\label{A:disk-isoperimetry}
 \operatorname{area}(D)\leq\operatorname{length}(\partial D).
\end{equation}
For a nonempty disk the inequality is strict, since its lifted boundary
is compact and the displayed vector field has $|V|<1$ at every finite
point. The later argument only requires the weak inequality for the
Cheeger infimum, but the strict version will exclude eigenvalue $1/4$.


For completeness we pass from disks to arbitrary domains, a step
which matters for the Cheeger constant. The disjoint boundary circles
of $U$ cut the underlying sphere into components; their dual graph is
a tree. Give each vertex its number of cone points, with total weight
$5$. Orient each edge toward the side with at least three cone points.
There is a unique sink: two different sinks would require disjoint
subtrees each containing at least three of the five points. Every
subtree separated from the sink by one edge is a disk with at most
two cone points, so its area is bounded by the length of that edge's
boundary circle, by \eqref{A:disk-isoperimetry}.
If the sink is outside $U$, choose for each connected cluster of
vertices belonging to $U$ its boundary edge nearest the sink. The
associated disk contains the entire cluster. The chosen edges are
distinct, so their inequalities give
$\operatorname{length}(\partial U)\geq\operatorname{area}(U)$.
If the sink is inside $U$, apply the same argument to the complement:
$\operatorname{length}(\partial U)\geq
\operatorname{area}(M\setminus U)\geq\operatorname{area}(U)$.
Thus $h(M)\geq1$.

We also include the analytic implication to specify the convention
and to avoid a hidden orbifold hypothesis. For a smooth mean-zero
real function $f$, choose a median $m$ and put
$f_+=(f-m)_+$, $f_-=(m-f)_+$. Each support has area at most half
the total area. The coarea formula and the preceding isoperimetric
inequality yield
\[
 \int f_\pm^2\leq\int|\nabla(f_\pm^2)|
 \leq2\Big(\int f_\pm^2\int|\nabla f_\pm|^2\Big)^{1/2}.
\]
Adding the resulting inequalities gives
\[
 \int|\nabla f|^2\geq\frac14\int(f-m)^2
 \geq\frac14\int f^2.
\]
Orbifold charts justify coarea and integration by parts by lifting
to disks; isolated cone points have area zero. The Rayleigh
principle gives the weak spectral bound. For a nonconstant
eigenfunction the positive or negative part is nonzero, and its regular
positive level sets have positive area for a set of levels of positive
measure. The strict disk inequality, and the same tree argument, make
coarea strict on those levels. Thus its Rayleigh quotient is strictly
greater than $1/4$. We claim no uniform strict margin over all moduli.
Statement A.2.1
verified that our actual group has exactly this signature.
\end{proof}

\subsection{The unrestricted relative count}
\paragraph{Statement A.3.1 (Known theorem applied with complete bookkeeping).}
For $X\geq2$,
\begin{equation}\label{A:unrestricted-asymptotic}
 S(X)=\sum_{1\leq z\leq X}\rho(3z-1)\rho(3z+1)
 =\frac{9L_0^2}{2\pi^2}X+O(X^{2/3}).
\end{equation}
The implied constant is fixed for the displayed lattice and region.
\begin{proof}
We use Theorem 1.2 of Lekkas--Petridis and its explicitly proved
two-geodesic extension, Remark 6 \cite{LekkasPetridis2026}.
Their hypotheses are a cocompact torsion-free Fuchsian group and
primitive cyclic geodesic stabilizers. Their conclusion is
\[
 N(Y;l,l')=\sum_{1/2<s_j\leq1}\frac2\pi\gamma_1(s_j)
  \Big(\int_lu_j\,ds\Big)\Big(\int_{l'}u_j\,ds\Big)Y^{s_j}
  +O(Y^{2/3}),
\]
where the $u_j$ are orthonormal Laplace eigenfunctions,
$\lambda_j=s_j(1-s_j)$, and
\[
 \gamma_1(s)=\frac\pi2
 \frac{\Gamma(s+1/2)\Gamma((s+1)/2)}{\Gamma(s/2)^3}
 \frac4{s(2s-1)},\qquad \gamma_1(1)=1.
\]
Intersecting axes contribute only finitely many double cosets.
The extension in Remark 6 replaces squared periods by products;
its proof uses the same transform and Cauchy--Schwarz for the two
period sums. We now verify the torsion removal and its exact effect,
instead of applying the torsion-free statement directly to $M$.

Set $K=\ker(\Gamma^+\to\operatorname{GL}_3(\mathbb F_5))$ and
$d=[\Gamma^+:K]$. It is normal of finite index and cocompact.
It is torsion-free: a nonidentity matrix congruent to $1$ modulo $5$
can be written $1+5^eB$ with $e\geq1$ and $B\not\equiv0\pmod5$.
If it has finite order, a power has prime order $p$. For $p\ne5$
the binomial expansion modulo $5^{e+1}$ contradicts order $p$;
for $p=5$ expansion modulo $5^{e+2}$ gives the same contradiction.
Put $H_i'=H_i\cap K$, $m_i=[H_i:H_i']$. These are the full
oriented cyclic stabilizers in $K$ and their primitive lengths are
$m_i l_i$. The values from Statement A.2.2 are
$l_0=L_0$, $l_1=2L_0$, $l_2=L_0$.

For $i=1,2$ no translate of the target axis equals the source axis:
such equality would make an integral norm-$-2$ normal proportional
to $(0,0,1)$, hence equal to $(0,0,\pm1/3)$, which is impossible.
Consequently $H_0\cap\gamma H_i\gamma^{-1}$ is trivial for every
$\gamma$. The natural map
\[
 H_0'\backslash\Gamma^+/H_i'\longrightarrow
 H_0\backslash\Gamma^+/H_i
\]
therefore has exactly $m_0m_i$ preimages of every double coset:
they are represented by the $m_0$ left and $m_i$ right residue
classes, and an equality between two of them would give a
nonidentity element of that intersection. Distance is unchanged.

Decompose $\Gamma^+$ into its $d$ cosets $K\gamma_r$.
The part in $K\gamma_r$ is the pair count in $K$ for $H_0'$ and
$\gamma_rH_i'\gamma_r^{-1}$. All these stabilizers are primitive.
There are no asymptotic pairs of distinct axes: two hyperbolic
elements of a discrete cocompact Fuchsian group sharing one boundary
fixed point share both, and then their axes coincide. To see the
claim directly, conjugate the common endpoint to infinity; distinct
second endpoints and conjugation by powers of one dilation produce
nonidentity group elements tending to the identity. Thus the
regularity condition $abcd\ne0$ in the cited construction holds
outside the finite intersection terms.

Apply the displayed theorem to each of the finitely many pairs in
$K\backslash\mathbb H_F$. Extra small eigenvalues of this cover do
not survive the sum. Indeed the deck group $\Gamma^+/K$ acts
unitarily on every Laplace eigenspace. The sum of the target periods
over $\gamma_r$ is $d$ times the period of the orthogonal projection
onto the deck-invariant subspace. We may choose an orthonormal
eigenbasis respecting that subspace. All other terms cancel exactly.
The invariant eigenfunctions are $d^{-1/2}$ times normalized
eigenfunctions on $M$. Their source and target periods acquire
factors $m_0/\sqrt d$ and $m_i/\sqrt d$, respectively. Summing $d$
cosets and dividing by $m_0m_i$ cancels these factors. By Statement
A.2.5 only the constant eigenfunction has $1/2<s_j\leq1$.
We obtain, for each $i$,
\[
 N_i(Y)=\frac{2l_0l_i}{\pi\operatorname{area}(M)}Y
       +O(Y^{2/3}).
\]
Since $\operatorname{area}(M)=\pi$, summing $i=1,2$, substituting
$Y=3X$, and using the exact identity
$N_1(3X)+N_2(3X)=4S(X)+2$ proves
\eqref{A:unrestricted-asymptotic}.

The theorem is written with a strict cutoff, whereas our count uses
a weak cutoff. Their difference at a height $z$ is at most
$4\rho(m_z)$. The elementary estimate $d(n)\ll_\epsilon n^\epsilon$
therefore makes this $O_\epsilon(X^\epsilon)$, which is absorbed.
For clarity that elementary estimate follows by separating the
finitely many primes $p<2^{1/\epsilon}$; for the remaining primes
$e+1\leq2^e\leq p^{\epsilon e}$, and the finitely many maxima of
$(e+1)p^{-\epsilon e}$ are finite.
\end{proof}

\paragraph{Statement A.3.2 (Proved reformulation and literature comparison).}
Writing $n=3z-1$ gives the exact shifted norm-coefficient sum
\[
 S(X)=\sum_{\substack{1\leq n\leq3X-1\\n\equiv2\pmod3}}
       (1*\chi_{12})(n)(1*\chi_{12})(n+2).
\]
These are the coefficients counting ideals in $\mathbb Q(\sqrt3)$:
$\zeta_{\mathbb Q(\sqrt3)}(s)=\zeta(s)L(s,\chi_{12})$.
This specifies the binary additive-divisor interpretation precisely:
the additive shift $2$ is in the integer norm variable, with the
progression $n\equiv2\pmod3$. It is not a sum over elements of the
quadratic field modulo an unstated unit convention.
For comparison, Motohashi's Theorem 1 and Corollary 1
\cite{Motohashi1994} give, for every fixed positive shift $h$,
\[
 \sum_{n\leq Y}d(n)d(n+h)=Y P_{2,h}(\log Y)
       +O_{h,\epsilon}(Y^{2/3+\epsilon}),
\]
where the quadratic polynomial has leading coefficient
$6\sigma_{-1}(h)/\pi^2$. In our case the main term is linear without a
logarithmic polynomial. The distinction is consistent with the simple
pole of the Dedekind zeta function, as opposed to the double pole of
$\zeta(s)^2$. The cited classical theorem concerns $d(n)$, so it is a
comparison and is not used as a theorem about the character-twisted
coefficients or the sharp ideal arc. Statement A.3.1 proves our formula
from the relative hyperbolic count instead.

\subsection{The sharp arc: an explicit power saving and the uniform transform}
\paragraph{Statement A.4.1 (Proved; weaker than the requested exponent).}
For every $\epsilon>0$,
\begin{equation}\label{A:sharp-7-8}
 R(X)=\frac{9L_0b_0}{2\pi^2}X+O_\epsilon(X^{7/8+\epsilon}).
\end{equation}
In particular this proves the proposed leading constant for the actual
sharp arc. It does not prove the requested exponent $2/3$.
\begin{proof}
Here is a direct application of the general lattice-counting theorem of
Gorodnik--Nevo \cite{GorodnikNevo2012}, with its domain regularity verified
for this arc. Work first in $G=\operatorname{PSL}_2(\mathbb R)$, identified
with the connected isometry group of $\mathbb H_F$. Conjugate the source
axis to the imaginary axis and put
\[
 A_u=\begin{pmatrix}e^{u/2}&0\\0&e^{-u/2}\end{pmatrix},\qquad
 B_r=\begin{pmatrix}\cosh(r/2)&\sinh(r/2)\\
                  \sinh(r/2)&\cosh(r/2)\end{pmatrix}.
\]
For $g=A_uB_rA_v=\begin{pmatrix}a&b\\c&d\end{pmatrix}$,
\[
 ad+bc=\cosh r,\qquad
 u=\frac12\log\left|\frac{ab}{cd}\right|,\qquad
 v=\frac12\log\left|\frac{ac}{bd}\right|.
\]
The parameter $u$ is arclength of the perpendicular foot on the source
axis, with the choice of origin and orientation in Statement A.2.3.
There are four regular cells, according to the sign of $ad+bc$ and the
side/orientation of the perpendicular; each is obtained from the displayed
cell by fixed matrices reversing axes and/or replacing $r$ by $-r$.
Their absolute Haar densities agree. Direct differentiation of
$g^{-1}dg$ gives
\begin{equation}\label{A:Haar-ABA}
 dg=\kappa\sinh r\,du\,dr\,dv\qquad(r>0),
\end{equation}
for a positive normalization constant $\kappa$. The exact symbolic check
is in \path{code/A_relative_coordinates.py}.

Choose $\tau_i\in G$ carrying the standard oriented normal to the
normal of the $i$th target geodesic, for $i=1,2$. Then
$H_i=\tau_i\langle A_{l_i}\rangle\tau_i^{-1}$. An element of
$H_0\backslash\Gamma^+/H_i$ has a unique representative in the coset
$\Gamma^+\tau_i$ with
\[
 0\leq u<L_0,\qquad 0\leq v<l_i,
\]
once its regular cell is specified. Indeed left and right translations
change $u$ and $v$ by the indicated primitive lengths, and their
intersection is trivial, as proved in Statement A.3.1. A single one of
the four cells is the positive-$s$, positive-$z$ ideal count. In that
cell the sharp arc is exactly $0\leq u\leq b_0$.
The lower bound $r\geq1$ removes no positive integral height, since
$\cosh r=3z\geq3$ and $\operatorname{arcosh}3>1$.
Thus, apart from the explicitly treated endpoints, its count is the
number of coset points in
\[
 D_i(X;I)=\{A_u B_r A_v:u\in I,\ 0\leq v<l_i,
              \ 1\leq r\leq\operatorname{arcosh}(3X)\}
\]
in the appropriate regular cell, where $I=[0,b_0]$ or $[0,L_0)$.
Changing which endpoint of a fundamental interval is included selects
one representative and introduces no multiplicity.

We verify Lipschitz well-roundedness, rather than assuming a sector
regularity theorem. On each such domain all four entries have absolute
value between fixed positive multiples of $e^{r/2}$: the ranges of
$u,v$ are fixed compact intervals and $r\geq1$. Left or right
multiplication by a matrix $1+O(\delta)$ changes each entry relatively
by $O(\delta)$, uniformly in $r$. The signs therefore remain fixed, the
two logarithmic formulas change $u,v$ by $O(\delta)$, and the formula
for $\cosh r$ changes $r$ by $O(\delta)$, since $\coth r$ is bounded
for $r\geq1$. The same estimates applied to inverse perturbations give
an inner comparison. Hence the outer and inner $\delta$-perturbations
are bounded by coordinate rectangles whose six faces have moved by
$O(\delta)$. Formula \eqref{A:Haar-ABA} shows
\[
 \operatorname{vol}(D_i^+(\delta))-
 \operatorname{vol}(D_i^-(\delta))
       =O(\delta\operatorname{vol}(D_i)),\qquad
 \operatorname{vol}(D_i)\asymp X.
\]
This is precisely Definition 1.1 of the cited paper with exponent
$a=1$. Small identity balls have upper local dimension $\rho=3$.

The representation on $L^2_0(\Gamma^+\backslash G)$ is $L^{2+}$.
For completeness, the unitary dual classification
\cite[Theorem (2.3)]{Moriyama1999} consists of the trivial representation,
principal series, discrete series and their limits, and complementary
series. Every complementary series has a $K$-fixed vector and corresponds
to a scalar Laplace eigenvalue in $(0,1/4)$, excluded by Statement A.2.5.
The other nontrivial representations are tempered: their $K$-finite
matrix coefficients are in $L^{2+}$ (in polar coordinates the principal
series bound is a constant times $(1+r)e^{-r/2}$ and Haar density is
$O(e^r)$; discrete series are square integrable, with their limits
satisfying the same $L^{2+}$ bound). Finite sums of such vectors are dense.
The constant representation is removed, and connectedness of the quotient
ensures there are no other invariant vectors.

To satisfy the algebraic simply-connected hypothesis in Gorodnik--Nevo,
lift the group, lattice and domains to $\operatorname{SL}_2(\mathbb R)$.
This is a real isotropic simply connected semisimple algebraic group;
the twofold central lift changes point counts and volumes by the same
factor. It does not change the representation on the quotient.
Theorem 4.5 and equation (3.1) of the cited paper, with $p=2$ and
$n_e(p)=1$, give the stable operator bound
$E(X)\ll_\eta\operatorname{vol}(D_i)^{-1/2+\eta}$.
The lattice is discrete and cocompact, so the identity and the finitely
many translates $\tau_i$ have positive injectivity radii. Theorem 1.9
and Corollary 1.11, with $a=1,\rho=3$, now give
\[
 \#(\Gamma^+\tau_i\cap D_i(X;I))
 =\frac{\operatorname{vol}(D_i(X;I))}
        {\operatorname{vol}(\Gamma^+\backslash G)}
    +O_\epsilon(X^{7/8+\epsilon}).
\]
This checks both the translated-coset and sharp-boundary hypotheses.

It remains to normalize the main coefficient. By
\eqref{A:Haar-ABA}, restriction from $[0,L_0)$ to $[0,b_0]$ multiplies
volume by $b_0/L_0$. The four regular cells have equal volume. Their
full count is $4S(X)+O(1)$ by Statement A.2.3, whose leading coefficient
was proved in Statement A.3.1. Consequently the positive-$s$, positive-$z$
cell has leading coefficient $9L_0^2/(2\pi^2)$, and the arc has coefficient
$9L_0b_0/(2\pi^2)$, as asserted.
The lower endpoint consists of the Pell family $t=0$, hence contributes
$O(\log X)$ if its convention is changed. The upper endpoint is empty:
$288s^2=289(18z^2-2)$ would imply $17\mid s$ and then
$144(s/17)^2=9z^2-1$, impossible modulo $9$.
Finally Statement A.1.3 changes the arc count by $O(\log X)$.
\end{proof}


\paragraph{Statement A.4.2 (Proved reduction; its needed uniform estimate remains undecided).}
For each integer $\ell$ define a unitary character on the nonzero ideals
of $\mathbb Z[\sqrt3]$ by
\[
 \psi_\ell((\alpha))=
 \left|\frac{\alpha}{\alpha'}\right|^{i\pi\ell/L_0},
 \qquad a_\ell(m)=\sum_{N\mathfrak a=m}\psi_\ell(\mathfrak a),
\]
where the prime denotes the nontrivial real embedding. Then
\begin{equation}\label{A:twisted-shift}
 T_\ell(X):=\sum_{z\leq X}\sum_{N\mathfrak a=m_z}
 e^{2\pi i\ell\theta(\mathfrak a)}
 =(-1)^\ell\sum_{z\leq X}a_\ell(3z-1)a_\ell(3z+1).
\end{equation}
For every integer $H\geq1$ the actual sharp-arc error satisfies
\begin{equation}\label{A:ET-reduction}
 |R(X)-C_0X|\ll X^{2/3}+\frac XH+
       \sum_{1\leq\ell\leq H}\frac{|T_\ell(X)|}{\ell}
       +\log(2X).
\end{equation}
In particular, the uniform bound
\begin{equation}\label{A:needed-twist}
 |T_\ell(X)|\ll_\epsilon X^{2/3+\epsilon}
           +X^{1/2+\epsilon}\ell^{1/2}
 \quad(1\leq\ell\leq X^{1/3})
\end{equation}
would prove the requested exponent $2/3+\epsilon$.
We have not proved \eqref{A:needed-twist}.
\begin{proof}
Changing the generator by $\pm\varepsilon_0^k$ changes the logarithm
of the absolute embedding ratio by $2kL_0$, and so leaves
$\psi_\ell$ unchanged. The character is multiplicative. On the
positive-norm generator of $m_z$ it is exactly
$e^{2\pi i\ell\theta}$. For the unique prime ideal above $2$,
$|(1+\sqrt3)/(1-\sqrt3)|=\varepsilon_0$, so its character is
$(-1)^\ell$. Away from $2$, the integers $3z-1$ and $3z+1$ are
coprime. At $2$, where they can share powers, there is just one prime
ideal, and its local coefficient is geometric. Combining these facts
proves \eqref{A:twisted-shift}, including the ramified factor.
Conjugation pairs inverse character values, so $a_\ell(m)$ is real.

Apply the Erd\H{o}s--Tur\'an inequality in the form of
\cite[equation (8.2)]{CarneiroVaaler2010} to the finite list of the
$S(X)$ ideal angles. That equation bounds the translate of the sawtooth
sum; the characteristic function of an interval minus its length is
the difference of two translates of that sawtooth function, away from
its endpoints. Thus the arc discrepancy is at most a fixed constant
times $S(X)/H+\sum_{\ell\leq H}|T_\ell(X)|/\ell$, plus endpoint
multiplicities. The upper endpoint is empty by the calculation in
Statement A.4.1. At the lower endpoint, $t=0$ gives
$s^2-2(3z)^2=-2$. Reduction in $\mathbb Z[\sqrt2]$ by its norm-one
unit $3+2\sqrt2$ puts a positive generator in
$[\sqrt2,(3+2\sqrt2)\sqrt2)$, where the equation of norm $-2$
has the sole seed $(s,3z)=(0,1)$. Thus the total number with
$z\leq X$ is $O(\log(2X))$; imposing divisibility of the second
coordinate by $3$ can only decrease it.
Statement A.3.1 gives $S(X)=cX+O(X^{2/3})$ and $S(X)\ll X$.
Statement A.1.3 now proves \eqref{A:ET-reduction}.
Under \eqref{A:needed-twist}, take $H=\lfloor X^{1/3}\rfloor$;
the two sums cost $O_\epsilon(X^{2/3+\epsilon}\log X)$ and
$O_\epsilon(X^{1/2+\epsilon}H^{1/2})$, respectively. Relabeling
$\epsilon$ proves the asserted conditional implication.
\end{proof}

\paragraph{Uniformity audit and current obstruction.}
A bound for each fixed $\ell$ with an unspecified constant depending on
$\ell$ cannot be inserted in \eqref{A:ET-reduction} with $H$ growing
with $X$. More quantitatively, a bound
$T_\ell(X)\ll_\epsilon \ell^A X^{\beta+\epsilon}$, $A>0$, leads by this
argument to exponent $(A+\beta)/(A+1)$, after balancing $X/H$ against
$X^\beta H^A$. Thus even a fixed-mode exponent $2/3$ with polynomial
loss in the mode does not establish the sharp-arc exponent $2/3$.
Theorems 1.2 and 1.5 of Lekkas--Petridis integrate the entire source
geodesic. Their Remark 6 changes the target geodesic, but does not
supply these growing-mode estimates. Voskou's four orientation cases
\cite{Voskou2024} are discrete sign cases, not intervals in the foot
parameter. The fixed compact angular modes in
\cite[Theorem 1.1]{BourgainKontorovichSarnak2010} likewise do not
supply the required uniform hyperbolic-foot estimate.
The direct domain method of Statement A.4.1 proves $7/8+\epsilon$.
Statement A.4.4 supplies a uniform averaged substitute by Mellin
inversion and proves the requested $2/3+\epsilon$ exponent. The
stronger individual-mode bound \eqref{A:needed-twist} is not needed
for that proof and is not asserted here.


\paragraph{Statement A.4.3 (Proved uniform spectral coefficient estimate).}
Let $M$ be a fixed compact hyperbolic surface, let $l_0,l_1$ be fixed
closed geodesics of lengths $L_0,L_1$, and let $u_j$ be an orthonormal
basis of scalar Laplace eigenfunctions with
$\lambda_j=1/4+t_j^2$, $t_j\geq0$. Set $s_j=1/2+it_j$ and
\[
 P_{j,\ell}=L_0^{-1/2}\int_{l_0}u_j(u)e^{2\pi i\ell u/L_0}\,du,
 \quad P_j'=L_1^{-1/2}\int_{l_1}u_j\,du,\quad
 \kappa_\ell=2\pi\ell/L_0,
\]
\[
 D_s(\kappa)=
 \frac{\Gamma((s+i\kappa)/2)\Gamma((s-i\kappa)/2)}{\Gamma(s)},
 \qquad
 Q_{j,H}=\sum_{|\ell|\leq H}b_\ell
             \frac{D_{s_j}(0)}{D_{s_j}(\kappa_\ell)}P_{j,\ell},
\]
where $|b_\ell|\leq B/(1+|\ell|)$ and $H\geq1$. Then, for $T\geq1$,
\begin{equation}\label{A:uniform-arc-period}
 \sum_{T\leq t_j<2T}|Q_{j,H}|^2
       \ll B^2T\log^2(2H),\qquad
 \sum_{T\leq t_j<2T}|P_j'Q_{j,H}|
       \ll BT\log(2H).
\end{equation}
The constants depend on the surface and geodesics, but not on $T,H$
or the displayed coefficients. Thus the angular sum itself can be
controlled without a power of $H$ in this normalized spectral expression.
The uniform inversion that turns this coefficient bound into a sharp-arc
count is proved in Statement A.4.4 below.
\begin{proof}
We first record an elementary broad-band restriction estimate, uniformly
in the real frequency $\kappa$:
\begin{equation}\label{A:broad-band-period}
 \sum_{t_j\leq2T}\left|\int_{l_0}u_j(u)e^{i\kappa u}\,du\right|^2
       \ll T.
\end{equation}
Indeed, in finitely many Fermi charts along the parametrized geodesic,
the one-dimensional fundamental theorem of calculus and Cauchy--Schwarz
give, for $0<\delta\ll1$,
\[
 \|f|_{l_0}\|_2^2\ll
 \delta^{-1}\|f\|_{L^2(M)}^2+
 \delta\|\nabla f\|_{L^2(M)}^2.
\]
The finite number of charts accounts for any self-intersections.
For a linear combination of eigenfunctions with $t_j\leq2T$,
$\|\nabla f\|_2^2\leq(4T^2+1/4)\|f\|_2^2$.
Take $\delta=T^{-1}$ (adjust a fixed constant for small $T$), and
apply the adjoint restriction operator to $e^{i\kappa u}$, whose
$L^2(l_0)$ norm is fixed. This proves \eqref{A:broad-band-period}.
The same estimate applies to $P_j'$.
We also use the elementary bound $\#\{t_j\leq2T\}\ll T^2$.
One proof partitions a fixed finite collection of coordinate patches
into $O(T^2)$ cells of diameter $c/T$ and uniformly bounded shape.
If the spectral space had larger dimension, a nonzero member would
have integral zero on every cell. The cell Poincar\'e inequalities would
give $\|f\|_2^2\leq C(c/T)^2\|\nabla f\|_2^2$, contradicting the
spectral bound when the fixed $c$ is sufficiently small.

Reznikov's mean-value estimate
\cite[equation (1.13), proved in Sections 2.3 and 3.1--3.2]{Reznikov2015}
states that the normalized hyperbolic coefficients $a_{j,\ell}$ satisfy
\begin{equation}\label{A:Reznikov-mean}
 \sum_{|\ell|\leq N}|a_{j,\ell}|^2\ll\max(N,1+t_j),
 \qquad P_{j,\ell}=a_{j,\ell}D_{s_j}(\kappa_\ell).
\end{equation}
A fixed scalar from the model functional in his equation (2.3) has
been absorbed into $a_{j,\ell}$. Explicitly, put his representation
parameter equal to $-2it_j$ and his diagonal-character parameter equal
to $2i\kappa_\ell$. His gamma arguments
$(1-\lambda_{\mathrm{Rep}}\pm s_{\mathrm{char}})/4$ then become
$(s_j\pm i\kappa_\ell)/2$, and his denominator argument
$(1-\lambda_{\mathrm{Rep}})/2$ is $s_j$. The scalar $2/\pi$
is independent of both parameters. The hypotheses are precisely a fixed
closed geodesic on a compact hyperbolic surface and an $L^2$-normalized
principal-series eigenfunction. His argument includes both components
of the diagonal group; retaining the component that gives scalar
periods only decreases its nonnegative sum. The assertion is uniform
in the eigenfunction as well as in $N$.

Split the angular modes at $|\kappa_\ell|=T/2$.
For $t_j\in[T,2T]$ and the lower modes, Stirling's formula in the
fixed vertical strip $\operatorname{Re}(s_j/2)=1/4$ gives
\[
 \left|\frac{D_{s_j}(0)}{D_{s_j}(\kappa_\ell)}\right|\ll1.
\]
For bounded $T$ this follows from continuity and the absence of gamma
zeros or poles; for large $T$ both $t_j\pm\kappa_\ell$ are comparable
to $T$ and the exponential factors cancel exactly.
Minkowski's inequality in the $j$ variable and
\eqref{A:broad-band-period} bound the square root of the lower-mode
sum by
\[
 \sum_{|\kappa_\ell|\leq T/2}|b_\ell|
 \left(\sum_{T\leq t_j<2T}|P_{j,\ell}|^2\right)^{1/2}
       \ll B\sqrt T\log(2H).
\]

For the higher modes divide the integers into dyadic blocks
$N\leq|\ell|<2N$, with $N\gg T$ up to a fixed geodesic-length factor.
On such a block \eqref{A:Reznikov-mean} and Cauchy--Schwarz give
\[
 \left|\sum_{N\leq|\ell|<2N}b_\ell a_{j,\ell}\right|
 \leq\left(\sum|a_{j,\ell}|^2\right)^{1/2}
       \left(\sum|b_\ell|^2\right)^{1/2}\ll B.
\]
Stirling's formula also gives $|D_{s_j}(0)|\ll T^{-1/2}$.
Since the gamma factors cancel exactly in the definition of $Q$,
the higher-mode contribution to each $Q_{j,H}$ is at most
$BT^{-1/2}\log(2H)$. There are $O(T^2)$ eigenvalues in the band,
so its squared sum is $O(B^2T\log^2(2H))$.
Combining the two ranges proves the first estimate in
\eqref{A:uniform-arc-period}; Cauchy--Schwarz with
$\sum|P_j'|^2\ll T$ proves the second.
\end{proof}

The estimate applies after torsion removal to the fixed cover used in
Statement A.3.1. The original angular frequencies are a fixed-index
subsequence of the cover's frequencies, and the constants in
\eqref{A:Reznikov-mean} change only by that fixed index. The next statement derives the corresponding uniform radial transform
and proves its inversion estimates explicitly.

\paragraph{Statement A.4.4 (Proved sharp-arc exponent).}
For every $\epsilon>0$, the actual count in the question satisfies
\begin{equation}\label{A:sharp-two-thirds}
 R(X)=C_0X+O_\epsilon(X^{2/3+\epsilon}).
\end{equation}
The proof uses an angularly uniform Mellin transform. In particular, it
does not deduce an arc estimate from a count over an entire geodesic.

\begin{proof}
We give the transform and all uniform estimates used to invert it.
Put $A_z=(9z^2-1)^{1/2}$ and $\kappa_\ell=2\pi\ell/L_0$.
In this proof periods are unnormalized; this changes the estimates of
Statement A.4.3 only by fixed geodesic-length factors. Choose real
orthonormal eigenfunctions $u_j$ on $\Gamma^+\backslash\mathbb H$ and put
\[
 p_{j,\ell}=\int_{l_0}u_j(u)e^{i\kappa_\ell u}\,du,
 \qquad p_j'=\sum_{i=1}^2\int_{l_i}u_j\,du.
\]
Apart from the constant eigenfunction, write $s_j=1/2+it_j$, with
$t_j>0$, using Statement A.2.5. The constants in the following estimates
may depend on this fixed orbifold and its three geodesics. The spectral
estimates on its torsion-free cover apply to these eigenfunctions:
their lifts are normalized by the square root of the cover degree,
and a period on a lifted geodesic is the original period multiplied by
its fixed covering degree. The angular frequencies form the corresponding
fixed-index subsequence. Thus no spectral gap is assumed on the cover.

\emph{1. An exact Mellin identity.}
For $\operatorname{Re}s>1$ form the absolutely convergent series
\[
 E_s(p)=\sum_{i=1}^2\sum_{H_i\backslash\Gamma^+}
       \cosh\bigl(d(p,\gamma^{-1}l_i)\bigr)^{-s}.
\]
Absolute local uniform convergence follows by placing fixed disjoint
balls around coset representatives in the hyperbolic cylinder:
the number at distance at most $r$ is $O(e^r)$, since a cylinder
of fixed neck length has area $O(e^r)$ in that range. Its derivatives
also converge locally after increasing $\operatorname{Re}s$; the
identities initially obtained there continue to $\operatorname{Re}s>1$.
Set
\[
 K_s(\kappa)=\int_{\mathbb R}(\cosh u)^{-s}e^{i\kappa u}\,du
 =\frac{2^{s-1}\Gamma((s+i\kappa)/2)\Gamma((s-i\kappa)/2)}{\Gamma(s)},
\]
\[
 M_\ell(s)=K_s(\kappa_\ell)^{-1}
                 \int_{l_0}E_s(u)e^{i\kappa_\ell u}\,du.
\]
The beta integral proves the displayed formula for $K_s$.
For a disjoint target axis, the signed distance of a point on $l_0$
from that axis has absolute hyperbolic sine
$A\cosh(u-u_0)$, where $A=\sinh d(l_0,l)$ and $u_0$ is the
perpendicular foot. Consequently its contribution to $M_\ell$ is
\[
 e^{i\kappa_\ell u_0}A^{-s}F_{s,\kappa_\ell}(A),\quad
 F_{s,\kappa}(A)=
 {}_2F_1\left(\frac{s+i\kappa}{2},\frac{s-i\kappa}{2};
                    \frac{s+1}{2};-A^{-2}\right).
\]
To verify this identity without an unquoted transform formula, expand
$(1+A^{-2}\operatorname{sech}^2u)^{-s/2}$ and integrate termwise.
The coefficient is
\[
 \frac{(-1)^k(s/2)_k}{k!A^{2k}}
      \frac{K_{s+2k}(\kappa)}{K_s(\kappa)}
 =\frac{(-1)^k((s+i\kappa)/2)_k((s-i\kappa)/2)_k}
          {k!((s+1)/2)_k A^{2k}},
\]
by the gamma duplication formula. Here $A=A_z\geq\sqrt8$, so the
series and the original integrals converge absolutely.

There are four copies of each positive-height ideal angle: changing
the normal's sign and changing the sign of its $z$ coordinate leave
the foot coordinate unchanged modulo $L_0$. These are the four
regular cells already counted in Statement A.2.3. The two exceptional
double cosets with $z=0$ meet $l_0$ orthogonally at $u=L_0/2$.
Indeed $r_2=(1,1,0)$ is orthogonal to $e_z$ and its intersection
point has $\tanh u=1/\sqrt3$. Their individual contributions are
$e^{i\kappa_\ell L_0/2}=(-1)^\ell$. Thus exactly
\begin{equation}\label{A:Mellin-geometry}
 M_\ell(s)=2(-1)^\ell+4D_\ell(s)+4H_\ell(s),
\end{equation}
where
\[
 D_\ell(s)=\sum_{z\geq1}\sum_{N\mathfrak a=m_z}
          e^{2\pi i\ell\theta(\mathfrak a)}A_z^{-s},
\]
\[
 H_\ell(s)=\sum_{z\geq1}\sum_{N\mathfrak a=m_z}
 e^{2\pi i\ell\theta(\mathfrak a)}A_z^{-s}
                    (F_{s,\kappa_\ell}(A_z)-1).
\]
All sums in \eqref{A:Mellin-geometry} are absolutely convergent
in their initial half-plane. The half-turns reversing an axis do not
belong to its cyclic subgroup $H_i$; they therefore give the two
oriented normal cosets, as required for this factor of four.

Unfolding the scalar spectral coefficient of $E_s$ along $l_i$ gives
\[
 \begin{aligned}
 &\left(\int_{l_i}u_j\right)
 \int_{\mathbb R}(1+x^2)^{-s/2}
 {}_2F_1\left(\frac{s_j}{2},\frac{1-s_j}{2};\frac12;-x^2\right)dx
 \\
 &\qquad=\left(\int_{l_i}u_j\right)
 \frac{\sqrt\pi\Gamma((s+s_j-1)/2)\Gamma((s-s_j)/2)}{\Gamma(s/2)^2}.
 \end{aligned}
\]
Here $x=\sinh r$ in Fermi coordinates. Integrating the eigenfunction
along the full target geodesic solves
$\xi''+\tanh r\,\xi'+\lambda_j\xi=0$; its odd initial-data solution
integrates to zero over $r\in\mathbb R$, and its even solution is
the displayed hypergeometric function. For the beta evaluation put
$y=x^2/(1+x^2)$, apply the power-series identity
\[
 {}_2F_1(s_j/2,(1-s_j)/2;1/2;-x^2)
 =(1+x^2)^{-s_j/2}{}_2F_1(s_j/2,s_j/2;1/2;y),
\]
and integrate the series against
$y^{-1/2}(1-y)^{(s+s_j-3)/2}$. The resulting Gauss beta sum
is the stated gamma quotient. All beta integrals are absolutely
convergent for $\operatorname{Re}s>1$.
It follows that, for $\ell\ne0$,
\begin{equation}\label{A:Mellin-spectral}
 M_\ell(s)=\sum_{j\geq1}p_j'p_{j,\ell}G(s,t_j,\kappa_\ell),
\end{equation}
\[
 G(s,t,\kappa)=
 \frac{\Gamma((s+1)/2)\Gamma((s+s_j-1)/2)\Gamma((s-s_j)/2)}
 {\Gamma(s/2)\Gamma((s+i\kappa)/2)\Gamma((s-i\kappa)/2)}.
\]
The constant eigenfunction vanishes because its nonzero Fourier
period on $l_0$ is zero. For each fixed $\ell$ the series in
\eqref{A:Mellin-spectral} converges normally on compact subsets away
from its poles, since its terms decay exponentially as
$t_j\to\infty$. It therefore supplies the needed meromorphic
continuation. In $\operatorname{Re}s>0$ its only possible poles are
$s=s_j$ and $s=1-s_j$.

\emph{2. A uniform bound for the geometric correction.}
Fix $0<\sigma<1/2$, write $s=\sigma+i\tau$, and put
$U=1+|\tau|$. With $\delta=\operatorname{arsinh}(1/A)$ one has
\begin{equation}\label{A:F-spherical-integral}
 F_{s,\kappa}(A)=
 \frac{\Gamma((s+1)/2)}{\sqrt\pi\Gamma(s/2)}
 \int_{-1}^1
 \bigl(\cosh(2\delta)+x\sinh(2\delta)\bigr)^{-(s+i\kappa)/2}
 (1-x^2)^{s/2-1}\,dx.
\end{equation}
For completeness, this identity follows by differentiating under
the integral and integrating once by parts in $x$: both sides,
as analytic even functions of $\delta$ near zero, solve
\[
 f''+2s\coth(2\delta)f'+(s^2+\kappa^2)f=0,
 \qquad f(0)=1,\quad f'(0)=0.
\]
Uniqueness follows coefficient by coefficient, since
$\operatorname{Re}s>0$; continuation along positive real $\delta$
then proves \eqref{A:F-spherical-integral}. Equivalently, it is
the quadratic hypergeometric transformation with parameters
$a+b=s$, $c=(s+1)/2$.
The base in the integral is positive. Its complex power has modulus
equal to the power $-\sigma/2$, independently of both $\tau$ and
$\kappa$; the weight has modulus $(1-x^2)^{\sigma/2-1}$.
Stirling's formula for the prefactor gives
\[
 |F_{s,\kappa}(A)|\ll_\sigma U^{1/2}.
\]
At $\delta=0$ the first derivative of the integrand is an odd
multiple of $x$, whose integral is zero. Its second derivative is
bounded by a constant times $(1+|s|+|\kappa|)^2$ times the same
integrable absolute weight, uniformly for
$0\leq\delta\leq\operatorname{arsinh}(1/\sqrt8)$.
Taylor's formula with integral remainder therefore gives
\begin{equation}\label{A:F-uniform}
 |F_{s,\kappa}(A)-1|
 \ll_\sigma U^{1/2}
       \min\{1,(U+|\kappa|)^2/A^2\}.
\end{equation}
Statement A.3.1 implies
$\#\{(z,\mathfrak a):A_z\leq V\}\ll V$.
Partial summation in \eqref{A:F-uniform} now proves normal
absolute convergence of $H_\ell(s)$ for $\operatorname{Re}s>0$
and, on our chosen line,
\begin{equation}\label{A:H-uniform}
 |H_\ell(\sigma+i\tau)|
       \ll_\sigma U^{1/2}(U+|\kappa_\ell|)^{1-\sigma}.
\end{equation}
The same convergence argument continues it to any closed strip
in $\operatorname{Re}s>-1$ avoiding the possible hypergeometric
parameter poles, but only the positive strip is needed here.

\emph{3. The remaining Mellin contour, uniformly in the angular cutoff.}
Let $T,H\geq2$. Let $w=w_T$ be a smooth function, equal to one
near zero, supported in $[0,2]$, with its transition confined to
$[1-1/T,1+1/T]$, $0\leq w\leq1$, and
$|w^{(k)}|\ll_kT^k$. Its Mellin transform satisfies, on every
fixed positive vertical strip and for every $N$,
\begin{equation}\label{A:w-Mellin-bound}
 |\widehat w(\sigma+i\tau)|
 \ll_{\sigma,N}U^{-1}(1+U/T)^{-N}.
\end{equation}
This follows by one integration by parts, whose derivative has
bounded $L^1$ norm, and repeated further integrations by parts in
the logarithmic variable. Let $b_\ell$, $1\leq\ell\leq H$, be
arbitrary complex numbers with $|b_\ell|\leq1/\ell$.
We prove
\begin{equation}\label{A:contour-uniform}
 \int_{\mathbb R}|\widehat w(\sigma+i\tau)|
 \left|\sum_{\ell=1}^H b_\ell M_\ell(\sigma+i\tau)\right|d\tau
       \ll_\sigma (T+H)\log(2H).
\end{equation}

Here are details of the gamma estimate and its summation, to make
the frequency uniformity explicit. Write
$p_{j,\ell}=a_{j,\ell}D_{s_j}(\kappa_\ell)$ as in
\eqref{A:Reznikov-mean}. Put
$a=3/4-\sigma/2$ and $b=1/2-\sigma/2$, so $1/2<a<3/4$.
For $\tau,t,\kappa\geq0$ Stirling's formula gives
\begin{align}\label{A:gamma-full-bound}
 |G(\sigma+i\tau,t,\kappa)D_{1/2+it}(\kappa)|
 &\ll_\sigma U^{1/2}
  \{(1+\tau+t)(1+|\tau-t|)\}^{-a}\\
 &\quad\cdot\{(1+\tau+\kappa)(1+|\tau-\kappa|)\}^{b}
 \\ &\quad\cdot\{(1+t+\kappa)(1+|t-\kappa|)\}^{-1/4}
 e^{-c\operatorname{dist}(t,[\min(\tau,\kappa),\max(\tau,\kappa)])}.
 \nonumber
\end{align}
One may take $c=\pi/2$. The exponential follows exactly from
\[
 t-\max(t,\kappa)-\max(t,\tau)+\max(\tau,\kappa)
 =-\operatorname{dist}(t,[\min(\tau,\kappa),\max(\tau,\kappa)]).
\]
The same bound holds with $|\tau|$, $|\kappa|$ for arbitrary signs.
The fixed real parts of the numerator gamma functions are
$\sigma/2-1/4$, which are negative but not nonpositive integers;
thus the estimate is bounded at $\tau=\pm t$.

In a dyadic spectral band $J\leq1+t_j<2J$, the elementary
estimates and Reznikov's bound used in Statement A.4.3 give
\begin{equation}\label{A:three-basic-sums}
 \sum_j|p_j'|^2\ll J,\quad
 \sum_j|p_j'|\ll J^{3/2},\quad
 \sum_j|p_j'p_{j,\ell}|\ll J,
 \quad\sum_{|\ell|\leq K}|a_{j,\ell}|^2\ll\max(K,J).
\end{equation}
The middle bound uses the elementary $O(J^2)$ eigenvalue count.
Changing between $\ell$ and $\kappa_\ell$ changes dyadic bounds
only by fixed constants.

First suppose $\kappa_\ell\leq t_j/2$. Then
$|D_{s_j}(\kappa_\ell)|\asymp J^{-1/2}$.
For $\tau\geq t_j/2$, \eqref{A:gamma-full-bound} implies
\[
 |G(\sigma+i\tau,t_j,\kappa_\ell)|
 \ll_\sigma
 \left(\frac{1+\tau}{1+|\tau-t_j|}\right)^a
 e^{-c(t_j-\tau)_+}.
\]
For $\tau<t_j/2$ there is instead exponential decay in $J$,
with an immaterial polynomial factor. Since $a<1$, integration
over a dyadic $\tau$ block of size $U\gg J$, with the factor
$U^{-1}(1+U/T)^{-N}$, is
$O((1+U/T)^{-N})$; the same assertion holds for $U\asymp J$
because the integrable singular factor is bounded by
$(J/(1+|\tau-t_j|))^a$.
Summing \eqref{A:three-basic-sums} over $\ell$ costs $\log(2H)$.
Summing over $U\gg J$ and then dyadic $J$ costs $O(T)$:
for $J\leq T$ the contribution is
$O(J\log(2T/J))$, whose dyadic sum is $O(T)$, and the remaining
bands decay faster than any power of $J/T$.

Second suppose $\kappa_\ell\geq2t_j$, in an angular block
$K\leq\kappa_\ell<2K$ with $K\gg J$.
Up to $\tau\leq\max(2t_j,\kappa_\ell/2)$,
\eqref{A:gamma-full-bound} is bounded by
\[
 K^{1/2-\sigma}U^{1/2}J^{-a}(1+|\tau-t_j|)^{-a}
 e^{-c(\tau-t_j)_+}.
\]
In the remaining range, the exponential factor is at most
$e^{-c'K}$. For $\tau\gg K$ the polynomial factor is bounded
independently of $\tau$ (the total $\tau$ power in
\eqref{A:gamma-full-bound} is zero). For $\tau\asymp K$ it is
at most a fixed power of $K$. After the coefficient bounds below,
the sum of all these tails is therefore $O(\log(2T))$:
the integral of $U^{-1}(1+U/T)^{-N}$ above a fixed positive
lower limit is $O(\log(2T))$, and the exponential in $K$
absorbs the polynomial factors. This accounts also for
$\tau\geq\kappa_\ell$, where the exponential is
$e^{-c(\kappa_\ell-t_j)}$, rather than
$e^{-c(\tau-t_j)}$.
Cauchy--Schwarz in $\ell$, using
$\sum|a_{j,\ell}|^2\ll K$ and
$\sum_{\ell\asymp K}|b_\ell|^2\ll K^{-1}$, introduces no
additional factor. The elementary beta-integral bound
\[
 \int_0^\infty U^{-1/2}(1+|\tau-t_j|)^{-a}
 e^{-c(\tau-t_j)_+}(1+U/T)^{-N}\,d\tau
 \ll_\sigma J^{-a}\min(J,T)^{1/2}
\]
therefore bounds the spectral band by
\[
 K^{1/2-\sigma}J^{3/2-2a}\min(J,T)^{1/2}
 =K^{1/2-\sigma}J^\sigma\min(J,T)^{1/2}.
\]
Its dyadic sum over $J\ll K$ is $O_\sigma(K)$ for $K\leq T$
and $O_\sigma(\sqrt{KT})$ for $K>T$. Summing $K\ll H$ is
$O_\sigma(T+H)$.

Finally suppose $t_j/2<\kappa_\ell<2t_j$, so $J\ll H$.
Cauchy--Schwarz in $\ell$ gives
\[
 \left|\sum b_\ell a_{j,\ell}
       G(\sigma+i\tau,t_j,\kappa_\ell)D_{s_j}(\kappa_\ell)\right|
 \ll J^{-1/2}
       \left(\sum_{\kappa_\ell\asymp J}|GD|^2\right)^{1/2}.
\]
Squaring \eqref{A:gamma-full-bound} and using
$\sum_{\kappa_\ell\asymp J}(1+|t_j-\kappa_\ell|)^{-1/2}
\ll J^{1/2}$ gives, in the three respective $\tau$ ranges,
\[
 \sum_{\kappa_\ell\asymp J}|GD|^2\ll_\sigma
 \begin{cases}
 U/J,&\tau\leq t_j/2,\\
 J^{2b+1/2}(1+|\tau-t_j|)^{-2a},&t_j/2<\tau<2t_j,\\
 1,&\tau\geq2t_j.
 \end{cases}
\]
After multiplication by $|\widehat w|$, integration in $\tau$,
and summation of $|p_j'|$ over the spectral band, these give
respectively
\[
 O_\sigma(\sqrt{J\min(J,T)}),\qquad
 O_\sigma(J(1+J/T)^{-N}),\qquad
 O_\sigma\bigl(J\log(2+T/J)\min(1,(T/J)^N)\bigr).
\]
For the middle range the exponent is exactly
$b-1/4-a=-1/2$ before the final $J^{3/2}$ sum.
Their dyadic sum over $J\ll H$ is $O_\sigma(T+H)$.
These three cases prove \eqref{A:contour-uniform}. The finitely
many $0<t_j<1$ are included in the band $J=1$; the strict spectral
gap ensures all constants there are finite.

\emph{4. Shifting the contour and bounding the residues.}
For $Y\geq2$ write
\[
 W_\ell(Y)=\sum_{z\geq1}\sum_{N\mathfrak a=m_z}
 e^{2\pi i\ell\theta(\mathfrak a)}w(A_z/Y).
\]
Mellin inversion of \eqref{A:Mellin-geometry} on a line
$\operatorname{Re}s>1$ gives $4W_\ell$ as the inverse Mellin
transform of $M_\ell-2(-1)^\ell-4H_\ell$.
For $H_\ell$ shift to $\operatorname{Re}s=\sigma$ without
crossing a pole. Equations \eqref{A:H-uniform} and
\eqref{A:w-Mellin-bound} show, uniformly in the coefficients,
\begin{equation}\label{A:correction-inversion}
 \left|\sum_{\ell=1}^Hb_\ell\frac1{2\pi i}
 \int_{(\sigma)}\widehat w(s)Y^sH_\ell(s)\,ds\right|
 \ll_\sigma Y^\sigma
     T^{1/2}(T+H)^{1-\sigma}\log(2H).
\end{equation}
Indeed the remaining integral is bounded by
$\int_0^\infty U^{-1/2}(U+H)^{1-\sigma}(1+U/T)^{-N}\,d\tau$.
The exceptional constant $2(-1)^\ell$ contributes
$2(-1)^\ell w(1/Y)$ and hence $O(\log(2H))$ in the angular sum.

Shift the $M_\ell$ contour to the same line. Its residue at $s=s_j$ is
\[
 p_j'p_{j,\ell}\,r_j\frac{D_{s_j}(0)}{D_{s_j}(\kappa_\ell)},\qquad
 r_j=\frac{2\Gamma((s_j+1)/2)\Gamma(s_j-1/2)}{\Gamma(s_j/2)^3},
 \qquad |r_j|\ll(1+t_j)^{1/2}.
\]
The residue at $1-s_j$ has the conjugate gamma factors and obeys
the same estimate. Statement A.4.3 and
\eqref{A:w-Mellin-bound} therefore bound all residues, after summing
the angular coefficients, by
\begin{equation}\label{A:residues-inversion}
 \ll Y^{1/2}\log(2H)
      \sum_{J\text{ dyadic}}J^{1/2}(1+J/T)^{-N}
 \ll Y^{1/2}T^{1/2}\log(2H).
\end{equation}
The lower-half-plane residues use the same proof with conjugated
gamma ratios. No positivity or common phase of the periods is used.
The remaining contour costs
$O_\sigma(Y^\sigma(T+H)\log(2H))$ by
\eqref{A:contour-uniform}.

For clarity, all contour shifts can first be made at finite heights.
The eigenvalue bound $N(V)\ll V^2$ permits heights tending to infinity
whose distance from every $\pm t_j$ is at least a fixed multiple of
$V^{-3}$: the excluded intervals in $[V,V+1]$ have total length
$O(V^{-1})$. On these horizontal segments the gamma formula and
the spectral coefficient bounds give polynomial growth for each
fixed $H$. Taking $N$ arbitrarily large in
\eqref{A:w-Mellin-bound} makes their integrals tend to zero.
Absolute convergence of the residue series follows from
\eqref{A:residues-inversion}. The estimates on the resulting vertical
line are uniform in $H,T$, so this limiting argument makes no
uncontrolled exchange of a growing angular cutoff with a contour limit.

Combining these estimates proves the uniform smoothed inequality
\begin{equation}\label{A:smoothed-angular-sum}
 \left|\sum_{\ell=1}^H b_\ell W_\ell(Y)\right|
 \ll_\sigma\log(2H)\left[
 Y^{1/2}T^{1/2}+Y^\sigma(T+H)
 +Y^\sigma T^{1/2}(T+H)^{1-\sigma}+1\right].
\end{equation}
Because the coefficients $b_\ell$ are arbitrary within
$|b_\ell|\leq1/\ell$, choose their phases to align $W_\ell(Y)$.
Then the left side is exactly
$\sum_{\ell\leq H}|W_\ell(Y)|/\ell$.
This averaged angular estimate is sufficient; it avoids asserting
the stronger, individual-mode estimate \eqref{A:needed-twist}.

\emph{5. Returning to both sharp boundaries.}
The Erd\H{o}s--Tur\'an inequality used in Statement A.4.2 applies
also to a finite list with nonnegative real weights. For rational
weights this follows by replication after clearing denominators;
arbitrary weights follow by approximation and continuity. Hence the
weighted arc discrepancy is at most
\[
 O\left(\frac{W_0(Y)}H+
       \sum_{\ell\leq H}\frac{|W_\ell(Y)|}{\ell}+\log(2Y)\right).
\]
The endpoint term is $O(\log(2Y))$ by the Pell calculation already
proved; the other arc endpoint is empty.
Since Statement A.3.1 gives
$\#\{(z,\mathfrak a):A_z\leq V\}=(c/3)V+O(V^{2/3})$,
partial summation gives
\[
 W_0(Y)=\frac c3Y\widehat w(1)+O(Y^{2/3}),
 \qquad c=\frac{9L_0^2}{2\pi^2},
\]
uniformly for the monotone smooth cutoffs under discussion, because
their total variation is one.

Choose two such cutoffs $w^-,w^+$ with
$0\leq w^-\leq\mathbf1_{[0,1]}\leq w^+\leq1$,
transition intervals respectively $[1-1/T,1]$ and $[1,1+1/T]$.
They bracket the sharp radial count with nonnegative weights and
satisfy $\widehat w^\pm(1)=1+O(1/T)$.
Put $T=H=\lfloor Y^{1/3}\rfloor$ in
\eqref{A:smoothed-angular-sum}. The three nonconstant error terms
are respectively
\[
 O(Y^{2/3}\log Y),\quad
 O_\sigma(Y^{1/3+\sigma}\log Y),\quad
 O_\sigma(Y^{1/2+2\sigma/3}\log Y).
\]
The angular and radial main-term bracketing errors are
$O(Y/H+Y/T)=O(Y^{2/3})$.
Take any fixed sufficiently small $\sigma>0$, and absorb logarithms
into $Y^\epsilon$. Finally put $Y=(9X^2-1)^{1/2}$.
The ideal-arc main term is $(\lambda_0c/3)Y=C_0X+O(X^{-1})$.
Statement A.1.3 changes this by only $O(\log(2X))$, proving
\eqref{A:sharp-two-thirds}.
\end{proof}


\paragraph{Statement A.4.5 (Proved endpoint-strip rigidity).}
Put $b=b_0$, $A=A_z=\sqrt{9z^2-1}$ and
$\delta=\operatorname{arsinh}(1/A)$. For an integral solution
$s^2-3t^2=2A^2$, $s>0$, write
$s=\sqrt2A\cosh u$, $t=\sqrt{2/3}A\sinh u$.
After choosing the lift of $u$ nearest the indicated endpoint,
\[
 |u|\leq\delta\ \Longrightarrow\ t=0,
 \qquad |u-b|\leq\delta\ \Longrightarrow\
 \bigl(4s-17z=0\text{ or }(s,t,z)=(4,0,1)\bigr).
\]
Consequently only $O(\log(2X))$ reduced ideals with $z\leq X$
can lie within these two endpoint strips. These are exactly the
possible endpoint-crossing terms in the spherical correction of
Statement A.4.4; this assertion concerns the geometric support of
that correction, not its complete amplitude.
\begin{proof}
The first implication follows from
$|t|\leq\sqrt{2/3}A\sinh\delta=\sqrt{2/3}<1$.
For the second use
\[
 \cosh b=\frac{17}{12\sqrt2},\qquad
 \sinh b=\frac1{12\sqrt2},\qquad
 \cosh\delta=\frac{3z}{A}.
\]
It follows that
\[
 \sqrt2A\cosh(b\pm\delta)=\frac{17z}{4}\pm\frac1{12}.
\]
Since $b>0$ and $\cosh$ is increasing on $[0,\infty)$, this
bounds $s$ when $\delta<b$, since the entire strip is positive.
In the remaining case $\delta\geq b$, we have $A\leq
1/\sinh b=12\sqrt2$, hence $z\leq5$. Directly for these five
heights, the only integral solutions in the second strip are
$(4,0,1)$ and $(17,1,4)$; equivalently one can note that the negative part lies
in the first strip and has $t=0$, and test $s^2=18z^2-2$ for
$z=1,2,3,4,5$. The only $t=0$ point in this finite range is $(4,0,1)$.
Outside this one stated exception, $|4s-17z|\leq1/3$, and integrality forces
$4s=17z$. The latter equation gives $z=4a$, $s=17a$ and
$a^2-3t^2=-2$, precisely the $n=0$ Pell family already counted.
The first equation gives $s^2-18z^2=-2$, another Pell equation;
its positive solutions form finitely many unit orbits, hence have
$O(\log(2X))$ members up to height $X$. Alternatively, ordinary
reduction by the fixed unit $17+12\sqrt2$ proves this bound directly.

Finally, in \eqref{A:F-spherical-integral} the multiplier of
$e^{i\kappa u}$ is an integral of
$e^{-i\kappa\log B(x)/2}$, with
$B(x)=\cosh(2\delta)+x\sinh(2\delta)$ and $-1\leq x\leq1$.
Since $e^{-2\delta}\leq B(x)\leq e^{2\delta}$, it only shifts
$u$ by a number in $[-\delta,\delta]$. An interval indicator
can therefore change only in the strips just considered. There
is also the scalar factor $B(x)^{-s/2}$: away from the strips
it gives $F_{s,0}(A)$, rather than $1$. No disappearance of that
amplitude correction is claimed.
\end{proof}

\paragraph{Statement A.4.6 (Proved existence of the sharp spectral arc functional).}
For each fixed nonconstant eigenfunction $u_j$, the limit
\[
 Q_{j,I}=\lim_{H\to\infty}
 \sum_{|\ell|\leq H}\widehat{\mathbf1_I}(\ell)
 \frac{D_{s_j}(0)}{D_{s_j}(\kappa_\ell)}p_{j,\ell}
\]
exists. Here $I=[0,b_0/L_0]$, the Fourier coefficient is with
respect to $\theta=u/L_0$, and the periods are unnormalized as
in Statement A.4.4. The limit is a well-defined interval value
of a distribution whose local primitive is H\"older continuous
with exponent $1/2$. This proves existence, not nonvanishing,
of the actual arc coefficient.
\begin{proof}
Conjugate the lift of the source axis to $\{ie^u:u\in\mathbb R\}$
in the upper half-plane. The Poisson--Helgason boundary
representation writes the lifted eigenfunction as
\[
 u_j(z)=\int_{\mathbb R}
       \left(\frac{\operatorname{Im}z}{|z-x|^2}\right)^{s_j}
       T_j(dx).
\]
This is a distributional integral, with the usual additional
chart at infinity. The precise theorem used is the boundary-value
representation in Anantharaman--Zelditch
\cite[Theorem 3.2]{AnantharamanZelditch2012}, together with their
Theorem 7.1 (Otal's boundary regularity theorem): a bounded
hyperbolic Laplace eigenfunction with bounded gradient and
parameter $s$, $\operatorname{Re}s\geq0$, has boundary values
which are the distributional derivative of a
$\operatorname{Re}s$-H\"older function. Our eigenfunction and
its gradient are bounded because the quotient is compact;
Statement A.2.5 gives $\operatorname{Re}s_j=1/2$. The source is
written for compact torsion-free quotients, so apply it to the
fixed cover of Statement A.3.1. The lifted eigenfunction is
unchanged as a function on the universal cover.

For $x=\pm e^v$, define the sum of two boundary distributions
\[
 \beta_j(dv)=(2e^v)^{-s_j}
       \bigl(T_j(d(e^v))+T_j(d(-e^v))\bigr),
\]
where $T_j(d(-e^v))$ denotes pushforward under
$x\mapsto\log|x|$ on the negative half-line, without an
orientation sign. Smooth changes of variable and multiplication
by a smooth function preserve the property of being the
derivative of a locally $1/2$-H\"older function. To see the
multiplication assertion explicitly, if $T=F'$ then
$aT=(aF)'-a'F$, and the second term has a continuously
differentiable primitive. The same formula after a smooth
coordinate change proves the other assertion.

The dilation $x\mapsto e^{L_0}x$ belongs to the source
stabilizer. Uniqueness of the Poisson boundary representation
implies $T_j(e^{L_0}dx)=e^{L_0s_j}T_j(dx)$, so $\beta_j$ is
$L_0$-periodic. Since
\[
 \left(\frac{e^u}{e^{2v}+e^{2u}}\right)^{s_j}
     =(2e^v)^{-s_j}\cosh(u-v)^{-s_j},
\]
its convolution with $\cosh(u)^{-s_j}$ is the restriction
$u_j(ie^u)$. The convolution is legitimate for the periodic
distribution because this kernel and each of its derivatives
decay exponentially. Taking its Fourier coefficients gives
\[
 p_{j,\ell}=K_{s_j}(\kappa_\ell)\,\beta_{j,\ell},
 \qquad K_s(\kappa)=2^{s-1}D_s(\kappa),\qquad
 \beta_{j,\ell}=\int_0^{L_0}e^{i\kappa_\ell v}\beta_j(dv).
\]
Therefore the displayed partial sum defining $Q_{j,I}$ equals
$K_{s_j}(0)$ times the pairing of $\beta_j$ with the symmetric
Fourier partial sum of $\mathbf1_{[0,b_0]}$.

Let $F_j$ be a continuous local primitive of $\beta_j$;
subtracting its affine mean gives a periodic $1/2$-H\"older
function. Its Fourier partial sums converge uniformly. For
completeness, convolving a H\"older function with the Fej\'er
kernel gives uniform error $O_j(H^{-1/2})$ by splitting the
kernel integral at $1/H$; composing the approximation with the
Dirichlet projection of degree $H$ and using its elementary
$L^1$ bound $O(\log(2H))$ gives an error
$O_j(H^{-1/2}\log(2H))$ for the ordinary partial sum (use a
Fej\'er polynomial of degree at most $H$). Integration by parts
then identifies the limit with
\[
 Q_{j,I}=K_{s_j}(0)\bigl(F_j(b_0)-F_j(0)\bigr).
\]
This value is independent of the primitive's additive constant,
and the convergence error is $O_j(H^{-1/2}\log(2H))$.
The implied constant here depends on the fixed eigenfunction;
no exchange with an unbounded spectral sum follows from this
fixed-eigenfunction assertion alone.
\end{proof}

\paragraph{Statement A.4.7 (Proved: exact smooth model for this particular sharp arc).}
This statement concerns the actual arc, and uses its arithmetic endpoints.
Put
\[
 a_*=17,\qquad B_*=12\sqrt2,\qquad
 R_*=\operatorname{arcosh}(17)=2\operatorname{arcosh}(3),
 \qquad b=\log\frac3{2\sqrt2}.
\]
Thus $a_*^2-B_*^2=1$ and $\coth R_*=\cosh b$.
For $\Re s>0$, define
\[
 W_s(u)=(a_*+B_*\cosh u)^{-s},\qquad
 H_s(\kappa)=\int_{\mathbb R}W_s(u)e^{i\kappa u}\,du,
 \qquad m_s(\kappa)=H_s(\kappa)/K_{s,\kappa}.
\]
Choose a real eigenbasis also diagonalizing the reflection involution.
Every mode that contributes to $p'_j$ is reflection-even. For those modes,
\begin{align}\label{A:U-smooth-boundary}
 \beta_j([0,b])
 &=\frac12\int_{\mathbb R}W_{s_j}(u)\,\beta_j(du),\\
 Q_{j,I}
 &=\frac{K_{s_j,0}}{2L_0}
       \sum_{\ell\in\mathbb Z}p_{j,\ell}\,
                    m_{s_j}(\kappa_\ell).
 \label{A:U-smooth-Q}
\end{align}
The first integral is the distributional pairing with a smooth,
exponentially decreasing function; the second series converges absolutely.
In particular no truncation at the sharp endpoints occurs in
\eqref{A:U-smooth-Q}.

The multiplier has the exact formula
\begin{align}\label{A:U-Legendre-multiplier}
 H_s(\kappa)
 &=\sqrt{\frac{2\pi}{B_*}}\,
   \frac{\Gamma(s+i\kappa)\Gamma(s-i\kappa)}{\Gamma(s)}
   P_{-1/2+i\kappa}^{\,1/2-s}(a_*/B_*),\\
 m_s(\kappa)
 &=\frac{2^{s-1/2}}{\sqrt{\pi B_*}}\,
   \Gamma\!\left(\frac{s+1+i\kappa}{2}\right)
   \Gamma\!\left(\frac{s+1-i\kappa}{2}\right)
   P_{-1/2+i\kappa}^{\,1/2-s}(a_*/B_*).
 \nonumber
\end{align}
For $s=1/2+it$, $t\geq0$, and $|\kappa|\geq t$,
\begin{equation}\label{A:U-forbidden-exponential}
 |m_s(\kappa)|\ll
 (1+t+|\kappa|)^{1/2}
       e^{-\pi(|\kappa|-t)/4}.
\end{equation}

For each fixed $\delta>0$ there is also a uniform stationary-phase
description on $|\rho|=|\kappa/t|\leq1-\delta$ as $t\to\infty$.
Set
\[
 q(\rho)=\sqrt{B_*^2+\rho^2},\qquad
 \xi(\rho)=
 \log\frac{q(\rho)+a_*\rho}{B_*\sqrt{1-\rho^2}},
\]
and
\[
 S(\rho)=\log(a_*+q(\rho))
       -\tfrac12\log(1-\rho^2)-\rho\,\xi(\rho).
\]
Then
\begin{equation}\label{A:U-stationary-multiplier}
 m_{1/2+it}(\kappa)
 =q(\rho)^{-1/2}(1-\rho^2)^{1/4}
       e^{-itS(\rho)}\{1+O_\delta(t^{-1})\}.
\end{equation}
Its canonical time and tangential displacement obey the exact identities
\begin{equation}\label{A:U-canonical-time}
 S'(\rho)=-\xi(\rho),\qquad
 S(\rho)-\rho S'(\rho)
   =\operatorname{arsinh}\!\bigl(B_*\cosh\xi(\rho)\bigr)
   \geq R_*.
\end{equation}
These formulas and the exponential forbidden-region estimate do not by
themselves assert nonvanishing of an eigenspace contraction.

\begin{proof}
All reflections represent the same nontrivial coset of
$W/\Gamma^+$. Thus an odd eigenfunction vanishes pointwise on every
reflection mirror, in particular on both target geodesics.
Its two target periods vanish. We may therefore restrict to
$W$-even eigenfunctions.

In the upper half-plane model with source axis $i e^u$, the reflections
in $z=0$ and $t=0$ act on its real boundary as $x\mapsto-x$ and
$x\mapsto1/x$, respectively. Consequently the boundary distribution
$T_j$ in Statement A.4.6 satisfies $T_j(-dx)=T_j(dx)$, and
$\beta_j$ is even in $u$. The element
\[
 U=S_4S_1,\qquad U(x)=\frac{3x+2\sqrt2}{2\sqrt2 x+3}
\]
belongs to $\Gamma^+$, fixes $x=1$, and maps $\infty$ to $e^b$.
The boundary transformation law, obtained directly from the Poisson
kernel and uniqueness of its distributional boundary value, is
$T_j(U\,dx)=|U'(x)|^{s_j}T_j(dx)$.
On the positive boundary, the two equal signs of $x$ give
$\beta_j(du)=2(2e^u)^{-s_j}T_j(d(e^u))$.
Changing variables $x=U(y)$ on the interval $[1,e^b]$ therefore gives
\[
 \beta_j([0,b])
 =\int_0^\infty
 \left(\frac{e^u}
 {(3e^u+2\sqrt2)(2\sqrt2e^u+3)}\right)^{s_j}
 \beta_j(du)
 =\int_0^\infty W_{s_j}(u)\,\beta_j(du).
\]
Evenness gives \eqref{A:U-smooth-boundary}. The change of variables
is first made away from the endpoints. The boundary primitive is
$1/2$-H\"older, so its interval masses are well-defined; the tail
of the transformed pairing tends to zero because $W_s$ and all its
derivatives decrease exponentially for fixed $\Re s>0$.
This justifies passage to the endpoints.
Periodize $W_s$ with period $L_0$ and use
$p_{j,\ell}=K_{s_j,\kappa_\ell}\beta_{j,\ell}$.
The smooth periodic kernel has rapidly decreasing Fourier coefficients,
which proves \eqref{A:U-smooth-Q}; the stronger convergence assertion
also follows from \eqref{A:U-forbidden-exponential}.

For the special-function identity, take
$x=a_*/B_*>1$, $\nu=-1/2+i\kappa$, and $\mu=s-1/2$
in DLMF \cite[equation (14.12.4)]{DLMF}. Its conditions
$\Re(\mu-\nu)>0$ and $\nu+\mu\notin\{-1,-2,\ldots\}$
hold for $\Re s>0$ and real $\kappa$.
Since $x^2-1=B_*^{-2}$, its integral is exactly one half
of $B_*^sH_s(\kappa)$. This gives the first formula.
The gamma duplication formula gives the second.
For $\kappa=0$, equation (14.12.6), with the conversion
from Olver's $\mathbf Q$ to the standard $Q$, gives
$H_s(0)=2Q_{s-1}(17)$. For example
$H_1(0)=2Q_0(17)=2b$, verifying the constant-mode normalization.

For the uniform bound, the analytic continuation of $W_s(u)$
has no singularity in $|\Im u|<\pi$: its first zeros are
at $\pm b\pm i\pi$.
Shift the contour to $\Im u=3\pi/4$ when $\kappa\geq0$,
and to $\Im u=-3\pi/4$ when $\kappa<0$.
For $v=3\pi/4$ and real $u$,
\[
 |a_*+B_*\cosh(u+iv)|^2
   =288\cosh^2u-408\cosh u+145
   \geq25\cosh^2u.
\]
Also $|\arg(a_*+B_*\cosh(u+iv))|\leq v$.
For $u\geq0$ the imaginary part is nonnegative, and
the imaginary part after multiplication by $e^{-iv}$ equals
$-\sin v\,(a_*+B_*e^{-u}\cos v)<0$;
the other signs follow by conjugation. Hence
\[
 |H_{1/2+it}(\kappa)|
 \leq 5^{-1/2}K_{1/2,0}
       e^{-3\pi(|\kappa|-t)/4}.
\]
The vertical connecting segments vanish by exponential decay in
$|\Re u|$. Uniform two-sided Stirling estimates give
\[
 |K_{1/2+it,\kappa}|
 \asymp
 [(1+|t+\kappa|)(1+|t-\kappa|)]^{-1/4}
 e^{-\frac\pi4(|t+\kappa|+|t-\kappa|-2t)}.
\]
Combining the two estimates proves
\eqref{A:U-forbidden-exponential}.

Finally apply stationary phase separately to $H$ and $K$.
The phases are $\rho u-\log(a_*+B_*\cosh u)$ and
$\rho u-\log\cosh u$.
Their unique stationary points are
\[
 u_W=\log\frac{q(\rho)+a_*\rho}{B_*(1-\rho)},
 \qquad u_K=\operatorname{artanh}\rho .
\]
On $|\rho|\leq1-\delta$ they stay in a fixed compact set and
are uniformly nondegenerate. Outside a larger fixed compact set
the phase derivatives are uniformly separated from zero;
integration by parts treats these exponentially decreasing tails.
At the stationary points, the amplitude divided by the square root
of the absolute second phase derivative equals $q(\rho)^{-1/2}$
for $H$, and $(1-\rho^2)^{-1/4}$ for $K$.
The two stationary-phase factors have the same sign and cancel
in the quotient. This proves \eqref{A:U-stationary-multiplier}.
The formulas for $S'$ and the canonical time follow by differentiation.
Explicitly
$\cosh\xi=q/(B_*\sqrt{1-\rho^2})$ and
$\sinh\xi=a_*\rho/(B_*\sqrt{1-\rho^2})$; substitution gives
\eqref{A:U-canonical-time}.
\end{proof}

\paragraph{Scope of this additional identity.}
The axis of $U$ is the mirror $t=0$, whose normal has norm $-3$.
It is not either norm-$-2$ target.
The two roots of the smooth kernel, $-e^{\pm b}$, are endpoints
of $U^{-1}g_0$, a translate of the source geodesic.
The positive boundary component used in the integral is not a
full period fundamental interval for that axis. Thus the identity
does not turn $Q_{j,I}$ into a positive scalar multiple of $p'_j$.
The nonzero complete-eigenspace coefficient required in
Statement A.5.6 is proved separately in Statement A.5.8, using
the outgoing physical kernel and the complete endpoint comparison.


\paragraph{Statement A.4.8 (Proved: exact outgoing physical kernel, including glancing).}
Retain $a_*,B_*,R_*$ from Statement A.4.7 and put
\[
 \tau(u)=\operatorname{arsinh}(B_*\cosh u),\qquad
 f_s(r)=2^{-s}(\cosh r)^{-s}
 {}_2F_1\!\left(\frac s2,\frac s2;s+\frac12;\operatorname{sech}^2r\right).
\]
For $\Re s>0$ the inverse Fourier kernel of $m_s$ is
\begin{equation}\label{A:exact-physical-kernel}
 k_s(u)=-\frac{\Gamma(s)}{\sqrt\pi\,\Gamma(s+1/2)}
          \frac{a_*}{\cosh\tau(u)}f_s'(\tau(u)).
\end{equation}
Here the convention is
$m_s(\kappa)=\int_{\mathbb R}k_s(u)e^{i\kappa u}\,du$.
For every contributing reflection-even eigenfunction the exact sharp
coefficient is therefore
\begin{equation}\label{A:Q-physical}
 Q_{j,I}=\int_{\mathbb R}u_j(i e^u)\,q_{t_j}(u)\,du,
 \qquad q_t(u)=\frac{K_{1/2+it,0}}2 k_{1/2+it}(u).
\end{equation}
The integral is absolutely convergent. As $t\to+\infty$,
\begin{equation}\label{A:physical-principal-symbol}
 q_t(u)=
 \frac{a_*}{2\cosh^{3/2}\tau(u)}
       e^{-it\tau(u)}\{1+O(t^{-1})\},
 \qquad u\in\mathbb R,
\end{equation}
uniformly over the entire real line.
The principal amplitude is positive and decreases like
$e^{-3|u|/2}$. Thus the physical formula remains uniform where
$|\kappa|/t$ tends to one in the Fourier description.

\begin{proof}
The hypergeometric equation gives directly
\[
 f_s''+\tanh r\,f_s'+s(1-s)f_s=0,\qquad
 f_s(r)\sim e^{-sr}\quad(r\to+\infty).
\]
We prove the Fourier identity first for real $s>1$ and then use
analytic continuation. In this proof temporarily allow
$a_*=\cosh R$, $B_*=\sinh R$ with arbitrary $R>0$.
The elementary function $W_s(R,u)$ satisfies
\[
 \left(\partial_R^2+\coth R\,\partial_R+s(1-s)
              -\frac{\partial_u^2}{\sinh^2R}\right)W_s(R,u)=0.
\]
To check the candidate kernel, set $y=\sinh r$ and write $f_s(r)=F_s(y)$.
Then
\[
 (1+y^2)F_s''+2yF_s'+s(1-s)F_s=0.
\]
Consequently $g=F_s'$ satisfies
\[
 (1+y^2)g''+4yg'+\{s(1-s)+2\}g=0.
\]
A direct differentiation shows that
$\cosh R\,g(\sinh R\cosh u)$ satisfies the same two-variable
equation as $W_s$. Its scalar multiple in
\eqref{A:exact-physical-kernel} therefore does so as well.
Fourier transformation in $u$ makes both transforms solutions of
\[
 h''+\coth R\,h'
       +\left(s(1-s)+\frac{\kappa^2}{\sinh^2R}\right)h=0.
\]

We compare their outgoing coefficients. Uniformly under the integrable
majorants needed for Fourier transformation,
\[
 W_s(R,u)\sim e^{-sR}\cosh^{-2s}(u/2),\qquad
 k_s(R,u)\sim
 \frac{s\Gamma(s)}{\sqrt\pi\,\Gamma(s+1/2)}
 e^{-sR}\cosh^{-s-1}u.
\]
The duplication formula gives the exact identity
\[
 \frac{K_{s,\kappa}K_{s+1,\kappa}}
            {2K_{2s,2\kappa}}
   =\frac{\sqrt\pi\,\Gamma(s+1/2)}{s\Gamma(s)}.
\]
It follows that the Fourier transform of $k_s$ and $H_s/K_s$
have the same leading $e^{-sR}$ coefficient.
For $s>1$, the other solution grows like $e^{(s-1)R}$;
the outgoing solution is unique.
Thus their transforms coincide.
All defining integrals are locally holomorphic for $\Re s>0$,
and $K_{s,\kappa}$ has neither zeros nor poles there.
The identity theorem proves \eqref{A:exact-physical-kernel}
throughout that half-plane. Statement A.4.7 and Fourier
periodization now give \eqref{A:Q-physical}.

For the uniform high-frequency assertion, write
$v_t(r)=\sqrt{\cosh r}\,f_{1/2+it}(r)$.
The equation becomes the elementary one-dimensional scattering equation
\[
 v_t''+\left(t^2-\frac14\operatorname{sech}^2r\right)v_t=0,
 \qquad v_t(r)\sim 2^{-1/2}e^{-itr}.
\]
On $r\geq R_*>0$, its outgoing Volterra equation has an
integrable smooth potential $V(r)=\operatorname{sech}^2r/4$.
Successive approximation, using the Green kernel
$t^{-1}\sin(t(r-v))$ and
$\int_r^\infty V(v)\,dv\ll e^{-2r}$, yields
\[
 v_t(r)=2^{-1/2}e^{-itr}
       \{1+O(t^{-1}e^{-2r})\},\qquad
 v_t'(r)=-it\,2^{-1/2}e^{-itr}+O(e^{-2r}).
\]
The constants are uniform for $t\geq1$ and $r\geq R_*$.
Therefore
\[
 f_{1/2+it}'(r)=
 -\frac{it}{\sqrt{2\cosh r}}e^{-itr}\{1+O(t^{-1})\}.
\]
If further finite-order symbol expansions are required, they follow
from the same Volterra equation: after removing $e^{-itr}$,
insert $\sum_{n=0}^N(2it)^{-n}a_n(r)$, with $a_0=1$ and
$a_{n+1}'=a_n''-Va_n$, $a_{n+1}(\infty)=0$.
The residual is $O(t^{-N}e^{-2r})$, and the outgoing Green kernel
gains a further factor $t^{-1}$. All coefficients and their fixed-order
$r$ derivatives are bounded, and the nonconstant coefficients decrease
exponentially. This also verifies that the preceding uniform
estimate is compatible with finite-order stationary-phase calculations.

Finally, Stirling's formula or direct stationary phase in $K$ gives
\[
 \frac{\Gamma(1/2+it)}{\Gamma(1+it)}
       =t^{-1/2}e^{-i\pi/4}\{1+O(t^{-1})\},\qquad
 K_{1/2+it,0}=
       \sqrt{\frac{2\pi}{t}}e^{-i\pi/4}\{1+O(t^{-1})\}.
\]
The factor $it$ from $f_s'$ cancels both phases.
Substitution in \eqref{A:exact-physical-kernel} proves
\eqref{A:physical-principal-symbol}. Since
$\cosh\tau(u)\asymp e^{|u|}$, it also proves the stated decay.
\end{proof}

\paragraph{Connection with the nonvanishing argument.}
The exact physical kernel makes the two pieces of a relative-wave
argument concrete. On bounded $u$ intervals, the phase is the signed
difference of distances to the target geodesic and to $U^{-1}g_0$;
its critical points represent common perpendiculars. At infinity,
these differences approach endpoint Busemann constants.
Those constants can accumulate, so their contribution requires a
quantitative summation before the shortest-perpendicular term
proves nonvanishing. Statement A.5.8 supplies that comparison:
the entire endpoint contribution is the Fourier transform of a
bounded-variation density with finite total variation, and is
of lower order than the positive stationary term.


\subsection{Exact census and square-root-scale evidence}
\paragraph{Statement A.5.1 (Proved by exact exhaustive computation).}
The census in Statement A.1.4 is exhaustive. Moreover the following
finite-range inequality is certified for every real $x$ in its stated
range:
\begin{equation}\label{A:finite-bound}
 |R(x)-C_0x|<\frac{43}{250}\sqrt{x}
 \qquad(10^4\leq x\leq10^9).
\end{equation}
This is a finite theorem, and asserts no bound as $x\to\infty$.
\begin{proof}
The principal census program \texttt{A\_census.cpp} uses independent
$n$-slices and verifies every output in $(s,t,z)$ coordinates. For
$n\geq1$, factor $18n^2-2=2(3n-1)(3n+1)$ with an exact smallest-prime-factor
sieve. In $\mathbb Z[\sqrt3]$, an inert prime of odd exponent rules out a
solution. An inert prime of exponent $2e$ contributes the scalar $p^e$;
the ramified prime $2$ contributes $(1+\sqrt3)^e$; and a split prime of
exponent $e$ contributes the $e+1$ generators
$\pi^k\bar\pi^{e-k}$, $0\leq k\leq e$.
Here $\pi$ is computed as the norm-Euclidean gcd of $p$ and
$r+\sqrt3$, where $r^2\equiv3\pmod p$ is found by the exact
Tonelli--Shanks algorithm. The norm of this gcd is asserted to be $\pm p$.
There is no factor at $3$, since $3n\pm1$ are prime to $3$.
The Euclidean division rounds the two rational coefficients; the strict
decrease of absolute norm is asserted at each division step.

Every product is reduced by exact powers of $2+\sqrt3$, with its real
sign decided by integer comparisons. Since the total norm is positive,
the resulting seeds are precisely
\[
 a+q\sqrt3\in[\sqrt{18n^2-2},\varepsilon_0\sqrt{18n^2-2}),
 \qquad a>2q\geq0.
\]
Unique factorization of ideals and Statement A.1.2 prove that there is
one seed for every unit orbit and no omitted orbit. Equality of distinct
computed seeds is also explicitly ruled out. From each seed, multiply
successively by $\varepsilon_0$ until $4a+17n>X$. These are all solutions
with $a>0,q\geq0$: they correspond to all generators at least
$\sqrt{18n^2-2}$, and the positive coefficient $a$ increases strictly
under this iteration. The $n=0$ family is generated separately from
$1+\sqrt3$ and the strict inequality rejects its first member.
The last possible $n$ is found by exact binary search using
$4\lceil\sqrt{18n^2-2}\rceil+17n\leq X$; at $X=10^9$ it is
$29\,437\,251$. Integer square roots are corrected by exact comparisons,
so floating-point initial guesses do not decide any predicate.

For every emitted point the code checks $F(s,t,z)=-2$,
$t+4z<s$, $4s\leq17z$, the inverse parity, and the original equation
$4(xy+xz+yz)-x^2-y^2-z^2=-1$ in signed 128-bit arithmetic.
Its sorted height list is compared literally with a direct scan of
$s$ at every $z\leq10^6$. A second program,
\texttt{A\_height\_census.cpp}, independently slices by every height
$z\leq10^7$, enumerates ideals of $18z^2-2$, tests the moving region
and the fixed arc $288s^2\leq289(18z^2-2)$, and verifies
\eqref{A:arc-exact} at every height. Its full ordered height list agrees
with the $n$-slice program. This is the required independent recount in
$(s,t,z)$ through $10^7$, and not an extrapolation of the literal scan.

The $10^9$ census contains $35\,363\,970$ solutions at
$23\,243\,639$ occupied heights. The maximum multiplicity is $52$,
first attained at $971\,256\,191$. The raw sorted little-endian
32-bit height receipt has SHA256
\begin{center}\small\ttfamily
29bc345c303f0f1885b13f459669411fc4202\\
a43d3ed95509844f8fac2c65a48
\end{center}
The lossless file \texttt{A\_census\_1e9.d16m8.gz} stores successive
height gaps as 16-bit integers and multiplicities as 8-bit integers;
the maximum gap is $601$. The attached pack/unpack program reconstructs
the raw file and verifies this hash.

To certify \eqref{A:finite-bound}, the Arb computation
\path{A_certify_constant.py}, at 256-bit precision, proves
\[
 \frac{35362047799}{10^{12}}<C_0<
 \frac{35362047800}{10^{12}}.
\]
For both rational endpoints $c/10^{12}$, at every jump height $z$ we
check both its right value and its left limit by the integer inequality
\[
 62500\big(10^{12}R-cz\big)^2
       <1849z\,10^{24}.
\]
The program checks a strict size guard before squaring, which keeps all
products inside signed 128-bit range. It also checks the endpoints
$10^4,10^9$. Convexity in $C_0$ gives the same bound for the enclosed
constant. Between consecutive jumps $R(x)=r$ is constant and
\[
 \frac{d}{dx}\frac{r-C_0x}{\sqrt x}
 =-\frac{r}{2x^{3/2}}-\frac{C_0}{2\sqrt x}<0.
\]
Thus the two one-sided jump values and the external endpoints control
the absolute value throughout every interval. All checks pass, proving
the assertion for real, as well as integer, $x$.
\end{proof}

\paragraph{Statement A.5.2 (Proved finite computations; descriptive evidence with its scope specified).}
On integer $X\leq10^8$ the extremes in Lead 5 are confirmed:
\[
 E(32\,461\,294)=-812.8300691488068897\ldots,
 \qquad E(77\,078\,404)=721.7934249571297570\ldots,
\]
where $E(X)=R(X)-C_0X$. On integer $X\leq10^9$ the extrema are
\[
 \begin{split}
 E(889\,576\,590)&=-2332.8971024648566233\ldots,\\
 E(586\,293\,324)&= 2668.4520483254967593\ldots.
 \end{split}
\]
For $10^4\leq X\leq10^9$ the largest absolute standardized discrepancy
at an integer is $0.1714001928088145597\ldots$, attained at
$X=867\,240$. These decimal evaluations use 80-digit arithmetic after
exactly locating the integer candidates; the finite bound itself is the
separate rational certificate in Statement A.5.1.

The deterministic sample $X=10^4+1000k\leq10^9$, with $999\,991$
observations, has mean $0.00067195$, standard deviation $0.04345364$,
skewness $0.21623$, and excess kurtosis $-0.71874$ for $E(X)/\sqrt X$.
Its median is $-0.00410075$, its $5\%$ and $95\%$ quantiles are
$-0.06359556$ and $0.07133799$, and its lag-one correlation is
$0.99986553$. Thus the sample has strong serial dependence; these
statistics do not constitute evidence of independent normal sampling,
and no limiting law is inferred.
The complete sample and its generating program are attached.

For the continuous mean-square integral, exact constancy of $R(x)$
between successive jumps reduces each interval of length $g$, beginning
at $a$ with $E(a)=e$, to
\[
 \int_a^{a+g}E(x)^2\,dx=e^2g-eC_0g^2+\frac{C_0^2g^3}{3}.
\]
Summation in long-double arithmetic gives the following descriptive
values; they are not interval-certified decimal values:
\[
\begin{array}{c|c}
 X&X^{-2}\int_X^{2X}|E(x)|^2\,dx\\\hline
 10^4&0.0021358981\\
 10^5&0.0057996192\\
 10^6&0.0048290724\\
 10^7&0.0091159753\\
 10^8&0.0027564219\\
 5\cdot10^8&0.0031277302
\end{array}
\]
The asserted finite computations are reproducible from
\path{A_statistics.py} and its output. No global mean-square,
$\Omega$, or $\Omega_\pm$ theorem follows from this finite table.

\paragraph{Statement A.5.3 (Proved weaker mean-square bound).}
For every $\epsilon>0$,
\[
 \int_X^{2X}|R(x)-C_0x|^2\,dx\ll_\epsilon X^{7/3+\epsilon}.
\]
\begin{proof}
Apply Statement A.4.4 with $\epsilon/2$ uniformly for
$x\in[X,2X]$, square it, and integrate over an interval of length $X$.
\end{proof}

\paragraph{Statement A.5.4 (Open: remaining pointwise conjecture and literature scope).}
The exact square-root conjecture for this count is
\[
 \text{for every }\epsilon>0\text{ there are }C_\epsilon,X_\epsilon>0
 \text{ such that }|R(X)-C_0X|\leq C_\epsilon X^{1/2+\epsilon}
 \quad(X\geq X_\epsilon).
\]
The requested mean-square bound is proved in Statement A.5.5.
Statement A.5.8 proves both
$R(X)-C_0X=\Omega_+(X^{1/2})$ and
$R(X)-C_0X=\Omega_-(X^{1/2})$, and hence the unsigned assertion.
The displayed pointwise square-root conjecture remains unproved,
and we have not established an implication from a numbered
published conjecture to it.
The numerical work above settles the finite census and tests the
standardized discrepancy; the infinite oscillation assertions
are settled by the separate analytic proof.

\emph{An exact sufficient arithmetic hypothesis.}
Retain the ideal characters and coefficients $a_\ell(m)$ of
Statement A.4.2, and put
\[
 U_\ell(X)=\sum_{1\leq z\leq X}a_\ell(18z^2-2).
\]
The following uniform estimate would imply the displayed
square-root conjecture:
\begin{equation}\label{A:sufficient-uniform-quadratic}
 \begin{aligned}
 \left|U_\ell(X)-\mathbf1_{\ell=0}\,cX\right|
 &\ll_\eta X^{1/2+\eta}(1+|\ell|)^\eta,\\
 &0\leq|\ell|\leq\lfloor X^{1/2}\rfloor,\qquad \eta>0.
 \end{aligned}
\end{equation}
The implied constant must be independent of both $X$ and $\ell$.
This is a precisely stated sufficient hypothesis, not a conjecture
attributed to the literature. To prove the implication, view the
ideals of norm $18z^2-2$, $z\leq X$, as a positive measure on
$\mathbb R/\mathbb Z$. Its mass is $U_0(X)$ and its $\ell$th
Fourier coefficient is exactly $U_\ell(X)$. Weighted
Erd\H{o}s--Tur\'an, or the Selberg majorants and minorants already
used in Statement A.4.4, gives for $H=\lfloor X^{1/2}\rfloor$
\[
 \left|\sum_{z\leq X}\#\{\mathfrak a:
 N\mathfrak a=18z^2-2,\ \theta(\mathfrak a)\in I\}
       -\lambda_0 U_0(X)\right|
 \ll \frac{U_0(X)}{H+1}
       +\sum_{\ell=1}^{H}\frac{|U_\ell(X)|}{\ell}.
\]
Endpoint conventions change only the already counted Pell
families, hence contribute $O(\log X)$.
Under \eqref{A:sufficient-uniform-quadratic} the right side is
$O_\eta(X^{1/2+3\eta/2}\log(2X))$.
Take $\eta=\epsilon/3$, absorb the logarithm into the remaining
power, and use the exact formula in Statement A.1.3 and
$C_0=\lambda_0c$. This proves the claimed implication. Equivalently,
the nonzero modes in \eqref{A:sufficient-uniform-quadratic} are
\[
 U_\ell(X)=(-1)^\ell
       \sum_{z\leq X}a_\ell(3z-1)a_\ell(3z+1);
\]
thus uniformity in a family of real-quadratic ideal characters
growing as fast as $|\ell|\leq X^{1/2}$ is essential.

The closest explicit shifted-convolution expectation located,
\cite[equation (2) and the paragraph immediately after it]
{NordentoftPetridisRisager2022}, concerns a fixed holomorphic
level-one form and expressly leaves its weight-uniform range
unspecified. It does not state
\eqref{A:sufficient-uniform-quadratic}.
The varying-form Conjectures 1.6 and 1.8 of
\cite{Nelson2022Quadratic} assume non-dihedral holomorphic forms
and irreducible quadratic polynomials, and ask for weaker
cancellation; Remark 1.9 says that variants for dihedral forms
or reducible polynomials still need to be formulated.
Our character family is dihedral and
$18z^2-2=2(3z-1)(3z+1)$ is reducible.
Finally, \cite[Conjecture $C_a(\epsilon)$]{Merikoski2025Primes}
has explicit epsilon dependence on the real-quadratic character
parameter, but its sum runs over the two-variable binary cubic
$c_1y_1^3-c_2y_2^3$. No specialization of that conjecture to
\eqref{A:sufficient-uniform-quadratic} is proved here.

A direct delta-method attempt using the theta-corrected
half-integral Kloosterman-sum Conjecture 1.4 of
Sun \cite{A5Sun2025} was tested on its full analytic weights,
including their inverse-phase amplitudes.  The calculation in
\path{records/B_A5_compensated_weight.tex}, with its exact and
Arb certificates, uses a cutoff supported near the actual quadric
and the dual vectors $(1,0,k)$.  On the fixed modulus block
$2X/3\leq q\leq3X/4$, their normalized weights satisfy
$\|W_k\|_\infty\asymp k^{-1}$ but
$\operatorname{Var}W_k=L+O(k^{-1})$ for an explicit $L>0$.
The variation summed over $k\leq H$ is therefore $\asymp H$,
also for smoothing of width $H^{-1}$.
Consequently, applying the proposed untwisted bound separately
to each frequency by Abel summation and then summing absolutely
has the scale $X^{1/2+\epsilon}H$ on this family.
Balancing this with a boundary cost $X/H$ gives a $3/4$
envelope for that procedure.  This is neither a full $3/4$
theorem nor a proof that Sun's conjecture cannot imply the
desired result by another argument.  It identifies the precise
remaining loss; cancellation between dual frequencies or an
arithmetic transformation retaining their joint phases is still
needed.  The multiplier, bad-prime factors, and complete theta
projection must also be matched before that conjecture is applied.

Completing the normal frequency before summing over moduli gives
a further exact reduction, recorded in
\path{records/A_normal_completion.tex}.
For the same family $(1,0,k)$, a finite quadratic Fourier
transform converts the Sali\'e sums to
$S(2-18h^2,-\overline4;q)$, and Chinese remaindering gives
\[
 S(2-18h^2,-\overline4;q)
       =-S(2-18h^2,-1;2q).
\]
The endpoint phase centers the new index at $h=X$ with width
$X/K$ for an original dyadic block $k\asymp K$.
The bounded-twist Linnik--Selberg conjecture proposed by
Browning--Kumaraswamy--Steiner and stated as
\cite[Conjecture 2]{A5Steiner2019} does apply to each such
completed modulus sum: its fixed parameters can be taken
as $s=24$ and $B=3$. Its resulting bound, summed absolutely
over the new index, is $O_\epsilon(X^{1+\epsilon}/K)$
for this top-face family when $X^\delta\leq K\leq X^{1/2}$,
with any fixed $\delta>0$. Before completion the proposed
Sun estimate gives $O_\epsilon(X^{1/2+\epsilon}K)$.
Taking the better bound still leaves the intermediate
block $K=X^{1/4}$ at scale $X^{3/4+\epsilon}$.
The two exact Fourier identities have been checked in
200 cyclotomic-ring computations, including composite
and prime-power moduli. This supplies a verified application
of a named conjecture to an isolated transformed family;
it supplies no implication to the full sharp square-root
bound. The missing estimate is joint cancellation over
$h,q$, with $2-18h^2$ running over a short quadratic sequence.

The exact scope of the closest sources matters. Theorem 1.5 of
Lekkas--Petridis \cite{LekkasPetridis2026} proves
$\int_X^{2X}|E(x,l)|^2\,dx\ll X^2\log^2X$ for the count integrating
an entire closed geodesic; its Conjecture 1.6 is the square-root
conjecture for that full-geodesic error. Neither statement includes a
fixed sharp foot arc. In \cite[Theorem 3]{Voskou2025Omega}, the
$\Omega_\delta(X^{1/2}(\log\log X)^{1/4-\delta})$ theorem concerns a
full geodesic paired with itself, with the eigenvalue-$1/4$ contribution
removed. Its squared periods supply a positivity absent from the
cross-geodesic count with a source arc. It does not assert
$\Omega_\pm$ for the present count. Statements A.4.4 and A.5.5 prove
the uniform spectral estimates and the requested mean square.
Statement A.5.6 further proves an absolutely convergent
almost-periodic expansion after fixed positive radial smoothing
of the actual sharp arc. It proves both oscillation signs conditional
on a precisely stated nonzero complete-eigenspace coefficient.
Statement A.5.8 proves that condition, indeed for infinitely many
spectral parameters, using the exact arithmetic endpoints,
the outgoing physical kernel, a summable endpoint density, and
the positive shortest-perpendicular contribution.
The square-root pointwise conjecture and its implication from
a numbered known conjecture remain open here.

Two recent results do not remove this gap.
Constantinescu--Nordentoft \cite[Theorems 1.3 and 1.5]{ConstantinescuNordentoft2025}
prove nonvanishing while varying the closed geodesic for a fixed
automorphic form; their general theorem assumes a cusp.
Chatzakos--Dimakis \cite[Theorem 1.2]{ChatzakosDimakis2026} prove
square-root cancellation, up to logarithms, for almost every
point paired with a fixed point on a compact quotient.
Their Section 1.6 explicitly excludes a conclusion for a
preassigned pair. Neither theorem supplies the fixed-cycle,
fixed-arc contraction considered here.


\paragraph{Statement A.5.5 (Proved sharp-arc mean square).}
For every $\epsilon>0$,
\begin{equation}\label{A:sharp-mean-square}
 \int_X^{2X}|R(x)-C_0x|^2\,dx\ll_\epsilon X^{2+\epsilon}.
\end{equation}
\begin{proof}
We extend the uniform transform of Statement A.4.4 to a negative
Mellin line and then apply the geodesic large sieve to its actual
arc coefficients. This checks the additional step needed to use that
large sieve here.

\emph{1. The geometric correction on a negative line.}
In \eqref{A:F-spherical-integral} put
$p=(s+i\kappa)/2$, $B(x)=\cosh(2\delta)+x\sinh(2\delta)$,
and $f(x)=B(x)^{-p}$, where $\delta=\operatorname{arsinh}(1/A)$.
For $\operatorname{Re}s>0$ the identity
\[
 (1-x^2)^{s/2-1}
 =\frac{s+1}{s}(1-x^2)^{s/2}
       -\frac1s\frac d{dx}\{x(1-x^2)^{s/2}\}
\]
and integration by parts give
\begin{equation}\label{A:F-regularized}
 F_{s,\kappa}(A)=
 \frac{\Gamma((s+3)/2)}{\sqrt\pi\Gamma((s+2)/2)}
 \int_{-1}^1(1-x^2)^{s/2}
       \left(f(x)+\frac{x f'(x)}{s+1}\right)dx.
\end{equation}
The right side converges absolutely and is analytic for
$\operatorname{Re}s>-2$, apart from its displayed parameter
singularities. Therefore it continues the left side to the range
$-1<\operatorname{Re}s<0$ used below.

Fix $-1<\sigma<0$, write $s=\sigma+i\tau$,
$U=1+|\tau|$ and $V=U+|\kappa|$. The gamma prefactor is
$O_\sigma(U^{1/2})$, $|s+1|^{-1}\ll_\sigma U^{-1}$,
and $|f'(x)|\ll_\sigma V/A$. Hence
\begin{equation}\label{A:F-negative-large}
 |F_{s,\kappa}(A)-1|
       \ll_\sigma U^{1/2}\left(1+\frac{V}{UA}\right).
\end{equation}
If $A\geq V$, differentiate the integrand in
\eqref{A:F-regularized} twice with respect to $\delta$ on
$0\leq\delta\leq\operatorname{arsinh}(1/A)$.
Its first derivative at zero is odd in $x$, so it integrates to
zero. Its second derivative is $O_\sigma(V^2)$ times the
integrable absolute weight: the possible third power of $p$
from $f'$ is multiplied by $\delta/(s+1)$ and is bounded by
$O_\sigma(V^2)$ because $V\delta\leq1$.
Taylor's integral remainder therefore gives
\begin{equation}\label{A:F-negative-small}
 |F_{s,\kappa}(A)-1|\ll_\sigma U^{1/2}V^2/A^2
 \qquad(A\geq V).
\end{equation}
Using the count $\#\{A_z\leq y\}\ll y$ with its ideal
multiplicities, sum \eqref{A:F-negative-large} over $A_z\leq V$
and \eqref{A:F-negative-small} over $A_z>V$.
The term $V/(UA)$ is summable with total
$O_\sigma(U^{-1/2}V^{1-\sigma})$ because $\sigma<0$;
the other two terms are
$O_\sigma(U^{1/2}V^{1-\sigma})$ because $\sigma>-1$.
Thus precisely the same estimate as \eqref{A:H-uniform} holds:
\begin{equation}\label{A:H-negative}
 |H_\ell(\sigma+i\tau)|
       \ll_\sigma U^{1/2}(U+|\kappa_\ell|)^{1-\sigma}.
\end{equation}
The series defining $H_\ell$ converges normally in this strip,
since its summands at large $A_z$ have two additional inverse
powers of $A_z$.

At $s=0$ the exact hypergeometric identity is
\[
 F_{0,\kappa}(A)=\cos\bigl(\kappa\operatorname{arsinh}(1/A)\bigr).
\]
It follows either by solving the differential equation in the proof
of \eqref{A:F-spherical-integral} at $s=0$, or by its power
series. In particular,
\begin{equation}\label{A:H-at-zero}
 H_0(0)=0,\qquad |H_\ell(0)|\ll1+|\ell|.
\end{equation}
For the second assertion sum
$\min(2,\kappa_\ell^2/A_z^2)$ using the same linear counting
bound. No estimate for the distribution of individual ideal angles
is being assumed in this step.

\emph{2. A uniform spectral expansion with square-root-size remainder.}
Fix a large base height $Y_0$, and take
$T=H=\lfloor\sqrt{Y_0}\rfloor$, fixed as $Y$ varies over
$[Y_0,2Y_0]$. Let $w=w_T$ be either smooth radial cutoff of
Statement A.4.4. Its Mellin transform continues meromorphically by
\[
 \widehat w(s)=-\frac1s\int_0^\infty w'(u)u^s\,du.
\]
It has a simple pole at zero of residue one and satisfies
\eqref{A:w-Mellin-bound} on any fixed nonzero real-part line.
Choose $\sigma=-1/2+\eta$, where $0<\eta<1/2$ is fixed.

The contour estimate \eqref{A:contour-uniform}, with $T=H$,
remains valid on this line. Here are the precise changes to its
proof. The gamma bound \eqref{A:gamma-full-bound} is unchanged,
with $a=1-\eta/2<1$ and $b=3/4-\eta/2$; its numerator gamma
functions still have fixed real parts strictly between $-1/2$
and $-1/4$, and so have no pole on the line. The low-angular and
comparable-angular integrations use only $a<1$ and remain
identical. In the high-angular range, $K\ll H=T$ and the sum is
\[
 K^{1/2-\sigma}\sum_{J\ll K}J^{\sigma+1/2}
       \ll_\eta K,
\]
because $\sigma+1/2=\eta>0$. The exponentially small tails
remain $O(\log(2T))$. These facts give
$O_\eta(T\log(2H))$ for the remaining contour before multiplication
by $Y^\sigma$. All constants are uniform in $Y\in[Y_0,2Y_0]$.
Allowing negative angular indices changes only a factor of two.

We also include angular index zero. Its nonconstant spectral
part obeys the same estimates, using the low-angular range.
The constant eigenfunction contributes
\[
 p'_0p_{0,0}
 \frac{\Gamma((s+1)/2)\Gamma((s-1)/2)}{\Gamma(s/2)^2}.
\]
On the negative contour this is $O_\eta(1)$ in $\tau$, and
its only pole in the traversed strip is at $s=1$.
Its residue in $M_0$ is $4c/3$, where
$c=9L_0^2/(2\pi^2)$. This follows also directly from
\eqref{A:Mellin-geometry} and the already proved
$\sum_{A_z\leq Y}\rho(m_z)=(c/3)Y+O(Y^{2/3})$;
$H_0$ is analytic at one.

There is no contribution at zero from $\widehat w(s)M_\ell(s)$:
the normally convergent gamma series shows $M_\ell(0)=0$
for every $\ell$. For $\ell\ne0$, the factor $1/\Gamma(s/2)$
gives this zero; for $\ell=0$ the nonconstant terms have three
such factors, and the displayed constant term has two.
In the geometric correction the pole of $\widehat w$ does give
$H_\ell(0)$, whose total against coefficients
$|b_\ell|\ll(1+|\ell|)^{-1}$ is $O(H)$ by
\eqref{A:H-at-zero}. The two exceptional double cosets contribute
only $O(\log(2H))$.
Equation \eqref{A:H-negative} bounds the negative-contour
correction by
\[
 O_\eta\left(Y_0^\sigma
       T^{1/2}(T+H)^{1-\sigma}\log(2H)\right)
 =O_\eta\left(Y_0^{1/2+\eta/2}\log Y_0\right).
\]
The remaining $M$ contour costs only
$O_\eta(Y_0^\sigma T\log(2H))$.
The finite-height contour justification from Statement A.4.4
continues without change; the next spectral poles have real part
$-3/2$, outside the traversed strip.

Consequently, for any real trigonometric polynomial
$B(\theta)=\sum_{|\ell|\leq H}b_\ell e^{2\pi i\ell\theta}$
with $|b_\ell|\ll(1+|\ell|)^{-1}$, the weighted count
\[
 \mathcal N_{w,B}(Y)=\sum_z\sum_{N\mathfrak a=m_z}
                  w(A_z/Y)B(\theta(\mathfrak a))
\]
has the expansion, uniformly on $[Y_0,2Y_0]$,
\begin{align}\label{A:sqrt-remainder-expansion}
 \mathcal N_{w,B}(Y)&=\frac{b_0c}{3}\widehat w(1)Y
 \\ &\quad+\frac14\sum_{j\geq1}p'_j
  \left(r_j\widehat w(s_j)Q_{j,H}Y^{s_j}
       +\overline{r_j\widehat w(s_j)Q_{j,H}}Y^{1-s_j}\right)
 \\ &\quad+O_\eta\left(Y_0^{1/2+\eta/2}\log Y_0\right).
 \nonumber
\end{align}
Here $Q_{j,H}=\sum b_\ell D_{s_j}(0)
p_{j,\ell}/D_{s_j}(\kappa_\ell)$ uses the unnormalized source
periods. The conjugation in the second term is valid because
the eigenfunctions are real and $b_{-\ell}=\overline{b_\ell}$.
The coefficients, including $\widehat w(s_j)$ and $Q_{j,H}$,
are fixed throughout this $Y$ interval.

\emph{3. The large sieve and its hypotheses.}
Theorem 3 of Lekkas--Voskou \cite{LekkasVoskou2024}, with weight
$m=0$, states that for a fixed cofinite Fuchsian group and a
fixed closed geodesic, arbitrary complex coefficients $d_j$, and
$\delta$-separated points $y_\nu\in[Y_0,2Y_0]$,
\[
 \sum_\nu\left|\sum_{|t_j|\leq V}d_jy_\nu^{it_j}
                    \int_lu_j\right|^2
 \ll (V+Y_0\delta^{-1})\sum_{|t_j|\leq V}|d_j|^2.
\]
The continuous-spectrum coefficients vanish on a compact quotient.
The same statement is Theorem 1.3 of \cite{LekkasPetridis2026}.
Apply it on the fixed compact torsion-free cover from Statement A.3.1,
with zero coefficients for all eigenfunctions not lifted from
$\Gamma^+$. The covering and period normalization factors are
fixed and were verified above. Apply it separately to the two
target geodesics and use Cauchy--Schwarz for their sum $p_j'$.
There is no requirement that the target equal the source geodesic
or that the coefficients $d_j$ be periods.

For a finite spectral sum, take equally spaced $y_\nu$ with
mesh $h$, choose $\delta=h/2$, multiply the inequality by $h$,
and let $h\to0$. The Riemann sums converge, giving
\begin{equation}\label{A:continuous-large-sieve}
 \int_{Y_0}^{2Y_0}
       \left|\sum_jd_jY^{it_j}p_j'\right|^2dY
       \ll Y_0\sum_j|d_j|^2.
\end{equation}
For the band $J\leq1+t_j<2J$, take
$d_j=r_j\widehat w(s_j)Q_{j,H}$.
By Statement A.4.3, Stirling's formula, and
\eqref{A:w-Mellin-bound},
\[
 \sum_{j\text{ in band}}|d_j|^2
 \ll J^{-1}(1+J/T)^{-2N}\sum_{j\text{ in band}}|Q_{j,H}|^2
 \ll\log^2(2H)(1+J/T)^{-2N}.
\]
The factor $Y^{1/2}$ in \eqref{A:sqrt-remainder-expansion}
therefore gives $L^2([Y_0,2Y_0])$ norm at most
$O(Y_0\log(2H)(1+J/T)^{-N})$ for this band.
Minkowski's inequality over the dyadic bands, and then passage
to the absolutely convergent spectral series, prove
\begin{equation}\label{A:smoothed-mean-square}
 \int_{Y_0}^{2Y_0}
 \left|\mathcal N_{w,B}(Y)-\frac{b_0c}{3}\widehat w(1)Y\right|^2dY
 \ll_\eta Y_0^2\log^4Y_0+Y_0^{2+\eta}\log^2Y_0.
\end{equation}

\emph{4. Sharp arc and sharp height.}
For the fixed arc $I=[0,\lambda_0]$ take its Selberg majorant
and minorant trigonometric polynomials $B_H^+,B_H^-$.
The exact elementary approximation used here is
\cite[equation (4.11)]{LesterMatomakiRadziwill2015}:
for every interval in $\mathbb R/\mathbb Z$ and integer $H\geq1$,
\[
 B_H^-\leq\mathbf1_I\leq B_H^+,\qquad
 b_0^\pm=|I|\pm\frac1{H+1},\qquad
 |b_\ell^\pm|\ll(1+|\ell|)^{-1}\quad(0<|\ell|\leq H).
\]
Endpoint half-value conventions can change the sum only by the
already proved $O(\log Y_0)$ Pell contribution.
These polynomials are real, have degree at most $H$, and are
fixed throughout the dyadic $Y$ interval. In particular $B_H^+\geq0$.
Use the radial cutoffs $w^\pm$ of Statement A.4.4 with the same
fixed $T=\lfloor\sqrt{Y_0}\rfloor$. Then
\[
 \mathcal N_{w^-,B_H^-}(Y)
 \leq\mathcal A(Y)
 \leq\mathcal N_{w^+,B_H^+}(Y)+O(\log Y_0),
\]
where $\mathcal A(Y)$ is the exact sharp ideal-arc count with
$A_z\leq Y$. The lower inequality is valid even where
$B_H^-$ is negative: if $w^-\ne0$ the radial sharp indicator
is one, and multiplication by $0\leq w^-\leq1$ still leaves
$w^-B_H^-\leq\mathbf1_I$. Outside that radial range it is zero.
The upper inequality follows from nonnegativity of $B_H^+$.

Both main terms differ from $\lambda_0cY/3$ by
$O(Y_0/H+Y_0/T)=O(\sqrt{Y_0})$.
Thus \eqref{A:smoothed-mean-square} for the two fixed pairs
$(w^-,B_H^-)$ and $(w^+,B_H^+)$ proves
\[
 \int_{Y_0}^{2Y_0}|\mathcal A(Y)-\lambda_0cY/3|^2dY
       \ll_\epsilon Y_0^{2+\epsilon},
\]
after choosing $\eta$ smaller than $\epsilon$ and absorbing
logarithms. Finally $Y=(9x^2-1)^{1/2}$ has derivative bounded
above and below by positive constants for $x\geq1$;
two dyadic $Y$ intervals cover the image of $[X,2X]$.
The difference between $\lambda_0cY/3$ and $C_0x$ is
$O(x^{-1})$, and Statement A.1.3 adds $O(\log(2x))$.
These terms have mean square $O(X\log^2X)$, proving
\eqref{A:sharp-mean-square}.
\end{proof}


\paragraph{Statement A.5.6 (Proved smoothed expansion and conditional two-sided oscillation).}
Let
\[
 \mathcal A(Y)=\sum_{A_z\leq Y}
       \#\{\mathfrak a:N\mathfrak a=m_z,\ \theta(\mathfrak a)\in I\},
 \qquad e(Y)=\frac{\mathcal A(Y)-\lambda_0cY/3}{\sqrt Y}.
\]
Fix a smooth nonincreasing function $w:[0,\infty)\to[0,1]$,
equal to $1$ near zero and to $0$ for sufficiently large argument.
Its derivative is supported in a compact subinterval of $(0,\infty)$.
Then, as $Y\to\infty$,
\begin{align}\label{A:fixed-smoothing-AP}
 \int_0^\infty e(Yv)\sqrt v\,(-w'(v))\,dv
 &=\frac14\sum_{j\geq1}p'_j
 \left(r_j\widehat w(s_j)Q_{j,I}Y^{it_j}
 +\overline{r_j\widehat w(s_j)Q_{j,I}}Y^{-it_j}\right)
 \nonumber\\
 &\quad+O_w(Y^{-1/25}).
\end{align}
The series converges absolutely and uniformly in $\log Y$.
It is a real uniformly almost-periodic function with mean zero.
The following is consequently a sufficient condition for the
requested two-sided oscillation:
\begin{equation}\label{A:exact-nonvanishing-condition}
 \text{there is a }t>0\text{ such that }\quad
 \mathcal C(t):=\sum_{j:t_j=t}p'_jQ_{j,I}\ne0.
\end{equation}
If \eqref{A:exact-nonvanishing-condition} holds, then
$R(X)-C_0X=\Omega_+(X^{1/2})$ and
$R(X)-C_0X=\Omega_-(X^{1/2})$.
The sum in \eqref{A:exact-nonvanishing-condition} is over the
complete eigenspace and is independent of its real orthonormal
basis. Nonvanishing is not asserted in this statement.

\begin{proof}
We give the additional estimates needed to pass to the sharp arc.
They also explain why the $O(H)$ zero-residue bound used in
Statement A.5.5 cannot simply be retained in this argument.

\emph{1. Uniform angular discrepancy.}
The proof of Statement A.4.4 gives the stronger intermediate fact
\begin{equation}\label{A:all-interval-discrepancy}
 D(V):=\sup_{J\subset\mathbb R/\mathbb Z}
 \left|\sum_{A_z\leq V}\sum_{N\mathfrak a=m_z}
       \bigl(\mathbf1_J(\theta(\mathfrak a))-|J|\bigr)\right|
 \ll_\epsilon V^{2/3+\epsilon}.
\end{equation}
Indeed its phase-aligned bound for
$\sum_{\ell\leq H}|W_\ell|/\ell$ is independent of the interval;
the weighted Erd\H{o}s--Tur\'an step and the two radial
brackets therefore apply to every $J$ with the same constants.
One can use half-open intervals, or control a point mass by
the same positive Fej\'er majorant. This avoids assuming any
uniform Pell estimate for a varying interval endpoint.
For every periodic absolutely continuous function $f$ of mean zero,
Stieltjes integration by parts consequently gives
\begin{equation}\label{A:BV-discrepancy}
 \left|\sum_{A_z\leq V}\sum_{\mathfrak a}f(\theta(\mathfrak a))\right|
 \leq D(V)\operatorname{Var}(f).
\end{equation}
The assertion follows first for step functions, by expressing
them through interval indicators, and then by approximation.

\emph{2. Canonical angular polynomials have bounded variation.}
We use a specified choice of the Selberg polynomials.
The sawtooth approximation in
\cite[Lemma 2.2]{WangZhang2023}, explicitly attributed there to
Vaaler's Theorem 18, states the following. For an integer $H\geq1$,
put
\[
 W(t)=\pi t(1-|t|)\cot(\pi t)+|t|\quad(0<|t|<1),
 \qquad W(0)=1,
\]
extend $W$ by zero outside $[-1,1]$, and set
\[
 \psi_H(x)=-\sum_{1\leq|\ell|\leq H}
       \frac{W(\ell/(H+1))}{2\pi i\ell}e^{2\pi i\ell x},
 \quad
 F_H(x)=\sum_{|\ell|\leq H}
       \left(1-\frac{|\ell|}{H+1}\right)e^{2\pi i\ell x}.
\]
Then $F_H\geq0$ and
$|\psi_H-\psi|\leq F_H/(2(H+1))$.
These hypotheses are purely trigonometric and impose no
condition on our lattice. For an interval $I=[a,b]$ take
\[
 B_H^\pm(x)=b-a+\psi_H(a-x)+\psi_H(x-b)
 \ \pm\frac{F_H(a-x)+F_H(x-b)}{2(H+1)}.
\]
With the usual half-value convention at the two endpoints,
these satisfy the pointwise brackets of Statement A.5.5.
The endpoint convention changes our actual count by
$O(\log(2Y))$, already proved in Statement A.1.3.

In addition to the previously used Fourier bounds, these
particular polynomials satisfy
\begin{equation}\label{A:Selberg-BV}
 \|B_H^\pm\|_\infty+\operatorname{Var}(B_H^\pm)\ll1,
 \qquad \|(B_H^\pm)^{(r)}\|_1\ll_r H^{r-1}\quad(r\geq1).
\end{equation}
Here is a direct verification of the extra assertion.
The extension of $W$ is $C^1$, with integrable piecewise second
derivative: its value and first derivative vanish at $\pm1$.
Twice summing its Fourier polynomial by parts gives
\[
 \left|\sum_{\ell\in\mathbb Z}W(\ell/(H+1))e^{2\pi i\ell x}\right|
 \ll \min\{H,(H\|x\|^2)^{-1}\}.
\]
Its $L^1$ norm is therefore bounded independently of $H$.
Since $\psi_H'$ equals one minus this polynomial, it has
bounded $L^1$ norm. The formula
$F_H(x)=\sin^2(\pi(H+1)x)/((H+1)\sin^2(\pi x))$
shows $\|F_H'\|_1\ll H$: split at $\|x\|=1/H$ and
differentiate, using $O(H^3|x|)$ in the first range and
$O(|x|^{-2}+H^{-1}|x|^{-3})$ in the second.
This proves the variation bound; bounded sup norm follows from
the sawtooth approximation and $F_H\leq H+1$.
Finally, the elementary $L^1$ Bernstein bound for band-limited
polynomials gives the derivative estimates in
\eqref{A:Selberg-BV}. It can be proved here without an
approximation theorem: choose a fixed smooth frequency cutoff
equal to one on $[-1,1]$ and supported in $[-2,2]$.
Its periodic kernel, differentiated $r-1$ times, has $L^1$
norm $O_r(H^{r-1})$, by periodizing its Schwartz inverse Fourier
transform. Convolve this kernel with $(B_H^\pm)'$.

\emph{3. Improvement of the zero residue.}
For either polynomial $B=B_H^\pm$, the exact identity at $s=0$
in Statement A.5.5 gives
\[
 H_B(0):=\sum_{|\ell|\leq H}b_\ell H_\ell(0)
   =\sum_z\sum_{\mathfrak a}f_{A_z}(\theta(\mathfrak a)),
\]
where, putting $\delta(A)=\operatorname{arsinh}(1/A)/L_0$,
\[
 f_A(x)=\tfrac12\{B(x+\delta(A))+B(x-\delta(A))-2B(x)\}.
\]
This series converges absolutely for each fixed $H$,
since its large-$A$ summand is $O_H(A^{-2})$ and
$S_A(V)\ll V$. Every $f_A$ has mean zero.
On a dyadic block $A\asymp Z$, \eqref{A:Selberg-BV} gives
\[
 \operatorname{Var}(f_A)\ll
      \min\{1,H^2/Z^2\},\qquad
 \operatorname{Var}(\partial_A f_A)
      \ll\min\{H/Z^2,H^2/Z^3\}.
\]
The second estimate follows from
$|\delta'(A)|\ll A^{-2}$ and the first or second difference
of $B''$; the first follows from the second difference of $B'$.
Partial summation in $A$, followed by
\eqref{A:BV-discrepancy}, bounds that dyadic block by
\[
 \begin{cases}
 O_\epsilon(HZ^{\vartheta-1}),&Z\leq H,\\
 O_\epsilon(H^2Z^{\vartheta-2}),&Z\geq H,
 \end{cases}
 \qquad \vartheta=2/3+\epsilon<1.
\]
For $A\leq V$ use instead $\|B\|_\infty\ll1$ and
$S_A(V)\ll V$. Summing the remaining dyadic blocks gives
\[
 |H_B(0)|\ll_\epsilon
 V+HV^{\vartheta-1}+H^\vartheta.
\]
Take $V=H^{1/(2-\vartheta)}$. Since $\epsilon$ can be
arbitrarily small, this proves
\begin{equation}\label{A:zero-residue-improved}
 H_B(0)\ll_\epsilon H^{3/4+\epsilon}.
\end{equation}

\emph{4. Fixed radial smoothing and passage to the sharp arc.}
Use the negative-line transform of Statement A.5.5 with the fixed
function $w$, $\sigma=-49/100$, and $H=\lfloor Y^{3/5}\rfloor$.
The parameter $T$ in the transform bounds is now a fixed
constant depending on $w$. The contour bound remains
$O_w((T+H)\log(2H))$ for $-1/2<\sigma<0$.
To verify the only extra range not used in Statement A.5.5,
take $K>T$ in its high-angular calculation. The dyadic sum is
\[
 K^{1/2-\sigma}\sum_{J\ll K}J^\sigma\min(J,T)^{1/2}
 \ll_\sigma K^{1/2-\sigma}T^{\sigma+1/2}\leq C_\sigma K.
\]
The low and comparable ranges use only $a<1$, exactly as before.
The spherical-correction inversion, the new zero-residue bound,
and the exceptional intersection terms therefore leave an error
\[
 O_w\bigl(Y^\sigma H\log(2H)
          +Y^\sigma H^{1-\sigma}\log(2H)
          +H^{3/4+1/100}+\log(2Y)\bigr)
   =O_w(Y^{23/50}).
\]
The zero Fourier mode has the same main term and pole at $s=1$
as in Statement A.5.5; all $M_\ell(0)$ vanish.

We also need a bound uniform in $H$ for the individual
$Q_{j,H}^\pm$. The quantitative assertion of
\cite[Theorem 7.1 and Section 7.1]{AnantharamanZelditch2012}
gives a polynomial bound in $t_j$ for the H\"older norm of
the boundary primitive. Its displayed constant is
$O((1+t_j)^{3/2})$ times
$\|u_j\|_\infty+\|\nabla u_j\|_\infty$.
On the fixed compact cover, the elliptic estimate and Sobolev
embedding give the deliberately coarse bound
$\|u_j\|_{C^1}\ll(1+t_j)^4$.
The fixed coordinate changes and multiplication by
$(2e^v)^{-s_j}$ in Statement A.4.6 cost at most a further fixed
power of $1+t_j$. Thus, for some absolute integer $M$ depending
only on this fixed geometric setup, the periodic primitive
$F_j$ of $\beta_j$ has
\[
 \|F_j\|_{C^{1/2}}\ll(1+t_j)^M.
\]
An affine mean is understood in this assertion, as in
Statement A.4.6.

Identify $u=L_0\theta$, absorbing only fixed coordinate factors.
The middle polynomial in $B_H^\pm$ is the convolution of the
interval indicator with the kernel having coefficients
$W(\ell/(H+1))$. Its bounded $L^1$ norm and the displayed
decay bound imply approximation of a $1/2$-H\"older primitive
with error $O(H^{-1/2})$. For the two Fej\'er terms,
integration by parts and the preceding bounds for $F_H'$ give
\[
 \frac1{H+1}\left|\int F_H(x-a)\,\beta_j(dx)\right|
 \ll (1+t_j)^M H^{-1/2};
\]
subtract $F_j(a)$ inside the derivative integral and use
$|F_j(x)-F_j(a)|\ll(1+t_j)^M\|x-a\|^{1/2}$.
Consequently
\begin{equation}\label{A:Selberg-Q-limit}
 |Q_{j,H}^\pm-Q_{j,I}|
       \ll(1+t_j)^M H^{-1/2},
 \qquad |Q_{j,H}^\pm|+|Q_{j,I}|\ll(1+t_j)^M,
\end{equation}
after increasing $M$ if necessary.
The rapid vertical decay of $\widehat w$, the polynomial gamma
bound on $r_j$, the broad-band period bound, and
$\#\{t_j\leq V\}\ll V^2$ make the residue series absolutely
convergent. Equation \eqref{A:Selberg-Q-limit} permits passage
to its sharp-arc limit with error $O_w(Y^{1/2}H^{-1/2})$.

Bracket the actual radially smoothed arc count by $B_H^-$ and
$B_H^+$. Their main terms differ from
$\lambda_0cY\widehat w(1)/3$ by $O_w(Y/H)$.
Both spectral terms converge to the same series, and all the
other errors just proved are $O_w(Y^{23/50})$.
This proves the corresponding unnormalized expansion.
Partial summation gives
\[
 \sum_z\sum_{\mathfrak a}w(A_z/Y)\mathbf1_I(\theta(\mathfrak a))
       -\frac{\lambda_0c}{3}Y\widehat w(1)
 =\sqrt Y\int_0^\infty e(Yv)\sqrt v\,(-w'(v))\,dv,
\]
and proves \eqref{A:fixed-smoothing-AP}.

\emph{5. The conditional two-sided conclusion.}
Group the absolutely convergent series by distinct $t$.
Its coefficient at frequency $t$ is
$r(t)\widehat w(1/2+it)\mathcal C(t)/4$.
Here $r(t)\ne0$ because the gamma function has no zeros.
If \eqref{A:exact-nonvanishing-condition} holds, choose $w$
with $\widehat w(1/2+it)\ne0$: smooth monotone approximations
to $\mathbf1_{[0,1]}$ have Mellin transforms tending to
$(1/2+it)^{-1}\ne0$.
The resulting real series $P(\log Y)$ is nonzero.
Termwise averaging, justified by absolute uniform convergence,
gives mean zero and strictly positive mean square.
Since it is bounded, these two facts imply
$\limsup P>0$ and $\liminf P<0$.
For example, if $\limsup P\leq0$, then for every $\varepsilon>0$
one has eventually $P\leq\varepsilon$ and $P\geq-C$.
The inequality
$P^2\leq(\varepsilon-C)P+\varepsilon C$ contradicts the
positive mean square after averaging and letting
$\varepsilon\to0$. The other sign is identical.

The kernel $\sqrt v\,(-w'(v))$ is nonnegative, has finite
positive integral, and is supported in a fixed compact
subinterval of $(0,\infty)$. Positive and negative values
bounded away from zero of its convolution with $e(Y)$
therefore force positive and negative values bounded away
from zero of $e$ at arbitrarily large arguments.
Finally $Y=\sqrt{9X^2-1}=3X+O(X^{-1})$ and
$R(X)=\mathcal A(Y)+O(\log(2X))$ transfer both assertions
to $R(X)-C_0X$. The eigenspace sum is basis-independent because
it is the contraction of the two linear period functionals
with the orthogonal projector onto that eigenspace.
\end{proof}
}{}
% Source audit for A.5, 15 November 2026.  Ready for root integration.
% No assertion that a corrected Sali\'e conjecture implies the actual arc count.
\subsection{An audit of a proposed square-root input}
\paragraph{Statement A.5.7 (Refuted conjectural input).}
The literal Conjecture~1 of Bir\'o \cite{A5Biro2026} cannot serve as a
sound conjectural input: it contradicts a proved first-moment formula.
We give the specialization, including its smoothness and normalization
checks.  This does not refute a version in which the unary-theta
contribution is subtracted.

For odd positive $c$, put
\[
 T(m,n;c)=\sum_{x\in(\Z/c\Z)^\times}
 \left(\frac{x}{c}\right)e\left(\frac{m\bar x+nx}{c}\right),
 \qquad e(u)=e^{2\pi i u},
 \qquad
 \epsilon_c=\begin{cases}1,&c\equiv1\pmod4,\\
                         i,&c\equiv3\pmod4.
            \end{cases}
\]
These are Bir\'o's conventions.  Louvel's $K_2(m,n;c)$ writes the two
arguments in the reverse order; replacing $x$ by $\bar x$ shows that
the two definitions agree, since
$\bigl(\frac{\bar x}{c}\bigr)=\bigl(\frac{x}{c}\bigr)$.

Choose a fixed nonnegative, nonzero function
$g\in C_c^\infty(\R)$ supported in $[1,2]$ and, for $C\geq1$, define
\[
 f_C(x)=\sqrt{C/x}\,g(x/C).
\]
This is smooth and supported in $[C,2C]$, and for every fixed
$j\geq0$ it satisfies
$x^j f_C^{(j)}(x)=O_{g,j}(1)$.
In Bir\'o's conjecture take
\[
 L=4,\quad K=1,\quad m=4,\quad n=1,
 \quad B=1,\quad\alpha=0,
 \quad r=1.
\]
All the hypotheses hold: $2\mid L$,
$\gcd(L,Kr)=1$, and $LK(1+B)=8\ll_\delta C^\delta$
for any $\delta>0$ once $C$ grows.  The stated derivative bounds hold
for that same $\delta$.  For each odd $c$, multiplication by $4$ is a
bijection of $(\Z/c\Z)^\times$.  Substitution $x=4y$ gives
\[
 T(4,\bar4;c)
 =\sum_y\left(\frac{4y}{c}\right)
 e\left(\frac{\bar y+y}{c}\right)
 =T(1,1;c).
\]
The choice $n=1$ also removes any possible notational ambiguity about
the placement of the overline in Bir\'o's second argument.
Thus his asserted estimate would give
\[
 \sum_{c\equiv1\pmod4}
 g(c/C)\frac{T(1,1;c)}{\sqrt c}
 =O_{g,\varepsilon}(C^{1/2+\varepsilon}).
\]

Theorem~1 of Louvel's published paper
\cite[p.~525]{A5Louvel2012} (Theorem~1.1 in its arXiv version)
has the following applicable specialization:
\begin{equation}\label{A5audit:louvel-specialization}
 A(U):=\sum_{\substack{0<c\leq U\\c\ \text{ odd}}}
 \frac{T(1,1;c)\overline{\epsilon_c}}{\sqrt c}
 =\frac8{\pi^2}U+O_\eta(U^{3/4+\eta})
 \qquad(\eta>0).
\end{equation}
For completeness, in his general theorem we take the odd, coprime
integers $D=f=1$, the even primitive character of conductor~$1$,
and positive $m=n=1$.  In his formula for the main coefficient take
$t=s=m'=n'=D_t=D_0=D_1=1$.
Then $t$ is squarefree and $1$ modulo~$4$, $s$ is supported by $t$,
$s^2t^3\mid D$, and all coprimality and equal-gcd conditions hold.
The Jacobi symbol with denominator~$1$ equals~$1$ and all prime
products are empty, so the coefficient is exactly $8/\pi^2$.
The conjugate on $\epsilon_c$ has been checked in both the author
preprint and the published formula; omitting it would change this
normalization.

The same theorem, with $m=n=-1$ and the other data unchanged, gives
\[
 A_-(U):=\sum_{\substack{0<c\leq U\\c\ \text{ odd}}}
 \frac{T(-1,-1;c)\overline{\epsilon_c}}{\sqrt c}
 =O_\eta(U^{3/4+\eta}),
\]
because its first case applies whenever either argument is negative.
The exact substitution $x=-y$ in the defining sum gives
\[
 T(-1,-1;c)=\left(\frac{-1}{c}\right)T(1,1;c).
\]
Consequently the single residue class $1$ modulo $4$ has the asymptotic
\begin{equation}\label{A5audit:single-class}
 A_1(U):=\sum_{\substack{0<c\leq U\\c\equiv1\pmod4}}
 \frac{T(1,1;c)}{\sqrt c}
 =\frac{A(U)+A_-(U)}2
 =\frac4{\pi^2}U+O_\eta(U^{3/4+\eta}).
\end{equation}

Abel summation applied to \eqref{A5audit:single-class} gives
\begin{align*}
 \sum_{c\equiv1\pmod4} g(c/C)
       \frac{T(1,1;c)}{\sqrt c}
 &=-\frac1C\int_C^{2C} A_1(u)g'(u/C)\,du\\
 &=\frac4{\pi^2}C\int_1^2g(v)\,dv
       +O_{g,\eta}(C^{3/4+\eta}).
\end{align*}
In the exact normalization of Bir\'o's left side, this states
\[
 \frac1C\sum_{c\equiv1\pmod4}f_C(c)T(4,\bar4;c)
 =\frac4{\pi^2}\sqrt C\int_1^2g(v)\,dv
 +O_{g,\eta}(C^{1/4+\eta}).
\]
The main coefficient is positive.  Taking, for example,
$\eta=1/16$ and $\varepsilon=1/8$ contradicts his asserted
$O(C^\varepsilon)$ bound as $C\to\infty$.
This is a single completely specified parameter family violating the
printed conjecture.

The program \path{code/B_A5_salie_audit.py} checks the change of
variables and the identity for $(-1,-1)$ exactly as equalities of integer coefficient vectors in
$\Z[X]/(X^c-1)$ for every odd $3\leq c\leq1001$.
It also checks the identity
\[
 T(1,1;c)=
 \left(\sum_{y\bmod c}e(y^2/c)\right)
 \sum_{v^2\equiv4\pmod c}e(v/c)
\]
exactly modulo the cyclotomic polynomial $\Phi_c$ for odd
$3\leq c\leq199$.
The accompanying first-moment evaluations use 192-bit Arb arithmetic
and the explicit bump
\[
 g(v)=\begin{cases}
 \exp\!\left(4-\dfrac1{(v-1)(2-v)}\right),&1<v<2,\\
 0,&\text{otherwise}.
 \end{cases}
\]
They are finite normalization checks, not the proof of the asymptotic
contradiction.  The latter is the preceding application of Louvel's
theorem.  On $[5/4,7/4]$ this bump is at least $e^{-4/3}$, so its
positive integral does not depend on a numerical quadrature.

For the literature status of this correction, on 15 November 2026 we
searched the exact paper identifier, its title together with
``correction'' and ``counterexample'', and Bir\'o--Louvel--Sali\'e
combinations.  The primary arXiv record still listed only version~1
of 13 April 2026.  We did not locate an earlier published correction
of this printed conjecture.  The mathematical input is Louvel's known
theorem, and the present deduction is an independently checked
specialization of it; we make no claim of priority.

\paragraph{The corrected candidate and the remaining reduction.}
Conjecture~1.4 of Sun \cite{A5Sun2025} explicitly subtracts
\[
 2Y^{1/2}\sum_{r_j=i/4}\tau_j(m,n)
\]
from the corresponding sum of half-integral-weight Kloosterman
sums divided by their moduli.  Its hypotheses include a weight
$k=\pm1/2$ multiplier on $\Gamma_0(N)$ and absence of Laplace
eigenvalues in $(3/16,1/4)$ for that multiplier space.
It gives an untwisted, uniformly parameterized modulus-sum estimate;
it does not state a bound for the angular-period coefficients of
our compact quaternionic surface.
We have not established a reduction from that conjecture to the
actual sharp-arc count.  In a direct ternary delta-method approach
one must specify the multiplier and its spectral condition, handle
the local factors at $2$ and $3$, identify and subtract every theta
term, and prove uniform bounds for the resulting modulus weights
as the physical angular and radial transition widths shrink to
$X^{-1/2}$.  Partial summation with an uncontrolled total-variation
factor does not prove the required exponent.

\paragraph{Statement A.5.8 (Proved: nonvanishing and both square-root oscillation signs).}
For the actual sharp arc $I$, there are infinitely many positive
spectral parameters $t$ such that
\[
 \mathcal C(t)=\sum_{j:t_j=t}p'_jQ_{j,I}\ne0.
\]
Consequently
\[
 R(X)-C_0X=\Omega_+(X^{1/2})
 \quad\hbox{and}\quad
 R(X)-C_0X=\Omega_-(X^{1/2}).
\]
In particular the requested unsigned $\Omega(X^{1/2})$ assertion holds.

\begin{proof}
We prove nonvanishing by a localized positive-frequency relative-wave
calculation, including its endpoint contribution. The oscillation
conclusion then follows from Statement A.5.6.
Throughout this proof set
\[
 a=17,\quad B=12\sqrt2,\quad R=\operatorname{arcosh}(17),
 \quad d_0=\operatorname{arcosh}(3),\quad
 \tau(u)=\operatorname{arsinh}(B\cosh u).
\]
Thus $R=2d_0$.
The period normalization throughout is
$p'_j=\sum_{i=1}^2\int_0^{L_i}u_j(\ell_i(v))\,dv$,
with the primitive cyclic translation stabilizers $H_i$ used in
A.3.1. If a larger axis normalizer identifies the two orientations,
unfolding with $H_i$ retains the resulting positive duplicate
multiplicity; it does not divide a period by an unrecorded factor.

\emph{1. The frequency window and the exact source functional.}
Choose a real even Schwartz function $\rho$ whose Fourier transform
\[
 \widehat\rho(y)=\int_{\mathbb R}\rho(t)e^{-ity}\,dt
\]
is smooth, nonnegative, supported in $(-\varepsilon,\varepsilon)$,
and positive at zero. The sufficiently small fixed $\varepsilon>0$
will be specified below.
Consider the absolutely convergent sum
\begin{equation}\label{A:nonvanishing-window}
 \mathscr S_T=
 \sum_{j\geq1}p'_jQ_{j,I}\rho(t_j-T)e^{it_jd_0},
 \qquad T\to+\infty.
\end{equation}
The polynomial individual bounds already proved in Statement A.5.6
justify this convergence. We shall prove
\begin{equation}\label{A:nonvanishing-wave-asymptotic}
 \mathscr S_T=C_*e^{-i\pi/4}T^{-1/2}+O(T^{-1}),
 \qquad C_*>0.
\end{equation}
If every complete-eigenspace contraction vanished, the left side
would be identically zero, so this asymptotic proves the claim.

Choose a real orthonormal eigenbasis diagonalizing the reflection $W$.
The reflection-odd functions have $p'_j=0$, as proved in A.4.7,
so only the even functions contribute to this sum.
For these functions Statement A.4.8 gives the exact representation
\[
 Q_{j,I}=\int_{\mathbb R}u_j(i e^u)\,q_{t_j}(u)\,du,
 \qquad
 q_t(u)=
 A(u)e^{-it\tau(u)}
   \left(1+\sum_{n=1}^{N}t^{-n}b_n(u)+r_{N,t}(u)\right),
\]
where
\[
 A(u)=\frac{a}{2\cosh^{3/2}\tau(u)},\qquad
 r_{N,t}(u)=O_N(t^{-N-1}).
\]
The coefficients and the needed fixed-order derivatives are bounded
after the common factor $A(u)$ is removed.
In particular $A(u)\ll e^{-3|u|/2}$.
The Volterra recurrence in A.4.8 permits $N$ to be arbitrarily large.

Here is a precise way to preserve finite propagation in using this
expansion. Expand each $t^{-n}$ to sufficiently many Taylor terms
about $t=T$ under $\rho(t-T)$.
Every retained term is a polynomial in $t-T$ times $\rho(t-T)$,
multiplied by a power of $T^{-1}$.
Its negative-frequency reflection makes it an even entire
Paley--Wiener multiplier. For the leading term this is
\[
 h_{T,u}(t)=
 \rho(t-T)e^{it(d_0-\tau(u))}
 +\rho(-t-T)e^{-it(d_0-\tau(u))}.
\]
On $t>0$ the second term is rapidly small in $T$.
Its Fourier transform is supported in the two intervals of radius
$\varepsilon$ about $\pm(d_0-\tau(u))$.
Polynomial factors merely differentiate that compactly supported
Fourier transform.
All discarded spectral terms can be bounded before geometric
unfolding: use the compact-source eigenfunction sup bound,
the polynomial individual target-period bound, Weyl's law, the
rapid decay of $\rho$, and $\int A(u)\,du<\infty$.
Choosing the number of expansion terms larger than the resulting
fixed polynomial loss makes their total $O(T^{-2})$.
This step does not assert finite propagation for a gamma quotient.

\emph{2. The point-to-geodesic spectral density on the cover.}
We record the local kernel, including its normalization.
In Fermi coordinates $(r,v)$ about a complete geodesic in
$\mathbb H^2$, the transverse measure is $\cosh r\,dr$.
The positive Laplacian acting on functions independent of $v$
is $-\partial_r^2-\tanh r\,\partial_r$.
Conjugation by $\sqrt{\cosh r}$, followed by subtraction of $1/4$,
gives $-\partial_r^2+\operatorname{sech}^2r/4\geq0$ on the full line.
There are no negative bound states; a zero bound state is excluded
by its quadratic form, and positive bound states are excluded by
the nondecaying incoming/outgoing asymptotics at infinity.
The explicit Jost resolvent has smooth boundary values for $t>0$,
so the following jump is the full spectral density of the line
source, with no omitted discrete transverse spectrum.
The resolvent of its line delta source is
\[
 G_s(r)=-\frac{f_s(|r|)}{2f_s'(0)},\qquad \Re s>1.
\]
Indeed it solves the homogeneous equation away from zero and has
derivative jump $-1$, as required by the positive Laplacian.
Its outgoing boundary values at $s=1/2\mp it$ determine the
spectral density in $t$ by the resolvent jump formula:
\begin{equation}\label{A:line-spectral-density}
 \mathcal P_t(r)=
 \frac{t}{\pi i}
 \left\{-\frac{f_{1/2-it}(|r|)}{2f_{1/2-it}'(0)}
       +\frac{f_{1/2+it}(|r|)}{2f_{1/2+it}'(0)}\right\}.
\end{equation}
The factor is fixed by the scalar identity
$(x-\lambda-i0)^{-1}-(x-\lambda+i0)^{-1}
 =2\pi i\,\delta(x-\lambda)$ in the spectral resolution,
and by $d\lambda=2t\,dt$ for $\lambda=1/4+t^2$.
The denominators do not vanish for $t>0$: otherwise the
Wronskian of the incoming and outgoing solutions would vanish
at zero, whereas its value at infinity is nonzero.
The outgoing Volterra calculation of A.4.8 applies also on $r\geq0$.
It yields
\begin{equation}\label{A:line-spectral-principal}
 \mathcal P_t(r)=
 \frac{e^{itr}+e^{-itr}}{2\pi\sqrt{\cosh r}}
 +\frac1{\sqrt{\cosh r}}
       \sum_{\sigma=\pm1}e^{i\sigma tr}
           \sum_{n=1}^{N}t^{-n}c_{\sigma,n}(r)
 +\mathcal R_N(t,r).
\end{equation}
The coefficients and their fixed-order derivatives are uniformly
bounded; after removing their limits at infinity their derivatives
decrease exponentially.
For every fixed $N$ the remainder can be made
$O_N(t^{-N-1}\cosh^{-1/2}r)$ by taking sufficiently many Volterra
terms. This is uniform for $r\geq0$.
Only $r\geq d_0-\varepsilon>0$ will contribute below.

The normalization can also be checked without an asymptotic formula.
The square-root connection coefficient of the hypergeometric function
at argument one gives
\[
 f_s'(0)=-\frac{2^{1-s}\sqrt\pi\,\Gamma(s+1/2)}
                         {\Gamma(s/2)^2}.
\]
Consequently the source kernel is exactly
\[
 q_t(u)=\frac{a}{2\cosh\tau(u)}
              \frac{f_{1/2+it}'(\tau(u))}{f_{1/2+it}'(0)}
        =-\frac{a}{\cosh\tau(u)}G_{1/2+it}'(\tau(u)).
\]
This identifies the source factor with the normal derivative of
the same outgoing line resolvent. In particular no unspecified
scattering phase remains in its positive principal amplitude.

For an even Paley--Wiener multiplier, the automorphic point-to-cycle
kernel is obtained by summing this complete-line kernel over
the translates of the two target axes.
This follows either by periodizing the wave kernel and integrating
over a primitive target period, or directly by unfolding that integral.
Finite propagation makes the geometric sum finite at each fixed
source point and fixed propagation time.
The construction is valid on our compact orbifold; equivalently,
one can perform it on the fixed torsion-free cover and project to
the invariant subspace with the normalization factors already
verified in A.3.1.
Unoriented axes can occur with fixed positive multiplicities from
the cyclic stabilizer. We retain these multiplicities throughout;
their positivity is all that is needed here.

We include the finite-propagation justification relevant to the
shifted operator. For a smooth solution of
$w_{vv}+(\Delta-1/4)w=0$ with compactly supported initial data,
use the local energy density
$\mathcal E=(|w_v|^2+|\nabla w|^2+|w|^2)/2$.
Its conservation identity has an additional term
$5\operatorname{Re}(\overline w\,w_v)/4$, bounded by
$5\mathcal E/4$.
Integrate over the complement of the expanding unit-speed
neighborhood of the initial support.
The moving-boundary flux is nonpositive by
$|\operatorname{Re}(\overline{w_v}\partial_n w)|
 \leq(|w_v|^2+|\partial_nw|^2)/2$.
The exterior energy therefore satisfies
$E'(v)\leq5E(v)/4$ and $E(0)=0$.
Gronwall's inequality gives zero exterior energy.
Approximation extends the support assertion to distributional
initial data and hence to the wave kernel.
The argument is local, so it holds on the fixed smooth cover
and descends to the orbifold invariant subspace.
Fourier inversion of an even Paley--Wiener multiplier expresses
it through this cosine wave kernel, proving the support assertion
used in unfolding.

These rearrangements remain absolutely integrable in $u$.
Indeed, since the source is closed, the number of target axes at
distance at most $D$ from $p(u)$ is $O(e^D)$ uniformly in $u$.
To see this, translate the nearest point on each axis into a fixed
primitive segment of length $L_i$. A fixed basepoint of that segment
then lies within $D+L_i$ of $p(u)$. The resulting orbit points have
uniformly bounded multiplicity and can be bounded by disjoint-ball
packing. The source parameter may be reduced modulo its fixed period,
so the constants are uniform in $u$.
The propagated kernel vanishes unless
$r\leq\tau(u)-d_0+\varepsilon$.
The amplitudes in \eqref{A:line-spectral-principal} therefore give
\[
 \sum_{\gamma:r_\gamma\leq\tau(u)+O(1)}
       (\cosh r_\gamma)^{-1/2}
       \ll e^{\tau(u)/2}.
\]
Multiplication by $A(u)\ll e^{-3\tau(u)/2}$ leaves an integrable
$O(e^{-\tau(u)})$ bound.
A polynomial in $\tau(u)$, arising in an auxiliary remainder
estimate, is harmless as well.

We make explicit the support of the integrated spectral remainder.
The exact $\mathcal P_t(r)$ is even in $t$.
After reflection to the full real line, expand its high-frequency
symbol about $T$ into polynomials in $t-T$ under $\rho(t-T)$.
This Taylor construction and even reflection are applied to
\emph{every} inverse power $t^{-n}$ appearing in the cover
expansion \eqref{A:line-spectral-principal}, as well as to every
such power in the source expansion of $q_t$.
We never apply finite propagation directly to a retained
nonentire factor $t^{-n}$.
Every retained term has Fourier support in
$|r-|d_0-\tau(u)||\leq\varepsilon$ (with its reflected branch);
in particular it vanishes for
$r>|d_0-\tau(u)|+\varepsilon$.
The exact even Paley--Wiener kernel vanishes there by finite
propagation as well. Hence the integrated remainder also vanishes
there, and its uniform
$O_N(T^{-N})\cosh^{-1/2}r$ bound is summed only over this finite
geometric set. The compact low-$t$ part satisfies the same bound:
the displayed gamma formula makes $f_s'(0)$ nonzero also at $t=0$,
and the explicit hypergeometric formula bounds
$\sqrt{\cosh r}\,f_{1/2+it}(r)$ uniformly for bounded $t$.
The rapid decay of $\rho(t-T)$ supplies every inverse power of $T$.
Thus no unrestricted image sum of an unsmoothed spectral remainder
has been interchanged.

The constant eigenfunction contributes only
$O_N(T^{-N})e^{\tau(u)/2}$ before this multiplication:
evaluate the even entire multiplier at $t=i/2$ and use the
rapid horizontal decay of $\rho$ in a fixed strip.
No nonconstant exceptional eigenvalues exist by A.2.5.

\emph{3. The geometric oscillatory integral.}
Let $r_\gamma(u)\geq0$ denote the distance from $p(u)$ to a
target translate, counted with the fixed positive multiplicity.
For the leading terms, Fourier inversion in the window is exact.
The branch with phase $d_0-\tau-r_\gamma$ vanishes because
$\tau\geq R=2d_0$ and $\varepsilon<d_0$.
The other branch gives
\begin{equation}\label{A:wave-leading-integral}
 \frac1{2\pi}\sum_\gamma
 \int_{\mathbb R}
 \frac{a\,\widehat\rho(\tau-r_\gamma-d_0)}
      {2\cosh^{3/2}\tau\,\cosh^{1/2}r_\gamma}
       e^{iT(d_0-\tau+r_\gamma)}\,du.
\end{equation}
Choose a compact interval
$J\subset(0,R)$ containing $d_0$ and the support in
$\ell=\tau-r_\gamma$ of this expression.
All subleading terms have the same form, with smooth bounded
amplitude factors and extra powers of $T^{-1}$, and derivatives
of $\widehat\rho$ from the polynomial Taylor factors.
These derivatives have the same compact support.
Their bounds will follow from the same pushforward and stationary
phase arguments as the leading term.

\emph{4. A summable endpoint density, without a spacing assumption.}
Use unit Lorentz normals
\[
 n=(N_0,N_1,N_2)
   =\left(\frac{s}{\sqrt2},\sqrt{\frac32}\,t,3z\right),
 \qquad N_0^2-N_1^2-N_2^2=-1.
\]
On each choice of sign making $r_\gamma>0$, put $v=\tanh u$ and
\[
 D(v)=\sqrt{a^2-v^2},\quad
 E(v,\ell)=D(v)\cosh\ell-B\sinh\ell.
\]
The equation $\ell=\tau-r_\gamma$ becomes
\begin{equation}\label{A:wave-phase-equation}
 F(v,\ell)=N_0-N_1v-B\cosh\ell+D(v)\sinh\ell=0.
\end{equation}
On $J\subset(0,R)$, both $E$ and
$B\cosh\ell-D\sinh\ell$ are bounded away from zero.
The principal measure has the exact simplification
\[
 \frac{a}{2\cosh^{3/2}\tau\cosh^{1/2}r_\gamma}\,du
       =\frac{a}{2D^{3/2}E^{1/2}}\,dv .
\]
Moreover
\[
 F_v=-N_1-\frac{v\sinh\ell}{D},\qquad
 F_\ell=E,\qquad
 \frac{dv}{d\ell}
     =\frac{E}{N_1+v\sinh\ell/D}.
\]
If $|N_1|$ exceeds a fixed constant depending only on $J$,
the map has no critical point on this phase window.
Its pushed-forward density, including the smooth window, is
\begin{equation}\label{A:wave-pushforward-density}
 h_n(\ell)=
 \widehat\rho(\ell-d_0)
 \frac{a\sqrt{E(v(\ell),\ell)}}
      {2D(v(\ell))^{3/2}
       |N_1+v(\ell)\sinh\ell/D(v(\ell))|}.
\end{equation}
It exists on an interval, possibly empty. Extending it by zero,
including the possible endpoint jumps, gives
\[
 \|h_n\|_{L^1}+\operatorname{Var}(h_n)\ll_J |N_1|^{-1}.
\]
For completeness, $|N_1+v\sinh\ell/D|\geq |N_1|/2$,
$|v'|\ll_J|N_1|^{-1}$, $D$ and $E$ are bounded above and below,
and direct differentiation of \eqref{A:wave-pushforward-density}
gives $\sup|h_n|+\sup|h_n'|\ll_J|N_1|^{-1}$.
There are at most two endpoint jumps.
The allowed interval follows from
$F(\pm1,\ell)=N_0\mp N_1-Be^{-\ell}$.
This proves the stated BV estimate without separating any
two endpoint Busemann constants.

This bound is summable over all but finitely many target translates.
For $|N_2|>1$ orient the normal by $N_0>0$ and put
$c=N_0-|N_1|$.
The phase equation forces $0<c\leq C_J$.
The normal relation gives
\[
 |N_1|=\frac{N_2^2-1-c^2}{2c},
 \qquad \|h_n\|_{BV}\ll_J\frac{c}{N_2^2}
\]
for sufficiently large $|N_2|$.
For a fixed reduced target axis, its source-stabilizer translates
have
\[
 c_m=\sqrt{N_2^2-1}\,e^{-|u_0+mL_0|},
 \qquad
 \sum_{m:c_m\leq C_J}c_m
       \leq\frac{2C_J}{1-e^{-L_0}}.
\]
Statement A.3.1 gives $O(Y)$ reduced axes with $|N_2|\leq Y$.
Thus dyadic summation of $N_2^{-2}$ proves convergence of the
total BV norm at large height.
There are finitely many reduced axes at bounded height, including
the intersecting axes with $N_2=0$.
For each, $|N_1(m)|\asymp e^{|m|L_0}$ in its two infinite tails;
a zero leading coefficient would mean a shared endpoint with
the source, excluded by discreteness and the distinct axes.
Their inverse $|N_1(m)|$ sums converge too.
Only finitely many actual translates with bounded $|N_1|$ remain.

Consequently all infinite tails of \eqref{A:wave-leading-integral}
are Fourier transforms of a compactly supported BV density with
finite total variation. Their contribution is $O(T^{-1})$.
The same proof applies to the subleading amplitudes:
on the phase surface $\cosh\tau=D/\sqrt{1-v^2}$ and
$\cosh r_\gamma=E/\sqrt{1-v^2}$.
The outgoing amplitude coefficients, and their first derivatives
in $(v,\ell)$, are uniformly bounded up to $v=\pm1$.
This follows from the exponentially decreasing derivatives of
the Volterra coefficients; the factor $(1-v^2)^{-1}$ from
differentiating $\tau$ is canceled by their $e^{-2\tau}$ decay.
The BV product estimate therefore applies term by term.

\emph{5. All remaining critical values and Hessian signs.}
Only finitely many actual translates remain.
At an interior critical point of $\ell(v)$,
\eqref{A:wave-phase-equation} and $F_v=0$ give
\[
 N_1=-\frac{v\sinh\ell}{D},\qquad
 N_0=B\cosh\ell-\frac{a^2\sinh\ell}{D}.
\]
Substitution in the normal equation gives
\[
 N_2=\pm a\left(\cosh\ell-\frac{B\sinh\ell}{D}\right).
\]
The expression in parentheses is positive for $\ell\in J$.
Choose
$n'=(B,0,\operatorname{sgn}(N_2)a)$.
It is the unit normal of one of the translates $U^{\pm1}g_0$,
and direct substitution gives
\[
 \langle n,n'\rangle=-\cosh\ell.
\]
In the original integral coordinates this inner product is
$12s\mp51z$, a multiple of $3$.
Its absolute value is $3|z'|$ for one of the integral transformed
coordinates $z'=17z\mp4s$.
It follows that every positive critical value satisfies
\[
 \cosh\ell=3|z'|,\qquad |z'|\geq1.
\]
Choose $J$ small enough around $d_0$ that
$\ell<\operatorname{arcosh}(6)$ on $J$.
Then every critical value in this window equals $d_0$.

All such critical points are nondegenerate and have the same
Hessian sign. Indeed
\[
 F_{vv}=-\frac{a^2\sinh\ell}{D^3},\qquad
 \ell''(v)=-\frac{F_{vv}}{F_\ell}
     =\frac{a^2\sinh\ell}{D^3E}>0.
\]
Thus the phase $d_0-\ell(v)$ always has negative second derivative.
There is at least one such point: the target normal
$(2\sqrt2,0,3)$, corresponding to $(s,t,z)=(4,0,1)$,
has $v=0$, $\ell=d_0$, and
\[
 E=3,\qquad \ell''(0)=\frac{2\sqrt2}{51}.
\]
Its principal amplitude is strictly positive.
In fact the stationary set is explicit. At $\ell=d_0$,
$|N_1|\leq\sinh(d_0)/B=1/6$.
Since $N_1=\sqrt{3/2}\,t$ with integral $t$, this forces $t=0$,
then $v=0$, $N_0=2\sqrt2$, and $N_2=\pm3$.
Thus the only two geometric axes are those with normals
$(4,0,\pm1)$ in the original coordinates, each with its positive
cyclic-stabilizer multiplicity.
For one unit of such multiplicity the coefficient of
$e^{-i\pi/4}T^{-1/2}$ is
$\widehat\rho(0)/(4\,2^{1/4}\sqrt\pi)$.

There is no endpoint of any of these finitely many branches
with $\ell=d_0$.
At $v=\pm1$ the phase equation would require
$N_0\mp N_1=B e^{-d_0}$, hence
\[
 s\mp\sqrt3\,t=72-48\sqrt2.
\]
The left side belongs to $\mathbb Q(\sqrt3)$ and the right side
is a nonrational element of $\mathbb Q(\sqrt2)$, which is impossible.
We can therefore shrink the fixed $\varepsilon$ further so that
the finite-translate windows are compactly supported away from
their endpoints.
Ordinary stationary phase applies to their finitely many critical
points, as can be seen directly in one dimension.
Near a critical point $v_0$ set
$y=\operatorname{sgn}(v-v_0)\sqrt{2(\ell(v)-d_0)}$.
Then $dy/dv|_{v_0}=\sqrt{\ell''(v_0)}>0$ and the phase becomes
$-y^2/2$.
The leading integral is the amplitude at zero multiplied by
\[
 \int_{\mathbb R}e^{-iTy^2/2}\,dy
       =\sqrt{\frac{2\pi}{T}}e^{-i\pi/4}.
\]
This identity follows by inserting $e^{-\eta y^2}$, evaluating
the ordinary Gaussian integral, and letting $\eta\downarrow0$.
Taylor expansion of the smooth amplitude gives the usual
$O(T^{-3/2})$ local remainder; outside these neighborhoods,
integration by parts with $(iT\Phi')^{-1}\partial_v$ gives
$O(T^{-1})$.
Every leading coefficient is a positive real multiple
of the same factor $e^{-i\pi/4}T^{-1/2}$.
The value $\widehat\rho(0)>0$, the positive axis multiplicities,
and the displayed example show that their sum is nonzero.
Nonstationary parts are $O(T^{-1})$, as are the infinite BV tails.
All subleading symbol terms have an additional inverse power of
$T$, and the pre-unfolding remainder was made $O(T^{-2})$.
This proves \eqref{A:nonvanishing-wave-asymptotic}.

Finally group \eqref{A:nonvanishing-window} by complete eigenspaces.
Its nonzero asymptotic proves that at least one $\mathcal C(t)$
is nonzero. In fact finitely many nonzero contractions would make
\eqref{A:nonvanishing-window} $O_N(T^{-N})$ for every $N$, by
the Schwartz decay of $\rho$; this contradicts the same asymptotic.
Thus infinitely many are nonzero.
Statement A.5.6 then gives both claimed signs of
square-root oscillation for the actual sharp count.
\end{proof}

\IfFileExists{sections/B.tex}{\section{The lattice intersection and its cubic symmetries}
\label{sec:B}

We use positive definite lattices with roots of squared length two. Thus
\(\det L\) is the determinant of a Gram matrix, and
\(\Theta_L(\tau)=\sum_{x\in L}q^{(x,x)/2}\), where
\(q=e^{2\pi i\tau}\). The quadratic discriminant form is
\(Q_L(x+L)=(x,x)/2\pmod{\Z}\). All matrices in our certificates
are rational or integral. The two different quadratic forms sometimes
called the discriminant form, \((x,x)\pmod{2\Z}\) and \(Q_L\),
therefore differ by the displayed factor two.

\subsection{An independent construction and the eight extensions}

Put
\[
 A_5=\{x\in\Z^6:\ \textstyle\sum x_i=0\},\qquad
 D_4=\{y\in\Z^4:\ \textstyle\sum y_i\in2\Z\},\qquad
 R=A_5^4\perp D_4.
\]
The scalar product is the ordinary Euclidean product on the ambient
\(28\)-space. Write \(\lambda=(5,-1,-1,-1,-1,-1)/6\), and
\(\omega_j=\sum_{i=1}^j e_i-(j/6)\sum_{i=1}^6e_i\).
The five glue vectors, with their first four entries denoting multiples
of \(\lambda\), are
\begin{equation}\label{eq:B-glue}
\begin{split}
 &(3,3,3,3;0),\quad(3,3,0,0;(1,0,0,0)),\\
 &(3,0,3,0;(1/2,1/2,1/2,1/2)),\\
 &(2,2,2,0;0),\quad(2,4,0,2;0).
\end{split}
\end{equation}
We define \(N\) to be their span together with \(R\), and put
\[
 \rho=((5,3,1,-1,-3,-5)/2)^4\,\oplus\,(3,2,1,0),
 \qquad \alpha=(e_1-e_2,0,0,0;0).
\]
Then \((\rho,\alpha)=1\), \(\rho^2=84\), and \(\rho\in N\).
For \(d\mid6\) define, without any choice of a lift modulo \(d^2\),
\begin{equation}\label{eq:B-neighbor}
 K_d=\{x\in N:(\rho,x)\in d\Z\},\qquad
 N_d=K_d+\Z\bigl(\rho/d-(6/d)\alpha\bigr).
\end{equation}
Let \(r\) cyclically shift the six coordinates in every \(A_5\) factor
by two places and fix the \(D_4\) factor. Set
\(\Lambda_i=r^iN_6\), \(i\in\Z/3\Z\), and
\(J=\Lambda_0\cap\Lambda_1\).

\begin{statement}[B.1 (Proved): the intersection and its discriminant form]
The lattice \(N\) is even unimodular, with root system \(A_5^4D_4\).
The three distinct lattices \(\Lambda_i\) are Leech lattices, every
pair has intersection \(J\), and
\[
 rJ=J,\qquad [\Lambda_i:J]=9,\qquad\det J=81.
\]
Put \(u=(\omega_2,\omega_4,0,\omega_2;0)\),
\(v=\rho/6-\alpha\), and \((e_1,e_2,e_3,e_4)=(u,v,ru,rv)\).
Their classes form a basis of \(J^\vee/J\) over \(\mathbb F_3\).
In this basis,
\begin{equation}\label{eq:B-disc}
 \bigl((e_i,e_j)\bigr)=
 \begin{pmatrix}
 4&2&-2&-4/3\\2&4&-2/3&-2/3\\
 -2&-2/3&4&2\\-4/3&-2/3&2&4
 \end{pmatrix},\qquad
 Q_J(a,b,c,d)=\frac{2ad+bc+bd}{3}\pmod{\Z}.
\end{equation}
In particular this discriminant form is split and nondegenerate.
\end{statement}

\begin{proof}
We give both the exact finite verification and the algebra that makes
it sufficient. The root lattice has determinant \(6^4\cdot4=5184\).
The classes in \eqref{eq:B-glue} span an isotropic subgroup of its
quadratic discriminant group of order \(72\): reduce the four
\(A_5\) coordinates modulo six and the \(D_4\) classes modulo its
root lattice. The orders of the independent primary components are
\(2^3\) and \(3^2\). Their pairings and half norms are integral.
Consequently \(N\) is even and has determinant \(5184/72^2=1\).

Here and below an assertion about a lattice span is verified by integer
Hermite normal form in the ambient coordinates multiplied by twelve.
The complete implementation is \path{code/B_exact_lattices.py};
the constructed bases and its output are in the correspondingly named
\path{records/B_*} files. A basis of \(K_d\) is obtained by taking
the Hermite normal form of
\[
 d\alpha,\qquad g-(\rho,g)\alpha
 \quad(g\text{ a generator of }N).
\]
The homomorphism \(x\mapsto(\rho,x)\pmod d\) is onto since it takes
\(\alpha\) to one, so \([N:K_d]=d\). The added vector \(w_d\)
in \eqref{eq:B-neighbor} pairs integrally with \(K_d\), has square
\(4(6/d)^2\), and has exact order \(d\) modulo \(K_d\).
These facts prove evenness and unimodularity of \(N_d\).

For clarity, intersection is not inferred from the coincidence of any
shells. The identity
\[
 A\cap B=(A^\vee+B^\vee)^\vee
\]
for full lattices is applied with exact rational Gram inverses and
Hermite normal forms. It gives determinant \(81\) for all three pairwise
intersections, identical bases for those intersections, and invariance
under \(r\). The determinant formula for an inclusion in a unimodular
lattice gives index nine. Direct reduction gives \(3e_i\in J\),
\((e_i,J)\subset\Z\), and
\(J+\sum_i\Z e_i=J^\vee\). The last lattice has index \(81\), so the
four classes are independent. Taking half the norm in the displayed
Gram matrix gives \eqref{eq:B-disc}. The plane spanned by its first two
basis vectors is totally isotropic. A nondegenerate four-dimensional
quadratic space containing such a plane is split.

We verify root systems by complete short-vector enumeration, not by a
random or a numerical search. The exact enumeration method is described
below. It gives \(144\) roots for \(N\), the roots already in \(R\),
and no roots in any \(\Lambda_i\). The uniqueness of the rootless even
unimodular positive definite lattice of rank 24 then identifies these
as Leech lattices. This is the rank-24 classification, in the form
stated and proved in Chapter II, Corollary 3.14, of
\citet{B_ChenevierLannes2019}; the rootless case is also established
by the construction and uniqueness argument of \citet{B_Borcherds1985}.
Every hypothesis here has been checked: rank 24, positive definiteness,
evenness, determinant one, and the complete absence of roots.
\end{proof}

\begin{statement}[B.2 (Proved): all eight even unimodular extensions]
There are exactly eight even unimodular lattices containing \(J\) in
\(J\otimes\Q\). In the coordinates of \eqref{eq:B-disc} they are
\(M_i=J+\Z\widetilde a_i+\Z\widetilde b_i\), where a tilde denotes the
linear combination of the \(e_j\) with the indicated coefficients:
\begin{center}\small
\begin{tabular}{ccllc}
\toprule
\(i\)&\(a_i\)&\(b_i\)&Root system&\(rM_i\)\\\midrule
0&0001&0010&none; \(\Lambda_1\)&\(M_7\)\\
1&0001&1100&none&\(M_3\)\\
2&0010&1000&\(A_2^{12}\); \(N_2\)&\(M_2\)\\
3&0011&1200&none&\(M_4\)\\
4&0012&0100&none&\(M_1\)\\
5&0100&1000&none; \(\Lambda_0\)&\(M_0\)\\
6&0112&1020&\(A_2^{12}\)&\(M_6\)\\
7&0121&1010&none; \(\Lambda_2\)&\(M_5\)\\\bottomrule
\end{tabular}
\end{center}
The two rooted extensions have exactly the same set of 72 roots in the
ambient Euclidean space. Moreover \(N_3\) has root system \(A_1^{24}\).
Thus Lead 6 holds, including its assertions about the action of \(r\).
\end{statement}

\begin{proof}
An intermediate lattice \(J\subset M\subset J^\vee\) is even exactly
when \(M/J\) is isotropic for \(Q_J\); it is unimodular exactly when
\((M/J)^\perp=M/J\). Indeed integrality and evenness are respectively
the vanishing of the induced bilinear form and of the quadratic form,
and the dual corresponds to the perpendicular subgroup. In our
nondegenerate group of order \(3^4\), these are precisely the isotropic
planes. The split four-space has 32 nonzero isotropic vectors and
16 isotropic lines. Every isotropic line belongs to two maximal
isotropic planes and each plane contains four lines, hence there are
\(16\cdot2/4=8\) planes. Alternatively substitution in
\(2ad+bc+bd=0\) gives exactly the eight spans in the table; the program
checks all 81 vectors and all pairs. The two rulings each have four
planes, as follows also by fixing one plane and using the parity of
the intersection dimension.

For each span the program computes an even unimodular basis and
exhausts every vector of square at most two. Six counts are zero;
the remaining counts are 72. Joining roots when their inner product
is nonzero partitions each of the latter root sets into twelve
rank-two components of size six, which are \(A_2\). Exact equality
of the two ambient root sets is checked. The same enumeration for
\(N_3\) gives 48 roots in 24 orthogonal rank-one components.
Applying \(r\) to the bases and reducing their spans gives the last
column. This verifies both fixed rooted lattices and the two cycles
of length three on the Leech lattices.
\end{proof}

\paragraph{The exact enumeration certificate.}
The program \path{code/B_enumerate.cpp} uses an integral,
fraction-free version of completing squares. After an integral
unimodular basis reduction, whose change-of-basis matrix and determinant
are checked, let \(G\) be the Gram matrix and \(D_i\) its first \(i\)
principal minors, with \(D_0=1\). Rational elimination gives integers
\(m_{ij}\) such that
\[
 x^tGx=\sum_{i=0}^{23}
 \frac{(D_{i+1}x_i+\sum_{j>i}m_{ij}x_j)^2}{D_iD_{i+1}}.
\]
All \(D_i\) are positive. Enumerate backwards from \(i=23\), using
integer square roots for the bounds. If the residual after the later
coordinates is \(R_{i+1}/D_{i+1}\), the new integral residual is
\[
 R_i=\frac{D_iR_{i+1}-(D_{i+1}x_i+\sum_{j>i}m_{ij}x_j)^2}{D_{i+1}}.
\]
The division is exact by the completed-square identity. Every permitted
integer in the closed interval is visited. This gives an induction
proof of completeness, including points on its boundary. The files
\texttt{*.enum} contain the exact minors and coefficients; files
\texttt{*.basis.json} identify the reduced bases. Root enumerations
also record every vector. No floating point decision enters the
enumeration. The largest minor in the construction of Statements
B.1--B.2 is \(2596\).

\subsection{Minimum, theta series, and the full automorphism group}

\begin{statement}[B.3 (Proved): the theta series of \(J\)]
The minimum of \(J\) is four, its kissing number is \(23058\), and
its theta series is a modular form of weight 12 on \(\Gamma_0(3)\)
with trivial character. With the usual normalizations of \(E_4\),
\(\Delta=\eta^{24}\), and \(\eta\), it is
\begin{equation}\label{eq:B-theta}
\begin{split}
 \Theta_J(\tau)={}&\frac{7381}{66430}E_4(\tau)^3
 +\frac{59049}{66430}E_4(3\tau)^3
 -\frac{531432}{6643}\Delta(\tau)\\
 &-\frac{4251528}{6643}\Delta(3\tau)
 +\frac{88938}{73}\eta(\tau)^{12}\eta(3\tau)^{12}.
\end{split}
\end{equation}
In particular its coefficients through \(q^8\) are
\begin{align*}
 1+23058q^2+1849032q^3+44291124q^4
 &+514246104q^5+3824355978q^6\\
 &+20831769864q^7+90541677312q^8+O(q^9).
\end{align*}
\end{statement}

\begin{proof}
The exact enumeration gives respectively
\(1,0,23058,1849032,44291124\) vectors of squares \(0,2,4,6,8\).
The norm-eight enumeration visits \(215418506\) nodes. Its complete
output and input are \path{records/B_J8.out} and
\path{records/B_J8.enum}; the separate norm-six run is retained.
Positive definiteness and the absence of norm-two vectors then prove
the first two assertions.

For the modular assertion we include the transformation argument.
Evenness gives \(\Theta_J(\tau+1)=\Theta_J(\tau)\). Poisson summation
on the rank-24 lattice gives
\[
 \Theta_J(-1/\tau)=\frac{\tau^{12}}9\Theta_{J^\vee}(\tau).
\]
Equation \eqref{eq:B-disc} gives \(3(x,x)/2\in\Z\) for every
\(x\in J^\vee\), so \(\Theta_{J^\vee}\) is invariant under
\(\tau\mapsto\tau+3\). The weight-12 slash action therefore makes
\(\Theta_J\) invariant under \(T\) and \(ST^3S^{-1}\), and hence
under \(\Gamma_0(3)\). To check the last generation assertion, write
the first column of a matrix as \((a,3c)^t\). Left multiplication by
these two parabolic generators permits \(a\mapsto a+3nc\) and
\(c\mapsto c+ma\). Nearest-integer reduction gives successively
\(|a|\le3|c|/2\) and \(|c|\le|a|/2\), reducing a nonzero
\(|c|\) by a factor at most \(3/4\). The terminal triangular matrix
is generated by \(T\) and \(-I\), the latter acting trivially in
weight 12. Holomorphy at the two cusps follows from the nonnegative
norms in the two Poisson expansions. The lattice level is exactly three, because \eqref{eq:B-disc} has
a value \(1/3\). The theta form itself is not a form for the full
modular group: its constant and linear coefficients would force
it to equal \(E_4^3-720\Delta\), whose quadratic coefficient is
196560 rather than 23058.

Each term on the right of \eqref{eq:B-theta} is holomorphic of weight
12 on \(\Gamma_0(3)\). For the eta product this follows directly from
\(\eta(\tau+1)=e^{\pi i/12}\eta(\tau)\) and
\(\eta(-1/\tau)=\sqrt{-i\tau}\eta(\tau)\): the product is invariant
under \(T\), and its Fricke transformation is
\(g(-1/(3\tau))=3^6\tau^{12}g(\tau)\), which gives the other
generator. It vanishes at both cusps.

Finally, a nonzero weight-12 holomorphic form on \(\Gamma_0(3)\)
cannot vanish to order more than four at infinity. One proof takes
the product of its four slash translates over
\(\Gamma_0(3)\backslash\mathrm{SL}_2(\Z)\); the product is a nonzero
holomorphic form of weight 48 for the full modular group. Its order
at infinity is at least that of the original form, and the usual
valence formula bounds that order by \(48/12=4\). The right side of
\eqref{eq:B-theta} has the five enumerated coefficients, so subtraction
and this bound prove the identity. The coefficient computation, using
integer products for eta and divisor sums for \(E_4\), is
\path{code/B_theta.py}; its five-by-five coefficient matrix has
nonzero determinant \(-387419760\). No extrapolation of shell counts
is used in the identity.
\end{proof}

The equality of the two root sets in Statement B.2 makes the group
calculation substantially smaller. Let \(R'\cong A_2^{12}\) be their
common span. In each component choose roots \(\beta_i,\gamma_i\)
with \((\beta_i,\gamma_i)=-1\), both having the same nonzero class
under the quotient homomorphism \(R'\to R'/(J\cap R')\cong\mathbb F_3\).
These choices, in ambient coordinates, are recorded in
\path{records/B_autJ_data.json}.

\begin{statement}[B.4 (Proved): the full group of \(J\)]
The full isometry group of \(J\) has order
\[
 |\operatorname{Aut}J|=12597120=2\cdot3^7\cdot2880
                      =2^7 3^9 5.
\]
An explicit structural description is
\[
 \operatorname{Aut}J\cong C_2\times
    \bigl((\mathbb F_3)^7\rtimes\operatorname{Aut}_{\rm mon}(D)\bigr),
\]
where \(D\) is the ternary code generated by
\begin{equation}\label{eq:B-code}
\begin{pmatrix}
1&1&0&0&0&0&0&0&2&2&2&2\\
0&0&1&0&0&2&0&2&2&2&0&2\\
0&0&0&1&0&2&0&0&2&1&2&1\\
0&0&0&0&1&2&0&2&1&0&1&1\\
0&0&0&0&0&0&1&2&1&1&2&2
\end{pmatrix}.
\end{equation}
Here \(\operatorname{Aut}_{\rm mon}\) means signed coordinate
permutations. Its order is 2880. Fourteen integral matrices generating
the full group, in the basis \path{B_J_basis.json}, are supplied in
\path{B_autJ_data.json} and verified by \path{B_autJ.py}.
Its action on the eight extensions has order 36 and has two orbits,
the six Leech extensions and the two rooted extensions.
\end{statement}

\begin{proof}
Every isometry of \(J\) permutes its even unimodular extensions and
preserves roots. It therefore preserves \(R'\), intrinsically obtained
from the two rooted extensions. This gives an upper bound by root
system isometries, and is the reason that the computation below
identifies the full group.

Put \(K=J\cap R'\). Exact reduction of the displayed root basis gives
\[
 K=\{\textstyle\sum_i(a_i\beta_i+b_i\gamma_i):
       a_i,b_i\in\Z,\ \sum_i(a_i+b_i)=0\pmod3\}.
\]
For \(x=\sum_i(a_i\beta_i+b_i\gamma_i)\in R'^\vee\), put
\(c_i=3b_i\pmod3\). Then \(a_i+b_i\in\Z\), and the homomorphism
\[
 R'^\vee/K\longrightarrow\mathbb F_3^{12}\oplus\mathbb F_3,
 \quad x\longmapsto
 \left(c,\sum_i(a_i+b_i)-\sum_i c_i\right)
\]
is an isomorphism; representatives of \(c_i\) can be changed by
multiples of three without changing the last coordinate. Reduction of
the 24 generators of \(J\) gives the graph of a functional \(f:D\to
\mathbb F_3\), with values \((0,1,1,2,1)\) on the rows of
\eqref{eq:B-code}. The code has weight enumerator
\(1+90z^6+140z^9+12z^{12}\).

In one component let \(r_i\) send \(\beta_i\) to \(\gamma_i\)
and \(\gamma_i\) to \(-\beta_i-\gamma_i\). Let \(t_i\) be the
involution with row-coordinate matrix
\(\left(\begin{smallmatrix}0&1\\1&0\end{smallmatrix}\right)
\left(\begin{smallmatrix}0&1\\-1&-1\end{smallmatrix}\right)\).
In the quotient coordinates above, \(r_i\) sends
\((c,h)\) to \((c,h-c_i)\), whereas \(t_i\) changes \(c_i\) to
\(-c_i\) and fixes \(h\). The product of the rotations with exponents
\[
 z_0=(0,0,1,1,2,0,1,0,0,0,0,0)
\]
changes the graph of \(f\) to \(D\oplus0\).

An isometry of \(A_2^{12}\) preserving \(K\) must act on its defining
nonzero linear functional by one common scalar \(\varepsilon=\pm1\).
The component root isometries are the dihedral groups of order twelve.
Those with scalar \(+1\) are generated by \(r_i,t_i\); the common
scalar \(-1\) is provided by the central global negation. After the
normalizing rotations, preservation of \(D\oplus0\) is therefore
equivalent to a monomial automorphism of \(D\), and a rotation vector
\(z\in D^\perp\). The latter space has dimension seven. This proves
the claimed group description, both its upper bound and existence.

We certify the finite code group twice. The first computation encodes
all 242 nonzero codewords and the 24 signed coordinates as a coloured
graph; the six returned signed coordinate generators are checked
against every codeword, and an exact permutation-group calculation
gives order 2880. A separate exhaustive computation in
\path{B_code_aut_exhaust.cpp} does not rely on the graph
algorithm. The five pivot columns of \eqref{eq:B-code} are the standard
basis of \(\mathbb F_3^5\). Its twelve columns form eleven projective
classes, one occurring twice. We exhaust all ordered images of the
five pivot columns among these eleven classes, with all \(2^5\) signs:
\(11\cdot10\cdot9\cdot8\cdot7\cdot32=1774080\) choices. Exactly
1440 linear transformations preserve the column multiset. Each has
two lifts to a signed coordinate permutation, arising from the repeated
class. The 2880 resulting permutations are distinct and all belong to
the group generated by the six displayed generators. This proves the
order without an unverified completeness assertion about a search.

The seven rotations from a basis of \(D^\perp\), the six monomial
generators, and global negation give fourteen generators. Conjugation
by the displayed normalizing rotations and by the exact root-to-\(J\)
basis matrix gives the recorded integral matrices. Each has determinant
\(\pm1\) and satisfies \(T G_JT^t=G_J\). Reducing their images of each
extension gives the stated permutation group and its two orbits.
\end{proof}

\begin{statement}[B.5 (Proved): the associated plane and its stabilizer]
Identify \(\Lambda=\Lambda_0=M_5\). The plane
\[
 P=3J^\vee/3\Lambda\ \subset\ \Lambda/3\Lambda
\]
is totally isotropic and has the two generators \(3ru,3rv\pmod{3\Lambda}\).
Its perpendicular preimage in \(\Lambda\) is \(J\). Its full
stabilizer in \(\mathrm{Co}_0=\operatorname{Aut}\Lambda\) has order
\(2099520\). Its \(\mathrm{Co}_0\)-orbit consists of
\(3960692736000\) isotropic planes. Intrinsically, it is the orbit
whose perpendicular lattice is isometric to the \(J\) above; the
choice of a Leech extension introduces no further orbit, because
\(\operatorname{Aut}J\) is transitive on all six such extensions.
Eleven integral generators in the displayed \(\Lambda_0\) basis are
given in \path{records/B_specific_plane_group.json}. This is orbit 4
in the complete classification below; the integral comparison
isometry is supplied in \path{B_specific_to_orbit4.json}.
\end{statement}

\begin{proof}
Since \(3J^\vee\subset J\subset\Lambda\), the quotient is defined.
The image has order nine: \([J^\vee:\Lambda]=9\). Pairings among
its vectors are divisible by three and its norms are divisible by
six, by \eqref{eq:B-disc}, proving isotropy for
\(x\mapsto x^2/2\pmod3\). The equality
\[
 x\in J\quad\Longleftrightarrow\quad
 (x,J^\vee)\subset\Z
 \quad\Longleftrightarrow\quad
 (x,3J^\vee)\subset3\Z
\]
proves the perpendicular assertion. The first two discriminant
basis classes generate \(\Lambda/J\), so the last two give the
specified basis of \(P\).

The stabilizer equals \(\operatorname{Aut}(J)_\Lambda\): in one
direction an isometry preserving \(P\) preserves its perpendicular
preimage; in the other it preserves \(3J^\vee/3\Lambda\). Statement
B.4 and orbit--stabilizer give \(12597120/6=2099520\). The order
\(8315553613086720000\) of \(\mathrm{Co}_0\), twice the listed
order of \(\mathrm{Co}_1\) in the primary ATLAS table
\citep{B_ATLASCo1}, then gives the displayed
orbit length. The isometry class characterization follows because
an isometry between two copies of \(J\) can be adjusted by an
automorphism of \(J\) to match any prescribed Leech extensions.
For the generator assertion we apply Schreier's lemma to the
six-point orbit of \(M_5\) under the fourteen matrices of B.4.
If \(t_i\) takes \(M_5\) to \(M_i\), the words
\(t_i g t_{g(i)}^{-1}\) generate its stabilizer. After change of
basis to \(\Lambda_0\), every retained matrix is integral,
unimodular, and preserves its Gram matrix. Their faithful action on
the 72 intrinsic roots has independently computed order 2099520;
\path{code/B_specific_plane_group.py} retains the words and checks.
\end{proof}

\subsubsection{Complete classification of isotropic planes modulo three}

For an isotropic line \(\ell\subset\Lambda/3\Lambda\), let
\(N(\ell)\) denote its unique even unimodular 3-neighbour. Explicitly,
choose \(v\in\Lambda\) spanning \(\ell\), replace \(v\) by
\(v+3a\) so that \(v^2\equiv0\pmod {18}\), and put
\[
 N(\ell)=\{x\in\Lambda:(x,v)\equiv0\pmod3\}+\Z(v/3).
\]
The replacement exists because the nonzero functional
\((v,{\cdot})\pmod3\) is onto. The resulting lattice is independent
of the replacement and sign of \(v\). Indeed the common index-three
intersection has split discriminant plane; its two isotropic lines
give precisely \(\Lambda\) and \(N(\ell)\). We call \((\ell,P)\),
where \(\ell\subset P\) is a totally isotropic plane, a marked plane.

\begin{statement}[B.9 (Proved): all 129 plane orbits and their full groups]
There are exactly 129 \(\mathrm{Co}_0\)-orbits of totally isotropic
planes in \(\Lambda/3\Lambda\). Of these, 106 have a line whose
3-neighbour is \(A_2^{12}\), another 20 have a line whose
3-neighbour is \(A_1^{24}\), and the remaining three have eight
Leech extensions. Table~\ref{tab:B129} gives every orbit, its full
stabilizer order, its orbit length, and the eight Niemeier extensions
of the perpendicular preimage \(J_P\).

For every row we supply a two-by-24 reduced matrix over \(\mathbb F_3\)
in the same integral Leech basis, not merely an abstract isometry
class. These matrices and all eight extension labels are in
\path{records/B_all_129_plane_orbits.json} and its CSV companion.
Full integral stabilizer generators in that same basis, with exact
order certificates, are in \path{records/B_plane_groups/001.json}
through \path{records/B_plane_groups/129.json}. The sum of the 129
orbit lengths is exactly
\[
 61549107892647466000,
 \qquad
 \sum_{[P]}\frac1{|\mathrm{Co}_{0,P}|}
   =\frac{14703561369481}{1986515435520}.
\]
The eight-extension multiset alone does not determine the orbit:
rows 1 and 12 both have four Leech and four \(A_2^{12}\) extensions,
but their stabilizers have orders 93312 and 103680.
\end{statement}

\begin{proof}
We give the finite reductions, their completeness arguments, and the
independent mass check. All quadratic forms in the finite calculation
are \(x\mapsto x^2/2\pmod3\).

The binary Golay code used in the common coordinate model is the
span of the following twelve hexadecimal masks; bit \(i\) means
coordinate \(i\), starting at zero:
\[
\begin{gathered}
\mathtt{fff0f0},\ \mathtt{ff0ff0},\ \mathtt{f0fff0},\
\mathtt{0ffff0},\ \mathtt{cccc00},\ \mathtt{aaaa00},\\
\mathtt{6ac0c0},\ \mathtt{c6a0a0},\ \mathtt{a6c00c},\
\mathtt{6ca00a},\ \mathtt{11111e},\ \mathtt{ffffff}.
\end{gathered}
\]
In \(\R^{24}\) with scalar product divided by eight, take the
integer span of \(8e_0\), the 23 vectors \(4(e_i-e_0)\), twice
these twelve binary code generators, and \((-3,1^{23})\).
Exact Hermite reduction gives the common basis; its Gram matrix is
even of determinant one, and exact norm-two enumeration is empty.
This identifies it as the Leech lattice. The ternary glue code in
the rooted marked model is the row span over \(\mathbb F_3\) of
\[
\begin{pmatrix}
1&0&0&0&0&1&0&2&2&1&1&0\\
0&1&0&0&0&2&0&1&0&1&1&2\\
0&0&1&0&0&2&0&2&2&2&0&2\\
0&0&0&1&0&2&0&0&2&1&2&1\\
0&0&0&0&1&2&0&2&1&0&1&1\\
0&0&0&0&0&0&1&2&1&1&2&2
\end{pmatrix}.
\]
Its self-duality and weight distribution are checked directly on
all 729 words. These displayed codes specify the inputs independently
of any software convention for Golay coordinates.

\paragraph{The four line orbits and full line stabilizers.}
There are four line types. The familiar labels \(3,6_{22},6_{23},9_{33}\)
refer to vector classes, not to their dimension. We use them only as
names for the explicit representatives in the supplied programs.
Their full line-stabilizer orders and orbit lengths are
\[
\begin{array}{c|r|r|c}
\text{type}&|\mathrm{Co}_{0,\ell}|&|\mathrm{Co}_0\ell|&N(\ell)\\\hline
3&991533312000&8386560&\Lambda\\
6_{22}&110361968640&75348000&\Lambda\\
6_{23}&489646080&16982784000&A_1^{24}\\
9_{33}&277136640&30005248000&A_2^{12}
\end{array}
\]
We justify fullness rather than inferring it from a generated subgroup.
For the last two rows, work in the indicated rooted neighbour \(N\).
The line stabilizer is exactly the group of root-system isometries
preserving both its glue and \(K=N\cap\Lambda\). Preservation of
\(N,K\) preserves the other unimodular extension \(\Lambda\), by
the two-line discriminant argument above.

For \(A_1^{24}\), in orthogonal root coordinates,
\(N=\Z^{24}+\tfrac12 C_{24}\), with the extended binary Golay
code, and \(K=\ker(\rho,{\cdot})\pmod3\),
\(\rho=\tfrac12\sum_i\alpha_i\). A signed permutation preserving
\(K\) must have a common sign on all roots. Thus the full group is
\(C_2\times M_{24}\), of order 489646080. The code graph and the
displayed coordinate permutations independently give the full
\(M_{24}\) order 244823040.
For \(A_2^{12}\), its glue is the self-dual ternary Golay code
\(C_{12}\) of dimension six. In a root basis
\(\beta_i,\gamma_i\) with \(\rho=\sum_i(\beta_i+\gamma_i)\),
the same quotient-functional argument used in B.4 gives
\[
 \mathrm{Co}_{0,\ell}
   =C_2\times\bigl(3^6\rtimes
           \operatorname{Aut}_{\rm mon}(C_{12})\bigr).
\]
The rotations are indexed by \(C_{12}^{\perp}=C_{12}\). The
monomial group, computed on the complete code graph, has order
190080, giving the displayed full order 277136640. All its lifted
24-dimensional matrices, including the stabilizers of the four
codeword strata below, are retained.

For the first two rows we also provide independent finite upper
bounds. Use the standard integral-coordinate Leech model, with
scalar product \(x\cdot y/8\), in
\path{B_standard_Leech_basis.json}. For type 3 take
\(v=(5,1^{23})\), of squared norm six. It is the unique vector of
norm six in its residue class: two such vectors differing by
\(3\Lambda\) would have squared distance at least 36 and at most
24. The 552 norm-four vectors \(w\) with \((v,w)=3\) span the
ambient space. Their Gram values off the diagonal are \(-1,1,2\);
the \(-1\) pairs are \(\{w,v-w\}\). We encode the inner-product-one
graph and separately mark those 276 pairs. Its full permutation
group has order 991533312000. Exact reconstruction of its linear
maps exhibits a generator that fails to preserve the Leech lattice.
The independently verified lattice subgroup of order 495766656000
has index two in that graph group. The lattice vector stabilizer
therefore has precisely that order, and adjoining global negation
gives the stated line stabilizer. The graph, maps, and failed
integrality witness are in \path{B_Co3_upper_bound.json}.

For type \(6_{22}\), take the three vectors whose nonzero coordinate
triples are the permutations of \((-8,4,4)\), in coordinate
positions 2,3,4. They have squared norm 12, sum zero, and lie in the
same nonzero class modulo \(3\Lambda\). These are all the shortest
representatives of that class. Distinct representatives of norm 12
must have inner product \(-6\): their difference has squared norm
at least 36, and the next possible squared norm in \(3\Lambda\)
is 54, larger than the maximum 48. A set of \(k\) such vectors has
Gram matrix \(18I_k-6\mathbf1\mathbf1^t\), whence \(k\le3\).
Their differences divided by three form a triangle of minimal
vectors. Fix two of these minimal vectors \(a,b\), with norms four
and inner product \(-2\). The 891 minimal vectors satisfying
\((w,a)=2,(w,b)=0\) span \(b^\perp\). Their Gram graph has full
group order 55180984320. A verified pointwise triangle subgroup has
order 9196830720. Of the six cosets of that subgroup in the graph
group, exact extension fixing \(b\) shows that just one preserves
the Leech lattice. Consequently the full pointwise triangle group
has order 9196830720. The displayed generators induce all six
permutations of the three norm-12 representatives, and global
negation doubles the class stabilizer to the line stabilizer.
This proves the second row. The complete six-coset test is
\path{B_U6_upper_bound.py}; no untested graph symmetry is counted as
a lattice isometry.

These calculations independently recover the customary names
\(\mathrm{Co}_3\) and \(U_6(2):S_3\). The corresponding known
stabilizer interpretations are recorded in
\citet[Section 2]{B_BrouwerLeech} and
\citet[Section 5, p.~10, statement preceding Corollary 7]{B_Fairbairn2010}.
Fairbairn's hypothesis is an unordered triangle of three minimal
vectors with sum zero; the difference triangle just constructed
has exactly that form. We use our graph and lattice calculations
for the numerical upper bounds. The first two types are distinct
because one has a norm-six representative and the other does not;
the latter assertion also follows from the minimum-distance bound
between a norm-six and a norm-12 vector in one class. The last two
types are distinguished by their rooted neighbours. Finally
\[
 8386560+75348000+16982784000+30005248000
 =47071766560
 =\frac{(3^{11}+1)(3^{12}-1)}2.
\]
This is the number of isotropic lines in the split quadratic
24-space. Thus the four verified line orbits exhaust all lines.

\paragraph{All planes with an \(A_2^{12}\) neighbour.}
For this line the quotient \(\ell^\perp/\ell\) is explicitly
\[
 W=K/(3N+\Z\rho),\qquad\dim_{\mathbb F_3}W=22.
\]
The comparison sends a vector of \(K\) to the same rational vector
in the Leech construction. Its kernel is \(3N+\Z\rho\); the
displayed matrices check the pairing, so it is an isometry of
quadratic spaces. Projection to the glue gives an equivariant
map \(W\to C_{12}\) with 16-dimensional kernel. The four codeword
orbits have weights \(0,6,9,12\) and sizes \(1,264,440,24\).
For a fixed representative of each, the full stabilizer acts
affinely on its 16-dimensional fibre. We give all affine matrices,
their integral lattice lifts, and the quadratic equation in
\path{B_A2_strata_setup.json}. The four complete enumerations give
\[
\begin{array}{c|r|r|r}
\text{codeword weight}&\text{code orbit}&
 \text{isotropic fibre vectors}&\text{fibre orbits}\\\hline
0&1&14467005&16\\
6&264&14348907&54\\
9&440&14348907&64\\
12&24&14348907&14
\end{array}
\]
Each enumeration examines all \(3^{16}=43046721\) vectors and
closes every retained orbit under every generator. Removing the
zero vector gives 147 marked-plane orbits. Since global negation
belongs to the full line group, a fibre orbit of size \(f\) over
a codeword orbit of size \(c\) gives a marked-plane orbit of size
\(fc/2\). Their sum is exactly
\[
 5230235650
   =\frac{(3^{10}+1)(3^{11}-1)}2,
\]
the full number of isotropic lines in \(W\). The dense orbit maps
are supplied losslessly compressed, together with their hashes.

For each of the 147 representatives we construct \(J_P\), its
split four-dimensional discriminant form, and all eight maximal
isotropic subspaces, hence all eight unimodular extensions. Complete
root enumeration identifies each extension. The four extensions
meeting the original Leech lattice in index three correspond
bijectively to the four lines of \(P\); the remaining three other
extensions have intersection index nine. This distinction is
checked by exact lattice intersections, not inferred from a root
label.

We change the marked line in every possible way that again gives
an \(A_2^{12}\) neighbour. In its root system choose
\(\beta_i,\gamma_i\) on which the intersection functional has
value one. Root-coordinate rotations normalize its value on the
glue, and an isomorphism of the complete Golay code graphs maps
the glue to the fixed model. The resulting rational map \(T\)
satisfies both the exact Gram identity and
\(\Lambda T=\Lambda\). Reducing the transported second vector in
the fibre tables gives the new marked orbit. These are actual
isometries, rather than matches of extension multisets. Their
connected components are exactly the unmarked plane orbits:
any isometry between two planes first sends the chosen mark to one
of the four possible marks, after which it lies in the full line
stabilizer already enumerated. The 147 marked orbits fuse to 106
unmarked orbits. All links and representatives are in
\path{B_A2_plane_orbits.json}.

\paragraph{All additional planes with an \(A_1^{24}\) neighbour.}
In the orthogonal root coordinates the corresponding quotient is
\[
 W=\{x\in\mathbb F_3^{24}:\sum_i x_i=0\}/\langle\mathbf1\rangle.
\]
An isotropic vector has all three colour-class sizes divisible by
three. Quotienting by constants and nonzero scalars makes these
three classes unordered. Apart from the zero vector, their sorted
sizes are
\[
\begin{gathered}
(0,3,21),(0,6,18),(0,9,15),(0,12,12),\\
(3,3,18),(3,6,15),(3,9,12),(6,6,12),(6,9,9).
\end{gathered}
\]
We first fix the smallest class. Its size is 0, 3 or 6; there are
four \(M_{24}\)-orbits, with lengths
\(1,2024,21252,113344\) and stabilizers of orders
\(244823040,120960,11520,2160\). The last two lengths sum to
\(\binom{24}{6}\). Exhaustive subset orbits of the second class,
followed by all allowed permutations of the three colours, give
114 marked-plane orbits. Their lengths again sum to 5230235650.
Every extension of every representative is constructed and its
roots are enumerated. Of these marked orbits, 83 have an
\(A_2^{12}\) mark and are mapped by the preceding explicit
isometry construction into the 106 existing classes. The other
31 marked orbits fuse, by all \(A_1^{24}\) changes of mark, to
20 further classes. The binary Golay graph gives the exact
coordinate permutation in every such change. The data and all
cross-identifications are in \path{B_A1_plane_orbits.json}.

\paragraph{The last three classes.}
The 126 classes just obtained contain
\(61549106976378112000\) planes. The residual count is
\(916269354000\). We find three further representatives in the
standard Leech coordinates; their second vectors are recorded by
the following unambiguous packing convention:
\[
 \operatorname{pack}(w)=
   \sum_{w_i=1}2^i+2^{24}\sum_{w_i=2}2^i,
 \quad w_i\in\mathbb F_3,\quad 0\le i<24.
\]
For type 3 the anchor is \((5,1^{23})\), and for type \(6_{22}\)
it is the vector with entries \((-8,4,4)\) in positions 2,3,4.
The table gives the complete orbit under the anchor-class group
(global negation acts trivially on marked planes).
\[
\begin{array}{c|r|r|r|r}
\text{anchor}&\operatorname{pack}(w)&\text{flag orbit}&
 |\mathrm{Co}_{0,P}|&|\mathrm{Co}_0P|\\\hline
3&9192064&37950&52254720&159134976000\\
6_{22}&416910284928&27720&15925248&522161640000\\
6_{22}&5174966569752&12474&35389440&234972738000
\end{array}
\]
The first plane has two type-3 lines and two type-\(6_{22}\) lines;
the other two have four type-\(6_{22}\) lines. Every eligible change
of mark stays in the same displayed flag orbit. All changes are
checked as integral Leech isometries, and their complete orbit
trees, parent edges and generators are retained. Thus these are
three distinct full plane orbits. Exact enumeration shows that
all 24 of their extensions are rootless. The seeded search used
to discover these representatives is not a completeness argument:
completeness follows from their closed group orbits and the exact
mass equality below.

\paragraph{Full group generators and order certificates.}
For clarity we state the constructive recipe used for every row.
In a rooted marked model let \(H=\mathrm{Co}_{0,\ell}\), and let
\(H_P\) be the stabilizer of the marked plane. In the \(A_2\)
model write \(H=C_2\times(R\rtimes M)\), where
\(R=C_{12}\cong3^6\) is the root-rotation group and
\(M=\operatorname{Aut}_{\rm mon}(C_{12})\). For a second quotient
vector \(y\), enumerate its at most 729 images under \(R\),
retaining a rotation \(r_x\) with \(yr_x=x\), and the rotation
kernel \(R_y\). For each signed monomial element \(m\) whose
image \(ym\) lies in that rotation orbit, retain
\(m r_{ym}^{-1}\). These elements and \(R_y\) generate the full
vector stabilizer by the semidirect decomposition. Adjoin global
negation to obtain \(H_P\). This gives a finite exact recipe using
only 729 rotations and 380160 signed monomial elements; it does
not enumerate a huge \(\mathrm{Co}_0\)-orbit. The program reduces
the generating set in its faithful action on the 72 roots and
checks that its order equals the independently known
\(|H|/|H\cdot P|\).

In the \(A_1\) model, adjoin to the 24 coordinate vertices and
759 octad vertices three vertices for the unordered colour classes.
The full coloured-graph group is exactly the \(M_{24}\) partition
stabilizer: each returned permutation is checked against every
octad and the three classes. Adjoin global negation. For either
model, for each mark whose transported flag lies in the original
\(H\)-orbit, compose the mark-changing isometry with an
\(H\)-transporter back to the representative. Adjoining these
matrices gives the full plane stabilizer, since its orbit on the
chosen mark has then been exhausted. Every generator is checked
to be integral in the common Leech basis, to preserve its Gram
matrix, and to preserve the two-dimensional plane modulo three.
Finally the union of the roots in all eight extensions spans the
ambient space and is intrinsic to \(J_P\). Its permutation action
is faithful; the orders of all 126 generated groups are independently
recomputed on that union. The complete program is
\path{B_plane_stabilizer_recipe.py}.

For the three remaining rows we use the saved flag-orbit trees.
If \(t_x\) is the tree word taking the representative to \(x\),
the words \(t_x g t_{xg}^{-1}\) generate its class-flag stabilizer.
Its order is computed in the faithful degree-276 or degree-891
class-group representation, giving respectively
\(13063680,1990656,4423680\). Adjoining global negation doubles
these orders. The exact changes of mark have anchor orbits of
sizes \(2,4,4\), giving precisely the three full orders displayed
above. All tree edges are rechecked by
\path{B_rare_plane_groups.py}; this proves both generation and
fullness without a guessed abstract group name.

\paragraph{Common coordinates and the final mass equality.}
Starting from the standard integral Leech basis, we construct its
type-\(9_{33}\) and type-\(6_{23}\) neighbours and use the root/glue
isometries above to identify them with the two marked models. The
exact maps preserve the forms and send the whole Leech lattice
onto the marked Leech lattice, as checked by integral Hermite
normal forms. Their inverses put every representative and every
group generator in one common basis. These comparison maps are
in \path{B_global_model_isometries.json}.

Independently of all orbit enumerations, an isotropic line has a
22-dimensional split perpendicular quotient, and a plane contains
four lines. Hence the total number of planes is
\[
 \frac14\,
 \frac{(3^{11}+1)(3^{12}-1)}2\,
 \frac{(3^{10}+1)(3^{11}-1)}2
 =61549107892647466000.
\]
The 126 rooted classes and the three new classes are disjoint by
the exhaustive changes of mark and their line types. Their exact
orbit lengths sum to this number: the last three lengths sum to
916269354000, precisely the residual count. Thus no plane remains.
Equivalently, division by
\(|\mathrm{Co}_0|=8315553613086720000\), as in
\citet[Theorem 1.1]{B_BrouwerLeech}, gives the stated rational mass.
The final stdout certificate \path{B_global_assembly.out} prints
both \texttt{RESIDUAL PLANES: 0} and \texttt{RESIDUAL MASS: 0}.
As a separate final check, \path{B_final_group_audit.py} rereads the
129 group files, verifies all 654 integral generator matrices, and
recomputes the same rational mass from their separately generated
full-group orders. Its stdout again gives both residuals as zero;
the per-file hashes are retained in \path{B_plane_group_manifest.json}.
\end{proof}

In Table~\ref{tab:B129}, the extension symbols have the following
fixed meanings; an exponent is the multiplicity of that entire
Niemeier lattice among the eight extensions, not a root multiplicity:
\[
\begin{array}{c|cccccc}
\text{symbol}&L&1&2&3&4&5\\\hline
\text{lattice}&\Lambda&A_1^{24}&A_2^{12}&A_3^8&A_4^6&A_5^4D_4
\end{array}
\qquad
\begin{array}{c|ccccc}
\text{symbol}&6&7&8&9&10\\\hline
\text{lattice}&D_4^6&A_6^4&A_7^2D_5^2&A_8^3&E_6^4
\end{array}
\]
The stratum and flag identifiers in the representative column refer
to the complete affine or subset orbit files described above. The
last three rows use the packed-coordinate convention in the proof.
The JSON and CSV also list all eight extensions separately, with
their exact bases and complete root sets referenced by filename.
{\scriptsize\setlength{\tabcolsep}{3pt}
\begin{longtable}{rllrrl}
\caption{All 129 isotropic-plane orbits. The representative is the stratum and flag identifier, or the packed rare vector; all planes are also supplied in one common integral Leech basis.}\label{tab:B129}\\
ID&mark&representative&$|\mathrm{Co}_{0,P}|$&$|\mathrm{Co}_0P|$&extensions\\\hline\endfirsthead
ID&mark&representative&$|\mathrm{Co}_{0,P}|$&$|\mathrm{Co}_0P|$&extensions\\\hline\endhead
1&A2&0,2&93312&89115586560000&$L^{4},2^{4}$\\
2&A2&0,3&23094720&360062976000&$L^{6},2^{2}$\\
3&A2&0,4&2592&3208161116160000&$L^{2},1^{3},2^{2},5$\\
4&A2&0,5&2099520&3960692736000&$L^{6},2^{2}$\\
5&A2&0,6&23328&356462346240000&$L^{2},2^{5},9$\\
6&A2&0,7&5038848&1650288640000&$L^{3},10,2^{4}$\\
7&A2&0,8&1920&4331017506816000&$L^{2},1^{2},2^{2},4^{2}$\\
8&A2&0,9&48&173240700272640000&$L,1,2^{5},4$\\
9&A2&0,10&12960&641632223232000&$L^{2},1^{3},2^{2},6$\\
10&A2&0,11&1296&6416322232320000&$L,1^{3},2^{4}$\\
11&A2&0,12&7776&1069387038720000&$L^{2},2^{6}$\\
12&A2&0,13&103680&80204027904000&$L^{4},2^{4}$\\
13&A2&0,14&648&12832644464640000&$L,2^{6},6$\\
14&A2&0,15&5184&1604080558080000&$L,2^{6},6$\\
15&A2&0,16&51840&160408055808000&$L^{2},2^{4},6^{2}$\\
16&A2&1,1&48&173240700272640000&$L,2^{6},5$\\
17&A2&1,2&4&2078888403271680000&$L,1,2^{5},4$\\
18&A2&1,3&12&692962801090560000&$L,2^{6},5$\\
19&A2&1,4&48&173240700272640000&$L,2^{5},3,4$\\
20&A2&1,5&4&2078888403271680000&$L,1,2^{4},3^{2}$\\
21&A2&1,6&4&2078888403271680000&$L,2^{5},3,4$\\
22&A2&1,7&4&2078888403271680000&$L,1^{2},2^{4},3$\\
23&A2&1,8&6&1385925602181120000&$L,2^{5},3,4$\\
24&A2&1,9&2&4157776806543360000&$L,1,2^{4},3^{2}$\\
25&A2&1,10&8&1039444201635840000&$L,1,2^{5},4$\\
26&A2&1,11&24&346481400545280000&$L,1,2^{5},4$\\
27&A2&1,12&4&2078888403271680000&$L,1,2^{4},3^{2}$\\
28&A2&1,13&4&2078888403271680000&$L,1^{2},2^{4},3$\\
29&A2&1,14&8&1039444201635840000&$L,1^{2},2^{4},3$\\
30&A2&1,15&36&230987600363520000&$L,2^{6},5$\\
31&A2&1,16&16&519722100817920000&$L,1^{2},2^{4},3$\\
32&A2&1,18&8&1039444201635840000&$L,1^{2},2^{4},3$\\
33&A2&1,19&48&173240700272640000&$L,1^{3},2^{4}$\\
34&A2&1,21&288&28873450045440000&$L^{2},2^{6}$\\
35&A2&1,23&12&692962801090560000&$L,2^{4},3^{3}$\\
36&A2&1,24&8&1039444201635840000&$L,1,2^{4},3^{2}$\\
37&A2&1,25&24&346481400545280000&$L,2^{4},3^{3}$\\
38&A2&1,27&48&173240700272640000&$L,2^{4},3^{3}$\\
39&A2&1,28&48&173240700272640000&$L,2^{4},3^{3}$\\
40&A2&1,29&32&259861050408960000&$L,1^{2},2^{4},3$\\
41&A2&1,30&48&173240700272640000&$L,1^{3},2^{4}$\\
42&A2&1,31&288&28873450045440000&$L^{2},2^{6}$\\
43&A2&1,32&24&346481400545280000&$L,1,2^{4},3,5$\\
44&A2&1,33&12&692962801090560000&$L,1^{2},2^{4},5$\\
45&A2&1,34&4&2078888403271680000&$L,1^{2},2^{3},3,4$\\
46&A2&1,35&16&519722100817920000&$L,1^{3},2^{3},4$\\
47&A2&1,37&16&519722100817920000&$L,1^{3},2^{2},3^{2}$\\
48&A2&1,38&64&129930525204480000&$L,1,2^{4},4^{2}$\\
49&A2&1,39&24&346481400545280000&$L,1^{3},2^{2},3^{2}$\\
50&A2&1,40&24&346481400545280000&$L,1^{3},2^{2},3^{2}$\\
51&A2&1,41&8&1039444201635840000&$L,1^{3},2^{2},3^{2}$\\
52&A2&1,42&16&519722100817920000&$L,1^{4},2^{2},3$\\
53&A2&1,43&48&173240700272640000&$L,1^{4},2^{2},3$\\
54&A2&1,44&192&43310175068160000&$L^{2},1^{2},2^{4}$\\
55&A2&1,46&8&1039444201635840000&$L,1^{3},2^{3},4$\\
56&A2&1,47&48&173240700272640000&$L,1,2^{5},7$\\
57&A2&1,48&480&17324070027264000&$L^{2},1^{2},2^{4}$\\
58&A2&1,49&24&346481400545280000&$L,1^{4},2^{2},3$\\
59&A2&1,50&384&21655087534080000&$L,1^{2},2^{4},8$\\
60&A2&1,51&120&69296280109056000&$L,1^{5},2,4$\\
61&A2&1,52&576&14436725022720000&$L^{2},1^{4},2^{2}$\\
62&A2&1,53&144&57746900090880000&$L,1^{4},2^{2},5$\\
63&A2&1,54&1440&5774690009088000&$L^{2},1^{4},2^{2}$\\
64&A2&2,1&4&2078888403271680000&$L,1^{2},2^{4},3$\\
65&A2&2,2&16&519722100817920000&$L,1^{4},2^{2},3$\\
66&A2&2,6&4&2078888403271680000&$L,1^{3},2^{4}$\\
67&A2&2,7&4&2078888403271680000&$L,2^{5},3,4$\\
68&A2&2,8&16&519722100817920000&$L,1^{3},2^{2},3^{2}$\\
69&A2&2,9&8&1039444201635840000&$L,1,2^{4},3^{2}$\\
70&A2&2,10&8&1039444201635840000&$L,1^{4},2^{2},3$\\
71&A2&2,11&8&1039444201635840000&$L,1,2^{4},3^{2}$\\
72&A2&2,14&4&2078888403271680000&$L,1,2^{5},4$\\
73&A2&2,15&4&2078888403271680000&$L,1,2^{4},3^{2}$\\
74&A2&2,16&8&1039444201635840000&$L,1,2^{4},3^{2}$\\
75&A2&2,18&4&2078888403271680000&$L,1^{2},2^{4},3$\\
76&A2&2,25&64&129930525204480000&$L^{2},1^{2},2^{4}$\\
77&A2&2,26&12&692962801090560000&$L,2^{4},3^{3}$\\
78&A2&2,28&4&2078888403271680000&$L,1^{2},2^{4},3$\\
79&A2&2,30&8&1039444201635840000&$L,1^{2},2^{4},3$\\
80&A2&2,36&48&173240700272640000&$L,2^{4},3^{3}$\\
81&A2&2,37&48&173240700272640000&$L,2^{4},3^{3}$\\
82&A2&2,38&12&692962801090560000&$L,1^{3},2^{4}$\\
83&A2&2,40&288&28873450045440000&$L^{2},2^{6}$\\
84&A2&2,41&12&692962801090560000&$L,1^{4},2^{2},3$\\
85&A2&2,42&864&9624483348480000&$L,1^{4},2^{2},6$\\
86&A2&2,44&24&346481400545280000&$L,1^{2},2^{4},3$\\
87&A2&2,46&72&115493800181760000&$L,2^{6},5$\\
88&A2&2,47&288&28873450045440000&$L,2^{6},6$\\
89&A2&2,50&24&346481400545280000&$L,1^{3},2^{3},4$\\
90&A2&2,53&288&28873450045440000&$L,2^{6},5$\\
91&A2&2,56&648&12832644464640000&$L^{2},2^{6}$\\
92&A2&2,58&1728&4812241674240000&$L^{2},1^{4},2^{2}$\\
93&A2&2,59&72&115493800181760000&$L,1^{3},2^{4}$\\
94&A2&2,60&72&115493800181760000&$L,1^{3},2^{4}$\\
95&A2&2,61&72&115493800181760000&$L,1^{3},2^{4}$\\
96&A2&2,62&648&12832644464640000&$L^{2},2^{6}$\\
97&A2&2,63&288&28873450045440000&$L^{2},1^{4},2^{2}$\\
98&A2&2,64&432&19248966696960000&$L^{2},1^{4},2^{2}$\\
99&A2&3,6&64&129930525204480000&$L,1^{2},2^{4},3$\\
100&A2&3,7&24&346481400545280000&$L,1^{4},2^{2},3$\\
101&A2&3,9&192&43310175068160000&$L,2^{4},3^{3}$\\
102&A2&3,10&240&34648140054528000&$L^{2},1^{4},2^{2}$\\
103&A2&3,11&48&173240700272640000&$L,1^{3},2^{4}$\\
104&A2&3,12&2304&3609181255680000&$L,2^{4},3^{3}$\\
105&A2&3,13&15840&524971819008000&$L^{3},1^{3},2^{2}$\\
106&A2&3,14&4320&1924896669696000&$L^{2},2^{6}$\\
107&A1&0,1&483840&17186577408000&$L^{4},1^{4}$\\
108&A1&0,2&46080&180459062784000&$L^{4},1^{4}$\\
109&A1&0,4&80640&103119464448000&$L^{4},1^{4}$\\
110&A1&0,5&2688&3093583933440000&$L^{2},1^{4},3^{2}$\\
111&A1&0,7&55296&150382552320000&$L^{4},1^{4}$\\
112&A1&0,8&3072&2706885941760000&$L^{2},1^{4},3^{2}$\\
113&A1&0,10&7920&1049943638016000&$L^{3},1^{3},2^{2}$\\
114&A1&0,11&760320&10936912896000&$L^{4},1^{4}$\\
115&A1&1,1&2304&3609181255680000&$L^{2},1^{6}$\\
116&A1&1,5&192&43310175068160000&$L,1^{6},3$\\
117&A1&1,10&96&86620350136320000&$L,1^{6},3$\\
118&A1&1,13&2880&2887345004544000&$L^{2},1^{6}$\\
119&A1&1,23&48&173240700272640000&$L,1^{5},2^{2}$\\
120&A1&1,24&24&346481400545280000&$L,1^{5},2^{2}$\\
121&A1&2,1&2304&3609181255680000&$L^{2},1^{6}$\\
122&A1&2,4&1152&7218362511360000&$L^{2},1^{6}$\\
123&A1&2,8&256&32482631301120000&$L,1^{6},3$\\
124&A1&2,10&2304&3609181255680000&$L^{2},1^{6}$\\
125&A1&2,12&80&103944420163584000&$L,1^{5},2^{2}$\\
126&A1&2,36&32&259861050408960000&$L,1^{5},2^{2}$\\
127&3&9192064&52254720&159134976000&$L^{8}$\\
128&622&416910284928&15925248&522161640000&$L^{8}$\\
129&622&5174966569752&35389440&234972738000&$L^{8}$\\
\end{longtable}
}

\subsection{Weyl-vector neighbours for all rooted Niemeier lattices}

\begin{statement}[B.6 (Refuted, with complete replacement table)]
The proposed rule that the Weyl \(d\)-neighbour has Coxeter number
\(h/d\), for every \(d\mid h\), is false. Of the 96 pairs consisting
of a rooted Niemeier lattice and a divisor of its Coxeter number,
exactly the following eleven violate the rule:
\begin{center}\small
\begin{tabular}{lrrll}
\toprule
Original root system&\(h\)&\(d\)&Neighbour root system&\(h'\ne h/d\)\\\midrule
\(D_{12}^2\)&22&2&\(D_6^4\)&\(10\ne11\)\\
\(D_8^3\)&14&2&\(D_4^6\)&\(6\ne7\)\\
\(E_8^3\)&30&2&\(D_8^3\)&\(14\ne15\)\\
\(E_8^3\)&30&3&\(A_8^3\)&\(9\ne10\)\\
\(E_8^3\)&30&5&\(A_4^6\)&\(5\ne6\)\\
\(D_6^4\)&10&2&\(A_3^8\)&\(4\ne5\)\\
\(E_6^4\)&12&3&\(A_2^{12}\)&\(3\ne4\)\\
\(D_4^6\)&6&2&\(A_1^{24}\)&\(2\ne3\)\\
\(D_{10}E_7^2\)&18&2&\(A_7^2D_5^2\)&\(8\ne9\)\\
\(D_{16}E_8\)&30&2&\(D_8^3\)&\(14\ne15\)\\
\(D_{24}\)&46&2&\(D_{12}^2\)&\(22\ne23\)\\\bottomrule
\end{tabular}
\end{center}
The complete list of all 96 neighbours, including root systems and
root counts, is \path{records/B_all_neighbours.csv}. Every
\(d=h\) neighbour is Leech.
\end{statement}

\begin{proof}
We specify the neighbour construction for arbitrary \(N\), so the
counterexamples are not an ambiguity of neighbour conventions. Choose
a positive system, let \(\rho\) be half its root sum, and let \(\alpha\)
be any simple root. For a rooted Niemeier lattice the components have
the same Coxeter number \(h\), \(\rho\in N\),
\(\rho^2=2h(h+1)\), and \((\rho,\alpha)=1\). These assertions are
respectively Lemmas 2.3, 2.6, and 2.4 of \citet{B_Borcherds1985};
in \citet{B_ChenevierLannes2019} they occur in Chapter III,
Propositions 4.2.1 and 4.2.5, with the common Coxeter number supplied
by the Niemeier root-system classification.
For \(d\mid h\) put
\[
 K_d=\{x\in N:(\rho,x)\in d\Z\},\qquad
 w_d=\rho/d-(h/d)\alpha,\qquad N_d=K_d+\Z w_d.
\]
The square of \(w_d\) is \(4h^2/d^2\), and \((w_d,K_d)\subset\Z\).
The vector \(dw_d=\rho-h\alpha\) lies in \(K_d\). Its order modulo
\(K_d\) is exactly \(d\): pairing a putative smaller multiple with
\(\alpha\), in the unimodular lattice \(N\), would require \(d\)
to divide that multiple because
\((\rho-h\alpha,\alpha)=1-2h\) is prime to \(d\).
Thus \(N_d\) is even and unimodular. Replacing \(\alpha\) by another
simple root changes \(w_d\) by an element of \(K_d\), so this
construction is independent of that choice. A Weyl change of positive
system gives an isometric neighbour.

For reproducibility the original integral Gram matrices for all 23
rooted lattices were retrieved from the catalogue of
\citet{B_NebeSloaneCatalogue}. We did not accept the labels of those
matrices as a certificate. The program \path{B_all_neighbours.py}
checks that each is positive definite, even and of determinant one,
enumerates all norm-two vectors, and verifies the stated root system.
Half the sum of the lexicographically positive roots gives an integral
\(\rho\), and the equalities above are independently checked in every
case. A root of pairing one supplies \(\alpha\). The kernel basis is
constructed by an elementary unimodular transformation of its pairing
row, followed by the single congruence condition modulo \(d\). An
exact rational span gives \(N_d\); evenness and determinant one are
checked again.

Every norm-two vector of each of the 96 resulting lattices is then
enumerated by the integer procedure already proved complete. The
connected components of the graph of nonorthogonal roots are
classified by their exact rank and number of roots. For an irreducible
simply laced positive definite root system these distinguish
\(A_n,D_n,E_6,E_7,E_8\), with the harmless equality \(A_3=D_3\).
The original and neighbour Gram matrices, Weyl vector, divisor,
and complete roots are retained for every pair. This is a finite
exhaustion of all divisors for all 23 possibilities, and yields the
table. The error in the proposed rule is the unwarranted assertion
that the quotient of a Coxeter number controls the roots surviving
or appearing in a smaller-divisor neighbour; the example
\(E_8^3,d=2\) already disproves it. Indeed there is not even an
allowed rank-24 Niemeier root system of common Coxeter number 15.

The final assertion also has a uniform proof: Chapter III,
Theorem 4.2.10(c), of \citet{B_ChenevierLannes2019} says that for an
even unimodular positive definite rank-24 lattice with roots and
Coxeter number \(h\), its Weyl \(h\)-neighbour is Leech.
Proposition 4.2.11 identifies this neighbour with the holy
construction. Our formula uses their correction vector of pairing
one and hence is precisely that neighbour. All 23 inputs satisfy
those hypotheses; the computation verifies the conclusion separately.
\end{proof}

\subsection{The Wilson cubic and an integral Niemeier embedding}

Write \(\mathbb O\) for the real division octonions with orthonormal
basis \(1,i_t\), \(t\in\Z/7\Z\), and multiplication
\(i_ti_{t+1}=i_{t+3}\), cyclically in every such triple. Put
\[
 s=(-1+i_0+\cdots+i_6)/2,
 \qquad L=D_8+\Z s,\qquad M=L\bar s,\quad H=Ls.
\]
Here \(L\) has 240 vectors of square two: the vectors
\(\pm e_i\pm e_j\), and the half-vectors with an odd number of minus
signs. Wilson's lattice is
\begin{equation}\label{eq:B-Wilson}
 \mathcal W=\{(x,y,z)\in L^3:
    x+y,y+z,z+x\in M,\ x+y+z\in H\},
\end{equation}
with squared length \((|x|^2+|y|^2+|z|^2)/2\). This is the
normalization in \citet{B_Wilson2009}. The cubic is
\(F(x,y,z)=\operatorname{Re}((xy)z)\); its polarized trilinear
expression is unambiguous because the octonion associator has real
part zero.

\begin{statement}[B.7 (Proved): the full Wilson cubic stabilizer]
The full stabilizer of \(F\) in \(\operatorname{Aut}\mathcal W\)
is
\[
 C_2^2\rtimes C_3\ \times\ (C_7\rtimes C_3)
 \ \cong\ A_4\times(C_7\rtimes C_3),
 \qquad\text{of order }252.
\]
Generators are the even sign changes of the three octonion
coordinates, their cyclic permutation, and the simultaneous Fano
permutations \(i_t\mapsto i_{t+1}\) and \(i_t\mapsto i_{2t}\).
Thus the subgroup exhibited in the workbench is the full stabilizer.
\end{statement}

\begin{proof}
First we reduce the continuous symmetry question to an explicitly
bounded finite calculation. The singular locus of \(F\) is the union
of the three coordinate octonion spaces. In fact its three derivatives
are, after octonion conjugation, the products of the other two
coordinates. The division property implies that at most one
coordinate is nonzero. A linear isometry preserving \(F\) must
therefore permute these three eight-dimensional subspaces.

The projections of \(\mathcal W\) to each block equal \(L\).
The projection of its intersection with a sum of two blocks to either
one of those blocks equals \(M\). To see this, \(M,H\subset L\)
contain \(2L\), their images modulo \(2L\) are complementary
four-dimensional subspaces, and the two-block section has
\(x,y\in M\), \(x+y\in2L\). These identities follow either directly
from the multiplication table or from the exact bases in
\path{B_wilson_stabilizer.py}. Consequently each block of a
lattice symmetry belongs to \(G=\operatorname{Aut}(L)_M\).

We determine this full projection group, rather than assuming that
familiar octonion maps generate it. The 240 roots of \(L\) are
explicit above. Choose the positive roots lexicographically and their
half sum; the eight roots pairing one with this half sum are simple.
Their root reflections generate the Weyl group of \(E_8\), which is
all \(\operatorname{Aut}(L)\), because the \(E_8\) Dynkin diagram
has no automorphisms. The exact permutation action on the roots gives
order \(696729600\). Its orbit on the subspace \(M/2L\) has size
270. The seven right multiplications \(R_{i_t}\), and the six maps
\[
 \frac12 R_{1+i_t}R_{1-i_0},\qquad t=1,\ldots,6,
\]
all preserve \(M\), have determinant \(+1\), and generate a group
of order \(2580480=696729600/270\). Orbit--stabilizer proves that
these generate the full \(G\). These checks, including the entire
finite subspace orbit, are \path{B_E8_projection.py}.

Consider now an ordered-block cubic symmetry \((A,B,C)\), with
\(A\in G\), and set \(D=\bar{\phantom{x}}\,C\,\bar{\phantom{x}}\),
where the bars denote the linear conjugation operator. The cubic
identity is equivalent to
\begin{equation}\label{eq:B-triality}
 (Ax)(By)=D(xy)\quad(x,y\in\mathbb O).
\end{equation}
For a fixed \(A\) there are at most two orthogonal solutions
\((B,C)\). Indeed the ratio of two solutions has first component
identity. Equation \eqref{eq:B-triality} with \(x=1\), then \(y=1\),
shows that its remaining maps are right multiplication by one unit
\(b\), and that \(x(yb)=(xy)b\) for all \(x,y\). The nucleus of
\(\mathbb O\) is \(\R\), hence \(b=\pm1\). In the supplied
multiplication table this last statement can also be verified by the
64 linear equations in the eight coordinates of \(b\); their
solution space is exactly \(\R\cdot1\).

For each of the thirteen generators of \(G\) we solve
\eqref{eq:B-triality} exactly. If \(a=A1\), then \(D=L_aB\), so
\((Ax)(By)=aB(xy)\) is a system of 512 rational linear equations in
the 64 entries of \(B\). In each case its nullspace is one dimensional
and admits a rational orthogonal normalization. The matrices are
recorded in \path{B_triality_lifts.json}. Adjoin the kernel
\((I,-I,-I)\). Their faithful permutation action on the roots in
the first and second blocks has order \(5160960=2|G|\).
This proves that we have exhausted all possible first projections
and their two lifts.

It remains to impose the lattice conditions on the other blocks.
The orbit of \(L\) under the third components of this lift group
has size 15, giving a stabilizer of order \(344064\). Its Schreier
generators are obtained by exact rational lattice Hermite normal
forms. On this subgroup all three components preserve \(L\).
Modulo \(2L^3\), the glue of \(\mathcal W\) is the twelve-dimensional
binary code
\[
 \{(m_1+n,m_2+n,m_1+m_2+n):
        m_1,m_2\in M/2L,\ n\in H/2L\}.
\]
The orbit of this code has size 4096. The ordered-block stabilizer
therefore has order
\(5160960/(15\cdot4096)=84\). Exact Schreier generators are retained
in \path{B_wilson_stabilizer.json}; their generated order is
independently checked in their permutation action. The code orbit
is exhaustive: a queue starts at the displayed code and closes under
every generator, comparing reduced binary echelon bases exactly.
Thus the order is an upper bound as well as a lower bound.

No odd permutation of blocks can occur. After a cyclic permutation
and composition with one of the even triality lifts having the same
first block, such a symmetry would have first block \(I\) and
interchange the other two. Its identity would force
\(x(yb)=(yx)b\) for all \(x,y\), for a unit \(b\). This is
impossible. For completeness the same exact linear calculation gives
a zero-dimensional solution space in the eight coordinates of \(b\),
whereas the non-reversed equation above gives dimension one. This is
a finite verification of the obstruction over \(\Q\), and hence
over \(\R\), since the coefficients are rational.

Finally the listed even sign changes, coordinate cycle and two Fano
permutations preserve \eqref{eq:B-Wilson} and the cubic. The first
two operations generate \(A_4\); the Fano permutations fix \(s\),
preserve the multiplication table, and generate \(C_7\rtimes C_3\).
The two groups commute and have trivial intersection. They give
\(4\cdot3\cdot21=252\) elements, precisely three times the
ordered-block upper bound 84. This proves both fullness and the
claimed isomorphism type.
\end{proof}

\begin{statement}[B.8 (Proved): a Niemeier embedding with integral twice cubic]
The Niemeier lattice \(E_8^3\), with its usual norm-two roots, has the
explicit isometric embedding \(L^3\subset\mathbb O^3\), where
\(L=D_8+\Z s\) as above and the Euclidean norm is
\(|x|^2+|y|^2+|z|^2\). Its cubic value group, meaning the additive
subgroup generated by all \(F(x,y,z)\), is
\[
 \langle F(L^3)\rangle_{\Z}=\tfrac12\Z.
\]
In particular \(2F\) is integer valued on this embedding. This
arithmetic property depends on the embedding: rotating the real and
\(i_0\) axes of the first block through \(\pi/4\) gives another
isometric embedding whose cubic value group is not \(\tfrac12\Z\).
More precisely, the second value group is
\(\tfrac12\Z+\tfrac{\sqrt2}{4}\Z\).
\end{statement}

\begin{proof}
Here the ambient metric is the ordinary octonion sum norm, the
coefficient of \(-\lambda\) in the Albert determinant
\(\lambda^3-\lambda(|x|^2+|y|^2+|z|^2)+2F(x,y,z)\).
Wilson's separately specified half-sum norm in B.7 is identified
with this metric by \((x,y,z)\mapsto(x,y,z)/\sqrt2\).
That change multiplies the cubic by \(1/(2\sqrt2)\) and leaves
its stabilizer unchanged; B.8 uses the ordinary metric throughout.
The lattice \(L\) is even unimodular: \(D_8\) has determinant four,
\(s^2=2\), \((s,D_8)\subset\Z\), and adjoining \(s\) has index
two. Thus its orthogonal cube is the stated Niemeier lattice.
The real trilinear expression is cyclic in its three arguments.
If all arguments belong to \(D_8\), its value is integral; if one
argument is \(s\), it lies in \(\tfrac12\Z\). If two arguments
are \(s\), cyclicity makes them adjacent and the coordinates of
\(s^2\) are half integers, so the same conclusion follows. Finally
\(\operatorname{Re}(s^3)=5/2\). Trilinearity in a generating set
therefore proves \(F(L^3)\subset\tfrac12\Z\). On the other hand
\((s,s,s)\) gives \(5/2\), and \((2,2,2)\) gives 8, so their
integer combinations generate \(\tfrac12\Z\).

Let \(T\) act orthogonally on the first octonion block by
\(T1=(1+i_0)/\sqrt2\), \(Ti_0=(-1+i_0)/\sqrt2\), fixing the
remaining six coordinates. The second embedding is
\((x,y,z)\mapsto(Tx,y,z)\). Its value at \((2,2,2)\) is
\(4\sqrt2\), which is not in \(\tfrac12\Z\). Thus the arithmetic
criterion offered in B.4 is satisfied by an explicit embedding, but
does not make the value group an invariant of the abstract Niemeier
lattice. No assertion of uniqueness or canonicality of that embedding
is implied by the word ``distinguished.''

We also compute the entire second group, rather than only an
irrational witness. Write \(e_0=1,e_{t+1}=i_t\), and take the
\(E_8\) basis
\(s,d_1,\ldots,d_7\), where
\(d_j=e_j-e_{j+1}\) for \(1\le j\le6\), and
\(d_7=e_6+e_7\). Its Gram determinant is one. Trilinearity reduces
the value group to the 512 triples of these basis vectors.
In \path{B_cubic_embeddings.py} every value is computed exactly
as a pair of rational coefficients of \(1,\sqrt2\); all have
first coefficient in \(\tfrac12\Z\) and second in
\(\tfrac14\Z\). Conversely the three exact values
\[
 F(Ts,s,d_4)=\tfrac12,\qquad
 F(Ts,d_7,s)=\tfrac{\sqrt2}{2},\qquad
 F(Ts,s,s)=\tfrac32+\tfrac{3\sqrt2}{4}
\]
generate \(\tfrac12\Z+\tfrac{\sqrt2}{4}\Z\). The complete
coefficient list and its two-dimensional integer Hermite reduction
are retained in \path{B_cubic_embeddings.json}. Thus both embeddings
and both distinct value groups have exact finite certificates.
\end{proof}

\paragraph{Literature scope.}
Borcherds' construction and the Weyl-neighbour results of Chenevier and
Lannes supply the rank-24 background and the endpoint \(d=h\); they
do not supply the eleven-exception table. Wilson's 2009 construction
supplies \eqref{eq:B-Wilson}, and his account of Conway's group supplies
large generating groups preserving that lattice. A generator of
\(\mathrm{Co}_0\) need not preserve the cubic, which is why the
separate upper-bound argument in B.7 is required. Searches for the
related Albert-algebra interpretation also locate
\citet[Corollary 2.3 and Section 4]{B_Rios2018}. Its equation (8)
uses the quadratic map \(v\mapsto vv^*\) into the full Albert
algebra. That construction does not identify the fixed linear
off-diagonal coordinates with a cubic-preserving action on
\(\mathbb O^3\); in particular \(v\) and \(-v\) have the same
quadratic image, whereas our nonzero cubic changes sign. We
therefore do not infer a cubic stabilizer from that interpretation.
Searches for the specific determinant-81 intersection, its shell counts 23058 and
1849032, and its group orders, together with searches for isotropic
planes modulo three, did not identify a published occurrence of this
configuration or a full plane-orbit table. We therefore make no
priority claim and do not label the configuration new. The exact
searches, sources opened, and theorem locations are retained in the
Part B research record. The full 129-orbit table in B.9 is an
independent exact classification; the failure to locate an earlier
table is not a claim of priority. The vector-residue table in
\citet[Proposition 4.2 and Table 4.2]{B_MartinetLatticesMod} and the
explicit vector transporters documented by \citet{B_mmgroup2026}
are useful comparisons, but neither supplies the plane fusion or
the 129 full stabilizers proved here. The known triangle stabilizer
in \citet[Section 5]{B_Fairbairn2010} is independently recovered by
our complete Gram-graph coset test. The Albert/Niemeier connection
itself is present in the work of \citet{B_ElkiesGross2001}; B.8
addresses the particular Euclidean embedding and cubic normalization
asked for here, and its elementary value-group calculation makes no
claim to introduce that broader connection.

The requested Conway--Sloane comparison is more specific than a
general reference to their book. Their Chapter 16 tables record the
rooted Niemeier systems and glue; the opened Chapter 17 material
classifies odd unimodular lattices through their two even neighbours,
an index-two construction. For Chapter 24 we opened the complete
original article \citep{B_ConwaySloane1982Holy}, whose Figure 3
includes both \(A_5^4D_4\) and \(A_2^{12}\). Its unnumbered holy
construction on p.~275, proved on pp.~281--283, constructs a Leech
lattice from each rooted Niemeier lattice, with the norm-two root
normalization used here. This is known background for B.1 and B.3.
The opened material does not state the common determinant-81
intersection, its six Leech and two rooted extensions, or the
129 plane orbits. We distinguish that precise search outcome from
an assertion that the configuration is absent from all literature.
}{}
\IfFileExists{sections/C_special.tex}{\section{C. Special values, local expansions and torsion}
\subsection{C.1. The normalized Green function and special arguments}
We write $\rho=(1+i\sqrt3)/2$, $T=\operatorname{Im}\tau$, and
\[
 U_\tau(z)=\log N(z,\tau)
 =\log\left|\frac{\vartheta _1(z,\tau)}{\eta(\tau)}\right|
       -\frac{\pi(\operatorname{Im}z)^2}{T}.
\]
Our theta variable has periods $1,\tau$; thus the standard theta series
is evaluated at $\pi z$ when its first argument has period $\pi$.
For a general lattice $\Lambda=\omega(\mathbb Z+\tau\mathbb Z)$ we set
$U_\Lambda(z)=U_\tau(z/\omega)$. The transformation laws of the theta
function show that this definition is independent of the oriented basis.

\begin{statement}[C.1.1; Proved: distribution under isogenies]
Let $\Lambda,\Lambda'\subset\mathbb C$ be lattices and let
$\alpha\ne0$ satisfy $\alpha\Lambda\subset\Lambda'$. Write
$\phi:\mathbb C/\Lambda\to\mathbb C/\Lambda'$ for $z\mapsto\alpha z$,
$K=\alpha^{-1}\Lambda'/\Lambda$, and
$d=[\Lambda':\alpha\Lambda]$. Then
\begin{equation}\label{eq:cs-isogeny}
 U_{\Lambda'}(\alpha z)=\sum_{v\in K}U_\Lambda(z+v).
\end{equation}
The equality holds away from the zeros and as an equality of locally
integrable functions at them. If $\Delta_\Lambda=g_2(\Lambda)^3-
27g_3(\Lambda)^2$, then
\[
 \prod_{0\ne v\in K}N_\Lambda(v)
 =|\alpha|\left|\frac{\Delta_{\Lambda'}}{\Delta_\Lambda}\right|^{1/12}.
\]
In particular, for every nonzero endomorphism of one elliptic curve,
$\prod_{v\in\ker\alpha,\ v\ne0}N(v)=|\alpha|$.
\end{statement}
\begin{proof}
The Jacobi product, with $q=e^{2\pi i\tau}$ and
$X=e^{2\pi iz}$, gives
\[
 \vartheta _1(z,\tau)/\eta(\tau)
 =i q^{1/12}X^{-1/2}(1-X)
     \prod_{m\ge1}(1-q^mX)(1-q^mX^{-1}),
\]
where a harmless fixed unit changes with the theta convention. For
$0<y<T$, integration of the logarithm over $z=x+iy$, $0\le x\le1$,
gives $-\pi T/6+\pi y-\pi y^2/T$. Each product factor has mean logarithm
zero by the absolutely convergent logarithmic series; the endpoint
singularities are integrable. Integrating in $y$ proves
$\int U_\Lambda\,d\mu_\Lambda=0$ for Haar probability measure.
Moreover
\[
 \Delta_{\rm Eucl}U_\Lambda
   =2\pi\delta_0-\frac{2\pi}{\operatorname{area}\Lambda},\qquad
 U_\Lambda(z)=\log|z|+\frac1{12}\log|\Delta_\Lambda|+o(1).
\]
The second formula follows from
$\vartheta _1'(0,\tau)=2\pi\eta(\tau)^3$ and
$\Delta_{\mathbb Z+\tau\mathbb Z}=(2\pi)^{12}\eta(\tau)^{24}$,
then scaling by $\omega$. The two sides of
\eqref{eq:cs-isogeny} have the same logarithmic masses at $K$ and the
same constant Laplacian, since
$d=|\alpha|^2\operatorname{area}\Lambda/
\operatorname{area}\Lambda'$. Their difference is a smooth harmonic
function on the compact torus, hence constant. Both means vanish:
translations preserve Haar measure and an isogeny sends Haar
probability measure to Haar probability measure. The constant is zero.
Taking $z\to0$ and cancelling the logarithmic singularity gives the
product. This also checks the precise normalization in checkpoint K1;
its kernel, Laplacian and mean-zero argument are all valid.
\end{proof}
\begin{statement}[C.1.2; Proved: the phases and the two periods]
With the branch continued from the specified cusp, the exact values are
\begin{align*}
 \frac{\vartheta _1((2-\rho)/3,\rho)}{\eta(\rho)}
   &=e^{\pi i/12}\,3^{1/4}e^{\pi/(6\sqrt3)},\\
 \frac{\vartheta _1((1-i)/2,i)}{\eta(i)}
   &=\sqrt2\,e^{\pi/4},\\
 \beta^{\rm part}(z_1)
   &=-\frac12+i\left(\frac1{\sqrt3}-\frac3{2\pi}\log48\right),\\
 \beta^{\rm part}(z_2)
   &=i\left(\frac32-\frac3\pi\log6\right).
\end{align*}
Consequently $N(z_1)=3^{1/4}$, $N(z_2)=\sqrt2$,
$D_{\rm part}(z_1)=-3\log48/(2\pi)$, and
$D_{\rm part}(z_2)=-3\log6/\pi$.
\end{statement}
\begin{proof}
We use the checked K1 calculation, rather than treat its modulus as an
open question: multiplication by $1+\rho$ has matrix
$\left(\begin{smallmatrix}1&-1\\1&2\end{smallmatrix}\right)$,
its kernel is $\{0,\mu,-\mu\}$, and $(1+\rho)\mu=1$ for
$\mu=(2-\rho)/3$. Statement C.1.1 gives $N(\mu)^2=\sqrt3$.
For the requested parallel calculation at $i$, multiplication by $1+i$
has kernel $\{0,(1-i)/2\}$ and absolute value $\sqrt2$; the same
identity gives $N((1-i)/2,i)=\sqrt2$ without a separate theta-constant
evaluation.

At $\rho$, $\mu=1/2-i/(2\sqrt3)$ and $q=-e^{-\pi\sqrt3}$. In
\[
 R(z,\tau)=2e^{\pi i\tau/6}\sin(\pi z)
 \prod_{m\ge1}\bigl(1-2q^m\cos(2\pi z)+q^{2m}\bigr),
\]
$\sin(\pi\mu)$ is positive real. The factors are
$(1+q^m e^{\pi/\sqrt3})(1+q^m e^{-\pi/\sqrt3})$, both positive:
when $m$ is odd the larger subtracted number is at most
$e^{-2\pi/\sqrt3}<1$. Thus the phase is exactly $\pi/12$ modulo
$2\pi$. There is no winding ambiguity in the continued logarithm.
Indeed continue from the cusp along $\tau=1/2+iT$,
$T\ge\sqrt3/2$. Real conjugation of the rhombic lattice keeps the
chosen zero on $\mu=1/2-iy$, $0<y<T$. It cannot leave that component:
a collision of the two zeros would require $g_3=0$, which does not
occur on this ray. Indeed, writing $q=-r$, $0<r<1/200$,
the tail after the first term of the $E_6$ series has absolute value
at most $504\cdot32r^2/((1-r^2)(1-8r))<1/2$, while its
first term is positive. Thus $E_6>1/2$. Hitting $y=0$ would again give a half-period zero, and
hitting $y=T$ would give a pole. Each displayed product factor remains
positive because $e^{-2\pi T}e^{2\pi y}<1$. The cusp phase tends to
$\pi/12$, so the continued argument has that value. On the imaginary
ray $\tau=iT$, $T\ge1$, the same branch is
$\mu=1/2-iy$ with $0<y\le T/2$; every product factor is positive,
so its argument at $i$ is zero. These checks supply the branch step
used in R3. The Gaussian factors are respectively
$e^{\pi/(6\sqrt3)}$ and $e^{\pi/4}$. Substitution in the definition of
$\beta^{\rm part}$ proves all four identities.
\end{proof}

\begin{statement}[C.1.3; Proved: numerical cross-check]
The independent analytic computation in
\path{code/Cspecial_torsion_numerics.py} gives, to fifty decimals,
\begin{align*}
 N(z_1)&=1.31607401295249246081921890179699905516006859020582\ldots,\\
 N(z_2)&=1.41421356237309504880168872420969807856967187537695\ldots,\\
 \operatorname{Im}\beta^{\rm part}(z_1)
  &=-1.27101206057537346286579673538412682778809507094949\ldots,\\
 \operatorname{Im}\beta^{\rm part}(z_2)
  &=-0.21100425815613411418178314271018934802718590587233\ldots.
\end{align*}
The theta/eta evaluations and exact expressions differ by less than
$10^{-100}$ at working precision $110$ decimal digits.
\end{statement}
\begin{proof}
The program directly evaluates the theta series and the eta product,
then compares the displayed closed forms. The eta product has $160$
factors. Since $|q|\le e^{-\pi\sqrt3}<1/200$, the omitted logarithm is
at most $\sum_{m>160}|q|^m/(1-|q|^m)<10^{-368}$.
The output records the measured floating-point discrepancies. The
exact values are established by C.1.2; this computation is their
independent numerical check, not an additional rounding certificate.
\end{proof}

\subsection{C.2. The Albanese coordinates and real structures}
\begin{statement}[C.2.1; Proved: exact nomes and continuation radius]
For $w=e^{-\pi ic}$, the fixed branches give
\begin{align*}
 \kappa_1&=c+\tfrac12-\frac{3i}{2\pi}\log48,&
 q_1&=-w^2/48^3,\\
 \kappa_2&=c+\tfrac32-\frac{3i}{\pi}\log6,&
 q_2&=iw/216,\qquad q_1=(27/64)q_2^2.
\end{align*}
The family is parameterized by $0<|q_2|<1$ and exists precisely in
that domain. Writing $J(q)=q^{-1}+744+\cdots$, one has
\[
 j(\operatorname{Alb}S_1)=J(-w^2/48^3).
\]
This function is holomorphic on $0<|w|<48^{3/2}$, has a pole of order
two at zero, and has a natural boundary at $|w|=48^{3/2}$.
In particular it continues holomorphically through every point of the
family boundary $|w|=216$.
\end{statement}
\begin{proof}
The precise Albanese lattices in R5 are verified by projecting all
translations on the product covers. At the order-three fibre the
projection is $(1/3)(\mathbb Z+\kappa_1\mathbb Z)$; multiplication by
three gives the first stated lattice. At the order-four fibre it is
$\tfrac12\mathbb Z+\tfrac14\kappa_2\mathbb Z$; multiplication by two
gives $\mathbb Z+(\kappa_2/2)\mathbb Z$. Thus the positive-imaginary
period representatives are $-\kappa_1$ and $-\kappa_2/2$.
Adding $2\mu_1=1-i/\sqrt3$ and
$3\mu_2=3/2-3i/2$ to C.1.2 proves the two $\kappa$ identities,
including their real parts. Exponentiation gives the signs in $q_1,q_2$;
the ratio is $216^2/48^3=27/64$.

We checked checkpoint K2's step giving $q_2=iw/216$: its factor
$1/2$ is exactly the Albanese projection just computed. Its
natural-boundary proof uses modular invariance at rational boundary
points and is valid. The existence criterion and exact $2\mathbb Z$
identification are R4 and R5; the latter is also Theorem~1.4
(label \texttt{p2:thm:exact-period}) of
\cite{DuGuoKusakabeXie2026v1}. Its hypotheses specify precisely the
retained period family, affine finite twists and cusp collar. They
coincide with the family in this prompt; the normalization differs
only by the fixed choice of the particular solution, which does not
alter parameter differences. Version 2 drops that theorem; we infer
no mathematical conclusion from this omission.

The Fourier expansion for $J$ converges for $0<|q|<1$ and has a simple
pole at $q=0$. The inverse image under $q=-w^2/48^3$ is the claimed
punctured disc and the pole has order two. The inequality
$216<48^{3/2}$ is equivalent to $27<64$. At every nonzero point on the
larger boundary the map to $q$ is locally biholomorphic. A continuation
there would continue $J$ through the unit circle, contradicting the
same rational-cusp argument checked in K2. The phrase ``up to the
radius'' therefore means the open punctured disc, not holomorphy on
its bounding circle or at its central pole.
\end{proof}

\begin{statement}[C.2.2; Proved: the real locus]
An antiholomorphic equivalence $X_c\to X_{c'}$ exists exactly when
$c'-(1-\bar c)\in2\mathbb Z$. Consequently $X_c$ admits an
antiholomorphic involution exactly when $q_2\in\mathbb R\setminus\{0\}$.
\end{statement}
\begin{proof}
We check the additional reflection argument, rather than infer a real
structure from the reality of one elliptic invariant. On the period
cover, reflection acts by
\[
 (\tau,\mu,\beta)\longmapsto(-\bar\tau,1-\bar\mu,-\bar\beta).
\]
The theta product and the cusp normalization prove this equality of
continued period functions. For the integral monodromy matrices
$A_1,A_2$ of the retained construction, the lattice comparison is
\[
 H=\begin{pmatrix}1&0&0&0\\0&1&0&0\\-6&0&-1&0\\0&-1&0&-1\end{pmatrix},
 \quad H^2=I,\quad HA_2H=A_2^{-1},\quad
 HA_1H=A_2^{-1}A_1^{-1}A_2,
\]
\[
 \Pi_{-\bar c}(\iota z)H=-\overline{\Pi_c(z)}.
\]
Thus $(z,\zeta)\mapsto(\iota z,-\bar\zeta)$ is an actual torus-family
map. Write $d=(0,0,0,1)^t$, $\varepsilon=(1,2,-4,0)^t$ and
$\varepsilon'=(1,3,-3,0)^t$. Reversing the finite deck generators
changes the target affine vectors to
\[
 v'_1=-\varepsilon+d,\qquad v'_2=\varepsilon'-3d.
\]
The second equality includes the sign of the inverse affine action.
The file \path{Cspecial_reflection.py} checks these matrix identities
over the integers.

Normalize this reflected datum to the retained vectors
$(\varepsilon,-\varepsilon')$. The linear fibre map is $-I_2$, the
parameter shift is $1$, and the lattice map is $-I+E$,
$E=d(1,0,0,0)$. The differences between target and transported affine
vectors are $2d$ and $-4d$. They lie in the images of the respective
norm operators. Explicit integral correcting cocycles are
$\ell_1=(0,0,-2,0)^t$ and $\ell_2=(0,2,0,0)^t$; the norm equations
are $N_1\ell_1=-2d$ and $N_2\ell_2=4d$. Rational solutions of
$(I-A_j)r_j=(v_j-(-I+E)v'_j)/m_j+\ell_j$ can be chosen with first
coordinate zero: $r_1=(0,-2/3,-2/3,0)^t$ and
$r_2=(0,1,1,0)^t$. These give translations holomorphic across both
root discs. Their cusp cocycle is
$\ell_0=-A_0(\ell_1+A_1\ell_2)=(0,0,4,-2)^t$; its coordinates are
integral, so the corresponding shear and inversion preserve the
periodic triangular cusp fan. They therefore extend across its toric
charts. The differences of the local corrections form a cocycle in
the rank-two vector bundle whose invariant line and quotient are
$\mathcal O$ and $\mathcal O(-1)$ on the base sphere. Both have zero
first cohomology, so the local maps glue. Equivalently one splits the
Laurent series of the cocycle into its nonnegative and negative
parts in these two line-bundle frames. This checks the precise
gluing step of S6\_02 B.4.2, including the finite affine translations
and cusp extension. The resulting map has target $X_{1-\bar c}$.

Every other antiholomorphic comparison composed with its inverse is
holomorphic, so R5 gives exactly the congruence in the statement.
A self-map exists iff $\operatorname{Re}c\in1/2+\mathbb Z$, iff
$w=-\bar w$, iff $q_2=iw/216$ is real. Its square is a holomorphic
translation $d_0\in\mathbb C/\mathbb Z$, since the automorphism group
is the $\mathbb C^*$ of Theorem~1.4 of
\cite{DuGuoKusakabeXie2026v1}. The constructed reflection acts on
this additive translation coordinate by conjugation. Commutation
with its square gives $d_0=\bar d_0$ in $\mathbb C/\mathbb Z$, hence
$d_0$ has a real representative. Composing the reflection with
translation by $-d_0/2$ makes its square the identity. Necessity and
sufficiency for an involution follow.
\end{proof}

\subsection{C.3. The next local terms}
Put $G=g_2(i)=\Gamma(1/4)^8/(16\pi^2)$ and
$D_*=-3\log6/\pi$.

\begin{statement}[C.3.1; Proved: fourth order at the maximum]
On either root branch at $\tau=i$, let
$h=\mu-(1-\tau)/2$ and $\delta=\tau-i$. Then
\begin{align*}
 \tau&=i-\frac{3\pi i}{2}h^2-\frac{9\pi^2i}{8}h^4+O(h^6),\\
 h^2&=\frac{2i}{3\pi}\delta+\frac1{3\pi}\delta^2+O(\delta^3),\\
 \mu&=\frac{1-\tau}{2}
  \mathbin{\pm}\left(\frac{2i\delta}{3\pi}\right)^{1/2}
          \left(1-\frac{i}{4}\delta+O(\delta^2)\right).
\end{align*}
The square root is taken holomorphically on the chosen root chart.
Furthermore
\begin{align}\label{eq:cs-Li}
 \log N={}&\tfrac12\log2-\frac\pi2|h|^2
   +\frac{3\pi^2}{8}|h|^4
   -\frac{3\pi^2}{4}|h|^2\operatorname{Re}(h^2)
   -\frac{5G}{192}\operatorname{Re}(h^4)+O(|h|^6),\\
 D_{\rm part}={}&D_* -\frac2\pi|\delta|
  +\frac1\pi\bigl(|\delta|^2+|\delta|\operatorname{Im}\delta\bigr)
  +\frac{5G}{72\pi^3}\operatorname{Re}(\delta^2)
  +O(|\delta|^3).\nonumber
\end{align}
In particular the leading terms of K3 and Lead 9 are correct.
\end{statement}
\begin{proof}
Write $e(\tau)=\wp_\tau((1-\tau)/2)$ and
$E=\pi^2E_2/3$. The half-period Taylor series is
\[
 \wp_\tau((1-\tau)/2+h)
 =e+(3e^2-g_2/4)h^2+(3e^3-g_2e/4)h^4+O(h^6).
\]
The exact coefficients checked in K3 are
$e(i)=0$, $e'(i)=iG/(6\pi)$. This proves the claimed identity
$3e'(i)=(2i/\pi)G/4$. For clarity its closed form follows from
$e=\pi^2(\vartheta_2^4-\vartheta_4^4)/3$ and the heat equation:
if $A=\vartheta_2(i)^4=\vartheta_4(i)^4$, then
$G=4\pi^4A^2$ and $e'(i)=2\pi^3iA^2/3$.
The classical beta integral
$K(1/2)=\Gamma(1/4)^2/(4\sqrt\pi)$ and
$\vartheta_3(i)^2=2K(1/2)/\pi$ give
$G=16K(1/2)^4=\Gamma(1/4)^8/(16\pi^2)$.
The beta evaluation follows by changing variable in
$\int_0^\infty(1+x^4)^{-1/2}\,dx
=\tfrac14 B(1/4,1/4)$; it is the lemniscatic complete period.

Ramanujan's identities give $g_2'(i)=2iG$. Differentiating
$e(-1/\tau)=\tau^2 e(\tau)$ twice at $i$ gives
$e''(i)=3i e'(i)=-G/(2\pi)$. Inserting
$\delta=A_2h^2+A_4h^4+O(h^6)$ into the displayed half-period series
therefore gives
\[
 A_2=-3\pi i/2,\qquad
 A_4=\frac{g_2'A_2/4-e''A_2^2/2}{e'}=-9\pi^2i/8.
\]
The implicit-function theorem applies because $e'(i)\ne0$;
evenness in $h$ gives the stated remainder. Series inversion gives
the formulas for $h^2$ and $\mu$.

Let $H(h,\tau)=\log(\vartheta_3(h,\tau)/\eta(\tau))$. The exact
identity from the half-period translation is
$\log N=\operatorname{Re}H-\pi(\operatorname{Im}h)^2/
\operatorname{Im}\tau$. At $(0,i)$ its relevant Taylor data are
\[
 \begin{gathered}
 H_0=\tfrac12\log2,\quad H_\tau=0,\quad
 [h^2]H=-\pi/2,\quad [h^4]H=G/48,\\
 \partial_\tau[h^2]H=-\pi i/4-iG/(16\pi),\qquad
 H_{\tau\tau}=-G/(24\pi^2).
 \end{gathered}
\]
Indeed $\partial_h^2\log\vartheta_3=-\wp(h+(1-\tau)/2)-E$;
the heat equation at $h=0$ gives $H_\tau=ie/(4\pi)$, and
$E'(i)=\pi i/2-iG/(24\pi)$. These identities prove every entry.
After substitution $\tau=\tau(h)$, the pure $h^4$ coefficient of
$H$ is $-3\pi^2/8-5G/192$. Expanding
$\operatorname{Im}\tau=1-(3\pi/2)\operatorname{Re}(h^2)+O(|h|^4)$
in the Gaussian gives \eqref{eq:cs-Li}. Finally
\[
 |h|^2=\frac{2}{3\pi}|\delta|
          \left(1+\tfrac12\operatorname{Im}\delta+O(|\delta|^2)\right)
\]
and substitution into the quartic expression, followed by
multiplication by $6/\pi$, gives the formula for $D_{\rm part}$.
The script \path{Cspecial_expansion_checks.py} independently checks
these coefficients against theta evaluations; its output is retained.
\end{proof}

\begin{statement}[C.3.2; Proved: the equianharmonic point is neither an extremum]
Set $T_0=\sqrt3/2$, $\mu_0=(2-\rho)/3$, and
$p=\wp'_\rho(\mu_0)$. Then $p\ne0$, $p^2=-g_3(\rho)$, and, for
$\delta=\tau-\rho$ on the order-three root chart,
\begin{align*}
 \mu&=\frac{2-\tau}{3}
      +\frac{3p}{16\pi^2}\delta^2
      +\frac{3ip}{32\pi^2T_0}\delta^3+O(\delta^4),\\
 t&=\frac{ip^2}{\pi^3}\delta^3
        \left(1+\frac{3i}{2T_0}\delta+O(\delta^2)\right),\\
 \log N&=\tfrac14\log3+\frac7{192}\operatorname{Re}t
       -\frac{9\pi}{512T_0|p|^{2/3}}|t|^{4/3}
       +O(|t|^{5/3}).
\end{align*}
Thus
\[
 D_{\rm part}(t)=-\frac3{2\pi}\log48
       +\frac7{32\pi}\operatorname{Re}t
       -\frac{27}{256T_0|p|^{2/3}}|t|^{4/3}
       +O(|t|^{5/3}).
\]
The point $t=0$ is neither a local maximum nor a local minimum.
\end{statement}
\begin{proof}
We use K4 after checking its differentiation: put
$\ell=\partial_z\log\vartheta_1(\mu,\tau)$,
$a=\operatorname{Im}\mu/T$, $Z=\ell+2\pi ia$, and
$R=\mu'-a$. The theta heat equation and
$\wp=-\partial_z\ell-E$ give on the zero locus
\[
 \mu'=\frac{i\ell}{2\pi}-\frac{ig_2}{6\pi\wp'},\qquad
 R=\frac{iZ}{2\pi}-\frac{ig_2}{6\pi\wp'},\qquad
 \partial_\tau\log N=\frac{Z}{8\pi i}(Z+4\pi iR).
\]
These equations are equivalent to K4's factored formula and verify
its normalizations. At $\rho$ the unit endomorphism fixing the
nonzero order-three point forces the gradient of $U_\rho$ there to
vanish, so $Z=0$. Explicitly the unit $\rho-1$ fixes $\mu_0$ modulo
the lattice and has no fixed nonzero real covector. Hence
$a=\mu'=-1/3$ and $\ell=2\pi i/3$.

Along the zero locus, a further heat-equation differentiation gives
\[
 \ell'=\frac{i\wp'}{4\pi}+\frac{iEg_2}{6\pi\wp'}.
\]
The Ramanujan derivatives at $\rho$ are
\[
 g_2'=\frac{3ip^2}{\pi},\quad
 g_2''=\frac{5i}{T_0}g_2',\quad
 \frac{(\wp')'}{\wp'}=\frac{3i}{2T_0},\quad E=\frac\pi{T_0}.
\]
They give $\mu''=3p/(8\pi^2)$,
$\mu'''=9ip/(16\pi^2T_0)$ and hence the asserted $\mu$ expansion.
Expansion of $t=g_2^3/(g_2^3-27g_3^2)$ gives the displayed $t$
formula. Since $\partial_\tau a=R/(2iT)$, one has
\[
 Z=\frac{ip}{4\pi}\delta+O(|\delta|^2),\qquad
 R=\frac{3p}{8\pi^2}\delta+O(|\delta|^2).
\]
The derivative formula now gives the harmonic cubic
$\operatorname{Re}(7ip^2\delta^3/(192\pi^3))=(7/192)
\operatorname{Re}t+O(|\delta|^4)$.

For the next term, direct differentiation of the Gaussian gives
\[
 \partial_\tau\partial_{\bar\tau}\log N
       =-\frac\pi{2T}|R|^2
       =-\frac{9|p|^2}{128\pi^3T_0}|\delta|^2+O(|\delta|^3).
\]
After the harmonic cubic has been written as a function of $t$, the
order-three stabilizer permits only a radial term in degree four:
in the linearizing coordinate
$(\tau-\rho)/(\tau-\bar\rho)$ invariant monomials have exponents
whose difference is divisible by three. The coefficient of
$|\delta|^4$ is therefore $-9|p|^2/(512\pi^3T_0)$, since
$\partial_\delta\partial_{\bar\delta}|\delta|^4=4|\delta|^2$.
Converting $|\delta|^4$ to $|t|^{4/3}$ proves the expansion, including
the remainder. Its nonzero linear term in $\operatorname{Re}t$
has both signs, which proves the last assertion. In particular
superharmonicity at this root point does not make it a local maximum.
\end{proof}


\subsection{C.4. Division identities and the entire torsion locus}
We use throughout this subsection the rational model
\[
 E_t:\quad y^2=x^3+A(t)x+B(t),\qquad
 A=-\tfrac34t^3(t-1),\quad B=\tfrac14t^4(t-1)^2,
 \quad P=(0,t^2(t-1)/2).
\]
Its discriminant is $\Delta(t)=27t^8(t-1)^3$.
The change of complex Weierstrass model does not change $N$.

\begin{statement}[C.4.1; Known, with derivation: Siegel functions and local heights]
For $a=(a_1,a_2)\in\mathbb R^2\setminus\mathbb Z^2$ and
$v=a_1\tau+a_2$, define
\[
 g_a(\tau)=-e^{\pi ia_2(a_1-1)}q^{B_2(a_1)/2}
 (1-e^{2\pi iv})\prod_{r\ge1}
 (1-q^re^{2\pi iv})(1-q^re^{-2\pi iv}),
 \quad B_2(x)=x^2-x+\tfrac16.
\]
Then
\[
 g_a(\tau)=i e^{\pi ia_1v}
       \frac{\vartheta_1(v,\tau)}{\eta(\tau)},\qquad
 N(v,\tau)=|g_a(\tau)|.
\]
When $a\in n^{-1}\mathbb Z^2\setminus\mathbb Z^2$, $n\ge2$,
$g_a^{12n}$ is a modular unit for $\Gamma(n)$ with Fourier coefficients
in $\mathbb Q(\zeta_n)$ in the variable $q^{1/n}$. On every complex
elliptic curve the model-independent Néron local function for $(O)$ is
\[
 \lambda_\infty(Q)=-\log N(Q).
\]
Thus $D_{\rm part}=-(6/\pi)(\lambda_\infty(P)+\log(2\sqrt3))$.
\end{statement}
\begin{proof}
Dividing the two product expressions gives the first identity exactly,
including the factor $i$. Its absolute value is the asserted norm.
This agrees with the convention in \cite{Shin2015}, equation (1).
For completeness the transformation argument needed here can be
checked directly. For $T:\tau\mapsto\tau+1$ and
$S:\tau\mapsto-1/\tau$,
\[
 g_a(T\tau)=e^{\pi i/6}g_{aT}(\tau),\qquad
 g_a(S\tau)=-i g_{aS}(\tau);
\]
the usual row action is used on $a$. For $b=(b_1,b_2)\in\mathbb Z^2$,
\[
 \frac{g_{a+b}}{g_a}
 =\exp\bigl(\pi i(a_1b_2-a_2b_1+b_1b_2+b_1+b_2)\bigr).
\]
These follow from the two Jacobi transformations and the lattice
translations. They are the cases of Proposition 2.3(ii)--(iii) of
\cite{Shin2015} used here. The multipliers on $\Gamma(n)$ are killed
by $12n$. The product gives a Laurent series in $q^{1/n}$ with
coefficients in $\mathbb Q(\zeta_n)$, because
$6nB_2(r/n)\in n^{-1}\mathbb Z$ and its leading unit is an $n$th
root of unity. The transformed products prove meromorphy at every
cusp. No factor has a zero or pole in the upper half-plane. This
proves precisely the modular-unit assertion, also recorded in
Remark 2.2 of that source; its hypotheses $n\ge2$ and a nonintegral
rational pair hold here.

For the local-height normalization, $-U_\Lambda$ has a logarithmic
pole at $O$, mean zero, and Laplacian
$-2\pi\delta_O+2\pi\,d\mu_\Lambda$. This is the unique normalized
Green function for the degree-one divisor $(O)$: two such functions
differ by a harmonic function on the compact curve and their means
remove the constant. Equivalently its Weierstrass expansion is
\[
 -U_\Lambda(z)-\tfrac12\log|x(z)|
             \longrightarrow-\tfrac1{12}\log|\Delta_\Lambda|.
\]
The transformation of $x$ and $\Delta$ under a Weierstrass homothety
cancels in this expression. This is the model-independent local
height normalization, rather than a local function with an arbitrary
additive constant.
\end{proof}

\begin{statement}[C.4.2; Proved: the full division identity]
Let $\psi_m$ be the standard division polynomial of a smooth
Weierstrass model, normalized by $\psi_1=1$, $\psi_2=2y$.
For $m\ge2$ and $Q,mQ\ne O$,
\begin{equation}\label{eq:cs-division}
 N(Q)^{m^2}
       =N(mQ)\frac{|\Delta_E|^{(m^2-1)/12}}{|\psi_m(Q)|}.
\end{equation}
For the displayed family this specializes to
\[
 N(P)^4=3^{3/4}|t-1|^{-1/4}N(2P),\qquad
 N(P)^9=16|t|^{-2/3}N(3P).
\]
If $P_t$ has exact order $n$ and $m\ge2$ satisfies
$m\equiv1$ or $-1\pmod n$, then
\begin{equation}\label{eq:cs-torsionvalue}
 N(P_t)^{12}
       =\frac{|\Delta(t)|}{|\psi_m(P_t)|^{12/(m^2-1)}}.
\end{equation}
The positive value on the right is independent of $m$.
\end{statement}
\begin{proof}
Use the complex uniformization $x=\wp_\Lambda(z)$,
$y=\wp'_\Lambda(z)/2$. The rational function $\psi_m$ has divisor
$\sum_{0\ne Q\in E[m]}(Q)-(m^2-1)(O)$ and leading Laurent term
$(-1)^{m-1}m z^{1-m^2}$. One way to verify both normalizations without
a model convention is to form
$(-1)^{m-1}\sigma_\Lambda(mz)/\sigma_\Lambda(z)^{m^2}$.
Its lattice factors cancel, its zeros are exactly the nonzero
$m$-division points, and its expansion at zero is the displayed
one. The addition formula for $\sigma$ gives the standard
division-polynomial recurrence with initial values $1,2y$, so it is
the stated rational function.

The function
\[
 m^2U_\Lambda(z)-U_\Lambda(mz)
                    +\log|\psi_m(Q(z))|
\]
is harmonic after its logarithmic singularities cancel. Its
Laplacian has zero constant term as well, so it is constant.
The expansion in C.1.1 and the leading term of $\psi_m$ make that
constant $(m^2-1)\log|\Delta_E|/12$. This proves
\eqref{eq:cs-division}. Direct substitution of
\[
 \psi_2(P)=t^2(t-1),\qquad
 \psi_3(P)=-9t^6(t-1)^2/16
\]
gives the two special identities. For $m\equiv\pm1\pmod n$,
$mP=\pm P$ and $N$ is even. Cancelling $N(P)$ in
\eqref{eq:cs-division} gives \eqref{eq:cs-torsionvalue}.
Every right side is the same positive number $N(P)^{12}$, which
proves independence of $m$, not merely numerical agreement.
\end{proof}

\begin{statement}[C.4.3; Proved: algebraicity and explicit fields for every order]
Let $t_0\in\mathbb C\setminus\{0,1\}$ and suppose $P_{t_0}$ has
exact order $n\ge2$. Then $t_0$ and $N(P_{t_0})$ are algebraic.
Put $K_0=\mathbb Q(t_0,\bar t_0)$,
\[
 W=\Delta(t_0)\Delta(\bar t_0)>0,\qquad
 V=\psi_{n+1}(P_{t_0})\psi_{n+1}(P_{\bar t_0})>0.
\]
With positive real radicals,
\[
 N(P_{t_0})^{12n}=W^{n/2}V^{-6/(n+2)}
 \ \in\ K_0\bigl(W^{1/2},V^{1/(n+2)}\bigr).
\]
This supplies an explicit field, without a claim that it is minimal.
More generally, if $m\equiv\pm1\pmod n$, then
\[
 \lambda_\infty(P_{t_0})
 =\frac{\log|\psi_m(P_{t_0})|}{m^2-1}
       -\frac1{12}\log|\Delta(t_0)|.
\]
These formulas determine the height and $N^{12}$ on the entire
torsion locus.
\end{statement}
\begin{proof}
The division polynomial $\psi_n(P_t)$ is a nonzero rational
polynomial in $t$. To check nonvanishing for every $n$, a hypothetical
torsion identity would express the elliptic logarithm as
$a\tau+b$ with constant rational $a,b$. Its bounded cusp limit
forces $a=0$, but that limit has nonzero imaginary part, a
contradiction. Thus $t_0$ is a zero of a nonzero polynomial over
$\mathbb Q$. Formula \eqref{eq:cs-torsionvalue} shows algebraicity
of $N$ using only positive rational powers of nonzero algebraic
absolute values. Taking $m=n+1$, raising to the $n$th power, and
using $|\Delta|=W^{1/2}$, $|\psi_{n+1}|=V^{1/2}$ gives the exact
field statement. The local-height formula is its logarithm.

Here is a polynomial relation involving only $t$ and $r=N^{12}$,
if one prefers that form to positive radicals. Remove the factors
at $0,1$ from $\psi_n(P_t)$ and divide successively by the
corresponding polynomials for proper divisors of $n$; let $F_n$ be
the resulting primitive polynomial. At a point of exact order
$d\mid n$, the formal group gives
$z([n/d]Q)=(n/d)z(Q)+O(z(Q)^2)$, so the vanishing multiplicities for
$\psi_d$ and $\psi_n$ agree. Hence this cancellation leaves exactly
the exact-order-$n$ divisor, with its multiplicities. Set
$m=n+1$ and $e=m^2-1$. Clearing rational denominators in
\[
 \operatorname{Res}_S\left(
 F_n(S),\
 r^{2e}\{\psi_m(P_T)\psi_m(P_S)\}^{12}
              -\{\Delta(T)\Delta(S)\}^{e}\right)
\]
gives a polynomial in $\mathbb Q[T,r]$ that vanishes at
$(t_0,N(P_{t_0})^{12})$, since $S=\bar t_0$ is a common root.
It is nonzero: at an exact-order-$n$ root no $\psi_m$ vanishes,
so the leading coefficient in $r$ of the resultant is nonzero.
This is an explicit recurrence-and-resultant prescription for
every $n$, including nonreal torsion parameters.
\end{proof}

\begin{statement}[C.4.4; Proved: the five values and all orders through eight]
The exact special values are
\[
 \begin{array}{c|c|c}
 t&\operatorname{ord}(P)&N(P)^{12}\\ \hline
 0&3&27\\
 1&2&64\\
 32/5&4&10\\
 64(47-27\sqrt3)/11&6&17+10\sqrt3\\
 64(47+27\sqrt3)/11&6&10\sqrt3-17 .
 \end{array}
\]
At $t=0,1$ the orders refer to the smooth models after the specified
base changes. On the smooth base, $F_2=F_3=1$ and
\[
 \begin{aligned}
 F_4&=5t-32,\\
 F_5&=89t^2+4736t-4096,\\
 F_6&=11t^2-6016t+8192,\\
 F_7&=881t^4+4745344t^3-12079104t^2\\
    &\hspace{13mm}+24641536t-16777216,\\
 F_8&=49t^4-2614336t^3-7864320t^2\\
    &\hspace{13mm}+19398656t-8388608.
 \end{aligned}
\]
All these polynomials are squarefree. At each root of exact order
$n=4,5,6,7,8$, the positive number $r=N^{12}$ is respectively
determined by
\[
 \begin{array}{c|l}
 n&\text{equation for }r\\ \hline
 4&r=64/|t|\\
 5&r^5=432^5/\bigl((1-t)(5t-32)^4\bigr)\\
 6&r^2=1728^2/F_5(t)\\
 7&r^{35}=2^{192}3^{81}/
       \bigl(|t|^8|t-1|^3|F_6(t)|^{12}\bigr)\\
 8&r^4=1728^4/|F_7(t)|.
 \end{array}
\]
The order-five roots are real and below one; the order-six roots
are real and above one. The two complex order-seven roots are
included in the formula.
\end{statement}
\begin{proof}
The division-polynomial recurrence, evaluated at $x=0$,
$y=t^2(t-1)/2$, gives
\[
 \begin{aligned}
 \psi_4(P)&=-t^{10}(t-1)^4F_4/32,\\
 \psi_5(P)&=t^{16}(t-1)^6F_5/4096,\\
 \psi_6(P)&=9t^{24}(t-1)^9F_6/65536,\\
 \psi_7(P)&=-t^{32}(t-1)^{12}F_7/16777216,\\
 \psi_8(P)&=-t^{42}(t-1)^{16}F_4F_8/268435456.
 \end{aligned}
\]
The recursion and all coefficients are checked over $\mathbb Q[t]$
by \path{Cspecial_division.py}. The same program computes
$\gcd(F_n,F_n')=1$ and excludes common factors for proper orders.
Thus no point or multiplicity is omitted. Applying
\eqref{eq:cs-torsionvalue} with $m=n-1$ gives every displayed
formula. The signs for $F_5,F_6$ follow directly from their quadratic
formula; in particular $F_5$ is positive at the two roots of $F_6$.
At order six, substituting
$t=64(47\mp27\sqrt3)/11$ gives the exact identities
\[
 F_5(t)(17\pm10\sqrt3)^2=1728^2.
\]
The positive roots are $17+10\sqrt3$ and $10\sqrt3-17$; their
product is $11$. The substitution is checked in
\path{Cspecial_relations.py}. The order-four value follows from
$64/(32/5)=10$, and the two central values are C.1.2.

For clarity, the recurrence used by the exact program is
\[
 \begin{aligned}
 \psi_{2m+1}&=\psi_{m+2}\psi_m^3-\psi_{m-1}\psi_{m+1}^3,\\
 \psi_{2m}&=\frac{\psi_m}{\psi_2}
   (\psi_{m+2}\psi_{m-1}^2-\psi_{m-2}\psi_{m+1}^2).
 \end{aligned}
\]
Its initial values at $P$ are $\psi_0=0,\psi_1=1,\psi_2=2y$,
$\psi_3=-A^2$, $\psi_4=4y(-8B^2-A^3)$. No approximate algebra
enters this calculation. The same code records $\psi_9$, and the
separate file \texttt{Cspecial\_psi16.out} records the factored
$\psi_{16}$ for the duplication ladder. A rational polynomial
relation for the nonreal order-seven values is explicitly computed
by \path{Cspecial_relations.py}, by the resultant prescription
of C.4.3.
\end{proof}

\begin{statement}[C.4.5; Proved: complete numerical verification through order eight]
All $13$ smooth exact-order roots for $2\le n\le8$ were evaluated
at $110$ decimal digits independently of the algebraic norm formula.
The discrepancy in $N^{12}$ is below $10^{-90}$ at every root;
the period coordinates also verify $nP=O$ to that tolerance.
There are $0,0,1,2,2,4,4$ roots for $n=2,\ldots,8$.
\end{statement}
\begin{proof}
The script \path{Cspecial_torsion_numerics.py} computes the
roots of $F_n$ to $115$ digits and constructs periods of
$y^2=x^3+Ax+B$ using its three roots $e_1,e_2,e_3$. If
$k=(e_2-e_3)/(e_1-e_3)$ and $K$ denotes the complete elliptic
integral with parameter $k$, it uses
\[
 \tau=iK(1-k)/K(k),\qquad
 v=\frac{F\bigl(\arcsin\sqrt{(e_1-e_3)/(-e_3)},k\bigr)}{2K(k)}.
\]
The identity
$\wp(u)=e_3+(e_1-e_3)/\operatorname{sn}^2
(\sqrt{e_1-e_3}\,u,k)$ verifies that this is a zero of $\wp$.
Permutations of the roots are tried until $\operatorname{Im}\tau>0$;
modular reduction and lattice translations improve numerical
conditioning and leave $N$ unchanged. The program then evaluates
the theta quotient directly. It checks both the norm formula and
that the coefficients of $nv$ in the basis $(\tau,1)$ are integral.
The full output gives at least $70$ digits for every $t,\tau,v,N$
and $N^{12}$, including both nonreal conjugates. The largest measured
norm discrepancy is below $2\cdot10^{-108}$. These are numerical
checks of the exact theorem, not numerical proofs of it.
\end{proof}

\subsection{C.5. CM and torsion}
\begin{statement}[C.5.1; Known, with derivation: no smooth CM-torsion parameter]
There is no $t\in\mathbb C\setminus\{0,1\}$ for which $E_t$ has
complex multiplication and $P_t$ is torsion. On the smooth models
at $t=0,1$ after base change of orders $3,4$, respectively, the
limiting points have orders $3,2$. Thus these are the only CM
torsion points on the completed family.
\end{statement}
\begin{proof}
The precise torsion theorem used is \cite{BourdonClarkStankewicz2017},
Theorem 4.15(a), on page 23 of arXiv:1411.2742v3:
for an elliptic curve $E/F$ with CM by any imaginary quadratic order,
\[
 F=\mathbb Q(j(E))\quad\Longrightarrow\quad
 E(F)_{\rm tors}\in
 \{0,\mathbb Z/2,\mathbb Z/3,\mathbb Z/4,
       \mathbb Z/6,(\mathbb Z/2)^2\}
\]
up to group isomorphism. The assertion allows nonmaximal orders.
The exceptional restrictions at discriminants $-3,-4$ belong only
to part (b), over the larger field containing the CM field; they
are not hypotheses of part (a).

For the present twist $F=\mathbb Q(t)=\mathbb Q(j)$ because
$j=1728t$. All Weierstrass coefficients and both coordinates of
$P$ belong to $F$, and $\Delta\ne0$ on the smooth base. Therefore
the theorem applies to the actual twist and the actual section,
not just to an isomorphic complex curve. The point $P$ is never
$O$, and $\psi_2(P),\psi_3(P)$ are nonzero there. The remaining
possible orders are $4$ and $6$, whose parameters were computed
exactly above. For order four, $j=55296/5$ is not an algebraic
integer. For order six the two $j$-values have monic minimal
polynomial
\[
 X^2-\frac{10395648}{11}X+\frac{24461180928}{11};
\]
it is irreducible over $\mathbb Q$, and its coefficients are not
integers. Neither root is an algebraic integer. By the integrality
theorem for singular moduli (for example, Theorem 20.12 and
Corollary 20.13 of \cite{Sutherland2023CM}, proved there for all
imaginary quadratic orders), this excludes complex multiplication
in every remaining case. The required hypothesis is simply that
the endomorphism ring is such an order, which is exactly the CM
assumption.

At zero put $t=s^3$, $x=s^4X$, $y=s^6Y$. The smooth equation is
\[
 Y^2=X^3-\tfrac34s(s^3-1)X+\tfrac14(s^3-1)^2,\qquad
 P_s=(0,(s^3-1)/2).
\]
Its central curve is $Y^2=X^3+1/4$ with endomorphism ring
$\mathbb Z[\rho]$, and the limiting point $(0,-1/2)$ has order
three: its horizontal tangent has triple contact, so $2P=-P$.
At one put $t=1+s^4$, $x=s^2X$, $y=s^3Y$. Then
\[
 Y^2=X^3-\tfrac34(1+s^4)^3X+
       \tfrac14(1+s^4)^4s^2,\qquad
 P_s=(0,(1+s^4)^2s/2).
\]
The central curve is $Y^2=X^3-3X/4$, with endomorphism ring
$\mathbb Z[i]$, and its limiting point $(0,0)$ has exact order two.
The displayed equations have nonzero central discriminants; hence
these are genuine smooth models, and their norm limits are those
already evaluated in C.1.2. This proves Lead 11, including its
twist and boundary qualifications.
\end{proof}
}{}
\IfFileExists{sections/C_heights.tex}{\section{C.6. Global heights and the eleven rational CM fibres}

We use the equation
\[
 E_t:Y^2=X^3-\frac34t^3(t-1)X+\frac14t^4(t-1)^2,
 \qquad P_t=(0,t^2(t-1)/2).
\]
Its discriminant is $27t^8(t-1)^3$. Our canonical height is
$\widehat h_O(P)=\frac12\lim_{n\to\infty}n^{-2}h(x(nP))$.
In particular the bilinear Mordell--Weil self-pairing used in the
supplied surface calculation is $2\widehat h_O$.
For a finite place $v$ above the rational prime $p$, we normalize
$|p|_v=p^{-1}$ and set $n_v=[K_v:\Q_p]$. At an archimedean
place $n_v$ is one or two, according as the place is real or complex.

\begin{statement}[C.6.1; Proved: the global formula]
For every number field $K$, every $t\in K\setminus\{0,1\}$,
and the section just specified,
\begin{equation}\label{eq:ch-global}
 \sum_{\sigma:K\hookrightarrow\C}\log N(P_{\sigma(t)})
 =\sum_{v\nmid\infty}n_v\lambda_v(P_t)
          -[K:\Q]\widehat h_O(P_t).
\end{equation}
Here the local heights are invariant under change of Weierstrass
model and satisfy, when $nP\ne O$,
\begin{equation}\label{eq:ch-local-mult}
 \lambda_v(nP)=n^2\lambda_v(P)-\log|\psi_n(P)|_v
                   +\frac{n^2-1}{12}\log|\Delta|_v.
\end{equation}
At a complex embedding $\lambda_v(P)=-\log N(P)$; complex
embeddings are counted separately in \eqref{eq:ch-global}.
For rational $t$, put $A_t=\exp(\sum_p\lambda_p(P_t))$.
Then $A_t$ is a positive real algebraic number and
\[
 N(P_t)=A_t e^{-\widehat h_O(P_t)},\qquad
 D_{\rm part}(t)=\frac6\pi
       \left(\log\frac{A_t}{2\sqrt3}-\widehat h_O(P_t)\right).
\]
Thus $N(P_t)$ is algebraic if and only if
$e^{\widehat h_O(P_t)}$ is algebraic.
\end{statement}
\begin{proof}
The sigma expression for the archimedean local height is
\[
 \lambda_v(P)=-\log|\sigma(u)|
       +\frac12\Re(u\eta(u))-\frac1{12}\log|\Delta|.
\]
The sigma--theta relation and the Legendre relation identify this
with $-\log N$ exactly, including its constant; this is also the
calculation in Statement C.7.1. The division identity
$\sigma(nu)=(-1)^{n-1}\psi_n(P)\sigma(u)^{n^2}$ gives
\eqref{eq:ch-local-mult}. At finite places the same identity is
the defining multiplication law of the canonical N\'eron function.
Its normalization is uniquely determined by continuity away from
$O$, its logarithmic pole at $O$, and the multiplication law.
These are the three properties printed in
\citet[Section 2.1]{demarco2017}; their duplication law includes
the discriminant term $\frac14\log|\Delta|_v$.
Under a coordinate scaling the changes in the sigma or formal
parameter term and in $\frac1{12}\log|\Delta|_v$ cancel.

Sum the duplication law over places. The product formula cancels
the division polynomial and discriminant terms, so that the sum
$H(P)=[K:\Q]^{-1}\sum_vn_v\lambda_v(P)$ satisfies
$H(2P)=4H(P)$. The pole normalization implies
$H(P)=\frac12h(x(P))+O_E(1)$. Iterating multiplication by two
therefore identifies $H$ with $\widehat h_O$. Moving the
archimedean terms to the other side gives \eqref{eq:ch-global}.
For a finite place, a multiple of $P$ lies in the identity component
of a minimal model. There the local height is a rational multiple
of $\log p$; the multiplication law shows the same for $P$.
Only finitely many places contribute. This proves the assertion
about $A_t$ and both of its consequences.
\end{proof}

\begin{statement}[C.6.2; Proved: exact finite local-height calculation]
Let $E/\Q$ be a global minimal Weierstrass model. If $mP$ is
in the identity component at $p$, and $mP\ne O$, then
\begin{equation}\label{eq:ch-local-exact}
 \frac{\lambda_p(P)}{\log p}
 =\frac{v_p(\Delta)}{12}
   +\frac{\frac12\max(0,-v_p(x(mP)))-v_p(\psi_m(P))}{m^2}.
\end{equation}
For $m=1$ we read $\psi_1=1$. This formula, with exact rational
coordinates and valuations, gives all the finite contributions
in the tables below.
\end{statement}
\begin{proof}
On the identity component of a minimal model the canonical
N\'eron function is
$\frac12\log\max(1,|x|_p)-\frac1{12}\log|\Delta|_p$.
This follows either from the formal parameter at $O$ followed
by the group law on the smooth locus, or from the uniqueness
properties used above. Substitute it for $\lambda_p(mP)$ in
\eqref{eq:ch-local-mult} and divide by $m^2\log p$.
The contribution $v_p(\Delta)/(12m^2)$ combines with the
discriminant term to give $v_p(\Delta)/12$.
Membership in the identity component is checked by reduction
to the nonsingular locus of the minimal cubic. Thus the test is
exact also at $2$ and $3$; it does not use a formula restricted
to residue characteristic at least five.
\end{proof}

\begin{statement}[C.6.3; Proved: the CM models and their height identities]
The eleven rational CM fibres have the global minimal models,
section coordinates, component classes, and finite local heights
in the following tables. On each model $a_1=a_2=0$, and we
write $[a_3,a_4,a_6]$ for
$y^2+a_3y=x^3+a_4x+a_6$. The displayed point is $(0,y_P)$.
The sign of $P$, which is immaterial to all these heights,
is chosen by the recorded change of variables.
The constants $A_t$ in increasing row order are
\[
 \sqrt6,\quad 2^{2/3}3^{1/4},\quad 2^{3/4},\quad
 2,\quad 2,\quad \sqrt3,\quad 1,\quad 3^{5/12},\quad
 1,\quad1,\quad1.
\]
\end{statement}
\begin{proof}
We substitute the eleven exact values $t=j/1728$, and perform
the integral minimal-model transformations and Tate reduction
algorithm. The transformations and resulting invariants are
printed in \path{code/C_cm_models.out}; the exact valuation
computations are in \path{code/C_cm_local_heights.out}.
For each bad prime the program computes $mP$ by the rational
group law, tests the two partial derivatives of the minimal
cubic, and evaluates $\psi_m$ by its exact recursion. Formula
\eqref{eq:ch-local-exact} then proves each entry. Outside the
listed primes the discriminant is a unit and $P$ is integral,
so the local height is zero. Multiplying $p^{\lambda_p/\log p}$
gives the listed constants $A_t$.

We specify the component labels without relying on an
unrecorded ordering of a resolution graph. The label $0$
means the identity component, and $c_P$ is the class represented
by the explicitly displayed section $P$. At an odd prime of
type $I_0^*$ this nonidentity class is the outer component
indexed by the root zero of the residual Tate cubic; the
scaling $x=pX$ identifies that root. When there are two
rational components, $c_P$ is their unique nonzero class.
The only remaining ambiguity is row 5 at $2$, where the
component group is $(\Z/2)^2$. Here let $T=(-14,0)$;
we label it with the basis $(c_P,c_T)$. Indeed
$P+T=(18,64)$ after the corresponding choice of sign,
and $P,T,P+T$ all reduce to the singular point, whereas
$2P=(25/4,99/8)$ reduces to the identity component.
Thus the three nonzero classes are explicitly separated.

The independent analytic check uses the exact quadratic CM
period ratio, solves $\wp(\mu)=0$ by its convergent Fourier
series, and evaluates the theta product at 270-digit working
precision. The value of $j$ agrees with its prescribed integer
to absolute error less than $10^{-245}$. For every row the
direct value of $\log N$ agrees with
$\log A_t-\widehat h_O(P_t)$ to more than 245 decimal places.
These numerical comparisons test the exact normalization;
the algebraic local-height entries were already proved by
integer and rational arithmetic. The full precision values
are in \path{records/C_cm_theta_pslq.json}.
\end{proof}

\begingroup\footnotesize\setlength{\tabcolsep}{4pt}
\begin{longtable}{r r r r r r}
\toprule Row & $j$ & $a_3$ & $a_4$ & $a_6$ & $y_P$\\\midrule\endhead
1 & 8000 & 0 & -5880 & 153664 & 392\\
2 & 54000 & 0 & -7260 & 234256 & 484\\
3 & 287496 & 0 & -2156 & 38416 & 196\\
4 & 16581375 & 0 & -7732620 & 8275904784 & 90972\\
5 & -3375 & 0 & -140 & 784 & -28\\
6 & -32768 & 1 & -3234 & 72630 & -270\\
7 & -884736 & 1 & -38 & 90 & -10\\
8 & -12288000 & 1 & -1920270 & 1024288020 & -32005\\
9 & -884736000 & 1 & -379260 & 89898842 & -9482\\
10 & -147197952000 & 1 & -347045930 & 2488447882842 & -1577482\\
11 & -262537412640768000 & 1 & -75065352465446420 & 7916029533155605844266842 & -2813543945482\\
\bottomrule\end{longtable}
\begin{longtable}{r r l c c l}
\toprule Row & $p$ & Type & Component & Order & $\lambda_p/\log p$\\\midrule\endhead
1 & 2 & $III^*$ & $c_P$ & 2 & $\frac{1}{2}$\\
1 & 3 & $I_0^*$ & $0$ & 1 & $\frac{1}{2}$\\
1 & 7 & $I_0^*$ & $c_P$ & 2 & $0$\\
2 & 2 & $I_{4}^*$ & $c_P$ & 2 & $\frac{2}{3}$\\
2 & 3 & $III$ & $0$ & 1 & $\frac{1}{4}$\\
2 & 11 & $I_0^*$ & $c_P$ & 2 & $0$\\
3 & 2 & $I_{5}^*$ & $c_P$ & 2 & $\frac{3}{4}$\\
3 & 7 & $I_0^*$ & $c_P$ & 2 & $0$\\
4 & 2 & $I_{8}^*$ & $c_P$ & 2 & $1$\\
4 & 3 & $I_0^*$ & $c_P$ & 2 & $0$\\
4 & 7 & $III$ & $c_P$ & 2 & $0$\\
4 & 19 & $I_0^*$ & $c_P$ & 2 & $0$\\
5 & 2 & $I_{8}^*$ & $c_P$ & 2 & $1$\\
5 & 7 & $III$ & $c_P$ & 2 & $0$\\
6 & 3 & $I_0^*$ & $0$ & 1 & $\frac{1}{2}$\\
6 & 7 & $I_0^*$ & $c_P$ & 2 & $0$\\
6 & 11 & $III$ & $c_P$ & 2 & $0$\\
7 & 19 & $III$ & $c_P$ & 2 & $0$\\
8 & 3 & $IV$ & $0$ & 1 & $\frac{5}{12}$\\
8 & 11 & $I_0^*$ & $c_P$ & 2 & $0$\\
8 & 23 & $I_0^*$ & $c_P$ & 2 & $0$\\
9 & 3 & $I_0^*$ & $c_P$ & 2 & $0$\\
9 & 7 & $I_0^*$ & $c_P$ & 2 & $0$\\
9 & 43 & $III$ & $c_P$ & 2 & $0$\\
10 & 7 & $I_0^*$ & $c_P$ & 2 & $0$\\
10 & 31 & $I_0^*$ & $c_P$ & 2 & $0$\\
10 & 67 & $III$ & $c_P$ & 2 & $0$\\
11 & 7 & $I_0^*$ & $c_P$ & 2 & $0$\\
11 & 11 & $I_0^*$ & $c_P$ & 2 & $0$\\
11 & 19 & $I_0^*$ & $c_P$ & 2 & $0$\\
11 & 127 & $I_0^*$ & $c_P$ & 2 & $0$\\
11 & 163 & $III$ & $c_P$ & 2 & $0$\\
\bottomrule\end{longtable}
\endgroup
\noindent The next table prints 52 significant digits. The last column is both the directly computed $\log N$ and the independently computed $\log A_t-\widehat h_O(P_t)$; the agreement is more precise than the digits printed.
\begingroup\small
\begin{longtable}{r l}
\toprule Row & Height and archimedean logarithm\\\midrule\endhead
1 & $\widehat h_O=0.68444758838339190721754413638128047178526153154571$\\
 & $\log N=0.2114321462306355931886945428090706645762338145457924$\\[3pt]
2 & $\widehat h_O=0.6417998032172362080432024005106271224631326545214062$\\
 & $\log N=0.09495138932308808775043032302545534908240674117445132$\\[3pt]
3 & $\widehat h_O=0.5335302122611885050129127004864161587766141055417204$\\
 & $\log N=-0.01366982684122952294998860939278373271998900477152897$\\[3pt]
4 & $\widehat h_O=0.9940668685687274300038154462317330832788241518412543$\\
 & $\log N=-0.300919688008782120586583324773556515203324017480999$\\[3pt]
5 & $\widehat h_O=0.4673656771814712800075202012628467836356203437033082$\\
 & $\log N=0.2257815033784740294097119201953297844398797906569471$\\[3pt]
6 & $\widehat h_O=0.4286091723813209970260194687110167110834518087280487$\\
 & $\log N=0.120696971952733848671603149750246141240293470183326$\\[3pt]
7 & $\widehat h_O=0.09064583701789963707894077069188164993051948122321821$\\
 & $\log N=-0.09064583701789963707894077069188164993051948122321821$\\[3pt]
8 & $\widehat h_O=0.7366208256267509672496883390906560396314066510934334$\\
 & $\log N=-0.2788657053483719291683361570396036626949522520006211$\\[3pt]
9 & $\widehat h_O=0.6018296198896332298783588845164397135769872973598109$\\
 & $\log N=-0.6018296198896332298783588845164397135769872973598109$\\[3pt]
10 & $\widehat h_O=1.002649286144114870320961646648539589847474352238484$\\
 & $\log N=-1.002649286144114870320961646648539589847474352238484$\\[3pt]
11 & $\widehat h_O=2.165489246796061077105970244492001901750597685962121$\\
 & $\log N=-2.165489246796061077105970244492001901750597685962121$\\[3pt]
\bottomrule\end{longtable}\endgroup


\begin{statement}[C.6.4; Proved: the maximum as a height inequality]
For $t\in K\setminus\{0,1\}$,
\begin{equation}\label{eq:ch-bound}
 \widehat h_O(P_t)\geq
 \frac1{[K:\Q]}\sum_{v\nmid\infty}n_v\lambda_v(P_t)
                    -\frac12\log2.
\end{equation}
The inequality is strict on this smooth base. Over $\Q$ it
is $\widehat h_O(P_t)\geq\log A_t-\frac12\log2$.
The generic height is $\widehat h_O(P)=1/12$, and the
variation divisor is $D_P=(\infty)/12$. Consequently
\[
 \widehat h_O(P_t)=\frac1{12}h(t)+O(1)
 \qquad(t\in\overline\Q\setminus\{0,1\}),
\]
where the constant is uniform over algebraic $t$ and is for
the fixed family and section.
\end{statement}
\begin{proof}
Apply the analytic maximum theorem C.8.5 at every embedding
and substitute into \eqref{eq:ch-global}. Equality would
require $j(\sigma(t))=1728$ for every embedding, hence $t=1$,
which has been excluded.

For the generic local heights at $0$ and $1$, the minimal
discriminant orders are $8$ and $3$, and the section meets
the nonidentity component of types $IV^*$ and $III$.
Their local corrections are $2/3$ and $1/4$, respectively,
so their normalized local heights are zero. At infinity,
put $s=1/t$, $X=s^{-2}x$, $Y=s^{-3}y$. The new equation is
\[
 y^2=x^3-\tfrac34(1-s)x+\tfrac14(1-s)^2,
 \qquad P=(0,(1-s)/2),\qquad\Delta=27s(1-s)^3.
\]
At $s=0$ the node is $(1/2,0)$, whereas $P=(0,1/2)$ is
nonsingular. Its height contribution is $1/12$.
At all other function-field places the point is integral
and the discriminant is a unit. This proves
$D_P=(\infty)/12$ and the generic height.

Theorem 1.1 of \citet{demarco2017} applies with $K=\Q$,
$B=\mathbb P^1$, the displayed elliptic surface and rational
section, since the generic height is positive. It equips
this divisor with continuous adelic semipositive metrics
whose associated height is exactly $\widehat h_O(P_t)$.
Comparing with the usual height for $(\infty)/12$ gives the
bounded difference. This is the specialization here of the
variation theorem of Tate, with the local refinements of
Silverman reproduced as Theorems 2.2 and 2.3 in that paper:
outside finitely many places the local correction is zero,
and at each remaining place it extends continuously after
subtracting its prescribed logarithmic singularity.
Our inequality is an explicit archimedean estimate with
constant $\frac12\log2$; the general bounded-variation
statement alone does not supply that constant.

We spell out the comparison with the opened primary author sources.
Theorem 4.1 of \citet[pp.~183--184]{callSilverman1993} concerns a
smooth projective base curve $T$, a degree-one base height $h_T$,
a section, and a fibrewise morphism whose pullback multiplies the
chosen divisor class by a real number $\alpha>1$. It gives
\[
 \lim_{h_T(t)\to\infty}
 \frac{\widehat h_t(P_t)}{h_T(t)}
       =\widehat h_{\Q(T)}(P).
\]
Here take $T=\mathbb P^1$, the class of $(O)$, and the generic
morphism $[2]$, with $[2]^*(O)\sim4(O)$ and $\alpha=4$.
The smooth fibres are precisely the complement of $0,1,\infty$;
the section extends on the regular projective elliptic surface.
All the hypotheses hold, and the limit is $1/12$.
The unnumbered remark immediately following that theorem states
Tate's stronger elliptic-family formula
$\widehat h_t(P_t)=h_{T,D_P}(t)+O_P(1)$, with
$\deg D_P=\widehat h_{\Q(T)}(P)$. This is an author-verified
reproduction of Tate's result, not a claim to have read the
inaccessible original scan \citep{tate1983variation}.
The same distinction applies to \citet{silverman1983specialization}:
the limit theorem is reproduced by the author on slide~22 of
\citet{silverman2006heights}, while slide~21 states the uniform
canonical-versus-Weil-height bound in terms of the base height.
For our explicit bounded remainder and local normalization, the
fully opened Theorems~1.1, 2.2 and 2.3 of
\citet{demarco2017}, applied above, remain the cited proof.
\end{proof}

\begin{statement}[C.6.5; Proved: rigorous high-precision height enclosures]
For every one of the eleven rational CM parameters in C.6.3,
the files \path{records/C_cm_height_arb.json} and
\path{code/C_cm_height_arb.out} give certified enclosures of
$\log N(P)$ and $\widehat h_O(P)$ of radius less than
$10^{-236}$. These enclosures use no floating-point height
value as a proof step.
\end{statement}
\begin{proof}
The CM discriminants, in the order of C.6.3, are
\[
 -8,-12,-16,-28,-7,-11,-19,-27,-43,-67,-163.
\]
Enumeration of primitive reduced forms
$(a,b,c)$, $|b|\le a\le c$, $b^2-4ac=D$, with the usual
nonnegative boundary convention, gives exactly one form of
each discriminant, always with $a=1$.
The range is finite and proved by $a\le\sqrt{|D|/3}$.
Thus the CM theorem used in C.5.1 makes the corresponding
singular modulus an integer. The Fourier enclosures below
place it within $1/4$ of the specified $j$, proving its
identity, rather than only numerical closeness. Exact class
polynomials were also computed independently in
\path{code/C_cm_class_polynomials.gp} and its output.

Write $\tau=iT$ for the first four entries and
$\tau=1/2+iT$ for the other seven, with
\[
 (T_1,T_2,T_3,T_4)=(\sqrt2,\sqrt3,2,\sqrt7),\qquad
 T_k=\sqrt{|D_k|}/2\quad(k\ge5).
\]
Put $Q=e^{-2\pi T}$, $q=Q$ or $-Q$ accordingly,
$M=350$, and $Y=2/5$. On $0\le y\le Y$,
\begin{equation}\label{eq:ch-wp-bracket}
 \frac{\wp(1/2-iy)}{\pi^2}
 =\frac1{\cosh^2\pi y}-\frac13+
 8\sum_{n\ge1}\frac{nq^n}{1-q^n}
                 \{1-(-1)^n\cosh(2\pi ny)\}.
\end{equation}
For $a=Qe^{2\pi Y}$, the absolute tail after $M$ in
$\wp$ is at most
\[
 \frac{16\pi^2}{1-Q}
 \frac{a^{M+1}\{M+1-Ma\}}{(1-a)^2}<10^{-875}.
\]
The proof is $1+\cosh(2\pi ny)\le2e^{2\pi nY}$ and
$\sum_{n>M}na^n=a^{M+1}(M+1-Ma)/(1-a)^2$.
Each previously computed decimal supplies only a rational
candidate centre $y_0$. At the two rational endpoints
$y_0\pm10^{-240}$ the truncated expression plus this
infinite tail is respectively strictly positive and strictly
negative. The code records all twenty-two strict signs.
Continuity produces a zero within each bracket. Every zero
of $\wp$ is one of the pair $\pm P$, and $N$ agrees on
that pair, so no unproved choice of root is needed.

The product formula on this bracket has positive real factors:
\begin{align*}
 \log N={}&\log(2\cosh\pi y)-\frac{\pi T}{6}
                      -\frac{\pi y^2}{T}\\
 &+\sum_{n\ge1}\{\log(1+q^ne^{2\pi y})
                         +\log(1+q^ne^{-2\pi y})\}.
\end{align*}
Since $|\log(1+x)|\le |x|/(1-|x|)$ for real $|x|<1$,
the omitted tail is bounded by
\[
 \frac{2e^{2\pi Y}Q^{M+1}}
 {(1-Q)(1-e^{2\pi Y}Q^{M+1})}<10^{-1260}.
\]
All calculations use Arb enclosures at $1152$ bits; the
radius comparisons and endpoint signs are exact interval
comparisons. The resulting radius is less than $10^{-236}$.
Adding the exact finite local logarithms of C.6.3 and
subtracting this enclosure gives the claimed height bound.

For completeness, the singular-modulus certificate uses
$E_4=1+240\sum n^3q^n/(1-q^n)$ and
$\Delta_0=q\prod(1-q^n)^{24}$. Their respective error
bounds, after truncation at $M$, are
\[
 \frac{240(M+1)^3Q^{M+1}(1+4Q+Q^2)}{(1-Q)^5},
 \qquad
 \frac{24Q^{M+1}}{(1-Q)(1-Q^{M+1})}
\]
for $E_4$ and for the logarithm of the product. The first
uses $(M+1+k)^3\le(M+1)^3(k+1)^3$ and the differentiated
geometric series; the second uses the same logarithm bound.
Thus $E_4^3/\Delta_0$ is rigorously enclosed. The complete
certificate is \path{code/C_cm_height_arb.py}, with its
output and rational bracket centres retained.
\end{proof}
}{}
\IfFileExists{sections/C_rank.tex}{\subsection{C.6. Supplement: ranks, saturation, and central derivatives}

We retain the order of the eleven rational CM values in the prompt, and
write $E_i$ for the global minimal models recorded in the preceding
tables. The real period in this subsection is
\[
 \Omega_i=\int_{E_i(\R)}|\omega_i|,
\]
including both connected components when $E_i(\R)$ is disconnected.
The canonical height $\widehat h$ is relative to $(O)$; PARI's
\texttt{ellheight} and the regulator convention used in the BSD formula
are relative to $2(O)$, and hence equal $2\widehat h$ in rank one.
This factor of two is essential in every comparison below.

\begin{statement}[C.6.R1; Proved by descent and full saturation]
The algebraic and analytic ranks agree for the eleven curves. They are
two for $j=8000$ and $j=-12288000$, and one for all nine other values.
For every rank-one curve the section $P$ is a primitive generator of
$E(\Q)/E(\Q)_{\rm tors}$. Its index is therefore one. For the two
rank-two curves the index of $\Z P$ in this free quotient is infinite.
\[
\begin{array}{r|r|r|r}
i&j&\operatorname{rank}E(\Q)&[E(\Q)/E(\Q)_{\rm tors}:\Z P]\\\hline
1&8000&2&\infty\\
2&54000&1&1\\
3&287496&1&1\\
4&16581375&1&1\\
5&-3375&1&1\\
6&-32768&1&1\\
7&-884736&1&1\\
8&-12288000&2&\infty\\
9&-884736000&1&1\\
10&-147197952000&1&1\\
11&-262537412640768000&1&1
\end{array}
\]
\end{statement}
\begin{proof}[Arithmetic verification]
Exact $2$-descent in PARI/GP returns equal lower and upper rank bounds
for all eleven curves; its entire output, including independent points,
is in \path{code/C_cm_rank_lvalues.out}. Independent eclib descent
and saturation were completed for the first nine curves. The separate
full-saturation calculation \path{code/Crank_saturation.cpp} treats
all nine rank-one points, using eclib's automatic height-based index
bound and checking every prime up to that bound. All nine returned
index one and an empty list of unsaturated primes.

For clarity, the rounded-down bounds for the indices in the subgroup
of points having good reduction everywhere are respectively
\[
 1,1,1,2,2,2,3,5,5
\]
in the rank-one order $i=2,3,4,5,6,7,9,10,11$; the algorithm also
checks the primes dividing component-group orders. The underlying
height lower bound is the one in
\citet[Theorem 1]{CrCremonaSiksek2006}, applied to a minimal
Weierstrass equation and the subgroup
$E(\Q)\cap E^0(\R)\cap\bigcap_pE^0(\Q_p)$.
For the largest curve the lower bound on the nonzero height in this
subgroup is $0.651940632595967949\ldots$, its regulator is
$17.323913974368488616\ldots$, and the resulting index bound is
$5.154885040175890631\ldots$. The finite-field saturation maps for
$2,3,5$ already have full rank at residue primes $17,17,59$.
These data certify primitivity; merely finding $P$ in the output of
\texttt{ellrank} would not do so. The complete height-bound and
saturation transcripts are retained, including the independent
descent run that was stopped after its ninth completed curve.

The analytic ranks follow from the explicit derivative sums and bounds
proved below. The two rank-two bases furnished by descent are
$(-7,441),(105,833)$ for $i=1$, and
$(0,32004),(22218,3305456)$ for $i=8$.
\end{proof}

\begin{statement}[C.6.R2; Proved: central derivatives with explicit tails]
The following are the values of $L'(E_i,1)/\Omega_i$; zeros in rows
one and eight are exact. More digits, the periods, derivatives and
tail bounds are in \path{results_Crank.csv}.
\begin{center}\small
\begin{tabular}{r|l}
$i$&$L'(E_i,1)/\Omega_i$\\\hline
1&$0$\\
2&$5.134398425737889664345619204085016979705061236171249722654226\ldots$\\
3&$2.134120849044754020051650801945664635106456422166881627967598\ldots$\\
4&$31.81013979419927776012209427941545866492237285892013668350005\ldots$\\
5&$1.869462708725885120030080805051387134542481374813232701178304\ldots$\\
6&$3.428873379050567976208155749688133688667614469824389504923957\ldots$\\
7&$0.362583348071598548315763082767526599722077924892872839203426\ldots$\\
8&$0$\\
9&$9.629273918234131678053742152263035417231796757756973886621556\ldots$\\
10&$16.04238857830583792513538634637663343755958963581574670348088\ldots$\\
11&$1247.321806154531180413038860827393095408344267114181914400926\ldots$
\end{tabular}
\end{center}
The two second derivatives, divided by the same real periods, are
\begin{align*}
 L''(E_1,1)/\Omega_1
 &=28.71408722758099118128474692322079946598060512912402172882525\ldots,\\
 L''(E_8,1)/\Omega_8
 &=78.03222041637032223732443782051140882655135939512581188900823\ldots.
\end{align*}
In particular the leading Taylor coefficients for these two curves
are half the displayed second derivatives.
\end{statement}
\begin{proof}[Derivative formulas and error control]
Let $N$ be the conductor, $c=2\pi/\sqrt N$, and
$L(E,s)=\sum_{n\geq1}a_nn^{-s}$. Put
\[
 E_1(x)=\int_x^\infty\frac{e^{-u}}u\,du,
 \qquad G(x)=\int_x^\infty e^{-u}\frac{\log(u/x)}u\,du.
\]
Splitting the Mellin integral at $1/\sqrt N$ and using the functional
equation gives, for root number $-1$,
\begin{equation}\label{eq:Cr-first}
 L'(E,1)=2\sum_{n\geq1}\frac{a_n}{n}E_1(cn).
\end{equation}
For root number $+1$ and $L(E,1)=0$, it gives
\begin{equation}\label{eq:Cr-second}
 L''(E,1)=4\sum_{n\geq1}\frac{a_n}{n}G(cn).
\end{equation}
One can verify the factors directly by writing the completed Mellin
integral as an integral over $[1,\infty)$ of
$(u^{s-1}+wu^{1-s})f(iu/\sqrt N)$: its first derivative has
the factor $2\log u$ when $w=-1$, and its second derivative
the factor $2\log^2u$ when $w=1$. Integration by parts gives
the two formulas above. The vanishing of the lower derivatives
removes all derivatives of the gamma and conductor factors.

Hasse's bound and multiplicativity imply
$|a_n|\leq d(n)\sqrt n\leq2n$. Also
$E_1(x)\leq e^{-x}/x$ and $G(x)\leq e^{-x}/x^2$.
The latter follows on writing $u=x+v$ and using
$\log(1+v/x)\leq v/x$. Consequently the absolute tails after
$n=M$ are at most
\begin{equation}\label{eq:Cr-tails}
 \frac{4e^{-c(M+1)}}{c(M+1)(1-e^{-c})},\qquad
 \frac{8e^{-c(M+1)}}{c^2(M+1)^2(1-e^{-c})},
\end{equation}
respectively. For the first ten curves we used
$M=\lceil180/c\rceil$, exact integer coefficients, and 1024-bit
MPFR arithmetic. For the rank-two kernel we used the convergent series
\[
 G(x)=\frac12(\log x+\gamma)^2+\frac{\pi^2}{12}
       +\sum_{k\geq1}\frac{(-x)^k}{k^2k!}.
\]
After $k>x$ its terms decrease in absolute value. The alternating
remainder was cut below $10^{-140}$. At most 740 terms were needed;
the largest intermediate term was less than $1.4\cdot10^{72}$.
The 1024-bit arithmetic therefore leaves far more than 100 digits
after this cancellation. The largest bound in
\eqref{eq:Cr-tails} was $3.46\cdot10^{-77}$. The source and complete
outputs are \path{code/Crank_coefficients.gp},
\path{code/Crank_derivatives.cpp}, and their corresponding outputs.

For $i=11$, a full vector of $819$ million coefficients would be
unnecessarily large. We instead evaluated the CM theta series in
bounded-memory segments. The curve is the quadratic twist, by
$d=185801=7\cdot11\cdot19\cdot127$, of the conductor-$163^2$
curve
\[
 y^2+y=x^3-2174420x+1234136692.
\]
Write $D=163$, $\chi_D(a)=(a/D)$, and $\chi_d(n)=(d/n)$.
The exact coefficient $a_n$ is the sum of
\[
 \chi_d(n)\chi_D(a)a \quad(n=a^2,\ a>0)
\]
and
\[
 \chi_d(n)\chi_D((D+1)u/2)u
 \quad\left(n=\frac{u^2+Dv^2}{4},\ u,v>0,
                           \ u\equiv v\pmod2\right).
\]
This is the Hecke-character theta series for
$\mathcal O_{\Q(\sqrt{-D})}$, grouping conjugate algebraic integers
and then applying the quadratic twist; see
\citet[Section 3]{CrDarmon2009}. The first 10000 coefficients agree
exactly with PARI for both the untwisted and twisted curve. The same
code at $D=67,d=217$ reproduces the independent PARI derivative.

To accelerate the sum, in a block $n=n_0+r$, $|r|\leq R$, form the
exact integer moments $M_k=\sum a_nr^k$. For $x=cn_0$ write
\[
 e_k=e^{-x}(-x)^k/k!,\quad g_0=e^{-x},\quad
 g_k=e_k-g_{k-1},\quad
 \phi_0=E_1(x),\quad\phi_k=-\phi_{k-1}-g_{k-1}/k.
\]
Then the block contribution to $L'/2$ is
$\sum_k\phi_k M_k/n_0^{k+1}$. This follows by Taylor expansion
of $E_1(x(1+s))/(1+s)$. On $|s|=1/2$ its modulus is at most
$2E_1(x/2)$, which is less than $14$ throughout the block range.
We take enough terms that $(2R/n_0)^{k+1}<10^{-85}$.
The Cauchy remainder, summed with the actual
$\sum|a_n|<2\cdot10^{12}$ and $n_0>7700$, is less than
$10^{-75}$. The values of $E_1$ at consecutive block centres
are continued by the analogous Taylor expansion and refreshed
from a correctly rounded MPFR evaluation every 1024 blocks.
The continuation bound is $10^{-110}$ per normalized expansion.
All moments are integers, not floating-point sums.

The 384-bit run used $R=128$ and cutoff $170$; a separate 512-bit
run used $R=96$ and cutoff $190$. Their values agree beyond 75
decimal places. The latter used $915818441$ coefficients in
segments, with an omitted tail less than $3.10\cdot10^{-78}$,
and gave
\[
 L'(E_{11},1)=
 2.296582198163380608965915774436446055185027904705664673866754\ldots.
\]
The elementary operation-rounding bounds at these precisions are
smaller than $10^{-90}$: in the recurrences above
$|e_k|\leq1$, $|g_k|\leq1$, and $|\phi_k|<20$;
there are fewer than $5\cdot10^6$ blocks and the recurrences are
refreshed before 1024 consecutive continuations. Bounding each
normalized moment by $\sum|a_n|(R/n_0)^k$ gives this estimate
without assuming cancellation in the total sum. The arithmetic
and truncation margins are amply sufficient for every displayed
50-digit value. Both executable builds and both transcripts are
retained as \path{code/Crank_cm_stream*}.

The exact root numbers are $+1$ for $i=1,8$ and $-1$ otherwise.
All points are nontorsion by Statement C.5. For a CM curve,
the Coates--Wiles theorem, in the form
\citet[Theorem 10.1]{CrRubin1999}, says that a nonzero central
Hecke $L$-value forces the Mordell--Weil group over the CM field
to be finite. For a nonmaximal CM order first pass to an isogenous
curve with maximal order; ranks and $L$-functions are unchanged.
Since $L(E/K,s)=L(E/\Q,s)^2$ for the CM field $K$, nontorsion
therefore forces $L(E,1)=0$. This proves the needed central
vanishing also in the two sign-$+1$ cases. The functional equation
then forces $L'(E,1)=0$ for those two curves. The positive first
or second derivatives just evaluated prove exactly the asserted
analytic ranks.
\end{proof}

\begin{statement}[C.6.R3; Proved: the consequences of Gross--Zagier and Kolyvagin]
For each of the nine rank-one curves there is a positive rational
number $q_i$ such that
\begin{equation}\label{eq:Cr-rational-factor}
 \frac{L'(E_i,1)}{\Omega_i}=q_i\widehat h(P_i).
\end{equation}
Their Tate--Shafarevich groups over $\Q$ are finite. Consequently
the corresponding $D_{\rm part}$ is an explicit logarithm minus
$(6/(\pi q_i))L'(E_i,1)/\Omega_i$. The two rank-two curves do
not satisfy a rank-one formula; there $L'(E_i,1)=0$ while
$\widehat h(P_i)>0$.
\end{statement}
\begin{proof}
Gross--Zagier's theorem applies because each curve is modular and
$L(E_i,1)=0$. CM modularity follows from its Hecke-character
theta series. Their Theorem I.7.3, in the opened original
preprint \citep{CrGrossZagier1986}, gives a rational point $R$ and
a nonzero rational factor relating $L'(E,1)$ to its canonical
height times the real period. Changing their period convention
to the full real period changes that factor only rationally.
As $L'(E_i,1)>0$, $R$ has infinite order. Primitivity gives
$R=mP_i+T$, $m\ne0$, $T$ torsion; its height is
$m^2\widehat h(P_i)$. This proves
\eqref{eq:Cr-rational-factor}, including positivity.

The rank-one Gross--Zagier--Kolyvagin implication asserts rank one
and finite Tate--Shafarevich group when the central zero is simple.
One can use Kolyvagin's Heegner-point theorem as stated in
\citet[Theorem 1.3]{CrGross1991}, together with the nonzero
Gross--Zagier Heegner height; a modern primary research statement
of the rank conclusion is \citet[Theorem 1.1]{CrJetchev2008}.
The hypotheses are satisfied separately by the positive derivative
and exact sign in each of the nine rows. The two remaining rows
have analytic rank two, and this implication makes no assertion
of a height formula or of finiteness of their Tate--Shafarevich
groups. In particular it cannot express the second derivative
in terms of the height of $P$ alone: the rank-two regulator involves
a second independent point and their height pairing.
\end{proof}

\paragraph{What the numerical BSD comparison does and does not determine.}
For a primitive rank-one point, the BSD prediction in these conventions is
\[
 \frac{L'(E,1)}{\Omega_E}
 =\frac{2\#\operatorname{Sha}(E/\Q)\prod_pc_p}
        {\#E(\Q)_{\rm tors}^{\,2}}\,\widehat h(P).
\]
The independently evaluated quotient after dividing by
$2\widehat h(P)\prod_pc_p/\#E(\Q)_{\rm tors}^2$ is one, to
more than 69 digits, for $i=2,3,4,5,6,7,9,10$, and nine for
$i=11$. Thus the numerically predicted rational factors $q_i$ are
\[
 8,4,32,4,8,4,16,16,576,
\]
in that order. This initial numerical comparison alone proves neither
these exact factors nor the orders $1$ or $9$ of $\operatorname{Sha}$.
The exact factors are subsequently proved in C.6.R4 and C.6.R5:
the latter supplies an independent Gross--Zagier denominator bound
and rigorous intervals for all remaining ratios. This determines
the analytic BSD quotients exactly, while the orders of
$\operatorname{Sha}$ are proved here only in the three cases of
C.6.R4. In particular primitivity alone does not determine
$\#\operatorname{Sha}$ for the last curve.

\begin{statement}[C.6.R4; Known with derivation: three exact leading factors]
For $i=3,5,7$, namely $j=287496,-3375,-884736$, one has
\[
 \operatorname{Sha}(E_i/\Q)=0,\qquad
 \frac{L'(E_i,1)}{\Omega_i}=4\widehat h_O(P_i).
\]
Thus the factors $q_3=q_5=q_7=4$ are exact, not numerical
rational reconstructions. For the smallest of these examples,
$E:y^2+y=x^3-38x+90$ and $P=(0,9)$, this gives the exact
reduction
\[
 N(P)=\exp\left(-\frac{L'(E,1)}{4\Omega_E}\right).
\]
\end{statement}
\begin{proof}
We use \citet[Theorem 1.2 and Table 9]{miller2011}. In the
notation of its Definition 1.1, the theorem proves
$\mathrm{BSD}(E/\Q,p)$ for every prime $p$ when the curve has
analytic rank at most one, conductor less than $5000$, and
its isogeny class is outside the eleven exceptions in Table 9.
The exception conductors are
\[
546,570,858,870,1050,1230,1938,1950,2370,2550,3270.
\]
The three conductors here are $3136,3136,361$, and their
analytic ranks are one by C.6.R2. Hence every hypothesis is
verified. The theorem is stated for all curves in the isogeny
class, using isogeny invariance, so no optimal-curve assumption
is being silently imposed on our minimal equations.

In all three rows
$2\prod_pc_p/\#E(\Q)_{\rm tors}^2=4$, and primitivity is
proved in C.6.R1. Therefore the quotient
$L'(E_i,1)/(4\Omega_i\widehat h_O(P_i))$ is the positive
integer $\#\operatorname{Sha}(E_i/\Q)$. To isolate this integer
without relying on an unbounded decimal approximation to a
canonical height, we give an elementary coarse height certificate.

After completing the square, write the duplication map on the
$x$-line as $[F(X,Z):G(X,Z)]$, with relatively prime integral
quartics. The program
\path{code/C_height_coarse_certificate.py} supplies integer
homogeneous cubics $A_0,B_0,A_1,B_1$ and $D\ne0$ with
\[
 A_0F+B_0G=DZ^7,\qquad A_1F+B_1G=DX^7.
\]
Let $C$ be the larger of the coefficient one-norms of $F,G$
and of the two pairs $(A_j,B_j)$. For coprime integers $a,b$,
the common divisor of $F(a,b),G(a,b)$ divides $D$, by these
identities. Taking the identity corresponding to the larger
of $|a|,|b|$ proves, together with the elementary upper bound,
\[
 |h(x(2Q))-4h(x(Q))|\leq\log C.
\]
Telescoping consequently gives
\[
 \left|\widehat h_O(P)-\frac{h(x(2^6P))}{2\cdot4^6}\right|
 \leq\frac{\log C}{6\cdot4^6}.
\]
The exact constants in row order $3,5,7$ are
\[
 C=26610617537591808,\quad188716953600,\quad285931133.
\]
Six rational doublings, the displayed identities, and directed
Arb logarithms give the quotient enclosures
$1\pm0.003$, $1\pm0.003$, and $1\pm0.009$, respectively,
after including the central-derivative error bound from C.6.R2.
Each interval contains only the integer one. All quartics,
Bezout coefficients, rational iterates and enclosures are
retained in the program's output and its JSON certificate.
Finally C.6.3 gives $A=1$ in row seven, proving the exponential
formula. Its algebraicity remains the separate question C.7(c).
\end{proof}

\paragraph{The scope of Rubin's theorem.}
The applicable ordinary CM results concern a prime $p$ at which
$E$ has good ordinary reduction, equivalently $p$ splits in the CM
field. One precise opened formulation of Rubin's rank-one formula
is \citet[Theorem 1]{CrBDP2012}: for a curve over $\Q$ with CM
by the maximal order, an ordinary prime prime to the conductor,
and a simple central zero, an appropriate Katz $p$-adic $L$-value
is, up to a nonzero algebraic factor, the square of a formal
logarithm of a nontorsion rational point divided by a $p$-adic
period. Its input is analytic rank one, already verified here.
It is a $p$-adic identity; it does not identify the real exponential
$\exp\widehat h(P)$ as algebraic.

For the $p$-part of the real leading-term formula, a modern precise
version is \citet[Theorem A]{CrCastella2025}. In addition to analytic
rank one it requires full CM, abelian torsion field over the CM
field, $p\nmid6h_K$, splitting of $p$, conductor prime to $p$,
and that the different of $K$ divide the Hecke conductor exactly.
For the maximal odd-discriminant cases here, the latter follows
from the exponent two of the ramified prime in the rational
conductor. The nonmaximal cases of discriminant $-12$ and $-28$
can be passed through a $2$-isogeny at $p>3$. The discriminant
$-16$ case requires care: the direct application of this version
has a different-conductor hypothesis at $2$ that fails, and we
do not claim that this particular formulation proves its full
leading-term formula. Rubin's original ordinary main-conjecture
paper is \citep{CrRubin1991}; its bibliographic record was checked,
but its complete theorem text was not available in the opened
publisher copy, so no additional hypothesis-free claim is based
on it here.

In particular these ordinary-prime statements do not prove that
the $3$-primary part of $\operatorname{Sha}$ has order $9$ for the
last curve: $3$ is inert in $\Q(\sqrt{-163})$ and also excluded
from the stated $p>3$ hypotheses. C.6.R5 proves the corresponding
analytic quotient to be exactly $9$ by a separate denominator
argument. The identification with the actual Tate--Shafarevich
group remains a separate arithmetic question in that row.
}{}
\subsection{C.6. Supplement: exact leading factors from a genus-character denominator bound}

We retain the full real period $\Omega_i$ and the absolute canonical
height $\widehat h_O$ of the preceding subsection. The argument below
uses the Gross--Zagier formula to bound a rational denominator; it does
not assume the leading-term Birch--Swinnerton-Dyer formula.

\begin{statement}[C.6.R5; Proved: all remaining rational leading factors]
For the six rank-one rows $i=2,4,6,9,10,11$, one has
\[
 \frac{L'(E_i,1)}{\Omega_i}
   =q_i\widehat h_O(P_i),\qquad
 (q_2,q_4,q_6,q_9,q_{10},q_{11})=(8,32,8,16,16,576).
\]
Together with C.6.R4, this determines every rank-one factor:
\[
\begin{array}{c|rrrrrrrrr}
 i&2&3&4&5&6&7&9&10&11\\\hline
 q_i&8&4&32&4&8&4&16&16&576.
\end{array}
\]
The certified auxiliary ratios and denominator bounds used to identify
the six additional factors are as follows. Here $D$ is the auxiliary
negative fundamental discriminant,
$r_i=L(E_i^D,1)/\Omega_{i,D}$, and the denominator of $q_i$ is at most
$B_i$.
\[
\begin{array}{c|r|r|r|r|r}
 i&D&r_i&B_i&q_i&|a_i|\\\hline
 2&-23&64&256&8&32\\
 4&-31&64&256&32&64\\
 6&-19&8&32&8&8\\
 9&-19&16&64&16&16\\
 10&-3&400&1600&16&80\\
 11&-3&64&256&576&192
\end{array}
\]
The last column is the absolute genus-trace integer in the proof. These
identities determine the analytic BSD quotients exactly: they are one
for $i=2,4,6,9,10$ and nine for $i=11$. They do not by themselves
identify those numbers with the orders of the Tate--Shafarevich groups.
\end{statement}

\begin{proof}
We give the arithmetic denominator bound first and then the interval
certificate that isolates the listed rational numbers.

\paragraph{1. A parametrization under quadratic twisting.}
Let $E_0/\Q$ have conductor $N_0$, and suppose a parametrization
$\phi:X_0(N_0)\to E_0$ sends $\infty$ to $O$ and satisfies
\[
 \phi^*\omega_0=c_0f_0(q)\,dq/q,
\]
where $\omega_0$ is a minimal differential and $f_0$ is the normalized
newform. Let $b$ be an odd fundamental discriminant coprime to $N_0$,
and suppose that the minimal twist has conductor $N_0b^2$ and exact
invariants
\[
 c_4(E_0^b)=b^2c_4(E_0),\qquad
 c_6(E_0^b)=b^3c_6(E_0).
\]
Then $E_0^b$ has a parametrization from $X_0(N_0b^2)$ with the same
Manin factor $c_0$.

Indeed, put $m=|b|$ and, using addition on $E_0$, define
\[
 F(z)=\sum_{a\bmod m}\chi_b(a)\phi(z+a/m).
\]
For $\gamma=\left(\begin{smallmatrix}A&B\\C&D_0\end{smallmatrix}\right)$
in $\Gamma_0(N_0m^2)$ choose $a'\equiv aD_0^2\pmod m$.
Direct multiplication gives
$T_{a/m}\gamma T_{a'/m}^{-1}\in\Gamma_0(N_0)$, where $T_t$
denotes translation by $t$. Because $\chi_b(D_0^2)=1$, the sum is
invariant. Each summand is an algebraic morphism on a common
congruence cover; invariance therefore descends the sum to
$X_0(N_0m^2)$ over $\Q(\zeta_m)$. Its expansion at $\infty$
gives $F^{\sigma_u}=\chi_b(u)F$ for
$\sigma_u(\zeta_m)=\zeta_m^u$. Quadratic descent gives the required
map to $E_0^b$ over $\Q$, with $\infty$ mapping to $O$.
On differentials, the Gauss-sum identity gives
\[
 F^*\omega_0=c_0\tau(\chi_b)f_b(q)\,dq/q,
 \qquad \tau(\chi_b)^2=b.
\]
For the twist isomorphism $\iota:E_0^b\to E_0$ the exact invariant
identities imply
$\iota^*\omega_0=\sqrt b\,\omega_b$ up to sign, with no additional
minimal-model scaling. Cancellation proves the assertion. The map
need not be optimal; optimality is not used below.

If $c_0=1$, let $t_b=\#E_0^b(\Q)_{\rm tors}$ and
$n_b=\#\pi_0(E_0^b(\R))$. Manin--Drinfeld says that the image of
the rational cuspidal divisor $[0]-[\infty]$ is torsion. Its elliptic
logarithm, along the imaginary axis, is $\pm L(E_0^b,1)$.
Consequently $t_bL(E_0^b,1)$ is an integral multiple of the positive
real lattice period $\Omega_b^+$, and
\begin{equation}\label{eq:Cgz-central-integrality}
 t_bn_b\,\frac{L(E_0^b,1)}{\Omega_b}\in\Z,
 \qquad \Omega_b=n_b\Omega_b^+.
\end{equation}

\paragraph{2. The six small base curves and Heegner hypotheses.}
Use the following data; $E_i=E_0^d$ and $K=\Q(\sqrt D)$.
\[
\begin{array}{c|c|r|r|r}
 i&E_0&N_0&d&D\\\hline
 2&576\mathrm e2&576&-11&-23\\
 4&3136\mathrm r2&3136&57&-31\\
 6&121\mathrm b1&121&21&-19\\
 9&1849\mathrm a1&1849&21&-19\\
 10&4489\mathrm a1&4489&217&-3\\
 11&26569\mathrm a1&26569&185801&-3
\end{array}
\]
All six base Manin factors are one. For the first five bases this
is independently returned by PARI's unconditional modular-symbol
algorithm. For $26569\mathrm a1$, the value one and modular degree
$59368$ are explicitly tabulated by
\citet[Appendix B, p.~102, unnumbered $26569\mathrm a1/26569\mathrm a2$
block]{CgzBelabas2019}; their Algorithm 1 describes the computation.
The first five completed computations are in
\path{code/B_C6_gz_setup.out}. The separate exact audit
\path{code/B_C6_gz_exact_metadata.gp} checks the minimal equations,
both $c_4,c_6$ scaling identities for $b=d,dD$, both conductor
identities, all torsion orders, and all splitting conditions. It
does not recognize any of these identities from floating-point
scaling factors.

In every row $d,D,dD$ are fundamental discriminants,
$\gcd(d,D)=\gcd(dD,N_0)=1$, and every prime dividing $N_0$ splits
in $K$. The genus character for $K(\sqrt d)/K$ factors through
$\operatorname{Pic}(\mathcal O_{|d|})$ and has conductor exactly
$|d|$. On ideals prime to $d$ it is
\[
 \chi(\mathfrak a)=\chi_d(N\mathfrak a).
\]
For a unit of the order locally at a prime $p\mid d$, its norm is
a square modulo $p$, so this character is trivial on the order's
units. It is nontrivial on the full local unit group: at a split
prime the norm is multiplication on two copies of $\Q_p^\times$,
and at an inert prime the residue-field norm onto $\mathbb F_p^\times$
is surjective. Since $p$ is odd and unramified in $K$, this also
proves that the conductor contains every $p\mid d$ exactly once.
Thus the primitive ring-class-character hypotheses of
\citet[Theorem 1.1, pp.~2524--2525]{CgzCai2014} hold. Their
$\mu(N_0,D)=0$ and
$u=[\mathcal O_{|d|}^{\times}:\Z^{\times}]=1$; the latter remains
true for $D=-3$ because $|d|>1$.

Let $P_{\mathfrak a}$ be the corresponding Heegner points and put
\[
 Y=\sum_{[\mathfrak a]}\chi(\mathfrak a)\phi(P_{\mathfrak a}).
\]
The coefficients are $\pm1$, so this is an actual algebraic point,
not merely an element of a tensor product with $\Q$. It is fixed
by $\ker\chi$ and anti-invariant under the other automorphism of
$K(\sqrt d)/K$. Twisting therefore gives $Y_d\in E_i(K)$.
Artin factorization gives
\[
 L(s,E_0,\chi)=L(E_i,s)L(E_i^D,s).
\]
The intervals below prove $L(E_i^D,1)>0$, hence the auxiliary
curve has rank zero. The rank-zero implication used here is the
same unconditional modular/CM rank implication used in C.6.R2--R3.
Thus $E_i(K)$ has rank one. C.6.R1 proves that $P_i$ is primitive
over $\Q$. Since $Y_d-\overline{Y_d}$ is torsion and
$Y_d+\overline{Y_d}$ is rational, there is an integer $a_i$ with
\[
 2Y_d=a_iP_i\pmod{E_i(K)_{\rm tors}},\qquad
 \widehat h_O(Y_d)=\frac{a_i^2}{4}\widehat h_O(P_i).
\]

\paragraph{3. The exact height and period normalization.}
Write $V_0=\operatorname{area}(\C/\Lambda_0)$ for the period area
of the minimal differential. Integrating $\phi^*\omega_0$ gives
\[
 4\pi^2(f_0,f_0)_{\Gamma_0(N_0)}
   =\frac{\deg\phi}{c_0^2}V_0.
\]
The degree in the denominator of CST's Theorem 1.1 cancels, giving
\begin{equation}\label{eq:Cgz-product}
 L'(E_i,1)L(E_i^D,1)
  =\frac{2V_0}{c_0^2|d|\sqrt{|D|}}\,\widehat h_K(Y).
\end{equation}
Here $\widehat h_K$ is the diagonal N\'eron--Tate pairing relative
to $K$. This convention is explicit in
\citet[\S7.1.1, pp.~218--219, Theorem 7.1.2]{CgzYZZ2011}:
their underlying height is normalized by $[L:K]^{-1}$ over a
field of definition $L$, and their pairing is
$\widehat h_{\mathcal L}(x+y)-\widehat h_{\mathcal L}(x)
-\widehat h_{\mathcal L}(y)$.

One can also see the factor geometrically. For
$\mathcal L=\mathcal O(O)$, pull the normalized Poincar\'e bundle
back under $1\times\lambda_{\mathcal L}$. The result is
$m^*\mathcal L\otimes p_1^*\mathcal L^{-1}\otimes
p_2^*\mathcal L^{-1}$. On the diagonal its canonical height is
that of $[2]^*\mathcal L\otimes\mathcal L^{-2}\simeq
\mathcal L^2$, namely $2\widehat h_O$ in the absolute convention.
Changing from the absolute convention to relative heights over
$K$ multiplies by $[K:\Q]=2$. Therefore
\[
 \widehat h_K(Y)=4\widehat h_O(Y_d)
                =a_i^2\widehat h_O(P_i).
\]
This factor comes from the Poincar\'e pairing and the base field,
independently of any numerical height routine.

A real-adapted basis of $\Lambda_0$ gives
\[
 \Omega_i\Omega_{i,D}=
 \begin{cases}
 2V_0/(|d|\sqrt{|D|}),&\Delta(E_0)<0,\\
 4V_0/(|d|\sqrt{|D|}),&\Delta(E_0)>0.
 \end{cases}
\]
Indeed, a triangular real lattice has one real component and its
purely imaginary primitive period has twice the height of its
nonreal basis vector. A rectangular lattice has two real
components. Twisting by the two opposite signs interchanges the
real and imaginary directions, and the exact minimal differential
scalings in step 2 supply $1/(|d|\sqrt{|D|})$.
Substituting into \eqref{eq:Cgz-product} gives
\begin{equation}\label{eq:Cgz-denominator}
 q_ir_i=
 \begin{cases}
 a_i^2/c_0^2,&\Delta(E_0)<0,\\
 a_i^2/(2c_0^2),&\Delta(E_0)>0.
 \end{cases}
\end{equation}
In particular, if $r_i=A_i/B_i'$ is reduced and $c_0=1$, the
denominator of $q_i$ divides $2|A_i|$. We use the weaker upper
bound $B_i=4|A_i|$ in the interval certificate.

\paragraph{4. Global identification of the CM coefficient series.}
For $s=11,43,67,163$, put
$F_s=\Q(\sqrt{-s})$ and $\mathfrak p=(\sqrt{-s})$.
These fields have class number one and units $\{\pm1\}$.
The character
\[
 \psi((\alpha))=
  \left(\frac{\alpha\bmod\mathfrak p}{s}\right)\alpha
\]
is well defined because $(-1/s)=-1$. It has infinity type $(1,0)$
and finite conductor $\mathfrak p$. Hecke's theta construction
places $\sum_{\mathfrak a}\psi(\mathfrak a)q^{N\mathfrak a}$ in
$S_2(\Gamma_0(s^2))$: the nebentypus is
$\chi_{-s}(n)(n/s)=1$. This exact special case is also written
in \citet[\S3, pp.~4--5]{CrDarmon2009}. Grouping conjugate
generators $\alpha=(u+v\sqrt{-s})/2$ gives the integral formula
used by the coefficient stream:
\[
 b_n=\sum_{a^2=n,\ a>0}\left(\frac as\right)a
  +\sum_{\substack{u,v>0,\ u\equiv v\ (2)\\u^2+sv^2=4n}}
       \left(\frac{(s+1)u/2}{s}\right)u.
\]
For all four bases, exact PARI elliptic-curve coefficients equal
$b_n$ for $1\leq n\leq10000$.

This finite comparison proves equality of the entire modular forms.
For completeness, if a nonzero form $g$ of weight $k$ on
$\Gamma_0(N)$ has trivial character, take the product of
$g|_k\gamma$ over right coset representatives in
$\Gamma_0(N)\backslash\mathrm{SL}_2(\Z)$. It is a nonzero
holomorphic full-level form of weight $kI$, where
$I=N\prod_{p\mid N}(1+1/p)$. Every factor has nonnegative
order at infinity, so the product's order is at least
$\operatorname{ord}_\infty g$. The full-level valence formula
therefore implies $\operatorname{ord}_\infty g\leq kI/12$.
Equivalently, agreement through $\lfloor kI/12\rfloor$ proves
equality. The stronger arithmetic congruence statement is
\citet[Theorem 9.18]{CgzStein2007}; only the elementary
characteristic-zero argument is needed here. For $k=2,N=s^2$
the four bounds are
\[
 \left\lfloor\frac{s^2+s}{6}\right\rfloor
     =22,315,759,4455,
\]
all below $10000$. Applying the exact quadratic twist character
then identifies every coefficient in both the target and auxiliary
streams, including coefficients past the tested range. The exact
metadata transcript prints the four bounds and their comparisons.

\paragraph{5. Directed enclosures of both $L$-sums.}
For $c=2\pi/\sqrt N$, the root-number $+1$ formula is
\[
 L(E,1)=2\sum_{n\geq1}a_n\frac{e^{-cn}}n,
\]
while for root number $-1$ the first derivative is
\eqref{eq:Cr-first}. The exact coefficients in the smaller cases
come from retained PARI files. The largest auxiliary stream uses
$s=163$, signed twist $-557403$, conductor
$8254937936042721$, and $M=867617471$; the target stream uses
$s=163,d=185801,M=289205824$. Both use the cutoff $cM\geq60$
up to rounding of the integer cutoff, with the actual $M$ used in
the tail bound. No precision conclusion depends on the nominal
cutoff.

Here are the finite-sum and rounding certificates implemented in
\path{code/B_C6_central_stream.cpp}. All intervals use 256-bit
MPFR with outward rounding at every operation. In a block centred
at $b$ with radius $R=128$, set
$S_k=\sum a_n(n-b)^k\in\Z$. For $g(z)=e^{-cz}/z$, put
\[
 C_0=1,\quad C_k=C_{k-1}c/k,\qquad
 T_0=1/b,\quad T_k=(T_{k-1}+C_k)/b.
\]
The coefficient of $(z-b)^k$ is $(-1)^ke^{-cb}T_k$.
For $h(z)=E_1(cz)/z$, put
\[
 e_0=E_1(cb),\quad e_k=(-1)^ke^{-cb}T_{k-1}/k\ (k\geq1),
 \qquad h_0=e_0/b,\quad h_k=(e_k-h_{k-1})/b.
\]
The block sums are the finite sums of these interval coefficients
times the exact integers $S_k$.

On $|z-b|=b/2$, the two kernels have bounds
\[
 |g(z)|\leq\frac{2e^{-cb/2}}b=:U_g,
 \qquad
 |h(z)|\leq\frac{4e^{-cb/2}}{cb^2}=:U_h.
\]
For the second bound, use
$E_1(cz)=\int_1^\infty e^{-czt}\,dt/t$ when $\Re z>0$.
Thus truncation at degree $k$ contributes at most
\[
 U_g\frac{(2R/b)^{k+1}}{1-2R/b}\sum|a_n|
 \quad\hbox{or}\quad
 U_h\frac{(2R/b)^{k+1}}{1-2R/b}\sum|a_n|.
\]
The degree is selected to make the geometric factor small, but
the actual remainder is separately evaluated with upward rounding.
For the derivative stream, the initial $E_1$ in each segment is
evaluated using correctly rounded $-\operatorname{Ei}(-x)$ and
monotonicity. Subsequent centres, spaced by $W=257$, use
\[
 E_1(c(b+W))=E_1(cb)-\int_0^Wg(b+x)\,dx.
\]
Integrating the same Taylor expansion gives the interval for this
advance. Its omitted part is bounded by
$WU_g(2W/b)^{k+1}/(1-2W/b)$ and added to both endpoints.
The first $10280$ terms are summed directly. Segment boundaries
are multiples of $257$, so no coefficient is omitted or repeated.

Hasse's bound and multiplicativity give $|a_n|\leq2n$.
The infinite tails, including the factor two, are consequently at
most
\[
 \frac{4e^{-c(M+1)}}{1-e^{-c}},\qquad
 \frac{4e^{-c(M+1)}}{c(M+1)(1-e^{-c})},
\]
respectively. The program evaluates both with the downward bound
for $c$ and upward bounds for the remaining positive expressions.
It finally adds the block remainders, interval arithmetic errors,
and infinite tail. For row eleven the central-value tail is below
$5.07\cdot10^{-19}$, and the derivative tail is below
$2.82\cdot10^{-21}$. These modest precisions already suffice for
rational separation. All twelve completed interval outputs are
retained, as are the exact coefficient files.

\paragraph{6. Periods, heights, and exact rational separation.}
The periods are enclosed independently with 512-bit Arb. Let
$\tau=i\sqrt s$ for $i=2,4$ and
$\tau=(1+i\sqrt s)/2$ for $i=6,9,10,11$; here
$s=3,7,11,43,67,163$, respectively. The exact CM identification
and reduced-form class-number check are part of C.6.5. Put
$T=\Im\tau$, $q_\tau=e^{2\pi i\tau}$, and
$r=|q_\tau|<10^{-4}$. We enclose
\[
 E_4(\tau)=1+240\sum_{n=1}^{100}\sigma_3(n)q_\tau^n
                   +\text{tail}.
\]
Since $\sigma_3(n)<2n^3$, with $A=101$ the absolute tail is at most
\[
 480r^A\left\{
 \frac{A^3}{1-r}+\frac{3A^2r}{(1-r)^2}
 +\frac{3Ar(1+r)}{(1-r)^3}
 +\frac{r(1+4r+r^2)}{(1-r)^4}\right\}.
\]
The exact invariants $c_4,c_6$ of each minimal equation give
$\kappa=2\pi(E_4(\tau)/c_4)^{1/4}$. The sign of $E_6(\tau)$
is positive: $\sigma_5(n)<2n^5$ and
$n^5\leq32^{n-1}$ give
$|E_6(\tau)-1|\leq1008r/(1-32r)<1$.
Thus the full real period is
\[
 \Omega=n_{\R}\kappa\begin{cases}
 1,&c_6>0,\\
 T,&c_6<0,\ \Delta>0,\\
 2T,&c_6<0,\ \Delta<0,
 \end{cases}
 \qquad n_{\R}=\begin{cases}2,&\Delta>0,\\1,&\Delta<0.\end{cases}
\]
These formulas apply separately to the target and auxiliary curve.
Every operation, the Fourier tail, and the fourth root is enclosed
by Arb. The height intervals are the independent certificates
\path{records/C_cm_height_arb.json} of C.6.5.

The final computation is
\path{code/B_C6_exact_factors.py}. First, the intervals for
$r_i$ isolate a unique integer after multiplication by the exact
$t_bn_b$ in \eqref{eq:Cgz-central-integrality}. This proves the
six auxiliary rational values, with no BSD input. Next form the
interval $L'(E_i,1)/(\Omega_i\widehat h_O(P_i))$. The enclosures
below are deliberately enlarged decimal summaries of the recorded
Arb intervals:
\[
\begin{array}{c|c|c}
 i&|r_i-r_i^{\rm listed}|\text{ is at most}
   &|q_i-q_i^{\rm listed}|\text{ is at most}\\\hline
 2&4\cdot10^{-22}&2\cdot10^{-53}\\
 4&4\cdot10^{-20}&3\cdot10^{-45}\\
 6&4\cdot10^{-22}&2\cdot10^{-53}\\
 9&5\cdot10^{-21}&4\cdot10^{-53}\\
 10&3\cdot10^{-20}&10^{-23}\\
 11&7\cdot10^{-15}&8\cdot10^{-19}
\end{array}
\]
Each last-column bound is strictly less than $1/B_i$. A rational
number with denominator at most $B_i$ differing from the listed
integer would differ by at least $1/B_i$. Equation
\eqref{eq:Cgz-denominator} therefore makes every listed $q_i$
exact. The same equation gives the listed $|a_i|$ by integer
square-root checks. The JSON certificate records all input
intervals, exact invariants, rational values, and positive separation
margins. Its stdout concludes that all six remaining factors are
exactly certified. This completes the proof.
\end{proof}

\IfFileExists{sections/C_torsion_heights.tex}{\subsection{C.6. Supplement: the three specified smooth torsion fibres}

\begin{statement}[C.6.T1; Proved, with exact local and descent computations]
At $t=32/5$ a global minimal model and the image of the section are
\[
 E:y^2=x^3-2x+1,\qquad P=(0,1),\qquad \Delta=80.
\]
The conductor is $40$, $E(\Q)=\Z/4\Z$ is generated by $P$,
and $\widehat h(P)=0$. The index of $\Z P$ in the quotient
modulo torsion is one. The complete bad-prime data are
\[
\begin{array}{c|c|c|c|c}
p&\operatorname{ord}_p\Delta&\text{type}&\text{component of }P&\lambda_p(P)\\\hline
2&4&\mathrm{III}&\text{unique nonzero class, order }2&\log2/12\\
5&1&\mathrm I_1&\text{identity}&\log5/12.
\end{array}
\]
All other finite local heights vanish. Consequently
\[
 \log N=\frac{\log10}{12},\qquad
 D_{\rm part}=\frac6\pi\log\frac{10^{1/12}}{2\sqrt3}.
\]
Numerically the two sides of the global height identity equal
\[
 0.191882091082837140334832621223697017300091790719064414669444\ldots.
\]
With the full real period $\Omega=5.9376498936895458116293755517\ldots$,
the central derivative is
\[
 \frac{L'(E,1)}\Omega=
 0.071331625531738774092617964261659879861112925567606085287841\ldots.
\]
The analytic rank is zero; a vanishing first derivative is therefore
not expected even though the point $P$ is torsion.
\end{statement}
\begin{proof}
The change of variables from the displayed family is
$x=(24/5)^2X$, $y=(24/5)^3Y$. Tate's algorithm gives the
listed minimal model, conductor and local data. Since $P$ has
order four, $3P=-P$. The local division identity consequently
gives
\[
 \lambda_p(P)=\left(\frac{v_p(\Delta)}{12}
                         -\frac{v_p(\psi_3(P))}{8}\right)\log p.
\]
Here $\psi_3(P)=-4$, proving both local coefficients exactly.
The exact descent output is $[0,0,0,[]]$, and the torsion
calculation gives $\Z/4\Z$. These computations, the direct
global equality and the analytic derivative are in
\path{code/Crank_torsion_local.gp} and
\path{code/Crank_torsion_final.gp}, with their outputs.
As a further check, the exact plus modular symbol from $0$ to
$\infty$ is $-1/4$ in the real-period lattice normalization;
the numerical ratio $L(E,1)/\Omega$ is $1/8$. The completed
functional equation then gives the checked expression
$(\gamma+\log(2\pi/\sqrt{40}))/8$ for the displayed derivative
ratio. The local-height proof does not use this $L$-value check.
\end{proof}

\begin{statement}[C.6.T2; Proved, with exact number-field local arithmetic]
Put $s=\sqrt3$ and
$t_\pm=(3008\pm1728s)/11$. Over $K=\Q(s)$, a global minimal
model for the fibre at $t_-$ is
\[
 E_-:y^2=x^3-(18+10s)x+21+12s,\qquad P_-=(0,3+2s).
\]
The model and point for $t_+$ are its conjugates. The
discriminant of this model is $1032912+596352s$. Define
\[
 \mathfrak p_2=(2,1+s),\qquad\mathfrak p_3=(s),\qquad
 \mathfrak p_{11,-}=(11,s+5),\quad
 \mathfrak p_{11,+}=(11,s-5).
\]
The bad-prime data of $E_-$ are
\[
\begin{array}{c|c|c|c|c|c}
v&e_v&v(\Delta)&\text{type}&\text{component order of }P_-&\lambda_v(P_-)\\\hline
\mathfrak p_2&2&8&\mathrm{IV}^*&3&0\\
\mathfrak p_3&2&3&\mathrm{III}&2&0\\
\mathfrak p_{11,-}&1&1&\mathrm I_1&1&\log11/12.
\end{array}
\]
For $E_+$ replace $\mathfrak p_{11,-}$ by
$\mathfrak p_{11,+}$. The nonzero component at
$\mathfrak p_3$ is unique; the two nonzero classes at
$\mathfrak p_2$ are interchanged by $P\mapsto-P$ and the
class of $P$ is a generator. All omitted local heights vanish.
Both points have order six, and
\begin{align*}
 \log N(P_-)&=\frac1{12}\log(17+10\sqrt3),\\
 \log N(P_+)&=\frac1{12}\log(10\sqrt3-17),\\
 \sum_{\sigma:K\hookrightarrow\R}\log N(P_\sigma)
 &=\frac{\log11}{12}
 =\sum_{v<\infty}n_v\lambda_v(P)-2\widehat h(P).
\end{align*}
The common value in the last line is
\[
 0.199824606066530878671828631497094108318475571161451431268214\ldots.
\]
In particular $A_\mp=(17+10\sqrt3)^{1/12}$ and
$(10\sqrt3-17)^{1/12}$ respectively give
$D_{\rm part}(t_\mp)=(6/\pi)\log(A_\mp/(2\sqrt3))$.

The natural Mordell--Weil group here is $E_\pm(K)$; the
notation $E_{t_\pm}(\Q)$ requested in the prompt does not
define an elliptic curve over $\Q$, since its $j$-invariant
is irrational. Over $K$, both curves have rank one and torsion
group $\Z/6\Z$. Thus the index of the torsion section in the
free quotient is infinite. For example $(1,-1-s)$ has infinite
order on $E_-$.
\end{statement}
\begin{proof}
For $t_-$ the change from the original family to this model
has scale $u=48(-64+37s)/11$, with translation parameters zero.
It is obtained by multiplying the first global minimalization
scale by the unit $(2+s)^{-3}$. Hence it preserves all finite
minimal valuations. The conjugate scale gives $E_+$.
Tate's algorithm yields the displayed table. Its discriminant
ideal is
$\mathfrak p_2^8\mathfrak p_3^3\mathfrak p_{11,-}$ and its
conductor ideal is
$\mathfrak p_2^2\mathfrak p_3^2\mathfrak p_{11,-}$.

As $5P=-P$, the normalized local heights are computed exactly by
\[
 \frac{\lambda_v(P)}{\log p}
 =\frac1{e_v}\left(\frac{v(\Delta)}{12}
                         -\frac{v(\psi_5(P))}{24}\right).
\]
The valuations of $\psi_5(P)$ at the three listed primes are
$16,6,0$, respectively. This proves the table, including zero
contribution at the two additive primes. Here $f_v=1$ at all
three primes, so weighting by $n_v=e_vf_v$ gives precisely
$\log11/12$. The exact torsion values proved in C.4 have product
$(17+10\sqrt3)(10\sqrt3-17)=11$, proving the archimedean side.

For the rank, translate the rational two-torsion point to zero.
The resulting equation is
\[
 y^2=x^3-(9+6s)x^2+(45+26s)x.
\]
The $2$-isogeny descent gives
$[E(K):\widehat\phi E'(K)]=2$ and
$[E'(K):\phi E(K)]=4$. Every surviving covering has a
displayed point; every excluded covering has a recorded local
obstruction at $\mathfrak p_2$ or $\mathfrak p_3$.
The quotient formula $2^{r+2}=2\cdot4$ proves $r=1$.
The field arithmetic was certified by \texttt{bnfcertify};
Simon descent was run with \texttt{LIMBIGPRIME=0}, so all local
tests were deterministic. The complete transcript is
\path{code/Crank_torsion_descent_sage.out}. The unmodified
Simon scripts maintained in Sage are retained with their source
URLs; the older upstream scripts and their unsuccessful
compatibility attempts are also retained. The torsion calculation
gives $\Z/6\Z$, generated by $P$, independently of the rank
calculation.
\end{proof}

\paragraph{The corresponding number-field $L$-derivative.}
For completeness, the natural replacement for the nonexistent
rational-curve quantity at $t_\pm$ can also be computed. Put
\[
 \Omega_K=\frac1{\sqrt{12}}
  \prod_{\sigma:K\hookrightarrow\R}
       \int_{E_\sigma(\R)}|\omega_\sigma|.
\]
The Hasse--Weil $L$-function over $K$ has analytic conductor
$12^2\cdot396=57024$, gamma parameters $(0,0,1,1)$, and sign
$-1$. Direct 100-digit evaluation gives
\begin{align*}
 L'(E/K,1)&=2.046708855992864213402012895411307398829979893338038716037006\ldots,\\
 L'(E/K,1)/\Omega_K
 &=0.260797680551692455136263550519918387218960539765929916318910\ldots.
\end{align*}
This is a value for the curve over $K$, common to the two
conjugates. If the period is instead defined as the product of
the two full real periods, omitting the discriminant factor,
the ratio is
$0.075285805535274833467509818428226475613981023404668273560620\ldots$.
Both conventions and all input periods are explicitly recorded
in \path{code/Crank_torsion_final.out}; neither value enters
the exact torsion height proof.
}{}
\IfFileExists{sections/C_analytic.tex}{\section{C.8. An analytic proof of the elliptic maximum}

Throughout this section $q=e^{2\pi i\tau}$, $T=\Im\tau$,
$\Delta_0=q\prod_{n\ge1}(1-q^n)^{24}$, and
$t=E_4^3/(1728\Delta_0)$. Thus $t-1=E_6^2/(1728\Delta_0)$.
The discriminant $\Delta_0$ is the modular discriminant and must not
be confused with the discriminant $27t^8(t-1)^3$ of the rational
Weierstrass equation. We use the division identity
\eqref{eq:cs-division}, already proved above, and the modular
invariance of $N$. No interval cover or numerical root computation
enters the following proof.

\begin{statement}[C.8.1; Proved: three uniform bounds on the elliptic norm]
For every $z\in\C$ and every $\tau$ in the standard modular
fundamental region, we have
\[
 \begin{array}{c|ccc}
 \text{restriction}&T\le2&T\le7/5&T\le11/10\\ \hline
 N(z,\tau)&<17/10&<3/2&<143/100.
 \end{array}
\]
If $\wp_\tau(z)=0$ and $T\ge2$, then $N(z,\tau)<541/500$.
\end{statement}
\begin{proof}
Taking absolute values in the theta series and completing the square
gives
\[
 e^{-\pi(\Im z)^2/T}|\vartheta_1(z,\tau)|
 \le\sum_{n\in\Z}
 e^{-\pi T(n+1/2+\Im z/T)^2}
 \le\sum_{n\in\Z}e^{-\pi Tn^2}.
\]
The last inequality follows from Poisson summation: the Fourier
coefficients of the periodized Gaussian are positive, so its maximum
is at an integer. The eta product and $n^2\ge1+3(n-1)$ for $n\ge1$
therefore imply, on $a\le T\le b$,
\begin{equation}\label{eq:ca-sup}
 N(z,\tau)\le
 \exp\!\left(\frac{r_a}{(1-r_a)^2}\right)
 \max\{F(a),F(b)\},\qquad
 F(T)=e^{\pi T/12}
 \left(1+\frac{2e^{-\pi T}}{1-e^{-3\pi T}}\right),
\end{equation}
where $r_a=e^{-2\pi a}$. Indeed,
$-\sum\log(1-r^n)\le r/(1-r)^2$. The function $F$ is convex,
as follows by expanding its denominator as a geometric series;
each summand is an exponential of a real linear function of $T$.
We use $6/7<\sqrt3/2$ and the following rational bounds:
\[
\begin{array}{c|@{\,}c@{\,}c@{\,}c@{\,}c@{\,}c}
T&6/7&1&11/10&7/5&2\\ \hline
F(T)&<14211/10000&<14116/10000&<14180/10000&<14782/10000&<16944/10000.
\end{array}
\]
The eta factors in \eqref{eq:ca-sup} are less than $10047/10000$
for $a=6/7$ and less than $10019/10000$ for $a=1$.
Multiplying these fractions proves the asserted bounds, using
$[6/7,1]$ and $[1,2]$ separately for the first one.

We specify elementary arithmetic for all the scalar estimates in
this proof. One may use $314159/100000<\pi<355/113$ and, for
$0\le x\le20$, the Taylor bounds
\begin{equation}\label{eq:ca-exp}
 S_{80}(x)\le e^x\le
 S_{80}(x)+\frac{x^{81}}{81!}\frac1{1-x/82},
 \qquad S_{80}(x)=\sum_{j=0}^{80}\frac{x^j}{j!}.
\end{equation}
Taking reciprocals bounds negative exponentials. The remainder
bound follows by comparison of consecutive terms, not by numerical
sampling. These inequalities verify each displayed rational bound;
much shorter Taylor sums also suffice. Thus the decimal diagnostics
in the accompanying programs are dispensable.

For the cusp we check directly the analytic step of supplied K6.
Choose $|\Im z|\le T/2$ by translation, put $b_0=43/100$ and
$r_0=e^{-4\pi}$. The Fourier series
\[
 \frac{\wp_\tau(z)}{\pi^2}
 =\csc^2(\pi z)-\frac13
   +8\sum_{n\ge1}\frac{nq^n}{1-q^n}
          \bigl(1-\cos(2\pi nz)\bigr)
\]
shows, if $|\Im z|\ge b_0$, that
\begin{align*}
 \left|\frac{\wp_\tau(z)}{\pi^2}+\frac13\right|
 &\le \frac1{\sinh^2(\pi b_0)}
 +8\left\{\frac{r_0}{(1-r_0)^3}
 +\frac1{2(1-r_0)}
 \left(\frac{e^{-2\pi}}{(1-e^{-2\pi})^2}
       +\frac{e^{-6\pi}}{(1-e^{-6\pi})^2}\right)\right\}
 <\frac13.
\end{align*}
The two terms in the last parentheses bound respectively
$r^n e^{2\pi n|\Im z|}$ and $r^n e^{-2\pi n|\Im z|}$.
This inequality excludes a zero. Hence a zero has
$|\Im z|<b_0$. The Jacobi product, $\log(1+x)\le x$, and
the monotonicity of $\cosh(\pi y)e^{-\pi y^2/T}$ on
$0\le y\le b_0$ give
\[
 N(z,\tau)\le
 2\cosh(\pi b_0)e^{-\pi/3-\pi b_0^2/2}
 \exp\!\left(
 \frac{2\cosh(2\pi b_0)r_0}{1-r_0}
 +\frac{r_0^2}{1-r_0^2}\right)<\frac{541}{500}.
\]
For the stated monotonicity, the concave function
$\tanh(\pi y)-y$ is zero at zero and positive at $b_0$.
For $T\ge2$ the right side before substituting $T=2$ decreases:
the derivative of $-\pi T/6-\pi b_0^2/T$ is negative and
both product errors decrease. The two strict scalar inequalities
follow from \eqref{eq:ca-exp}. This proves the cusp assertion
without using any certificate from K6.
\end{proof}

\begin{statement}[C.8.2; Proved: an explicit local maximum on a larger disc]
If $|\tau-i|\le1/50$ and $\wp_\tau(z)=0$, then
$N(z,\tau)\le\sqrt2$, with equality exactly when $\tau=i$.
\end{statement}
\begin{proof}
Write $\delta=\tau-i$, $h=z-(1-\tau)/2$ and
\[
 G(h,\tau)=\wp_\tau((1-\tau)/2+h),\qquad
 H(h,\tau)=\log\frac{\vartheta_3(h,\tau)}{\eta(\tau)}.
\]
We choose the holomorphic logarithm whose value at $(0,i)$ is
$\log\sqrt2$. The half-period theta identity gives
\begin{equation}\label{eq:ca-H}
 \log N(z,\tau)=\Re H(h,\tau)-\frac{\pi(\Im h)^2}{T}.
\end{equation}
Both $G$ and $H$ are even in $h$. On the polydisc
$|h|\le1/4$, $|\delta|\le1/10$, they satisfy
\begin{equation}\label{eq:ca-GHbounds}
 |G|<30,\qquad |H|<2/3.
\end{equation}
Here are sufficient series estimates, including their margins.
Put
\[
 A_k=\sum_{n\ge1}n^k e^{-9\pi n^2/10}\cosh(\pi n/2),
 \quad a_0=2A_0,\quad a_1=4\pi A_1,\quad a_2=8\pi^2A_2.
\]
Then $a_0<3/10$, $a_1<19/10$, $a_2<12$.
The first two summands suffice to verify these inequalities: for
$n\ge3$, each of the two exponentials in the hyperbolic cosine
has successive-term ratio at most $e^{-29\pi/5}$, and the
factor $n^k$ increases this ratio by at most $(4/3)^2$.
We use \eqref{eq:ca-exp} for the first terms and this geometric
bound for the tail. With $r=e^{-9\pi/5}$, the identities
\[
 G=-\partial_h^2\log\vartheta_3-\frac{\pi^2}{3}E_2,
 \qquad
 |E_2|\le1+\frac{24r}{(1-r)^3}
\]
give
\[
 |G|\le\frac{a_2}{1-a_0}+\frac{a_1^2}{(1-a_0)^2}
       +\frac{\pi^2}{3}\left(1+\frac{24r}{(1-r)^3}\right)<30.
\]
Also
\[
 |H|\le-\log(1-a_0)+\frac{11\pi}{120}
                    +\frac{r}{(1-r)^2}<\frac23.
\]
The logarithm is well defined since $|\vartheta_3-1|<3/10$;
the displayed eta product fixes its branch. A geometric bound
for $-\log(1-x)=\sum_{n\ge1}x^n/n$ verifies the last estimate.

Let
\[
 e(\tau)=\frac{\pi^2}{3}(\vartheta_2(0,\tau)^4-
                              \vartheta_4(0,\tau)^4),\qquad
 E(\tau)=\frac{\pi^2}{3}E_2(\tau),\qquad
 g=\frac{g_2(i)}4.
\]
On $|\delta|\le1/50$ the following loose bounds hold:
\begin{equation}\label{eq:ca-derivatives}
 |e'|<14,\quad |e''|<60,\quad |g_2'|<500,
 \quad |E'|<2,\quad |E_6''|<100,
 \qquad 47<g<95/2.
\end{equation}
For completeness these are scalar Fourier-series bounds, not
interval computations. At $T\ge49/50$, upper bounds for
$|\vartheta_j|,|\vartheta_j'|,|\vartheta_j''|$ are
\[
\begin{array}{c|ccc}
j&0\text{ derivatives}&1\text{ derivative}&2\text{ derivatives}\\ \hline
2&929/1000&742/1000&670/1000\\
4&1093/1000&290/1000&911/1000.
\end{array}
\]
They follow by differentiating the absolutely convergent theta
series, taking absolute values, retaining two terms, and using
the same successive-term bound. Differentiating the fourth
powers proves the first two bounds in \eqref{eq:ca-derivatives}.
For $r\le1/470$, the remaining derivative majorants are
\begin{align*}
 |g_2'|&\le640\pi^5\frac{r(1+11r+11r^2+r^3)}{(1-r)^7}<500,\\
 |E'|&\le16\pi^3\frac{r(1+r)}{(1-r)^5}<2,\\
 |E_6''|&\le2016\pi^2
 \frac{r(1+r)(1+120r+1191r^2+2416r^3+1191r^4+120r^5+r^6)}
 {(1-r)^{11}}<100.
\end{align*}
These follow by summing $\sum n^k r^n$ and bounding
$(1-r^n)^{-1}$ by $(1-r)^{-1}$. Substitution of $r=1/470$
and $\pi<355/113$ verifies the rational inequalities.
Finally $g=\pi^4E_4(i)/3$ and
$E_4=1+240\sum n^3q^n/(1-q^n)$ give $47<g<95/2$.
The identity of K3, checked by differentiating $e$ and using
the theta heat equation, reads
\[
 e(i)=0,\qquad e'(i)=id,\qquad d=\frac{2g}{3\pi}>99/10.
\]

The differential equation $\wp''=6\wp^2-g_2/2$ determines the
first coefficients in the even expansion:
\[
 G(h,\tau)=e+a h^2+ea h^4+\mathcal R_G,
 \qquad a=3e^2-g_2/4.
\]
Consequently, on $|\delta|\le1/50$,
\[
 |e|\le7/25,\quad |a+g|\le\frac52+3(7/25)^2=\frac{3419}{1250},
 \quad |a|\le\frac{62794}{1250}.
\]
Cauchy's estimate from \eqref{eq:ca-GHbounds} yields
\[
 |\mathcal R_G|\le
 \frac{30(4|h|)^6}{1-16|h|^2}.
\]
On $|h|=1/10$ the term $ah^2$ has modulus at least
$55331/125000>0.4426$, whereas the sum of the other terms is at most
\[
 \frac7{25}+\frac7{25}\frac{62794}{1250}\frac1{10000}
       +\frac{30(4/25)^3}{1-4/25}<0.4277.
\]
Rouch\'e's theorem puts two zeros, with multiplicity, in
$|h|<1/10$. The disc injects into the torus and $\wp$ has exactly
two zeros on the torus, so it contains both of them.

Write $G=e+h^2Q$. The preceding estimates give, for $|h|\le1/10$,
\[
 |Q+g|\le
 \frac{3419}{1250}+\frac7{25}\frac{62794}{1250}\frac1{100}
       +\frac{30(4/25)^3}{(1-4/25)/100}<\frac{88}{5}.
\]
Taylor's formula and $|e''|<60$ give
$|e|\ge(99/10-3/5)|\delta|=(93/10)|\delta|$.
At either zero, therefore,
\begin{equation}\label{eq:ca-delta-h}
 |\delta|<7|h|^2,
\end{equation}
because $(95/2+88/5)/(93/10)=7$.

Expand
$H=H_0(\tau)+H_2(\tau)h^2+H_4(\tau)h^4+\mathcal R_H$.
The heat equation and the identity for $G$ give
\[
 H_0'=\frac{ie}{4\pi},\quad
 H_2=-\frac{e+E}{2},\quad H_4=-\frac a{12}.
\]
Thus $H_0(i)=\log\sqrt2$, $H_0'(i)=0$,
$|H_0''|<6/5$, $H_2(i)=-\pi/2$, $|H_2'|<8$, and
$|H_4|<21/5$. Cauchy's estimate gives
\[
 |\mathcal R_H|\le
 \frac{(2/3)(4|h|)^6}{1-16|h|^2}
 \le\frac{2048}{63}|h|^4
 \quad (|h|\le1/10).
\]
Combining \eqref{eq:ca-H} and \eqref{eq:ca-delta-h}, we obtain
\begin{align*}
 \log N-\log\sqrt2
 &\le-\frac\pi2\Re(h^2)-\pi(\Im h)^2+C|h|^4\\
 &=-\frac\pi2|h|^2+C|h|^4,\\
 C&=\frac{147}{5}+56+\frac{21}{5}
                 +\frac{2048}{63}+\frac{1100}{49}<150.
\end{align*}
The final term of $C$ bounds the change of $1/T$ from $1$:
$7\pi/(1-1/50)<1100/49$. Since $150|h|^2<3/2<\pi/2$,
the result is strict unless $h=0$. Equation
\eqref{eq:ca-delta-h} then gives $\delta=0$, where equality
holds by the checked isogeny computation at $i$.
\end{proof}

\begin{statement}[C.8.3; Proved: the comparison of the two local discs]
Every complex number $t_0$ with $|t_0-1|\le1/500$ has a
representative $\tau$ with $|\tau-i|<1/50$ and $j(\tau)=1728t_0$.
Consequently every zero of $\wp$ for that modulus satisfies
$N\le\sqrt2$, with equality only for $t_0=1$.
\end{statement}
\begin{proof}
At $i$, $E_6=0$ and the Ramanujan differential identity gives
$|E_6'(i)|=\pi E_4(i)^2>13/2$, since $E_4(i)>29/20$.
On $|\delta|=1/50$, \eqref{eq:ca-derivatives} gives
\[
 |E_6(\tau)|>\frac{13}{100}-\frac1{50}=\frac{11}{100}.
\]
On the whole disc, $|q|<1/470$ and
$|\Delta_0|\le |q|\exp(24|q|/(1-|q|))<1/400$.
It follows on the boundary that
\[
 |t-1|=\frac{|E_6|^2}{1728|\Delta_0|}
 >\frac{121}{43200}>\frac1{500}.
\]
Rouch\'e applied to $E_6'(i)\delta$ shows that $E_6$ has
exactly one zero in the disc. Hence $t-1$ has two zeros there,
with multiplicity, and another application of Rouch\'e shows
that $t-t_0$ has two for each $|t_0-1|\le1/500$.
This argument proves coverage of the $t$-disc, rather than
inferring coverage merely from a boundary bound. Elliptic
curves over $\C$ with the same $j$ are isomorphic. In terms of
lattices the isomorphism is a homothety; it carries the two
zeros of $\wp$ to the two zeros of the scaled $\wp$ and
preserves $N$. Statement C.8.2 applies to either zero.
\end{proof}

\begin{statement}[C.8.4; Proved: the first three normalized duplication steps]
Put $u=t-1$. For $n=2^k$, $k=1,2,3$, and $nP\ne O$,
\begin{equation}\label{eq:ca-ladder}
 \left(\frac{N(P)}{\sqrt2}\right)^{n^2}
 =\frac{3^{3/4}}{2n}\,
       \frac{N(nP)|u|^{-1/4}}{|A_k(u)|},
\end{equation}
where
\begin{align*}
 A_1(u)&=1,\qquad A_2(u)=1-5u/27,\\
 A_3(u)&=1-\frac{217}{27}u-\frac{20486}{729}u^2
       +\frac{10910}{19683}u^3
       +\frac{161383}{177147}u^4
       -\frac{245}{14348907}u^5.
\end{align*}
On $|u|\le1/28$ we have $|A_3(u)|\ge1-10|u|$.
For every $k\ge1$, not just these first three steps, define
\[
 E_k(t)=\frac{|\Delta_E|^{(4^k-1)/12}|t-1|^{1/4}}
 {2^{4^k/2}|\psi_{2^k}(P)|}.
\]
Then $(N(P)/\sqrt2)^{4^k}=E_k(t)N(2^kP)|t-1|^{-1/4}$,
and the removable limiting value is
\[
 \lim_{t\to1}E_k(t)=\frac{3^{3/4}}{2^{k+1}}
                 =2^{1-k}\frac{3^{3/4}}4.
\]
\end{statement}
\begin{proof}
Substitute the exact division polynomials from Statement C.4.2
in \eqref{eq:cs-division}. Equivalently, remove
$t^{2(n^2-1)/3}(t-1)^{n^2/4}$ from $\psi_n(P)$ and divide
the resulting polynomial by its value at $t=1$. For $n=2,4,8$
this gives the displayed polynomials and the constant
$3^{3/4}/(2n)$. All operations are rational polynomial operations;
\path{code/C_ladder_algebra.py} reproduces them and records
the next two steps as well. The general identity is C.4.2.
For its limit write $\mu=(1-\tau)/2+h$ as in C.3.1. Since
$n=2^k$ is even, $n\mu$ is congruent to $nh$ modulo the
period lattice. Hence
\[
 N(nP)=2\pi|\eta(i)|^2 n|h|(1+o(1)),\qquad
 t-1=\tfrac94\pi^4E_4(i)h^4(1+o(1)).
\]
The identity $E_4(i)^3=1728\eta(i)^{24}$ then gives
$N(nP)|t-1|^{-1/4}\to2n/3^{3/4}$.
Since the left side of the general ladder tends to one,
this proves the stated limit for every $k$.
The triangle inequality gives
$|A_3(u)|\ge1-|u|M$, where
\[
 M\le\frac{81}{10}+\frac{29}{28}
                +\frac1{28^2}+\frac1{28^3}+\frac1{28^4}<10.
\]
Here we bounded respectively the five absolute coefficients by
$81/10,29,1,1,1$. This is a rational inequality on one disc,
not a subdivision or a root search.
\end{proof}

\begin{statement}[C.8.5; Proved: R4 by explicit analytic estimates]
For every lattice $\Z+\tau\Z$ and every zero $z$ of its
Weierstrass function,
\[
 N(z,\tau)\le\sqrt2.
\]
Equality holds exactly for the modular orbit of $i$.
In the parameter of the family this says
$\max D_{\mathrm{part}}=-(3/\pi)\log6$, only at $t=1$.
\end{statement}
\begin{proof}
Reduce to the standard modular region. The portion $T\ge2$
is covered by Statement C.8.1. Suppose $T\le2$ and put
$r=|t-1|$. For $r\ge9/10$, the step $n=2$ in
\eqref{eq:ca-ladder} and $N(2P)<17/10$ give strict inequality,
because
\[
 \frac{27(17/10)^4}{256}<\frac9{10}.
\]
There is no $2$-torsion point of the section on the smooth base.

If $r\le9/10$, then $T<7/5$. Indeed, $T\ge7/5$ would give
$|q|<1/6500$ and
\[
 |E_4|>24/25,\qquad |\Delta_0|<1/6400,
 \qquad |t|>\frac{6400(24/25)^3}{1728}>19/10,
\]
contrary to $|t|\le19/10$. The estimates follow from
\[
 \sum_{n\ge1}\sigma_3(n)r^n
 \le\frac{r(1+4r+r^2)}{(1-r)^5},\qquad
 |\Delta_0|\le r\exp\frac{24r}{1-r}.
\]
For $1/28\le r\le9/10$, the $n=4$ step and
$N(4P)<3/2$ now apply. The function
$(1-5r/27)r^{1/4}$ is increasing up to $r=27/25$.
At its left endpoint its fourth power is
\[
 \frac1{28}\left(\frac{751}{756}\right)^4
 >\frac{2187}{65536}
 =\left(\frac{3^{3/4}}8\frac32\right)^4.
\]
This proves strict inequality on the whole annulus. The lower
bound for $A_2$ is positive there, so no division by a zero has
occurred.

If $r\le1/28$, then $T<11/10$. For otherwise $|q|<1/1000$,
and the Lambert series with
\[
 \sum_{n\ge1}\sigma_5(n)r^n
 \le\frac{r(1+26r+66r^2+26r^3+r^4)}{(1-r)^7}
\]
gives $|E_6|>9/20$, $|\Delta_0|<1/900$, and hence
$|t-1|>900(9/20)^2/1728>1/28$.
On $1/500\le r\le1/28$, the $n=8$ step,
$N(8P)<143/100$, and $|A_3|\ge1-10r$ prove the claim.
Indeed $(1-10r)r^{1/4}$ has a single critical point in this
interval, a maximum at $r=1/50$, so its minimum is at an endpoint.
At both endpoints its fourth power exceeds $9/5000$, whereas
\[
 \left(\frac{3^{3/4}}{16}\frac{143}{100}\right)^4
 <\frac7{4000}<\frac9{5000}.
\]
Again the displayed positive bound excludes a division-polynomial
zero. The remaining disc $r\le1/500$ is covered by Statement
C.8.3. The smooth limiting fibre at $t=0$ has
$N=3^{1/4}<\sqrt2$; the formula also extends there by
continuity, so the singular rational model causes no omitted
modulus. The only equality occurs at $t=1$.
\end{proof}

The supplied R4 and K6 already establish this inequality by
computer-assisted proofs. Statements C.8.1--C.8.5 give another
proof with a three-step division ladder and the explicit radius
$1/50$ in the $\tau$-plane. The programs
\path{code/C_ladder_estimates.py} and
\path{code/C_local_analytic_bounds.py} record numerical diagnostics
of the constants; the proof uses the displayed rational estimates
and their series remainder bounds alone.
}{}
\IfFileExists{sections/C_integral.tex}{
\subsection{Integral and Gaussian descriptions}

\begin{statement}[C.8.6; Proved: the Mahler-measure identities]
Let $\Lambda\subset\C$ be a lattice of area $A$, let $v$ be a zero
of $\wp_\Lambda$, and let $d\mu=A^{-1}dA_z$ be Haar probability
measure on $\C/\Lambda$. Put $U(z)=\log N(z)$ and
$L=U(v)$. Then
\begin{align}
 6L&=\int_{\C/\Lambda}\log\left|\frac{\wp'(z)^2}{\wp(z)^3}\right|d\mu(z),
 \label{eq:ca-mahler}\\
 N(v)^2&=|\Delta(\Lambda)|^{1/6}
      \exp\left(-\int_{\C/\Lambda}\log|\wp(z)|d\mu(z)\right).
 \label{eq:ca-wpmean}
\end{align}
If $C(u)=1-(g_2/4)u^2-(g_3/4)u^3$, the pushforward of $d\mu$
under $u=1/\wp(z)$ is
\begin{equation}\label{eq:ca-nu}
 d\nu(u)=\frac1{2A}\frac{dA_u}{|u|\,|C(u)|}.
\end{equation}
In particular its unnormalized integral is exactly $2A$, and
\[
 N(v)^6=4\exp\left(\int_{\C}\log|C(u)|d\nu(u)\right).
\]
Writing $e_1,e_2,e_3$ for the three branch values and
$p(x)=\int\log|x-\wp(z)|d\mu(z)$, the maximum theorem is equivalently
\[
 \int\log|C|d\nu\leq\log2,\qquad
 p(0)\geq\frac{p(e_1)+p(e_2)+p(e_3)}3-\frac{\log2}3.
\]
Equality in either inequality occurs exactly when $g_3=0$.
\end{statement}
\begin{proof}
The mean of $U$ is zero by the theta product calculation in C.1,
and near zero
$U(z)=\log|z|+\frac1{12}\log|\Delta|+o(1)$.
The divisor of $\wp$ is $(v)+(-v)-2(O)$, with multiplicity
two when $v=-v$. Subtract its logarithmic divisor from the
corresponding combination of $U$'s. The result is harmonic
on the compact torus and hence constant. Its expansion at
zero gives
\[
 \log|\wp(z)|=U(z-v)+U(z+v)-2U(z)
                  +\frac16\log|\Delta|-2L.
\]
Integration proves \eqref{eq:ca-wpmean}.
Similarly the divisor of $\wp'$ is the sum of the three
nonzero half-periods minus $3(O)$. The multiplication-by-two
identity of C.1 gives the product of $N$ at those half-periods
as $2$. Since $\wp'(z)\sim-2z^{-3}$, comparison at zero yields
\[
 \log|\wp'(z)|=\sum_{w\in E[2]\setminus\{O\}}U(z-w)-3U(z)
                         +\frac14\log|\Delta|.
\]
Its mean is $\frac14\log|\Delta|$. Twice this identity minus
three times the preceding mean proves \eqref{eq:ca-mahler}.
All logarithmic singularities are locally integrable, so these
integrations also cover a zero at a half-period.

The map $z\mapsto x=\wp(z)$ has degree two. Its two inverse
sheets each have Jacobian $|\wp'(z)|^{-2}$; branch values form
a measure-zero set. Thus its pushforward density is
$2/(A|4x^3-g_2x-g_3|)$. Under $x=1/u$ the Jacobian is
$|u|^{-4}$ and the denominator is $4|u|^{-3}|C(u)|$.
This proves \eqref{eq:ca-nu} and its normalization. The identity
$\wp'^2/\wp^3=4C(u)$ gives the next formula.

Replacing $v$ above by a half-period gives
$p(e_j)=\frac16\log|\Delta|-2U(w_j)$; summing and using
$\sum_jU(w_j)=\log2$ proves the centroid formulation.
The equality characterization is that of C.8.5,
since $j=1728$ is equivalent to $g_3=0$ for a nonsingular lattice.
\end{proof}

\begin{statement}[C.8.7; Proved: the two complex beta integrals]
For the square and equianharmonic lattices respectively, the
integrals in C.8.6 are
\[
 \int\log|C|\,d\nu=\log2,\qquad
 \int\log|C|\,d\nu=\frac32\log3-2\log2.
\]
They recover $N^2=2$ and $N^4=3$ with the same normalization.
\end{statement}
\begin{proof}
By homogeneity it suffices to take $C(u)=1-u^m$, with $m=2$
or $m=3$. For $-1<\Re s<1-1/m$ put
\[
 I_m(s)=\int_{\C} |u|^{-1}|1-u^m|^{s-1}dA_u.
\]
The substitution $w=u^m$, with all $m$ sheets counted, gives
\[
 I_m(s)=\frac1m B_{\C}\left(\frac1{2m},\frac{s+1}2\right),
 \quad
 B_{\C}(a,b)=\pi\frac{\Gamma(a)\Gamma(b)\Gamma(1-a-b)}
                   {\Gamma(1-a)\Gamma(1-b)\Gamma(a+b)}.
\]
We include a proof of the integral evaluation. For
$0<a,b$ and $a+b<1$, insert
\[
 |w|^{2a-2}=\frac1{\Gamma(1-a)}
       \int_0^\infty t^{-a}e^{-t|w|^2}\,dt
\]
and the corresponding formula for $|1-w|^{2b-2}$.
Tonelli's theorem applies to the positive integrand. The
Gaussian integral over $w$ is
$\pi(t+s)^{-1}\exp(-ts/(t+s))$.
Putting $t=rv$, $s=r(1-v)$ reduces the remaining integrals
to $\Gamma(1-a-b)$ and the beta integral $B(a,b)$, proving
the displayed formula. The same integrable bounds justify
differentiation in a neighbourhood of $s=0$.

Consequently, with $a=1/(2m)$ and $\psi=\Gamma'/\Gamma$,
\[
 \frac{I_m'(0)}{I_m(0)}
 =\frac12\left(2\psi(1/2)-\psi(1/2-a)-\psi(1/2+a)\right).
\]
The gamma duplication and triplication formulas, differentiated
logarithmically, evaluate this as $\log2$ for $m=2$ and
$\frac32\log3-2\log2$ for $m=3$. These can also be checked
directly from
$\psi(1/2)=-\gamma-2\log2$,
$\psi(1/4)+\psi(3/4)=-2\gamma-6\log2$ and
$\psi(1/3)+\psi(2/3)=-2\gamma-3\log3$.
Equation \eqref{eq:ca-mahler} completes the normalization test.
\end{proof}

\begin{statement}[C.8.8; Proved: the signed Gaussian identity]
For $\Lambda=\Z+\tau\Z$, $T=\Im\tau$, and every $z\in\C$,
\[
 \sqrt{2T}\,|\eta(\tau)|^2N(z,\tau)^2
 =\sum_{m\in\frac12\Lambda}\epsilon(m)
                         e^{-2\pi|z-m|^2/T},
\]
where $\epsilon=-1$ on $\Lambda$ and $+1$ on the three other
cosets of $\Lambda$ in $\frac12\Lambda$.
Both sides and all their derivatives converge locally uniformly.
\end{statement}
\begin{proof}
Write $z=b+a\tau$ and expand in the Fourier characters
$e^{2\pi i(kb+la)}$ of $(a,b)\in(\R/\Z)^2$.
For $G(z)=\sum_{\lambda\in\Lambda}e^{-2\pi|z-\lambda|^2/T}$,
Poisson summation gives the coefficient
\[
 \frac12\exp\left(-\frac{\pi|k\tau-l|^2}{2T}\right).
\]
The right side is the sum of the four half-period translates
of $G$, minus $2G$. Its coefficient is therefore
\[
 (-1)^{k+l+kl}
       \exp\left(-\frac{\pi|k\tau-l|^2}{2T}\right).
\]
To check the sign, it is $+1$ when both $k,l$ are even and
$-1$ otherwise.

On the left use
\[
 \vartheta_1(z,\tau)
 =-i\sum_{n\in\Z}(-1)^n
        e^{\pi i(n+1/2)^2\tau+2\pi i(n+1/2)z}.
\]
Multiply by its conjugate and by
$\sqrt{2T}e^{-2\pi(\Im z)^2/T}$.
In the double sum put $k=n-m$ and apply one-dimensional
Poisson summation to the remaining Gaussian in $m$.
Completing the square gives precisely the coefficient just
displayed. This proves equality of the Fourier series.
Positive definiteness of $|k\tau-l|^2$ and Gaussian decay
justify the changes of order, Poisson summation, and every
claimed differentiation.
\end{proof}

}{}
\IfFileExists{sections/C_measure.tex}{
\subsection{The Betti measure and torsion equidistribution}

We take $dd^c=(i/\pi)\partial\bar\partial$, so that
$dd^c\log|w|=\delta_0$. For a local complex coordinate $w$,
$dA_w=(i/2)dw\wedge d\bar w$.

\begin{statement}[C.8.9; Refuted sign and normalization, with the corrected current identity proved]
On a local period chart write $\mu=a\tau+b$, with $a,b$ real,
and put $T=\Im\tau$. The positive Betti measure is
\[
 m_P=db\wedge da
 =\frac{|\mu_w-a\tau_w|^2}{T}\,dA_w.
\]
Thus the form $da\wedge db$ in Lead 15(b) has the wrong sign.
The globally single-valued function $L(t)=\log N(P_t)$ satisfies
on $\mathbb P^1(\C)$
\begin{equation}\label{eq:ca-current}
 dd^cL=\frac1{12}\delta_\infty-m_P,\qquad
               \int_{\mathbb P^1}m_P=\frac1{12}.
\end{equation}
There are no atoms of $m_P$ at $0$, $1$ or $\infty$.
Let $T_0=\sqrt3/2$, $\rho=(1+i\sqrt3)/2$, and
$p=\wp'_\rho((2-\rho)/3)$. The density in the ordinary
$t$ coordinate satisfies
\begin{align}
 \frac{dm_P}{dA_t}
 &=\frac{|t|^{-2/3}}{64T_0|p|^{2/3}}
                      \bigl(1+O(|t|^{1/3})\bigr),
                         &&t\longrightarrow0,\label{eq:ca-density0}\\
 \frac{dm_P}{dA_t}
 &=\frac{|t-1|^{-3/2}}{24\pi^2\sqrt{E_4(i)}}
                      \bigl(1+O(|t-1|^{1/2})\bigr),
                         &&t\longrightarrow1.\label{eq:ca-density1}
\end{align}
\end{statement}
\begin{proof}
On a chart avoiding a zero of the section, the theta quotient
is holomorphic and nonzero. Its log modulus is pluriharmonic.
Only $-\pi(\Im\mu)^2/T$ contributes to the curvature.
Direct differentiation of this real quotient gives
\[
 \partial_w\partial_{\bar w}\frac{(\Im\mu)^2}{T}
       =\frac{|\mu_w-a\tau_w|^2}{2T}.
\]
Also
\[
 a_w=\frac{\mu_w-a\tau_w}{2iT},\qquad
 b_w=-\bar\tau\frac{\mu_w-a\tau_w}{2iT}.
\]
Taking the wedge product proves the formula for $m_P$ and
$dd^cL=-m_P$ on the chart. The transformation of $(a,b)$
under an integral change of period basis has determinant one,
so the measure patches. Translation of $\mu$ by a locally
constant integral period changes neither the formula nor
the measure. The formula remains valid on a root chart
at an elliptic modulus.

At $t=0$ the expansion proved in C.3 is
\[
 L(t)=\frac14\log3+\frac7{192}\Re t
       -B|t|^{4/3}+O(|t|^{5/3}),\qquad
 B=\frac{9\pi}{512T_0|p|^{2/3}}.
\]
The expansion is a convergent real-analytic expansion on a
root chart, so it can be differentiated termwise with the
indicated orders. The Euclidean identity
$\Delta_{\R^2}|t|^\alpha=\alpha^2|t|^{\alpha-2}$ gives
\eqref{eq:ca-density0}.

At $t=1$ put $h=\mu-(1-\tau)/2$. The checked K3 expansion
and the modular $j$ expansion give
\[
 t-1=A_1h^4(1+O(h^2)),\qquad
 A_1=\frac94\pi^4E_4(i)>0,\qquad
 L=\log\sqrt2-\frac\pi2|h|^2+O(|h|^4).
\]
Therefore
$L=\log\sqrt2-\pi|t-1|^{1/2}/(2\sqrt{A_1})+O(|t-1|)$.
Termwise differentiation on the root chart proves
\eqref{eq:ca-density1}. Both densities are locally integrable,
and the continuous potentials at $0$ and $1$ have no atoms.

For infinity take $u=1/t$, and set $\ell=\log(5+2\sqrt6)$.
The cusp branch is
$\mu=1/2-i\ell/(2\pi)+O(q)$, $q=e^{2\pi i\tau}$.
The theta and eta products, with the nonzero factor
$|2\sin\pi\mu|=2\sqrt3+O(q)$, give
\begin{equation}\label{eq:ca-cuspexp}
 L=\frac1{12}\log|u|+\log(2\sqrt3)-\frac1{12}\log1728
          -\frac{\ell^2}{2\log(1728/|u|)}
          +O\left(\frac{|u|}{|\log|u||}+|u|\right).
\end{equation}
Here $q=u/1728+O(u^2)$; the estimate follows directly by
bounding the convergent products. The nonlogarithmic
correction is continuous at $u=0$ and has no point mass.
Equivalently its Betti density is
\[
 \frac{dm_P}{dA_u}
  =\frac{\ell^2}{2\pi|u|^2\log(1728/|u|)^3}
       \bigl(1+o(1)\bigr),
\]
which is integrable. The single logarithmic term in
\eqref{eq:ca-cuspexp} contributes $\delta_\infty/12$.
This proves the global current identity; pairing it with
the constant function one gives the total mass.
\end{proof}

\begin{statement}[C.8.10; Proved: the correctly normalized equidistribution]
Let $D_n$ be the divisor, with intersection multiplicity, of
parameters $t\in\mathbb P^1\setminus\{0,1,\infty\}$ for which
$nP_t=O$. Then
\[
 \deg D_n=\frac{n^2}{12}+O(1),\qquad
 \frac1{n^2}[D_n]\ \longrightarrow\ m_P,\qquad
 \frac1{\deg D_n}[D_n]\ \longrightarrow\ 12m_P
\]
weakly as $n\to\infty$ through indices for which $D_n\ne0$.
For any nonrepeating sequence of individual algebraic torsion
parameters, their Galois-orbit probability measures likewise
converge to $12m_P$. Thus the probability-limit assertion
in Lead 15(b) requires the factor $12$.
\end{statement}
\begin{proof}
We apply \citet[Theorem 1.1 and Corollary 1.2]{demarco2017}.
Their hypotheses are a number field, a smooth projective
base curve, an elliptic surface and a rational section of
nonzero generic canonical height. Here they are $\Q$,
$\mathbb P^1$, the displayed rational elliptic surface,
and $P$, whose generic height for $(O)$ is $1/12$ by C.6.4.
At a torsion specialization the canonical height is zero.
Corollary 1.2 consequently gives convergence of every
nonrepeating sequence of Galois-orbit probability measures.

The limiting measure is the curvature of the height metric
divided by the degree of its line bundle. In our normalization,
C.8.9 computes the curvature as $m_P$, and C.6.4 computes
that degree as $1/12$. Hence the probability measure is
$12m_P$. This division by the degree is essential: a measure
of mass $1/12$ cannot be a weak limit of probability measures.
The source's informal description of the pullback current
is interpreted with this normalization.

For completeness, the passage from individual orbits to
the full division divisor does not require uniformity of
the equidistribution theorem in $n$. The surface height
formula, with our half-self-pairing convention, reads
\[
 \widehat h_O(Q)=\chi(\mathcal O_{\mathcal E})+(Q\cdot O)
                -\frac12\sum_v\operatorname{contr}_v(Q).
\]
Here $\chi=1$, and the finitely many fibre corrections are
bounded as $Q=nP$ varies, because they depend only on its
component. Since $\widehat h_O(nP)=n^2/12$, it follows
that $(nP\cdot O)=n^2/12+O(1)$. Removing the singular fibres
changes this by $O(1)$: the local intersection at each
such fibre is bounded as $n$ varies in characteristic zero.
One can see this in the formal group after a fixed
multiple enters it; multiplication by a nonzero integer
is a local analytic isomorphism over the residue field
$\C$. Thus $\deg D_n=n^2/12+O(1)$.

At any fixed smooth torsion parameter $t_0$, of exact order
$d$, the multiplicity in $D_n$ is zero unless $d\mid n$,
and is constant when $d\mid n$. Indeed the formal parameter
of $nP=[n/d](dP)$ has first term $(n/d)$ times that of $dP$,
with nonzero constant coefficient. Fix a continuous test
function and $\varepsilon>0$. Corollary 1.2 implies that
all but finitely many distinct Galois orbits of torsion
parameters have average within $\varepsilon$ of the limiting
integral; otherwise a countersequence would violate it.
The total weighted contribution of the finitely many
exceptional orbits is bounded independently of $n$, by
the multiplicity observation. Divide by $\deg D_n$ and
let $n\to\infty$. This proves the probability convergence
for $D_n$, and the degree asymptotic proves the unnormalized
version.
\end{proof}

\begin{statement}[C.8.11; Proved: the two first-derivative identities]
On an ordinary modulus chart put
\[
 a=\frac{\Im\mu}{T},\qquad
 W=\frac{\vartheta_1'(\mu,\tau)}{\vartheta_1(\mu,\tau)}
                         +2\pi ia,\qquad
 R=\frac{d\mu}{d\tau}-a.
\]
Then
\[
 L_\tau=\frac{W}{8\pi i}(W+4\pi iR),\qquad
 \det\frac{d}{d\tau}
 \begin{pmatrix}\tau&\mu\\\mu&\beta^{\rm part}/6\end{pmatrix}
       =\frac{(W+2\pi iR)^2}{4\pi^2}.
\]
Moreover
\[
 R=\frac{g_2}{6\pi i\wp'(\mu)}-\frac{W}{2\pi i},
 \qquad
 L_\tau=\frac{iW(3W\wp'(\mu)-2g_2)}
                         {24\pi\wp'(\mu)}.
\]
The critical-point factors in K4 follow with their stated
normalization after multiplying by $6/\pi$.
\end{statement}
\begin{proof}
Put $f=\vartheta_1'/\vartheta_1$ and
$E=\pi^2E_2/3$. The theta heat equation and
$\eta'/\eta=\pi iE_2/12$ give
\[
 \partial_\tau\log(\vartheta_1/\eta)
       =\frac{f_z+f^2+E}{4\pi i}.
\]
At $\wp(\mu)=0$, the identity $f_z=-\wp-E$ reduces
this to $f^2/(4\pi i)$. Thus, along the section,
\[
 \frac{d}{d\tau}\log(\vartheta_1(\mu,\tau)/\eta)
       =f\mu'+\frac{f^2}{4\pi i}.
\]
Differentiate the Gaussian $-\pi(\Im\mu)^2/T$ using
Wirtinger derivatives, retaining the factor $1/2$ in
the derivative of the real part of a holomorphic logarithm.
Substitution of $f=W-2\pi ia$ and $\mu'=R+a$ gives the
first formula. Since
$\beta^{\rm part}/6=(i/\pi)\log(\vartheta_1/(2\sqrt3\eta))$,
its derivative is $if\mu'/\pi+f^2/(4\pi^2)$.
Subtracting $\mu'^2$ gives
$(f+2\pi i\mu')^2/(4\pi^2)$, the second formula.

Finally differentiate $\wp_\tau(\mu(\tau))=0$ by the theta
heat equation. After substituting $\wp(\mu)=0$,
$\wp''(\mu)=-g_2/2$, and $f=W-2\pi ia$, the result is
$6\pi i\wp'R=g_2-3W\wp'$. This is the formula for $R$.
Inserting it into the first identity gives the last one.
These computations also independently verify the K4 step
used in C.3; its apparent singularity at a half-period is
handled by the root chart rather than division by zero.
\end{proof}

}{}
\IfFileExists{sections/C_critical.tex}{\section{C.8. Critical points and the global level curves}

We continue to write $L(t)=\log N(\mu(\tau),\tau)$, where
$t=j(\tau)/1728$ and $\wp_\tau(\mu)=0$. The sign and lattice
translation of $\mu$ do not affect any expression below. We use
Eichler and Zagier's unnumbered theorem on p.~399 of
\cite{EichlerZagier1982}, also stated as equation (6) and the
unnumbered theorem in Section~2 of \cite{DukeImamoglu2008}.
For the cusp branch in this problem their theorem says
\begin{equation}\label{eq:cc-EZ}
 \begin{aligned}
 \mu&=\frac12-i\beta-
       \frac{36i\sqrt6}{\pi}\sum_{n\ge1}\frac{a_nq^n}{n^2},
 &\beta&=\frac{\log(5+2\sqrt6)}{2\pi},\\
 \sum_{n\ge1}a_nq^n&=\frac{\Delta_0}{E_6^{3/2}},
 &\mu''&=144\pi i\sqrt6\frac{\Delta_0}{E_6^{3/2}}.
 \end{aligned}
\end{equation}
The branch of $E_6^{3/2}$ tends to one at the cusp. Thus the
formula applies directly on $\Im\tau>1$ and then by analytic
continuation away from the half-period branch points. Our later
use at $i$ is on the square-root chart specified explicitly below.
The original theorem and the hypergeometric reformulation concern
the logarithm $\mu$; neither asserts the global inequality for
$L$ proved here. Searches through these papers and their cited
work, and through the torus Green-function literature, did not
locate the radial monotonicity statement below. We make no
priority claim.

\begin{statement}[C.8.12; Proved: a radial derivative identity]
Put $p=\wp'_\tau(\mu)$ and
\[
 X=\frac{3Wp}{g_2},\qquad F=1-X,
 \qquad W=\frac{\vartheta_1'(\mu,\tau)}{\vartheta_1(\mu,\tau)}
              +\frac{2\pi i\Im\mu}{\Im\tau}.
\]
These are single-valued modular invariants wherever $g_2\ne0$.
They have continuous extensions at $t=0,1$, with
$X(0)=1/4$ and $X(1)=0$. At every $t\ne0,1$,
\begin{equation}\label{eq:cc-radial}
 (t-1)L_t=\frac{X(X-2)}{24},\qquad
 r\frac{d}{dr}L(1+re^{i\theta})
       =\frac{\Re(X(X-2))}{12}.
\end{equation}
Also, if $a=\Im\mu/\Im\tau$ and $R=\mu'-a$, then
\begin{equation}\label{eq:cc-FR}
 F=\frac{6\pi ip}{g_2}R.
\end{equation}
Consequently $|F|<1$ implies strict decrease along the ray
from $1$ through $t$.
\end{statement}
\begin{proof}
A lattice translation leaves $W$ unchanged, and replacing $\mu$
by $-\mu$ changes the signs of both $W$ and $p$. Under a modular
change, $W,p,g_2$ have respective weights $1,3,4$, so $X$ is
invariant. The two identities of C.8.11 give
\[
 L_\tau=\frac{iW(3Wp-2g_2)}{24\pi p},\qquad
 R=\frac{g_2-3Wp}{6\pi ip}.
\]
Ramanujan's identities and $p^2=-g_3$ give
\[
 t'_\tau=-2\pi it\frac{E_6}{E_4},\qquad
 \frac{t-1}{t'_\tau}=-\frac{E_6}{2\pi iE_4}.
\]
Substitution, using $g_2=4\pi^4E_4/3$ and
$g_3=8\pi^6E_6/27$, proves \eqref{eq:cc-radial} and
\eqref{eq:cc-FR}. Since $X(X-2)=F^2-1$, $|F|<1$ makes
its real part strictly negative.

At $i$, $W=p=0$ and $g_2\ne0$, so $X\to0$. At $\rho$, put $\delta=\tau-\rho$ and $p_0=\wp'_\rho(\mu_0)$.
The expansion in C.3.2 and Ramanujan's identities give
\[
 R=\frac{3p_0}{8\pi^2}\delta+O(|\delta|^2),\qquad
 g_2=\frac{3ip_0^2}{\pi}\delta+O(\delta^2).
\]
Substitution into \eqref{eq:cc-FR} gives $F(\rho)=3/4$.
Equivalently, the two possible limits furnished by
\eqref{eq:cc-radial} and $L_t(0)=7/384$ are $X=1/4,7/4$;
the expansion of $R$ selects $1/4$. The same expansions show
continuity on the coarse $t$-line.
\end{proof}

\begin{statement}[C.8.13; Proved: the cusp exclusion]
If $\Im\tau\ge7/5$, then $|F|<17/20$.
\end{statement}
\begin{proof}
Here $|q|<1/6600$. On the auxiliary circle $|q|=1/2000$,
the Lambert series and eta product give
\[
 |E_6-1|\le
 504r\frac{1+26r+66r^2+26r^3+r^4}{(1-r)^7},\qquad
 |\Delta_0|\le r\exp\!\frac{24r}{(1-r)^2},
 \quad r=\frac1{2000}.
\]
The right sides imply $|\Delta_0/E_6^{3/2}|<1/1250$.
In particular $E_6$ has no zero on this disc and the specified
branch is holomorphic there. By Cauchy's inequality, with
$s=10/33$,
\[
 \sum_{n\ge1}\frac{|a_nq^n|}{n}
 \le\frac{-\log(1-s)}{1250},\qquad
 \sum_{n\ge1}\frac{|a_nq^n|}{n^2}
 \le\frac1{1250}\left(s+\frac{s^2}{4(1-s)}\right).
\]
Equation \eqref{eq:cc-EZ} therefore yields
\[
 |\mu-(1/2-i\beta)|<\frac1{125},\qquad
 |\mu'|<\frac{13}{250},\qquad \beta<\frac{365}{1000}.
\]
For $r=1/6600$, the corresponding weight-four and weight-six
Lambert bounds are
\[
 240r\frac{1+4r+r^2}{(1-r)^5}<\frac{37}{1000},\qquad
 504r\frac{1+26r+66r^2+26r^3+r^4}{(1-r)^7}
       <\frac{77}{1000}.
\]
Thus $|E_4|>963/1000$ and $|E_6|^{1/2}<1038/1000$.
Since $6\pi|p/g_2|=\sqrt6|E_6|^{1/2}/|E_4|$,
\begin{equation}\label{eq:cc-cuspbound}
 |F|<\sqrt6\frac{1038}{963}
 \left\{\frac{365/1000+1/125}{7/5}+\frac{13}{250}\right\}
 <\frac{17}{20}.
\end{equation}
All scalar comparisons here follow from geometric series and
the exponential Taylor bounds \eqref{eq:ca-exp}; for the logarithm
use $-\log(1-s)=\sum s^n/n$ with its geometric remainder.
For reference, the upper quantities before the last rounding are
less than $789/10^6$, $755/10^5$, $51/1000$, and $841/1000$
for $|\Delta_0/E_6^{3/2}|$, the error in $\mu$, the bound on
$\mu'$, and the final bound on $|F|$, respectively.
\end{proof}

\begin{statement}[C.8.14; Proved: the compact exclusion by one Cayley disc]
Let $w=(\tau-i)/(\tau+i)$. On
$0<|w|\le2-\sqrt3$, the invariant $F$ satisfies $|F|<1$.
At the two equianharmonic points on the boundary its continuous
value is $3/4$.
\end{statement}
\begin{proof}
Set
\[
 \begin{gathered}
 r_0=2-\sqrt3,\qquad A=E_4(i),\qquad k=\pi^2A,\\
 v=\pi\sqrt A\,w,\qquad z=v^2=kw^2,\\
 z_0=kr_0^2,\qquad c_0=1-\frac1{\sqrt3}.
 \end{gathered}
\]
In this proof $z$ is a coordinate on the modular disc, not a
point of the elliptic curve. Define analytic functions $f_2,f_4,f_6$
by
\[
 \begin{gathered}
 \sqrt A f_2(v)=\frac{E_2(\tau)}{(1-w)^2}
                         -\frac3{\pi(1-w)},\\
 A f_4(v)=\frac{E_4(\tau)}{(1-w)^4},\qquad
 A^{3/2} f_6(v)=\frac{E_6(\tau)}{(1-w)^6}.
 \end{gathered}
\]
Ramanujan's differential identities become the rational system
\begin{equation}\label{eq:cc-ode}
 \begin{gathered}
 f_2'=\frac{f_4-f_2^2}{3},\qquad
 f_4'=\frac43(f_6-f_2f_4),\qquad
 f_6'=2(f_4^2-f_2f_6),\\
 (f_2,f_4,f_6)(0)=(0,1,0).
 \end{gathered}
\end{equation}
Uniqueness makes $f_4$ even and $f_2,f_6$ odd. Write
\[
 B(z)=f_4(v),\quad J(z)=\frac{f_6(v)}{2v},\quad
 D(z)=\frac{B(z)^3-4zJ(z)^2}{J(z)^{3/2}},\qquad
 H(z)=\sum_{n\ge0}\frac{-[z^n]D(z)}{16n^2-1}z^n.
\]
The branch has $J(0)^{1/2}=1$. The estimates below establish
its regularity on a larger disc, so these definitions are
unambiguous on $|z|\le z_0$.

Write $h=(1-w)(\mu-(1-\tau)/2)$. The Möbius chain rule and
\eqref{eq:cc-EZ} give
\[
 h''(w)=-\frac{4\mu''(\tau)}{(1-w)^3}
       =-\frac{i}{2\sqrt{3\pi}}w^{-3/2}D(kw^2).
\]
The root expansion at $i$ gives the branch
\begin{equation}\label{eq:cc-h}
 h=\frac{2i}{\sqrt{3\pi}}\sqrt w\,H(z).
\end{equation}
Indeed differentiating this series twice gives precisely the
coefficient denominator $16n^2-1$ above. The two possible affine
homogeneous terms are absent because the half-period difference
has a Puiseux expansion consisting only of half-integral powers.
Changing the sign of the square root changes the signs of $h,p,R$
and leaves $F$ unchanged.

Put $r=|w|$. The definitions of $a$ and $R$ give exactly
\[
 R=\frac{1-w}{2i}
       \left(h'+\frac{\overline w h+\overline h}{1-r^2}\right),
 \qquad
 \frac p{g_2}=-\frac{i\sqrt w}{\sqrt{3\pi}(1-w)}
                         \frac{\sqrt J}{B}.
\]
Consequently
\begin{equation}\label{eq:cc-Fseries}
 F=\frac{\sqrt J}{B}
 \left\{H+4zH'+\frac{2r}{1-r^2}(rH-\overline H)\right\}.
\end{equation}
All conjugates in this display are values at the same point;
$H$ has real coefficients, so $\overline{H(z)}=H(\overline z)$.

The apparent pole at the equianharmonic points must be removed
before estimating this expression. We have $B(-z_0)=0$.
At these points $R=0$, since the expansion at $\rho$ gives
$\mu'=-1/3=a$. Thus, if
\[
 S=\sqrt J H,\qquad T=\sqrt J(H+4zH'),
\]
then $T(-z_0)=c_0S(-z_0)$. Define the analytic quotients
\[
 C(z)=\frac{B(z)}{1+z/z_0},\qquad
 U(z)=\frac{T(z)-c_0S(z)}{(1+z/z_0)C(z)},\qquad
 V(z)=\frac{\sqrt{J(z)}}{C(z)}.
\]
The estimates below prove $C\ne0$ on the needed larger disc.
For $r<r_0$, equation \eqref{eq:cc-Fseries} is now
\begin{equation}\label{eq:cc-Ffactored}
 F=U+V\frac{dH+q(H-\overline H)}{1+z/z_0},
 \qquad
 d=\frac{2(r_0-r)}{(1+r_0)(1+r)},\quad q=\frac{2r}{1-r^2}.
\end{equation}
Here $q$ is only a real abbreviation within this display and
its estimates, not the Fourier nome.

We prove the following four coefficient bounds. For an analytic
power series $K=\sum K_nz^n$, put
$\|K\|_s=\sum|K_n|s^n$. Then, for $0\le s\le z_0$,
\begin{equation}\label{eq:cc-norms}
 \begin{aligned}
 \|U\|_s&\le\frac1{\sqrt3}+\frac{13}{200}s,&
 \|V\|_s&\le1+\frac7{50}s,\\
 \|H\|_s&\le1+\frac{113}{1000}s,&
 \|H'\|_s&\le\frac{13}{100}.
 \end{aligned}
\end{equation}
We give the coefficient arithmetic and all tails, rather than
leaving these inequalities to sampling a circle. In
\eqref{eq:cc-ode}, write $f_2=\sum a_nv^n$,
$f_4=\sum b_nv^n$, $f_6=\sum c_nv^n$. Starting with
$(a_0,b_0,c_0)=(0,1,0)$ (only in this recurrence the symbol
$c_0$ denotes the zero initial coefficient), compute
\begin{equation}\label{eq:cc-recurrence}
 \begin{aligned}
 a_{n+1}&=\frac{b_n-\sum_{j=0}^na_ja_{n-j}}{3(n+1)},\\
 b_{n+1}&=\frac{4(c_n-\sum_{j=0}^na_jb_{n-j})}{3(n+1)},\\
 c_{n+1}&=\frac{2(\sum_{j=0}^nb_jb_{n-j}-
                          \sum_{j=0}^na_jc_{n-j})}{n+1}.
 \end{aligned}
\end{equation}
Thus $B_n=b_{2n}$ and $J_n=c_{2n+1}/2$. For powers of a series
$K$ with $K_0=1$, the equally elementary recurrence is
\begin{equation}\label{eq:cc-power}
 [z^n]K^\alpha=\frac1n\sum_{j=1}^n
                 ((\alpha+1)j-n)K_j[z^{n-j}]K^\alpha.
\end{equation}
Together with $H_n=-D_n/(16n^2-1)$, these determine every
coefficient used. For example,
\[
 \begin{array}{c|rrrr}
 n&0&1&2&3\\ \hline
 B_n&1&10/9&5/27&4/81\\
 J_n&1&14/27&28/135&14/405\\
 H_n&1&13/135&-1999/153090&29299/28146690.
 \end{array}
\]
Division by $1+z/z_0$ uses $C_0=1$ and
$C_n=B_n-C_{n-1}/z_0$; the same recurrence applies to
$(T-c_0S)/(1+z/z_0)$. Division by $C$ uses the usual
convolution recurrence with constant coefficient one.

Here are sufficient rational enclosures for the only nonrational
constants in this arithmetic:
\begin{equation}\label{eq:cc-constants}
 \begin{gathered}
 14<k<72/5,\qquad
 1.0315619043558156493568263343177356527388253980203
 <z_0\\
 <1.0315619043558156493568263343177356527388253980205.
 \end{gathered}
\end{equation}
They follow without special-function evaluation from
$\pi=16\arctan(1/5)-4\arctan(1/239)$, the first 100 terms
of the alternating arctangent series, the exponential series
through degree 160, and
\[
 A=1+240\sum_{n=1}^{49}\frac{n^3e^{-2\pi n}}{1-e^{-2\pi n}}
       +\varepsilon_A,\qquad
 0<\varepsilon_A<
 \frac{240(4e^{-2\pi})^{50}}
 {(1-4e^{-2\pi})(1-e^{-2\pi})}.
\]
The last bound uses $n^3\le4^n$. Bounds for $\sqrt3$ follow
by squaring rational endpoints. All ensuing coefficient
comparisons are rational arithmetic using these enclosures.

To bound infinite tails, first consider $|z|=12/5$.
Here $|w|<5/12$, $\Im\tau>7/17$ and
$|e^{2\pi i\tau}|<1/13$. Lambert series give
\begin{equation}\label{eq:cc-outer}
 |B|<400,\qquad |J|<3000\qquad (|z|=12/5).
\end{equation}
For clarity the respective upper bounds before rounding are
\[
 \left(\frac{12}{7}\right)^4
 \left(1+240u\frac{1+4u+u^2}{(1-u)^5}\right)<400,
 \quad u=1/13,
\]
\[
 \frac13\left(\frac{12}{7}\right)^6
 \left(1+504u\frac{1+26u+66u^2+26u^3+u^4}{(1-u)^7}\right)<3000.
\]
We have used $A>1$ and $2\sqrt{12/5}>3$.
For $R=3/2$, let $\xi=5/8$ and retain coefficients through
degree 59. The tails of $B,J,C$ are respectively at most
\begin{equation}\label{eq:cc-tail60}
 \frac{400\xi^{60}}{1-\xi},\qquad
 \frac{3000\xi^{60}}{1-\xi},\qquad
 \frac{400z_0}{12/5-z_0}\frac{\xi^{60}}{1-\xi}.
\end{equation}
For the last bound, $B(-z_0)=0$ implies
$|C_n|\le400z_0(12/5)^{-n}/(12/5-z_0)$.
The finite sums from \eqref{eq:cc-recurrence}--\eqref{eq:cc-power}
with these tails give
\[
 \begin{array}{c|cccc}
 &\|B\|_{3/2}&\|J\|_{3/2}&\|C-1\|_{3/2}
                   &\min_{|z|\le3/2}\Re J(z)\\ \hline
 \text{bound}&<3276/1000&<2411/1000&<335/1000&>2049/10000.
 \end{array}
\]
For the last column, the lower bound used is exactly
\[
 1-J_2R^2-\frac{J_1^2}{8J_2}
       -\sum_{n=3}^{59}|J_n|R^n-
                   \frac{3000\xi^{60}}{1-\xi}.
\]
This is obtained by minimizing the quadratic
$\Re(1+J_1z+J_2z^2)$ on $|z|=R$; the minimum occurs at
$\cos\arg z=-J_1/(4J_2R)$. In particular, on $|z|\le3/2$,
\[
 |B|<33/10,\quad |J|<5/2,\quad |C|>3/5,\quad
 \Re J>1/5,\quad
 |D|\le\frac{(33/10)^3+6(5/2)^2}{(1/5)^{3/2}}<1000.
\]
Cauchy's inequality and the denominators in $H_n$ now give
\[
 |H|\le1+1000\sum_{n\ge1}\frac1{16n^2-1}
       =1+1000\left(\frac12-\frac\pi8\right)<110
                      \quad (|z|\le3/2).
\]
On $|z|\le7/5$, the same argument gives
\[
 |H+4zH'|\le1+1000\sum_{n\ge1}
                 \frac{(14/15)^n}{4n-1}
       \le1+\frac{1000}{3}\log15<910.
\]
It follows on $|z|=7/5$ that $|U|<10000$ and $|V|<3$;
indeed $|(1+z/z_0)C|>(7/(5z_0)-1)3/5$, and
$|\sqrt J|<8/5$.

Finally retain coefficients through degree 99 and put
$\eta=5z_0/7$. The tails in the coefficient norms for
$U,V,H,H'$ at $z_0$ are bounded respectively by
\begin{equation}\label{eq:cc-tail100}
 \frac{10000\eta^{100}}{1-\eta},\quad
 \frac{3\eta^{100}}{1-\eta},\quad
 \frac{110\eta^{100}}{1-\eta},\quad
 \frac{110\eta^{100}}{z_0}
 \left(\frac{100}{1-\eta}+\frac{\eta}{(1-\eta)^2}\right).
\end{equation}
The finite coefficient sums plus precisely these tails give
\begin{equation}\label{eq:cc-finite-table}
 \begin{array}{c|cccc}
 & (\|U\|_{z_0}-1/\sqrt3)/z_0
 & (\|V\|_{z_0}-1)/z_0
 & (\|H\|_{z_0}-1)/z_0 & \|H'\|_{z_0}\\ \hline
 \text{upper bound}&635/10000&1368/10000&1115/10000&1293/10000.
 \end{array}
\end{equation}
The recurrences, cutoffs and remainders above specify each finite
comparison completely. The attached exact integer/rational
script reproduces them with outward rounding on denominator
$10^{100}$ and independently checks every rational ODE
coefficient; no point of the modular domain is enumerated.
Equations \eqref{eq:cc-finite-table} imply
\eqref{eq:cc-norms}, with room in each inequality.

We finish with a single polynomial inequality. For $s=|z|=kr^2$,
\[
 |H-\overline H|\le2|\Im z|\,\|H'\|_s,\qquad
 \frac{|\Im z|}{|1+z/z_0|}\le s.
\]
The second inequality follows by writing $z/z_0=ue^{i\phi}$
and using $(\cos\phi+u)^2\ge0$. The useful cancellation is
\[
 |dH+q(H-\overline H)|^2
      =d^2|H|^2+4q(d+q)(\Im H)^2.
\]
Also
\[
 \frac d{|1+z/z_0|}\le
 e(r):=\frac{c_0}{(1+r)(1+r/r_0)},\qquad
 q\le\kappa r,\quad d+q\le c_0+\kappa r^2,
 \quad\kappa=1+\frac2{\sqrt3}.
\]
The last bounds use $(1-r_0^2)^{-1}=1/2+1/\sqrt3$.
It follows from \eqref{eq:cc-Ffactored}--\eqref{eq:cc-norms},
using $k<72/5$, that
\begin{equation}\label{eq:cc-scalar}
 |F|\le U_0(r)+V_0(r)
 \sqrt{e(r)^2H_0(r)^2+
       4\kappa r(c_0+\kappa r^2)(72/5)^2r^4(13/100)^2},
\end{equation}
where
\[
 U_0=\frac1{\sqrt3}+\frac{13}{200}\frac{72}{5}r^2,
 \quad V_0=1+\frac7{50}\frac{72}{5}r^2,
 \quad H_0=1+\frac{113}{1000}\frac{72}{5}r^2.
\]
To verify that the right side is less than one, let
\[
 M(r)=(1+r)^2(1+r/r_0)^2,
\]
\begin{equation}\label{eq:cc-P}
 \begin{aligned}
 P(r)={}&(1-U_0)^2M(r)\\
 &-V_0^2\left\{c_0^2H_0^2+
 4\kappa r(c_0+\kappa r^2)(72/5)^2r^4(13/100)^2M(r)\right\}.
 \end{aligned}
\end{equation}
This polynomial has $P(0)=0$. Write
\[
 \frac{P(r_0x)}{r_0x}
     =\sum_{j=0}^{14}b_j\binom{14}{j}x^j(1-x)^{14-j}.
\]
All coefficients belong to $\Q(\sqrt3)$. Explicit lower bounds,
obtained by expanding the displayed factored polynomial, are
\[
 \begin{array}{c|rrrrrrrr}
 j&0&1&2&3&4&5&6&7\\ \hline
 100b_j>&169&175&181&187&192&197&202&205\\[2pt]
 j&8&9&10&11&12&13&14&\\ \hline
 100b_j>&207&207&204&197&184&160&118&.
 \end{array}
\]
For an exact independent recipe, if $P(r)/r=\sum q_nr^n$,
then
\[
 b_j=\sum_{n=0}^j q_nr_0^n
                  \frac{\binom jn}{\binom{14}n}.
\]
For example $b_0=4-4\sqrt3/3$ and
$b_1=25357/13125-453\sqrt3/4375$. The rational enclosure
\[
 \begin{gathered}
 1.73205080756887729352744634150587236694280525381038<\sqrt3,\\
 \sqrt3<1.73205080756887729352744634150587236694280525381039.
 \end{gathered}
\]
verifies every entry by integer multiplication; its endpoints
are checked by squaring. The complete exact coefficient list is
also retained with the code. Every Bernstein basis function is
nonnegative on $[0,1]$, so $P(r)>0$ for $0<r\le r_0$.
Since $U_0<1$ there, squaring \eqref{eq:cc-scalar} is legitimate
and proves $|F|<1$. Formula \eqref{eq:cc-Ffactored} applies
at $r=r_0$ except at $z=-z_0$ by continuity; at that remaining
point C.8.12 gives $F=3/4$. This proves the assertion.
\end{proof}

\begin{statement}[C.8.15; Proved: every critical point and every level set]
For every zero $\mu$ of $\wp_\tau$, the zero sets of the two
critical factors are exactly
\[
 \begin{aligned}
 W=0&\quad\Longleftrightarrow\quad
       \tau\in\operatorname{SL}_2(\Z)\{i,\rho\},\\
 3W\wp'_\tau(\mu)-2g_2=0&\quad\Longleftrightarrow\quad
       \tau\in\operatorname{SL}_2(\Z)\rho.
 \end{aligned}
\]
On the coarse $t$-plane the function $L$ has no critical point
away from $t=1$. The point $t=0$ is a noncritical $C^1$ point,
with $L_t(0)=7/384$. The point $t=1$ is the unique global
maximum, of value $\log\sqrt2$, and has the fractional-power
local maximum described in C.3. In particular there is no
saddle in $0<t<1$. On the root charts, $h=0$ over $t=1$
is a nondegenerate maximum, while $\delta=0$ over $t=0$
is a degenerate critical point with nonzero harmonic cubic
leading term and three alternating ascending and descending
sectors. These are all critical points on the root charts.

For every angle $\theta$, $r\mapsto L(1+re^{i\theta})$ is
strictly decreasing for $r>0$, from $\log\sqrt2$ to $-\infty$.
For each $c<\log\sqrt2$, the level set $\{L=c\}$ is one
Jordan curve, strictly star-shaped about $1$. Its bounded
complement is $\{L>c\}$; its unbounded complement is
$\{L<c\}$. The curve is real analytic except possibly at
$t=0$ when $c=\frac14\log3$, where it is $C^1$ with nonzero
tangent. The closed superlevel sets form nested topological
discs. At the maximum the level set is the single point $1$,
and above it the level set is empty. The corresponding
statements for $D_{\rm part}$ follow by its increasing affine
change from $L$.
\end{statement}
\begin{proof}
Take a representative in the closed standard modular fundamental
region. If $T\ge7/5$, C.8.13 applies. Otherwise
$|\Re\tau|\le1/2$ and $\sqrt3/2\le T\le7/5$. For fixed
$\Re\tau$, the function
\[
 |w|^2=\frac{(\Re\tau)^2+(T-1)^2}
                    {(\Re\tau)^2+(T+1)^2}
\]
has its maximum on this interval at an endpoint and increases
with $(\Re\tau)^2$. The two endpoint bounds are $r_0^2$
and $41/601<r_0^2$. Thus C.8.14 covers the entire remaining
region. We conclude $|1-X|<1$ for every $t\ne1$.
Equations \eqref{eq:cc-radial} now prove strict radial
monotonicity, including $t=0$ by continuity. They also exclude
both critical equations $W=0$ and $3Wp-2g_2=0$ on the
ordinary base. The elliptic points must be checked on their
root charts: at $\rho$, the nonzero coarse derivative already
given rules out a critical point; at $i$, $W=p=0$ but $3Wp-2g_2=-2g_2\ne0$, and C.3 gives
the strict maximum. At $\rho$, $W=g_2=0$ and $p\ne0$,
so both displayed factors vanish. This proves their exact zero
sets, including the elliptic orbits. In particular the vanishing of both factors at
$\rho$ in the $\tau$ coordinate is a ramification effect,
not a saddle on the coarse base. More explicitly, on the root
chart at $\rho$,
\[
 L=L(0)+\frac7{192}\Re\left(\frac{ip_0^2}{\pi^3}\delta^3\right)
                   +O(|\delta|^4),
\]
with $p_0\ne0$, whereas on the chart at $i$,
$L=\log\sqrt2-\pi|h|^2/2+O(|h|^4)$.
These expansions give the asserted root-chart critical types.

At the cusp, the product formula and \eqref{eq:cc-EZ} give,
uniformly in $\Re\tau$,
\[
 L=-\frac{\pi T}{6}+\log(2\sqrt3)-
                      \frac{\pi\beta^2}{T}+O(e^{-2\pi T}),
 \qquad |t|=\frac{e^{2\pi T}}{1728}(1+o(1)).
\]
Hence $L\to-\infty$ uniformly as $|t|\to\infty$.
Continuity at $t=1$ supplies the other endpoint of each ray.
The intermediate value theorem and strict monotonicity give
one and only one radius $r_c(\theta)>0$ at each level $c$.
Properness and strict monotonicity imply continuity in $\theta$;
the nonvanishing radial derivative gives the stated local
regularity by the implicit-function theorem. Thus
$\theta\mapsto1+r_c(\theta)e^{i\theta}$ is a continuous
injective parametrization of a Jordan curve and has exactly
the asserted complementary regions.

The proposed route in the prompt that posits one saddle on
$(0,1)$ is therefore false. The corrected statement is the
strict radial foliation above. It supplies, in addition to
the separate duplication proof, a global proof of R4 and
resolves the critical-point and global-level-set question
left as D.4.6 in the supplied checkpoint.
\end{proof}
}{}
\IfFileExists{sections/C_periods.tex}{\section{C.7. Periods, compact subgroups, and the transcendence question}

We use a Weierstrass equation over $\overline\Q$ and its invariant
differential. Its period lattice is $\Lambda=\Z\omega_1+\Z\omega_2$,
with $\Im(\omega_2/\omega_1)>0$. The quasi-periods are defined by
$\zeta(z+\omega_j)=\zeta(z)+\eta_j$. We extend
$\eta:\Lambda\to\C$ real-linearly to $\C$. The discriminant in
this section is the discriminant of this equation, equivalently of
the lattice in the specified differential. The quantity $N$ is
independent of changing this algebraic equation and differential.

\begin{statement}[C.7.1; Proved: sigma expression and the third-kind period]
Let $P\in E(\overline\Q)$, with $P\ne O$ and $2P\ne O$, and
let $u$ be an elliptic logarithm of $P$. Then
\begin{equation}\label{eq:cp-sigma}
 2\log N(P)=\frac16\log|\Delta|
       +\Re\{\Lambda_1+u(\zeta(u)-\eta(u))\},
 \qquad \Lambda_1=2\log\sigma(u)-u\zeta(u).
\end{equation}
The choices of logarithms change the expression inside the real
part by purely imaginary integral multiples of $2\pi$.
There is an extension $G_P$ of $E$ by $\mathbb G_m$, defined
over $\overline\Q$ and corresponding, under the principal
polarization, to $P$, for which $\Lambda_1$, modulo a logarithm
of an algebraic number, is the negative of a relative period
of $[\Z\to G_P]$. Changing an algebraic differential of the
third kind adds an algebraic multiple of $u$ and algebraic
endpoint terms.
\end{statement}
\begin{proof}
The sigma--theta product, or its divisor and automorphy law
together with its derivative at zero, gives
\begin{equation}\label{eq:cp-klein}
 N(P)=|\Delta|^{1/12}|\sigma(u)|
                   \exp\!\left(-\frac12\Re(u\eta(u))\right).
\end{equation}
To check the constant, for $\Lambda=\Z+\tau\Z$ one uses
$\sigma(z)=e^{\eta_1z^2/2}\vartheta_1(z,\tau)/(2\pi\eta(\tau)^3)$,
$\Delta=(2\pi)^{12}\eta(\tau)^{24}$, and
$\eta(z)=\eta_1z-2\pi i\Im z/T$. Substitution gives exactly
the defining Gaussian in $N$. Homogeneity gives the general
lattice formula. Taking logarithms proves \eqref{eq:cp-sigma}.

We make the period identification explicit. Put
\[
 f_u(z)=\frac{\sigma(z+u)}{\sigma(z)\sigma(u)}e^{-\zeta(u)z}.
\]
Its logarithmic differential is
\[
 \frac{df_u}{f_u}
 =\frac12\frac{\wp'(z)-\wp'(u)}{\wp(z)-\wp(u)}\,dz.
\]
This is the pullback of an algebraic differential with residues
$+1$ at $-P$ and $-1$ at $O$. The apparent $0/0$ at $z=u$
has a removable singularity when $P$ is not $2$-torsion.
The rational cocycle $f_u(z+z')/(f_u(z)f_u(z'))$ defines
the algebraic extension; in its rational trivialization the
exponential map has the form
\[
 \exp_{G_P}(z,v)=(\exp_E z,e^v f_u(z)).
\]
Its periods include
\begin{equation}\label{eq:cp-periods}
 (0,2\pi i),\qquad
 (\omega_j,\zeta(u)\omega_j-\eta_j u),\quad j=1,2.
\end{equation}
These signs follow directly from the sigma automorphy law.
The identity
$\sigma(2u)=-\wp'(u)\sigma(u)^4$ implies
\[
 \log f_u(u)=\Lambda_1+\log(-\wp'(u)).
\]
For an algebraic lift $R=(P,r)$, $r\in\overline\Q^{\times}$,
the second coordinate of $\log_{G_P}R$ is
\begin{equation}\label{eq:cp-v}
 v=\log r-\Lambda_1-\log(-\wp'(u)).
\end{equation}
It is a period of the relative homology defining the indicated
$1$-motive. This also avoids treating an integral with an endpoint
at a residue pole as an unregularized ordinary integral.
The description agrees with the generalized logarithm and
Serre function in \citet[Definitions 2.3--2.4 and
Section 6]{bertolin2022}; the signs in
\eqref{eq:cp-periods} are for the exponential rather than the
quasi-periods of $f_u$ itself.
\end{proof}

\begin{statement}[C.7.2; Proved: the compact-subgroup criterion and the real formula]
For $P$ as in Statement C.7.1, $N(P)$ is algebraic if and only
if the fibre of $G_P$ above $P$ contains an algebraic point of
the maximal compact subgroup $G_P(\C)^c$.
For the family in the prompt, if $t>1$ is real, choose
$\tau=iT$, $T>1$, and the zero $\mu=1/2-iy$, $0<y<T/2$.
With $x=2y/T$, one has
\[
 N(P)^4=3^{3/4}(t-1)^{-1/4}
       \left|\frac{\vartheta_1(x,i/T)}{\eta(i/T)}\right|.
\]
\end{statement}
\begin{proof}
On the analytic universal cover of $G_P$ define the real
homomorphism
\[
 \chi(z,v)=\Re\{v-\zeta(u)z+u\eta(z)\}.
\]
It vanishes on the periods \eqref{eq:cp-periods}, hence descends
to $G_P(\C)$. Its restriction to the torus is $\log|\cdot|$.
The kernel is compact: in the universal cover it is a real
three-dimensional vector space modulo the rank-three period
lattice. Every compact subgroup has zero image in $\R$, so
this kernel is exactly $G_P(\C)^c$.
The identity
$\Re(u\eta(z))=\Re(z\eta(u))$ follows from the Legendre relation;
equivalently $\eta(z)=sz+\pi\bar z/\operatorname{area}(\Lambda)$
for a complex constant $s$. Substituting
\eqref{eq:cp-v} at $z=u$ gives
\begin{equation}\label{eq:cp-chi}
 \chi(R)=\log|r|-\log|\wp'(u)|
            -2\log N(P)+\frac16\log|\Delta|.
\end{equation}
If $r$ is algebraic and $\chi(R)=0$, this equation makes
$N(P)^2$ algebraic. Conversely, if $N(P)$ is algebraic,
choose the positive real algebraic number
$r=|\wp'(u)|\,|\Delta|^{-1/6}N(P)^2$.
Then \eqref{eq:cp-chi} vanishes. This proves both directions
with the chosen rational trivialization of the fibre.
The compact interpretation is the one discussed by
\citet[slides 5--8]{bertrand2017}; here the calculation fixes
its normalization on the diagonal $P=P$.

For the real formula, the rectangular elliptic curve has its
two zeros of $\wp$ on $1/2+i\R$; along that edge $\wp$ runs
between the positive root $e_1$ and the negative root $e_2$.
Thus the indicated representative exists. Translate $2\mu$
by $-1$ and apply the modular substitution $\tau\mapsto-1/\tau$.
The resulting elliptic argument is $-2y/T$, which is real.
The Gaussian is therefore one. Evenness of $N$ and the
duplication identity prove the displayed formula. This removes
the Gaussian from this particular expression but makes no
algebraicity assertion about its theta value.
\end{proof}

\subsection{The conditional theorem and its exact conjectural input}

We use Grothendieck's period conjecture in the form stated as
Conjecture 1.1 of \citet{bertolin2022}: for a motive $M$ over
$\overline\Q$, the field generated by its periods over
$\overline\Q$ has transcendence degree $\dim G_{\mathrm{mot}}(M)$.
Only $1$-motives occur below. We also use their Construction 3.4
and Lemma 3.5: if $B$ is the smallest abelian subvariety
generated by the point and extension data, and $Z(1)$ is the
central torus constructed from the Weil-pairing bracket and
the remaining algebraic lift, then
\begin{equation}\label{eq:cp-MTdim}
 \dim G_{\mathrm{mot}}(M)
 =\dim G_{\mathrm{mot}}(A)+2\dim B+\dim Z(1).
\end{equation}
The bracket is described explicitly in their Section 6.1.
Their Table 2 corrects the table in the earlier paper; the
dependent nondeficient case has unipotent dimension three.
We now verify which case occurs here.

\begin{statement}[C.7.3; Proved: the motivic groups needed on the diagonal]
Let $P$ be nontorsion and let $R$ be any algebraic lift above
$P$ in $G_P$. Put $M_R=[\Z\to G_P]$, $1\mapsto R$.
Its connected motivic group is, after a choice of Levi subgroup,
the semidirect product of $G_{\mathrm{mot}}(E)$ and the
three-dimensional Heisenberg group attached to $H_1(E)$ and
its principal polarization. Its dimension is $5$ for CM and
$7$ without CM.

For a non-real embedding put
$\mathcal M=M_R\oplus\overline{M_R}$.
If $E$ and $\bar E$ are not isogenous, its dimension is $13$.
If they are isogenous, choose an isogeny
$\phi:\bar E\to E$ of degree $d$ and write
$Q=\phi(\bar P)$. Let $r$ be the dimension over
$F=\End(E)\otimes\Q$ of the span of $P,Q$, modulo torsion.
Then
\[
\begin{array}{c|c|c|c}
\text{case}&\dim B&\dim Z&\dim G_{\mathrm{mot}}(\mathcal M)\\ \hline
r=2,\ E\text{ CM}&2&2&8\\
r=2,\ E\text{ not CM}&2&2&10\\
r=1,\ E\text{ CM}&1&1+e&5+e\\
r=1,\ E\text{ not CM}&1&1+e&7+e.
\end{array}
\]
Here $e$ is either zero or one: it is zero precisely when the
algebraic lift in the one-dimensional quotient of the two
central tori by the bracket torus is torsion. This criterion is
independent of the choices made modulo isogeny. In the real
case we use the single motive $M_R$; no extra quotient torus
is necessary.
\end{statement}
\begin{proof}
All curves, polarizations, points, and lifts are algebraic over
$\overline\Q$, as required by Construction 3.4 and Lemma 3.5.
For $M_R$ the smallest abelian subvariety containing $(P,P)$ is
the diagonal $B=E$, since $P$ is nontorsion. If $\psi$ is the
polarization pairing, the bracket on this diagonal is
\[
 [(x,x),(y,y)]=\psi(x,y)-\psi(y,x)=2\psi(x,y).
\]
It is nonzero and spans $\Q(1)$. Hence $Z'=Z=\mathbb G_m$;
there is no deficient or Ribet exception, including for CM.
The pure group is $\mathrm{GL}_2$ without CM and
$\operatorname{Res}_{F/\Q}\mathbb G_m$ for the imaginary
quadratic CM field. Their dimensions are four and two.
This proves the single-motive assertion, including its Lie
bracket and Levi action: the action on the centre is the
polarization multiplier.

For the direct sum only the two diagonal central coordinates
occur. If the two elliptic curves are not isogenous, they
cannot be CM conjugates. The pure group is
$\{(g_1,g_2)\in\mathrm{GL}_2^2:\det g_1=\det g_2\}$,
of dimension seven. The point data generate $E\times\bar E$,
and the two brackets are independent, giving $7+4+2=13$.

If the curves are isogenous, transport the second polarized
factor through $\phi$. Its polarization pairing becomes
$\psi/d$. The smallest abelian subvariety has dimension $r$.
For $r=2$, independent variation in its two factors makes the
two diagonal brackets span both central coordinates; thus
$Z'=Z=\mathbb G_m^2$. The unipotent group is the product
of two Heisenberg groups, with the same pure group acting
on both standard factors.
For $r=1$, after multiplying the points by a common integer
we may write $Q=aP$, $a\in F^\times$. The bracket image is
the one-dimensional line with proportional central coordinates
\begin{equation}\label{eq:cp-positive-centre}
 (1,c),\qquad c=\frac{a a^\dagger}{d}>0,
\end{equation}
where $\dagger$ is the Rosati involution. On an elliptic curve
$aa^\dagger$ is a positive rational number: a square in the
non-CM case and the positive field norm in the CM case.
Construction 3.4(3) adds a further central direction exactly
when the image of the algebraic lift in the quotient torus
is nontorsion. This is the integer $e$ above.
The unipotent Lie algebra is the Heisenberg algebra, with one
additional central Tate line when $e=1$. Formula
\eqref{eq:cp-MTdim} proves every entry of the table. This
description computes the group even when the lift contributes
a Kummer factor; specifying only the span of $P,Q$ would
not determine that last factor.
\end{proof}

\begin{statement}[C.7.4; Proved conditionally: transcendence of the local norm]
Assume Grothendieck's period conjecture for $1$-motives over
$\overline\Q$. For every algebraic parameter $t\ne0,1$ at
which $P_t$ is nontorsion, $N(P_t)$ is transcendental.
More generally the conclusion holds for every nontorsion
algebraic point on an elliptic curve over $\overline\Q$.
The weaker named Relative Semi-abelian Conjecture 4.7 of
\citet{bertolin2022}, restricted to the groups and algebraic
lifts used in this proof, already suffices.
\end{statement}
\begin{proof}
Suppose $N(P)$ is algebraic. Statement C.7.2 supplies an
algebraic lift $R=(P,r)$ in $G_P(\C)^c$, and its conjugate
supplies a lift $\bar R$ in $G_{\bar P}(\C)^c$.
Let $v$ denote the logarithmic coordinate in
\eqref{eq:cp-v}. Put
\[
 K_0=\overline\Q\bigl(\text{entries of the full de Rham--Betti
                    period matrix of }G_P\times G_{\bar P}\bigr).
\]
Here ``full'' includes the integrals of differentials of the second
kind, hence the quasi-periods $\eta_j$, as well as the periods
$\omega_j$, the third-kind periods $\kappa_j$, and $2\pi i$.
It is the matrix $\Omega_G$ in their Definition 2.4, not merely
the lattice in the Lie algebra of the semi-abelian variety.
This field contains $u,\zeta(u)$ and their complex conjugates.
Indeed it contains the pure periods, $2\pi i$, and
$\kappa_j=\zeta(u)\omega_j-\eta_j u$; the determinant of
the two equations for $u,\zeta(u)$ is, up to sign,
$\omega_1\eta_2-\omega_2\eta_1= -2\pi i$.
This is also the observation of \citet[Remark 4.13]{bertolin2022}.
The real-linear expression for $\eta(u)$ is
\begin{equation}\label{eq:cp-real-linear}
 \eta(u)=
 \frac{(u\bar\omega_2-\bar u\omega_2)\eta_1
       +(\bar u\omega_1-u\bar\omega_1)\eta_2}
      {\omega_1\bar\omega_2-\bar\omega_1\omega_2}.
\end{equation}
The denominator is nonzero. Thus $u\eta(u)$ and its conjugate
belong to $K_0$. The compact equation $\chi(R)=0$ becomes
the algebraic relation over $K_0$
\begin{equation}\label{eq:cp-relation}
 v+\bar v
 =u\zeta(u)+\bar u\overline{\zeta(u)}
                -u\eta(u)-\bar u\overline{\eta(u)}.
\end{equation}
Complex conjugates here are coordinates of the conjugate
$1$-motive. They are not being treated as rational functions
of un-conjugated periods without justification.

Consider the period torsor of $\mathcal M$ over the period
torsor of $G=G_P\times G_{\bar P}$. The kernel of the
corresponding motivic groups is its weight $-2$ vector group
$\operatorname{Lie}Z$, of dimension $\dim Z$. There is no
weight $-1$ kernel: the point coordinates are exactly the
extension coordinates under the polarization, so $B_Q=0$.
More explicitly, projection of the smallest abelian variety $B$
to the variety generated by the extension coordinates is an
isogeny, since $B$ is the graph of the same point data in the
two factors. It therefore induces an isomorphism on the
weight $-1$ graded Lie algebras. The motive $[0\to G]$ has
only weights $-1$ and $-2$, and its unipotent Lie algebra has
no weight $-2$ graded piece. Thus the entire kernel is exactly
$W_{-2}\operatorname{Lie}G_{\rm mot}(\mathcal M)$; no additional
abelian direction or unaccounted logarithm remains in the fibre.
After evaluating the base periods, the fibre is an affine
torsor over $K_0$ under that vector group. Translation in
its two central coordinates changes $(v,\bar v)$ by
$2\pi i(c_1,c_2)$, with the signs fixed by using a positive
Tate loop for each motive. The linear functional
$c_1+c_2$ is nonzero on this central group. It is nonzero
on the whole plane when $\dim B=2$, and when $\dim B=1$
it is already nonzero on the bracket line
\eqref{eq:cp-positive-centre}, because $1+c>0$.

Consequently \eqref{eq:cp-relation} cuts a proper affine
hyperplane in the period-torsor fibre. It forces the relative
transcendence degree of the actual periods over $K_0$ to be
at most $\dim Z-1$. The pure-base period torsor has dimension
$\dim G_{\mathrm{mot}}(G)$, so the total transcendence degree
is at most $\dim G_{\mathrm{mot}}(\mathcal M)-1$.
This contradicts the asserted period conjecture. For real
data the single motive suffices: conjugation acts on its
Betti realization over $\Q$, and the same equation has a
nonzero coefficient on its one central direction. One can
also double the point and lift into the identity component,
choose $u,\omega_1$ real, and read the relation directly as
\[
 \omega_1 v-\omega_1u\zeta(u)+\eta_1u^2=0.
\]
Doubling preserves the equivalence of algebraicity of $N$
by the division formula and its algebraic nonzero factors.

For the weaker conjectural input, Conjecture 4.7 asserts
relative transcendence degree at least $2\dim B_Q+\dim Z$.
All $R_i$ used here are algebraic, $B_Q=0$, and the
generalized logarithms add exactly the central coordinates
to $K_0$. Thus it contradicts the same hyperplane relation
directly. Proposition 4.12 of that paper derives Conjecture
4.7 from the generalized period conjecture; in our case the
field of definition is algebraic, so the ordinary period
conjecture is enough. The weaker Conjecture 4.9, formulated
using arbitrary algebraic translates, gives a zero lower
bound at an algebraic point and does not by itself yield
this conclusion. Linear relations among $1$-periods are
also insufficient: after clearing the denominator in
\eqref{eq:cp-real-linear}, the prohibited relation is
polynomial, not a linear relation with algebraic coefficients.
\end{proof}

\begin{statement}[C.7.5; Proved computationally: the specified integer-relation tests]
At all eleven rational CM parameters, an independent theta
calculation and the global height calculation agree to more
than $245$ decimal places. At $200$ decimal digits, PSLQ
with tolerance $10^{-185}$, maximum coefficient $10^8$ and
$10000$ steps returns no relation in each prescribed list
consisting of $\log N$, logarithms of the bad primes together
with $\log2,\log3$, $\pi$, and independent real period
components. At $j=287496$ the known identity
$\omega_1=2|\Im\omega_2|$ is removed first; repeated prime
logarithms and zero period components are likewise omitted.
An additional direct algebraicity search tests both $N$ and $N^{12}$
at all eleven points, using the vectors $(1,x,\ldots,x^d)$ for
$d=4,8,12$, working precision $235$ decimal digits, tolerance
$10^{-210}$, maximum coefficient $10^{12}$ and $20000$ steps.
All $66$ calls return no relation.
\end{statement}
\begin{proof}[Computational verification]
The complete inputs, $255$-digit values, comparison errors,
PSLQ labels and returns are in
\path{code/C_cm_250.gp}, \path{code/C_cm_250.out},
\path{code/C_cm_theta_pslq.py}, and its output. The theta
calculation uses the exact CM period ratios and solves the
Fourier equation for $\wp(1/2-iy)=0$ independently of PARI's
height implementation. It verifies the singular modulus
as well. The additional polynomial search is reproduced by
\path{code/C7_norm_polynomial_search.py}; its output and JSON
record include every parameter, source hash and returned vector.
These tests prove no transcendence assertion. A null PSLQ return
is not an exhaustive certificate excluding all polynomials within
the stated coefficient and degree limits.
\end{proof}

\subsection{The unconditional problem remains under investigation}

Statements C.7.1--C.7.5 do not decide the unconditional
arithmetic nature of $N$ at a nontorsion point. An explicit
small instance is
\[
 E:y^2=x^3-140x+784,\qquad P=(0,28),\qquad j=-3375,
 \qquad N(P)=2e^{-\widehat h(P)}.
\]
Here $P$ is a primitive generator modulo the rational
$2$-torsion and
\[
 \widehat h(P)=
 0.46736567718147128000752020126284678363562034370331\ldots.
\]
The exact unresolved assertion is that $2e^{-\widehat h(P)}$
is transcendental. The corresponding example of conductor
$361$, $y^2+y=x^3-38x+90$, $P=(0,9)$, has
$N=e^{-\widehat h(P)}$ and no finite local contribution.
The relative-period proof isolates its obstruction: an
unconditional theorem excluding
\eqref{eq:cp-relation} for the positive diagonal central
direction. A theorem about linear $1$-period relations
does not supply that exclusion. Further attempts and the
CM theta comparison are recorded below as they are tested.
}{}
\IfFileExists{sections/C_cm_transcendence.tex}{\subsection{CM theta normalization and further unconditional deductions}

The holomorphic theta normalization in Lead 16(e) differs from the
norm occurring in the question. We make that distinction explicit,
including the finite-character ambiguity in the multiplication law.

\begin{statement}[C.7.6; Proved: the exact CM multiplication law]
Let $E/\overline{\Q}$ have complex multiplication by an order
$\mathcal O$, and use an algebraic invariant differential. Write
$A=\operatorname{area}(\Lambda)$ and
\[
 \eta(z)=s_2z+\frac{\pi}{A}\bar z,\qquad
 \theta(z)=\sigma(z)e^{-s_2z^2/2}.
\]
Then $s_2\in\overline{\Q}$. For every nonzero
$\alpha\in\mathcal O$, with $d=\alpha\bar\alpha$, there is
$F_\alpha\in\overline{\Q}(E)$ such that
\begin{equation}\label{eq:ct-CMmult}
 \theta(\alpha z)^2=F_\alpha(\exp_E z)\theta(z)^{2d}.
\end{equation}
Its divisor and leading Laurent coefficient are
\[
 \operatorname{div}(F_\alpha)=2[\alpha]^*(O)-2d(O),\qquad
 F_\alpha(\exp_Ez)=\alpha^2z^{2-2d}(1+O(z^2)).
\]
Thus the unsquared quotient is an algebraic function on $E$, but
it need not be a rational function on $E$. For a nontorsion
$P=\exp_Eu\in E(\overline{\Q})$, algebraicity of $\theta(u)$
would imply algebraicity of $\theta(\alpha u)$ for all
$\alpha\in\mathcal O$. The additive group $\mathcal O u$ has
rank two, not more than four.
\end{statement}
\begin{proof}
The Legendre relation gives the displayed real-linear decomposition
of $\eta$. The sigma automorphy formula gives
\[
 \theta(z+\omega)=\epsilon(\omega)
  \exp\left(\frac{\pi\bar\omega}{A}
                         (z+\omega/2)\right)\theta(z),
 \qquad \epsilon(\omega)\in\{1,-1\}.
\]
For $\alpha\Lambda\subset\Lambda$, the exponential multipliers
of $\theta(\alpha z)/\theta(z)^d$ cancel because
$\alpha\bar\alpha=d$. Its remaining multiplier is a sign.
Squaring therefore gives a meromorphic function on $E$, with the
stated divisor. This divisor is defined over $\overline{\Q}$.
The existence of a rational function with this divisor over $\C$
implies its existence over $\overline{\Q}$: the relevant
one-dimensional Riemann--Roch space is obtained by extension of
scalars. Normalize its leading Laurent coefficient, in the formal
parameter of the algebraic differential, to $\alpha^2$. This
identifies it with the squared analytic quotient. The action of
$\alpha$ on this differential is algebraic, so $\alpha^2$ is an
allowed normalization.

Choose a nonreal $\alpha\in\mathcal O$. Since the normalized
sigma series is $\sigma(z)=z+O(z^5)$ with algebraic Laurent
coefficients, the next coefficient of the squared quotient is
\[
 \alpha^2z^{2-2d}
       \{1-s_2(\alpha^2-d)z^2+O(z^4)\}.
\]
The coefficients of the rational function in this algebraic formal
parameter are algebraic. As $\alpha^2-d\ne0$, this proves
$s_2\in\overline{\Q}$ without an unnormalized appeal to a CM
weight-two form. It also proves that the functions in
\eqref{eq:ct-CMmult} are algebraic over the chosen field.
At nontorsion $u$, none of the arguments $\alpha u$ is a lattice
point; evaluating the identity proves the algebraicity implication.
The injective map $\mathcal O\to\C$, $\alpha\mapsto\alpha u$,
proves the rank assertion.

The qualification about rationality is necessary. For a square
lattice, $s_2=0$ and $\alpha=1+i$ give
$F_{1+i}=2i\wp$. The unsquared quotient has a simple zero at
the nonzero point in $\ker(1+i)$ and a simple pole at $O$,
which is not a principal divisor on $E$; its square is the
rational function just displayed.
\end{proof}

\begin{statement}[C.7.7; Known with derivation: two unconditional normalized values]
Let $E/\overline{\Q}$ have an algebraic invariant differential,
and let $u$ be an elliptic logarithm of an algebraic nontorsion
point $P$. Then
\[
 \sigma(u)\exp\left(-\frac12u\zeta(u)\right)
\]
is transcendental. If $E$ has CM and $s_2$ is normalized as in
C.7.6, then $\zeta(u)-s_2u$ is also transcendental. The first
assertion concerns the displayed holomorphic normalization;
it does not assert transcendence of either $N(P)$ or
$\theta(u)=\sigma(u)e^{-s_2u^2/2}$.
\end{statement}
\begin{proof}
We use \citet[Theorem 1 and Corollary 1]{waldschmidt1979sigma},
whose complete proof was opened in the translation hosted by its
author. Its Theorem 1 states that, for algebraic invariants
$g_2,g_3$, algebraic elliptic points $\exp_Eu$, $\exp_Ev$, and
$\beta\in\overline{\Q}$, the number
\[
 \frac{\sigma(u+v)}{\sigma(u)\sigma(v)}
                      \exp((\beta-\zeta(v))u)
\]
is transcendental if $\exp_Ev$ is nontorsion and $u,u+v$
are not lattice points. Here ``algebraic elliptic point''
means that the Weierstrass coordinates are algebraic; it does
not mean that its logarithm is algebraic. Take $v=u$ and
$\beta=0$. All hypotheses hold because $P$ is nontorsion.
The sigma duplication formula identifies the displayed number
with
\[
 -\wp'(u)\sigma(u)^2e^{-u\zeta(u)}.
\]
The factor $-\wp'(u)$ is nonzero and algebraic. Therefore
$\sigma(u)e^{-u\zeta(u)/2}$ is transcendental, as also stated
explicitly after Corollary 1 of the cited paper.

For the second assertion we apply the analytic subgroup theorem
in \citet[Theorem 6.1]{huber2022}. In the algebraic universal
vector extension $\widetilde E$ a logarithm of an algebraic lift
of $P$ has coordinates $(u,\zeta(u)-b)$, where $b$ is an algebraic
fibre coordinate, with a possible algebraic linear change of the
second coordinate. If $b=\zeta(u)-s_2u$ were algebraic, this
logarithm would belong to the algebraic line
$\{(z,s_2z):z\in\C\}$. The theorem applies to the connected
commutative algebraic group $\widetilde E$ over
$\overline{\Q}$ and this line, defined over $\overline{\Q}$ by
C.7.6; its exponential contains a nonzero algebraic lift.
It would provide a connected one-dimensional algebraic subgroup
whose Lie algebra is that line. Its projection to $E$ is
nontrivial. A one-dimensional affine algebraic group has no
nonconstant algebraic-group homomorphism to an elliptic curve,
so the subgroup is an elliptic curve isogenous to $E$.
Its inclusion would split the universal vector extension after
that isogeny. This is impossible: the nonzero extension class in
$\operatorname{Ext}^1(E,\mathbb G_a)=H^1(E,\mathcal O_E)$
remains nonzero under isogeny, since pullback is invertible in
characteristic zero.
\end{proof}

\paragraph{A theta criterion requiring its original normalization.}
An earlier draft applied the pointwise criterion printed in
\citet[final slide, comment (c)]{bertrand2017} to
$g(u,v)-s_2uv$, where
$g(u,v)=\log(\sigma(u+v)/(\sigma(u)\sigma(v)))$.
It inferred transcendence of $\theta(u)$ by taking $v=u$.
We withdraw that inference pending verification of the original
lift conventions in \citet[Corollary 3]{bertrand1984}.
Indeed replacing $v$ by $v+\omega$, $\omega\in\Lambda$,
changes $g(u,v)-s_2uv$, modulo $2\pi i\Z$, by
\[
 (\eta(\omega)-s_2\omega)u=\frac{\pi\bar\omega u}{A},
\]
whereas the condition $N\exp_Ev=f(\exp_Eu)$ with $f$
Rosati-antisymmetric depends only on the elliptic points.
Thus the two-sided statement cannot be used with unrestricted
lifts without checking its normalization. The author's slide
and the original paper's first two pages were accessible;
the full original Corollary 3 was not. We use neither direction
of that criterion here. The withdrawn text and the lift audit
are retained in the records.

A subsequent source check resolves the normalization of the
one-variable holomorphic assertion itself:
\citet[p.~15-03]{waldschmidt1988theta} explicitly states that
$\sigma(u)e^{-s_2u^2/2}$ is transcendental when the invariants
are algebraic, $E$ has CM, and $\exp_Eu$ is an algebraic
nontorsion point.  That source defines $s_2$ by the regularized
lattice sum and attributes the result to Bertrand's 1982
paper \emph{Valeurs de fonctions th\^eta et hauteurs p-adiques}.
We have checked this complete normalized statement in
Waldschmidt's published exposition; the full original 1982
proof has not been retrieved.  This source clarification
does not validate the unrestricted-lift inference withdrawn
above, and the theorem concerns the holomorphic value rather
than the metric norm.  C.7.7 retains the two assertions whose
proofs have been fully checked here.

\paragraph{The archimedean term still requiring control.}
The exact comparison is
\begin{equation}\label{eq:ct-Gaussian}
 N(P)=|\Delta|^{1/12}|\theta(u)|
                         \exp\left(-\frac{\pi|u|^2}{2A}\right).
\end{equation}
It follows from C.7.1 and the real-linear decomposition of
$\eta$. The normalization proved transcendental in C.7.7 instead
gives the exact identity
\[
 N(P)=|\Delta|^{1/12}
 \left|\sigma(u)e^{-u\zeta(u)/2}\right|
 \exp\left(\frac12\Re\{u(\zeta(u)-\eta(u))\}\right).
\]
Neither formula excludes algebraicity of $N(P)$: transcendence
of one complex factor does not control its absolute value or
its product with the remaining exponential.
The generic finite-order interpolation approach in Lead 16(e)
supplies only the rank-two set $\mathcal O u$; it does not
meet the proposed rank-greater-than-four criterion. The
derivative approach would need the algebraicity of
$\theta'(u)/\theta(u)=\zeta(u)-s_2u$, which C.7.7 excludes.
Replacing the norm by the holomorphic theta value would
therefore change the problem precisely by the factor in
\eqref{eq:ct-Gaussian}.

The related independence theorem of
\citet[Corollary~2]{reyssat1982weierstrass} says that two of
$e^{i\pi u/\omega_1}$, $s(u)$ and $s'(u)$ are algebraically
independent, where $s(z)=\sigma(z)e^{-\eta_1z^2/(2\omega_1)}$.
For real $u$, hypothetical algebraicity of $N(P)$ makes
$s(u)$ algebraic, so the conclusion instead forces independence
of $e^{i\pi u/\omega_1}$ and
$\zeta(u)-\eta_1u/\omega_1$; it yields no contradiction.
The periodic-functional-equation theorem reported in
\citet[p.~15-05]{waldschmidt1988theta} assumes algebraic
$q=e^{i\pi\tau}$, whereas this $q$ is transcendental for
imaginary quadratic $\tau$ by Gelfond--Schneider.  These
specific checked applications therefore do not decide the
metric norm; their source scopes and the larger CM
independence tuples are recorded separately.

We also tested an additive-extension route.
\citet[Theorem 3.1.2]{bost2009} proves that an algebraic point
of the universal vector extension of an abelian variety which
lies in its maximal compact subgroup at one complex embedding
is torsion. The theorem's hypotheses require an algebraic
point of the vector extension, equivalently an algebraic line
bundle with an algebraic connection having unitary monodromy.
C.7.2 supplies only an algebraic point in one fibre of a
multiplicative extension. It does not supply that algebraic
connection. Differentiating a metric value at one algebraic
point cannot establish algebraicity of the connection there
or globally. This is the exact missing implication in that
attempt.

Another route uses the Latt\`es map $f_E$ defined by
$x(2Q)=f_E(x(Q))$. Its canonical height is
$\widehat h_{f_E}(x(P))=2\widehat h_O(P)$. For rational
parameters C.6 therefore makes algebraicity of $N(P)$
equivalent to algebraicity of
$\exp(\widehat h_{f_E}(x(P)))$.
\citet[Theorem 1.4 and Corollary 1.6]{nguyen2023} settle the
corresponding multiplicative-height problem for nonintegrable
polynomials using B\"ottcher coordinates in the basin of
infinity. They do not apply to $f_E$. Indeed an elliptic
duplication Latt\`es map has no exceptional point: the inverse
images of any point contain the images of all successive
division points of an elliptic point, an infinite set. A
polynomial has the exceptional point $\infty$, so no conjugate
of $f_E$ is polynomial. The Julia set is the whole sphere
and the required basin of infinity is absent. This checks
the failed hypothesis, rather than treating a theorem for
polynomials as a theorem for all rational maps.

We also checked the arithmetic algebraicity criterion of
\citet[Theorems 1.3 and 2.3]{xie2024algebraicity}.
Besides sufficiently rapid height-normalized approximation by
points in a fixed number field, it requires the analytic germ
to be formally defined over that field. This hypothesis fails
for the compact leaf in C.7.9: its tangent has the proved
transcendental slope
$c=\zeta(u)-\eta_1u/\omega_1$.
A smooth formal germ defined over $\overline{\Q}$ has an
algebraic tangent; algebraic translations, coordinate changes
and isogenies preserve this obstruction. Assuming $N(P)$
algebraic supplies an algebraic point on the leaf and its
integer multiples, but not an algebraic formal germ. Attempting
to force convergence by successive division also loses the
fixed-number-field hypothesis: if $[2]^nQ_n=Q$ with all
$Q_n\in E(K)$ for one number field, then
$\widehat h_O(Q_n)=4^{-n}\widehat h_O(Q)>0$.
The $Q_n$ are distinct with bounded ample height, contrary to
Northcott finiteness. The full hypothesis audit and the opened
preprint are retained in the records.

\paragraph{A current source whose theorem text remains unchecked.}
The publisher's abstract for \citet{pham2025sigma} announces
complex and $p$-adic linear-independence results for semiabelian
varieties and applications to Weierstrass sigma functions.
We opened its metadata, abstract and reference list, but both
the full-HTML link and the publisher's citation-PDF link returned
the abstract/access page. No accessible author manuscript was
located. Consequently we could not check any theorem statement
or its hypotheses. This is a potentially relevant unverified
source: we make no claim either that it settles the Gaussian
norm question or that its results are inapplicable to it.

}{}
\IfFileExists{sections/C_padic.tex}{\subsection{The precise nonarchimedean comparison}
\begin{statement}[C.7.8; Known with derivation: a nonzero ordinary height]
Let $E/\Q$ have complex multiplication by the full ring of integers
of an imaginary quadratic field $K$ of class number one. Assume
that its functional-equation sign is $-1$, and let $p>3$ be a
prime of good ordinary reduction, $p\mathcal O_K=\mathfrak p
\mathfrak p^*$. For a nontorsion $P\in E(\Q)$ the
$\mathfrak p$-adic height and the cyclotomic $p$-adic height
of $P$ over $K$, in the normalization of
\citet[Appendix A]{agboola2006}, are nonzero.
In particular this applies to
\[
 E:y^2+y=x^3-38x+90,\qquad P=(0,9),\qquad p=5.
\]
\end{statement}
\begin{proof}
We use the exact theorem reproduced as Theorem A.1 in Rubin's
appendix to \citet{agboola2006}, attributed there to
\citet{bertrand1984}: for a number field $F$ and a point of
infinite order in $E(F)$, its height $h_{F,\mathfrak p}$ is
nonzero. We retain the standing hypotheses of that appendix,
including the sign assumption, rather than infer a wider theorem
from the abbreviated statement. The height takes values in
the one-dimensional $\Q_p$-space attached to the unique
$\Z_p$-extension unramified outside $\mathfrak p$.
Taking $F=K$ proves its nonvanishing. The polarized height of
a rational point has zero anticyclotomic projection: complex
conjugation fixes the point and acts as $-1$ on that projection.
Lemma A.2 of the same appendix states that if two of the
cyclotomic, anticyclotomic and $\mathfrak p$ projections
vanish, then the third vanishes. It follows that the cyclotomic
projection is nonzero as well. This argument asserts nonvanishing
on this point, not nondegeneracy of a height pairing on an
arbitrary Mordell--Weil space tensored with $\Q_p$.

For the displayed curve, C.6.R1--R2 prove nontorsion and
analytic rank one, hence sign $-1$. Its CM order has
discriminant $-19$, which is maximal, and its class number is
one: the reduced positive binary forms of that discriminant
give only $(1,1,5)$ up to the usual boundary convention.
The discriminant of the minimal curve is $-19^3$, so reduction
at $5$ is good. Counting directly over $\F_5$ gives seven
points, hence $a_5=-1$ and ordinary reduction; equivalently
$-19$ is a nonzero square modulo $5$. Every hypothesis of the
cited theorem therefore holds in this concrete unresolved
archimedean example.
\end{proof}

The comparison explains the missing step in an attempted
archimedean adaptation. The ordinary $p$-adic sigma construction
uses a holomorphic quadratic correction. In the CM normalization
the corresponding complex holomorphic function is $\theta$
of C.7.6, whereas the real metric also contains
$\exp(-\pi|u|^2/(2A))$. Theorem A.1 supplies no identity
between the ordinary $p$-adic height and that real factor.
Transferring the conclusion would require precisely such a
new argument; replacing the real factor by the holomorphic
CM correction does not preserve the norm being asked about.

There is also a definite limitation on the maximal-compact
approach. \citet[Theorem 2.1]{bertrand1995} constructs
nontrivial multiplicative extensions of polarized abelian
varieties with nontorsion algebraic points in their maximal
real compact subgroup. Its hypotheses include an endomorphism
not fixed by the Rosati involution and a nontorsion rational
point over a real number field. The construction uses a base
point related to the extension class by a Rosati-antisymmetric
endomorphism. It cannot be replaced by a theorem saying that
every algebraic compact point of every multiplicative
extension is torsion. For our diagonal class, an equality
$NP=f(P)$ with $N>0$ and $f+f^\dagger=0$ would force
$P$ torsion: $N-f$ has no zero eigenvalue, since
$(N-f)^\dagger(N-f)=N^2-f^2$ is positive. In the elliptic
case this is simply a nonzero isogeny. The diagonal positivity
used in C.7.3--C.7.4 is therefore essential, but these
counterexamples neither prove nor refute the desired
unconditional diagonal statement. The additive result in
\citet[Theorem 3.1]{bertrand1995}, and its higher-dimensional
extension in \citet{bost2009}, requires a point in a vector
extension, a hypothesis not furnished by C.7.2.
}{}
\subsection{The fixed semiabelian extension and the diagonal subgroup attempt}

\begin{statement}[C.7.9; Proved: the endomorphism and slope obstructions]
Let $E/\overline\Q$ be an elliptic curve, let
$q\in E^\vee(\overline\Q)$ be nontorsion, and let
\[
 1\longrightarrow\mathbb G_m\longrightarrow G_q
       \longrightarrow E\longrightarrow0
\]
be the extension of class $q$. Then
$\operatorname{End}_{\overline\Q}(G_q)=\Z$, including when
$E$ has complex multiplication. In particular a noninteger
CM endomorphism of $E$ does not lift to an endomorphism of this
fixed extension.

For
\[
 E:y^2+y=x^3-38x+90,\qquad P=(0,9),\qquad Q=2P=(4,-2),
\]
one has
\[
 N(Q)=19^{1/4}N(P)^4.
\]
Use the differential $dx/(2y+1)$, choose the positive real
period $\omega_1$, and choose the real logarithm
$u\in(0,\omega_1/2)$ of $Q$. Put
\begin{equation}\label{eq:cd-slope}
 c=\zeta(u)-\frac{\eta_1}{\omega_1}u,
 \qquad \eta_1=\eta(\omega_1).
\end{equation}
Then $c$ is a real transcendental number. If $N(P)$ were
algebraic, the extension $G_Q$ of C.7.1 would have a
nontorsion algebraic point whose logarithm is $(u,cu)$.
The complex line
$\ell_c=\{(z,cz):z\in\C\}$ is not defined over
$\overline\Q$ in the algebraic logarithm coordinates of
C.7.1. Thus the analytic subgroup theorem cannot be applied
to that line using the hypothesized compact point.
\end{statement}
\begin{proof}
The torus is the maximal connected affine subgroup of $G_q$.
Consequently every endomorphism of $G_q$ induces multiplication
by some $n\in\Z$ on $\mathbb G_m$ and an endomorphism
$\alpha$ on $E$. Functoriality of extension classes gives
\begin{equation}\label{eq:cd-compatibility}
 \alpha^*q=nq\quad\text{in }E^\vee.
\end{equation}
Here $\alpha^*$ is the dual endomorphism, not $\alpha$
acting on an arbitrarily identified copy of $E$. To make the
convention explicit, choose a principal polarization
$\lambda:E\to E^\vee$, put $Q_0=\lambda^{-1}q$, and
write $\alpha^\dagger=\lambda^{-1}\alpha^*\lambda$.
Equation \eqref{eq:cd-compatibility} says
$(\alpha^\dagger-[n])Q_0=O$. A nonzero endomorphism of an
elliptic curve is an isogeny and has finite kernel. Since
$Q_0$ is nontorsion, $\alpha^\dagger=[n]$, hence
$\alpha=[n]$ as well.

The difference between the given endomorphism and $[n]$ has
image in the torus and is trivial on the torus. It therefore
factors through an algebraic-group homomorphism
$E\to\mathbb G_m$. This homomorphism is zero, since a
morphism from the proper connected curve $E$ to the affine
group $\mathbb G_m$ is constant and a group homomorphism
preserves the identity. Thus there is no additional ambiguity
from $\operatorname{Hom}(E,\mathbb G_m)$, and the
endomorphism is exactly $[n]$. Conversely all integer
multiplications exist. This proves the first assertion.

For the concrete curve, the tangent slope at $P$ is $-2$,
giving $2P=(4,-2)$. Its discriminant is $-19^3$ and
$\psi_2(P)=2\cdot9+1=19$. The division identity in C.4
therefore gives the asserted formula for $N(Q)$, so replacing
$P$ by $Q$ preserves algebraicity of the norm. This doubling
also places the base point in the identity component of the
real elliptic group without any component assumption. In the
chosen differential the Weierstrass coordinates are
$(\wp,\wp')=(x,2y+1)$, with $g_2=152$, $g_3=-361$.
The negative discriminant means that the real cubic has one
real root. Along $(0,\omega_1/2)$, $\wp$ decreases from
$+\infty$ to that root and $\wp'<0$. Since
$\wp'(u)=2(-2)+1=-3$, the indicated real logarithm of $Q$
exists.

The period vectors of $G_Q$ in C.7.1 include
\[
 (\omega_1,\zeta(u)\omega_1-\eta_1u)
       =\omega_1(1,c).
\]
These are algebraic logarithm coordinates: $dz$ and the
invariant differential $d\log r-df_u/f_u$ are algebraic,
because $df_u/f_u$ is the algebraic third-kind differential
displayed in C.7.1. Real invariants and real $u,\omega_1$
make $c$ real.

Suppose $c$ were algebraic. Then $\ell_c$ would be an
algebraic Lie subspace. Dividing the nonzero period
$\omega_1(1,c)$ by a sufficiently large positive integer
gives a vector in $\ell_c$ whose exponential is a nonzero
torsion point of $G_Q$, hence an algebraic point. The
analytic subgroup theorem
\citep[Theorem 6.1]{huber2022} supplies a positive-dimensional
connected algebraic subgroup with Lie algebra contained in
$\ell_c$. Its dimension is one, and its projection to $E$
is nontrivial. Such a subgroup must be an elliptic curve
$E'$ projecting isogenously to $E$: a connected
one-dimensional affine algebraic group cannot map
nontrivially to an elliptic curve. Its inclusion splits the
pullback of $G_Q$ to $E'$. The dual of that isogeny would
therefore kill the extension class $\lambda(Q)$, forcing
$Q$ torsion. This contradicts C.6.R1. We conclude that $c$
is transcendental. In particular the line is not defined over
$\overline\Q$. This argument uses the period vector alone;
it does not assume algebraicity of $N(P)$.

Finally assume, only for the attempted contradiction, that
$N(P)$ is algebraic. The positive real number
\[
 r=3\,19^{-1/2}N(Q)^2
\]
is algebraic, and C.7.2 places the point $R=(Q,r)$ of
$G_Q$ in its maximal compact subgroup. The chosen real
logarithm has
\[
 f_u(u)=-\wp'(u)\sigma(u)^2e^{-u\zeta(u)}
        =3\sigma(u)^2e^{-u\zeta(u)}>0.
\]
Consequently $v=\log(r/f_u(u))$ is real and
$(u,v)$ is a logarithm of $R$, with no imaginary logarithm
ambiguity. Substituting this real logarithm into the exact
compactness equation $\chi=0$ in C.7.2 gives
\[
 v=\zeta(u)u-u\eta(u)
   =u\left(\zeta(u)-\frac{\eta_1}{\omega_1}u\right)=cu.
\]
Its projection $Q$ is nontorsion, so $R$ is nontorsion.
The input missing from this proposed use of the analytic
subgroup theorem is a proper Lie subspace defined over
$\overline\Q$; the actual compact slope line is proved
not to have that property. The CM decomposition makes the
same obstruction visible as
\[
 c=\zeta(u)-s_2u-\frac{\pi}{A}u.
\]
The last term is the derivative of the real Gaussian
correction. It cannot be discarded while retaining the
compactness equation.
\end{proof}

\paragraph{What functional transcendence gives in this construction.}
The functional theorem can be applied fully, but its conclusion
does not contradict this hypothetical algebraic compact point.
Indeed consider
\[
 f:\C\longrightarrow\operatorname{Lie}(G_Q),\qquad
 f(z)=(z,cz).
\]
No proper connected algebraic subgroup of $G_Q$ projects
onto $E$: as in the proof above, such a subgroup would split
the nontorsion extension after an isogeny. Every proper
connected subgroup is therefore contained in the torus,
whereas $\exp_{G_Q}f$ has nonconstant elliptic projection.
Thus $f$ is $G_Q$-nondegenerate in Definition 1.2 of
\citet{noguchi2024}. Their Theorem 1.3 gives
\[
 \dim_{\C}\overline{\{(\exp_{G_Q}f(z),f(z)):z\in\C\}}^{\rm Zar}
       \ge \dim G_Q+1=3.
\]
The graph is contained in $G_Q\times\ell_c$, of dimension
three, so equality holds. This is a statement over $\C$
about a family of meromorphic functions. It remains valid
whether or not the single value $\exp_{G_Q}f(u)$ is
algebraic. Under the norm-algebraicity hypothesis the integer
values $\exp_{G_Q}f(nu)=[n]R$ are algebraic, but they do
not give an algebraic relation holding on the entire graph.
There is no contradiction with its full dimension.
One would need an additional arithmetic specialization
argument, or a way to obtain a further functional relation
from that single metric value. Neither is supplied by the
checked functional theorem. The endomorphism calculation
above also prevents obtaining extra CM multiplication
directions inside this fixed $G_Q$ by lifting the
noninteger endomorphisms of $E$.

\subsection{A CM-stable enlargement and its Gaussian eigencoordinate}

The fixed-extension calculation in C.7.9 does not prevent
enlarging the torus. We test that alternative explicitly.
Retain its curve, its point $Q=2P$, and its real logarithm
$u$. Write $\beta=\sqrt{-19}=i\sqrt{19}$, $d=19$,
$D=\zeta(u)-s_2u$, and $k=\pi/A$. All endomorphisms and
extensions in what follows are defined over $\overline\Q$.

\begin{statement}[C.7.10; Proved: the enlarged CM action does not isolate the Gaussian]
Let
\[
 \mathcal G=G_Q\mathbin{\times_E}G_{\beta Q},\qquad
 1\longrightarrow\mathbb G_m^2\longrightarrow\mathcal G
       \longrightarrow E\longrightarrow0.
\]
There is an endomorphism $B$ of $\mathcal G$ inducing
$\beta$ on $E$ and
$(t_1,t_2)\mapsto(t_2^{-1},t_1^d)$ on the torus.
It satisfies $B^2=[-d]$. Moreover
\[
 \operatorname{Hom}(\mathcal G,\mathbb G_m)=0,
 \qquad
 \operatorname{End}(\mathcal G)\otimes\Q=\Q(\beta).
\]
Thus $\mathcal G$ has no algebraic $\mathbb G_m$ quotient,
including after isogeny.

There is an algebraic compact point in the second extension
above $Q$ whose logarithm is
\[
 (u,v_2),\qquad v_2=-\beta uD+au,
 \qquad
 a=\frac1{2\beta}
       \frac{d}{du}\log F_\beta(Q)\in\overline\Q.
\]
Here the derivative means differentiation of the rational
function $F_\beta(\exp_E z)$ with respect to the algebraic
invariant differential, evaluated at $z=u$.
If $N(P)$ is algebraic, C.7.9 combines this point with an
algebraic compact point in the first extension to give
\[
 U=(u,v_1,v_2),\qquad v_1=uD-ku^2,
 \qquad \exp_{\mathcal G}U\in\mathcal G(\overline\Q).
\]
In these algebraic logarithm coordinates the differential of
$B$ and its $-\beta$ eigenprojection are
\[
 dB=\begin{pmatrix}
 \beta&0&0\\ a&0&-1\\ \beta a&d&0
 \end{pmatrix},\qquad
 \Pi_-:=\frac12\left(I-\frac{dB}{\beta}\right),
\]
\[
 \Pi_-U=-\frac{\pi u^2}{2A}(0,1,\beta).
\]
The first torus coordinate of $\exp(\Pi_-U)$ is exactly
the Gaussian $\exp(-\pi u^2/(2A))$. The projector is not
the differential of an algebraic quasi-endomorphism. No
sequence of rational polynomials in $B$, even allowing
inverses of isogenies, can converge to it as complex linear
operators. Algebraicity of $\exp(U)$ does not supply
algebraicity of $\exp(\Pi_-U)$.
\end{statement}
\begin{proof}
Use the principal polarization to identify the two extension
classes with $Q$ and $\beta Q$. They are linearly independent
over $\Q$ modulo torsion: any nonzero $m+n\beta$ is an
isogeny, so $(m+n\beta)Q$ cannot be torsion. A character
of $\mathbb G_m^2$ given by $(m,n)\in\Z^2$ extends to
$\mathcal G$ exactly when the pushout extension has zero
class, namely when $mQ+n\beta Q=O$. Hence the restriction
of every homomorphism $\mathcal G\to\mathbb G_m$ to
the torus is zero. Such a homomorphism then factors through
$E$ and is zero. This proves the asserted absence of
characters. An isogenous semiabelian variety cannot acquire
a nonzero character, because composing it with the isogeny
would give a nonzero character of $\mathcal G$.

The Rosati adjoint of $\beta$ is $-\beta$. Pullback by
$\beta$ sends the ordered extension classes to
$(-\beta Q,dQ)$, exactly as does pushout by the displayed
torus matrix. Equality of these extension classes provides
the endomorphism $B$. It is unique because
$\operatorname{Hom}(E,\mathbb G_m^2)=0$. Its square
induces $[-d]$ on the torus and on $E$, so the same
uniqueness gives $B^2=[-d]$. The map
\[
 \operatorname{End}(\mathcal G)\otimes\Q
     \longrightarrow\operatorname{End}(E)\otimes\Q
          =\Q(\beta)
\]
is injective: an endomorphism inducing zero on $E$ has image
in the torus, and the absence of characters makes it zero.
Its image contains $\beta$, hence equals $\Q(\beta)$.

Next construct the compact point without assuming anything
about $N(P)$. The identities in C.7.6, applied at $\beta$
and $1+\beta$, give the algebraic nonzero number
\begin{equation}\label{eq:ce-ribet-r}
 r_2=\frac{\theta((1+\beta)u)}{\theta(u)\theta(\beta u)},
 \qquad r_2^2=\frac{F_{1+\beta}(Q)}{F_\beta(Q)}.
\end{equation}
Indeed $(1+\beta)(1+\bar\beta)=1+d$, so the powers
of $\theta(u)$ cancel. None of the arguments is a lattice
point since $Q$ is nontorsion. Written in sigma functions,
\[
 r_2=\frac{\sigma((1+\beta)u)}{\sigma(u)\sigma(\beta u)}
                         e^{-s_2\beta u^2}.
\]
The point $(Q,r_2)$ in the rational trivialization of
$G_{\beta Q}$ therefore has logarithm
\[
 (u,v_2),\qquad
 v_2=u\{\zeta(\beta u)-s_2\beta u\}.
\]
Differentiating C.7.6 shows that
\begin{equation}\label{eq:ce-zeta}
 \zeta(\beta u)-s_2\beta u=-\beta D+a,
 \qquad a\in\overline\Q,
\end{equation}
with the value of $a$ in the statement. The compactness
character of this extension vanishes, since
\[
 \Re\{v_2-\zeta(\beta u)u+\beta u\eta(u)\}
    =\Re\left(\frac{\beta\pi|u|^2}{A}\right)=0.
\]
This also gives a direct, fully normalized construction of
the antisymmetric compact lift; no unrestricted theta
criterion is being used.

For completeness, the matrix $dB$ can be recovered without
choosing a transcendental change of coordinates. Its torus
block is the stated integer matrix and its first entry is
$\beta$, so initially write its first column as
$(\beta,a_1,a_2)^t$, with $a_1,a_2$ algebraic.
The identity $B^2=[-d]$ gives $a_2=\beta a_1$.
For $\omega\in\Lambda$, set
\[
 \kappa_1(\omega)=\zeta(u)\omega-\eta(\omega)u,
 \quad
 \kappa_2(\omega)=\zeta(\beta u)\omega-\eta(\omega)\beta u.
\]
Equation \eqref{eq:ce-zeta} and the real-linear formula for
$\eta$ imply
\[
 \kappa_1(\beta\omega)=a\omega-\kappa_2(\omega).
\]
Comparing the second coordinates of the corresponding period
vectors under $B$ gives
$(a_1-a)\omega_j\in2\pi i\Z$ for $j=1,2$.
Since $\omega_2/\omega_1$ is nonreal, both integers must
be zero, so $a_1=a$. This proves the matrix formula.

Under the hypothetical algebraicity of $N(P)$, C.7.9 gives
the first compact logarithm $v_1=uD-ku^2$. Pairing the
two algebraic lifts over the same base point gives $U$.
Substitution into the displayed matrix yields
\[
 v_2+\beta v_1-au=-\beta ku^2,
 \qquad
 \Pi_-U=-\frac{ku^2}{2}(0,1,\beta),
\]
which proves the exact Gaussian identity.

The two eigenspaces of $dB$ have eigenvalues $\beta$ and
$-\beta$. Every algebraic quasi-endomorphism has differential
$a_0I+b_0dB$ for rational $a_0,b_0$, with eigenvalues
$a_0+b_0\beta$ and $a_0-b_0\beta$. The projector
$\Pi_-$ would require these values to be $0$ and $1$,
which is impossible. There is also a uniform obstruction to
approximation. For arbitrary real $a_0,b_0$, put
$x=a_0+b_0\beta$, so that its other eigenvalue is
$\bar x$. Then
\begin{equation}\label{eq:ce-gap}
 \max\{|x|,|\bar x-1|\}\ge\frac12,
\end{equation}
because $1\le|\bar x|+|\bar x-1|=|x|+|\bar x-1|$.
These are exactly the two eigenvalue errors in approximating
$\Pi_-$. The spectral radius is bounded by every induced
operator norm, so no such sequence can converge to the
projector. Compositions and inverses of isogenies still belong
to the same endomorphism algebra and satisfy this bound.

Finally the image of $\Pi_-$ is the toric Lie line
$\C(0,1,\beta)$. It is not the Lie algebra of an algebraic
subtorus: an algebraic one-dimensional subtorus of
$\mathbb G_m^2$ has rational slope. Since $u\ne0$,
the projected vector is nonzero. The analytic subgroup
theorem in fact shows that its exponential cannot be an
algebraic point of both torus coordinates. If the projection
preserved algebraic exponentials here, that would give a
contradiction; it is precisely the extra implication that
is unavailable. The construction isolates the Gaussian in
a Lie eigencoordinate, without producing an algebraic
group map that evaluates it.
\end{proof}

The spectral comparison with the ordinary prime $5$ is
explicit. The CM number $\alpha=(1+\beta)/2$ satisfies
$\alpha^2-\alpha+5=0$. Its two archimedean conjugates
both have absolute value $\sqrt5$. At the two places over
$5$, their normalized valuations are $0$ and $1$, because
the polynomial reduces to $X(X-1)$ and has product $5$.
The separation available in that nonarchimedean linear
algebra is absent in \eqref{eq:ce-gap}. This comparison
does not turn a limiting operation on complex Lie spaces
into an algebraic-group endomorphism or transfer a
$p$-adic height conclusion to the real Gaussian.

\paragraph{An exact compact witness on the conductor-$361$ curve.}
In the preceding normalization the constants and the second
fibre coordinate are explicitly
\begin{equation}\label{eq:ce-explicit-witness}
 s_2=-2,\qquad a=-\frac{21\beta}{703},\qquad
 r_2=-\frac{165}{74}+\frac{1839\beta}{1406}.
\end{equation}
We give polynomial certificates for these values, so they
do not depend on recognition from numerical approximations.
Put $f=4x^3-152x+361$ and
\[
 \begin{aligned}
 h={}&x^9-38x^8+437x^7-1444x^6-7942x^5\\
     &+82308x^4-274360x^3+390963x^2-130321x-130321,\\
 N={}&(-19x+76)h^2+f(h''h-(h')^2)+(6x^2-76)h'h,\\
 K={}&N'h-2Nh'.
 \end{aligned}
\]
Direct expansion gives
\[
 -K^2f=4N^3-152\cdot19^2Nh^4+361\cdot19^3h^6,
 \qquad \gcd(N,h)=1.
\]
Hence, in coordinates $Y=2y+1$, the map
\[
 (x,Y)\longmapsto
 \left(\frac{N}{19h^2},\frac{KY}{19\beta h^3}\right)
\]
is a degree-$19$ endomorphism fixing $O$, with pullback of
$dx/Y$ equal to $\beta\,dx/Y$. It is therefore the CM
endomorphism $\beta$. Divisors and the leading coefficient
at $O$ give $F_\beta=\beta^2h^2$. Its Laurent coefficient
of relative degree two is $-76$. Comparison with C.7.6
gives $38s_2=-76$, proving $s_2=-2$.

For the other kernel put $h_+=(H_0+\beta H_1)/10$, where
\[
 \begin{aligned}
 H_0={}&10x^8-380x^7+4864x^6-28880x^5+72561x^4\\
       &+54872x^3-795644x^2+1824494x-1433531,\\
 H_1={}&20x^7-304x^6+760x^5+10469x^4-83752x^3\\
       &+246924x^2-315514x+130321.
 \end{aligned}
\]
Writing $\bar h_+$ for replacement of $\beta$ by $-\beta$,
the finite polynomial identities and coprimality checks are
\[
 \begin{gathered}
 19xh^2-N=5f h_+\bar h_+,\qquad
 h_+\mid K+19\beta h^3,\\
 \gcd(h_+,fh)=\gcd(h_+,h_+')=1.
 \end{gathered}
\]
At a root of $h_+$ the displayed CM map sends $(x,Y)$ to
$(x,-Y)$, so these are precisely the eight pairs of
non-$2$-torsion points in $\ker(1+\beta)$. The other three
nonzero kernel points are $E[2]\setminus\{O\}$, since
$1+\beta=2\alpha$ in the CM order. Its norm is $20$,
which verifies the count. Thus
\[
 F_{1+\beta}=\frac{(1+\beta)^2}{4}f h_+^2.
\]
Taking the square ratio in \eqref{eq:ce-ribet-r} and fixing
the sign by its expansion at $u=0$ gives the exact rational
function
\[
 r_2=-\frac{1+\beta}{2\beta}\frac{\wp'(u)h_+(\wp(u))}
                                             {h(\wp(u))}.
\]
At $Q$, the needed evaluations are
\[
 h(4)=-37,\qquad h'(4)=7,\qquad
 h_+(4)=\frac{829-279\beta}{10},\qquad\wp'(u)=-3.
\]
They give exactly \eqref{eq:ce-explicit-witness}. In
particular $(Q,r_2)$ is an explicit nontorsion algebraic
point in the maximal compact subgroup of $G_{\beta Q}$.
Its torus-coordinate norm is
$r_2\bar r_2=974799/26011=4-\wp(\beta u)$.

As an independent numerical check, evaluation at $110$
decimal digits gives
\[
 \begin{aligned}
 c&=1.14793408195385447403276278660663356664777041472753\ldots,\\
 e^{-\pi u^2/(2A)}
  &=0.68133647525194864004469591167703191986051023921164\ldots.
 \end{aligned}
\]
The sigma quotient and the zeta expression for $r_2$ and $a$
agree with their exact values to more than $110$ digits.
The exact polynomial identities, evaluations and numerical
outputs are retained in \path{code/C7_cm_explicit.gp} and
its output. The polynomial construction proves the values;
the numerical comparison only checks the analytic branch.

\subsection{Complexifying the Gaussian: values, jets, and the rank criterion}

The antiholomorphic Gaussian can be made holomorphic by introducing a
second variable. This gives more algebraic values under the hypothesis
in C.7(c), but the derivative and growth hypotheses of the resulting
transcendence arguments must still be checked. We do this on the
conductor-$361$ model, where all constants and conjugations are fixed.

\begin{statement}[C.7.11; Proved: the complexified route and its exact obstructions]
Use the curve, point and real logarithm of C.7.9:
\[
 E:y^2+y=x^3-38x+90,\qquad Q=(4,-2)=\exp_Eu,
 \qquad 0<u<\omega_1/2.
\]
Put $\beta=i\sqrt{19}$, $k=\pi/A$, $s_2=-2$, and
\[
 \theta(z)=\sigma(z)e^{-s_2z^2/2},\qquad
 D(z)=\zeta(z)-s_2z,\qquad
 \mathcal F(z,w)=\sqrt{19}\,\theta(z)\theta(w)e^{-kzw},
 \qquad H=\mathcal F^2.
\]
Then $\mathcal F(z,\bar z)=N(z)^2$, and $\mathcal F$ is an
entire function of two complex variables of order at most two.
It is periodic under $(z,w)\mapsto(z+\omega,w+\bar\omega)$,
$\omega\in\Lambda$. The subgroup
\[
 \Gamma=\{(\alpha u+\omega,\bar\alpha u+\bar\omega):
                   \alpha\in\Z[\beta],\ \omega\in\Lambda\}
 \subset\C^2
\]
has rank four over $\Z$ and lies in the real plane
$\{(z,\bar z):z\in\C\}$. If $N(Q)$ is algebraic, then the
values of $H$ on $\Gamma$ lie in one number field.

For every nonzero $\alpha\in\Z[\beta]$ there is an algebraic
number $a_\alpha$ such that
\[
 D(\alpha u)=\bar\alpha D(u)+a_\alpha.
\]
At $p_\alpha=(\alpha u,\bar\alpha u)$ the algebraic constant
direction $V_\alpha=\alpha\partial_z-\bar\alpha\partial_w$
satisfies
\begin{align*}
 V_\alpha\log\mathcal F(p_\alpha)
    &=\alpha a_\alpha-\bar\alpha\bar a_\alpha
                         \in\overline{\Q},\\
 V_\alpha^2\log\mathcal F(p_\alpha)
    &=-\alpha^2\{\wp(\alpha u)+s_2\}
      -\bar\alpha^2\{\wp(\bar\alpha u)+s_2\}
      +2\alpha\bar\alpha k.
\end{align*}
The number $k$ is transcendental. Consequently the second
directional derivative of $H$ at $p_\alpha$ is transcendental
whenever $H(p_\alpha)$ is algebraic. In particular,
\[
 (\partial_z-\partial_w)H(u,u)=0,\qquad
 \frac{(\partial_z-\partial_w)^2H(u,u)}{H(u,u)}=4k-8.
\]
Thus the cancelling first derivative does not give algebraic jets
of all orders. Moreover the rank-four set does not satisfy the
explicit derivative-free finite-order criterion audited below.
These are obstructions to the stated applications of those
criteria; they do not decide the algebraicity of $N(Q)$.
\end{statement}

\begin{proof}
The real model and its algebraic differential give real sigma
coefficients and a conjugation-stable lattice. In particular
$\overline{\theta(z)}=\theta(\bar z)$. The discriminant is
$-19^3$, so C.7.1 gives the asserted diagonal identity.
The sigma function and the exponential of the quadratic
polynomial $-kzw$ have order at most two. From C.7.6,
\[
 \theta(z+\omega)=\epsilon(\omega)
         e^{k\bar\omega(z+\omega/2)}\theta(z),
 \qquad \epsilon(\bar\omega)=\epsilon(\omega)\in\{1,-1\}.
\]
Multiplying this formula and its conjugate version, the two
signs multiply to one and the exponent cancels the change of
$-kzw$. This proves periodicity on all of $\C^2$, not only
on the real diagonal.

If $\alpha u+\omega=0$, then $[\alpha]Q=O$.
Nontorsion of $Q$ and the fact that a nonzero elliptic
endomorphism is an isogeny imply $\alpha=0$, and then
$\omega=0$. Thus the rank-two module $\Z[\beta]u$ and the
rank-two period lattice are independent over $\Z$.
This proves the rank assertion, with no extra algebraic
independence assumption on $u$ and the periods.

For $d=\alpha\bar\alpha$, the squared CM law C.7.6 gives
the exact fixed-field identity
\begin{equation}\label{eq:cc-fixedfield}
 H(\alpha u,\bar\alpha u)
   =19^{1-d}F_\alpha(Q)\overline{F_\alpha(Q)}H(u,u)^d.
\end{equation}
All the endomorphisms and their normalized rational functions
$F_\alpha$ are defined over a single number field: choose
a field containing the coefficients of the finitely many
generators of the CM order, and close it under conjugation.
The uniqueness of the Laurent normalization in C.7.6 gives
the same field of definition for every $F_\alpha$.
Adjoining the hypothesized algebraic number $H(u,u)$ gives
one field for all the values in \eqref{eq:cc-fixedfield}.
The value for $\alpha=0$ is zero, and period shifts do not
change it. Squaring was included because algebraicity of the
individual positive square roots alone would not establish
their membership in one fixed number field.

Logarithmically differentiating the squared CM law at $u$
gives
\[
 a_\alpha=\frac{1}{2\alpha}
     \frac{(d/dz)F_\alpha(\exp_Ez)|_{z=u}}{F_\alpha(Q)}
                  \in\overline{\Q}.
\]
The denominator is nonzero since $Q$ is nontorsion. The
derivation $d/dz$ is the algebraic invariant derivation on
$E$, so the displayed derivative has algebraic value.
Since $D(u)$ is real, conjugation gives the corresponding
formula at $\bar\alpha u$. On a neighborhood where
$\mathcal F\ne0$, its logarithmic gradient and Hessian are
\[
 (\partial_z\log\mathcal F,\partial_w\log\mathcal F)
       =(D(z)-kw,D(w)-kz),\qquad
 \operatorname{Hess}\log\mathcal F=
 \begin{pmatrix}-\wp(z)-s_2&-k\\-k&-\wp(w)-s_2\end{pmatrix}.
\]
Evaluation on $(\alpha,-\bar\alpha)$ proves both identities
in the statement. In particular the Gaussian cancels in the
first derivative and reappears in the second derivative.

To prove $k$ transcendental, recall the real period formula
$\eta_1=(s_2+k)\omega_1$. If $k$ were algebraic, the
nonzero period vector $(\omega_1,\eta_1)$ of the algebraic
universal vector extension of $E$ would lie on the algebraic
line of slope $s_2+k$. Dividing this period by a sufficiently
large integer supplies a nonzero algebraic torsion exponential.
The analytic subgroup theorem
\citep[Theorem 6.1]{huber2022} would give a one-dimensional
algebraic subgroup with that Lie line. Its projection to $E$
is nonzero, so it is an elliptic curve isogenous to $E$ and
splits the vector extension after that isogeny. This is
impossible: its nonzero class in $H^1(E,\mathcal O_E)$ stays
nonzero under an isogeny, whose pullback has a left inverse
after dividing the trace by the degree. These are the same
algebraic logarithm coordinates and extension as in C.7.7;
a fixed algebraic linear change of coordinates leaves the
argument unchanged.

Write the two logarithmic derivatives as
$V_\alpha\log\mathcal F=B_\alpha$ and
$V_\alpha^2\log\mathcal F=A_\alpha+2dk$, where
$A_\alpha,B_\alpha$ are algebraic. Since $H=\mathcal F^2$,
\[
 \frac{V_\alpha^2H(p_\alpha)}{H(p_\alpha)}
       =4B_\alpha^2+2A_\alpha+4dk.
\]
The right side is transcendental, and the asserted conclusion
follows when $H(p_\alpha)\ne0$ is algebraic.
At $\alpha=1$, $a_1=0$, $\wp(u)=4$ and $s_2=-2$,
which give the last displayed formula in the statement.
\end{proof}

\paragraph{The precise several-variable criteria that were checked.}
The differential criterion of
\citet[Theorems 3.6 and 3.7]{waldschmidt2009auxiliary}
requires a finite algebraically independent family whose
partial derivatives belong to the polynomial ring generated
by that family over one number field. Algebraic values of
the full family then give algebraic values of all iterated
partial derivatives. The second-jet calculation above proves
that a family containing $H$ cannot satisfy both requirements
at the hypothesized algebraic orbit points. The issue persists
although the first cancelling direction is algebraic; the
direction also varies with $\alpha$.

For a criterion without differential equations we opened
\citet[Theorem 1, p.~24]{waldschmidt1975dimension}, including
its proof and a visual check of the displayed inequalities.
For $d$ meromorphic functions on $\C^n$ with growth and
arithmetic-size exponents $\rho_1,\ldots,\rho_d$, a set with
cardinality $N^{\ell-o(1)}$ in a ball of radius $O(N)$ and
separation at least $N^{-\lambda-o(1)}$ yields algebraic
dependence under the strict condition
\begin{equation}\label{eq:cc-rankcriterion}
 (d-1)\ell>\sum_{j=1}^d\rho_j
                     +2d(n-1)(\lambda+1).
\end{equation}
The theorem also requires all function values to belong to
one number field and bounds on denominators in the entire
representations of the meromorphic functions. We do not
silently discard these hypotheses.

The natural independent triple here is
$\wp(z),\wp(w),H(z,w)$, all of order at most two.
The first two are algebraically independent. If $H$ satisfied
a polynomial over $\C(\wp(z),\wp(w))$, fix a generic pair
$(z,w)$ and translate only $z$ by $m\omega$, with
$0\ne\omega\in\Lambda$. The coefficients of the polynomial
remain constant, whereas the absolute values of its proposed
roots $H(z+m\omega,w)$ grow as
$\exp(k|\omega|^2m^2+O(|m|))$. This contradicts the finite
set of roots of a nonzero polynomial. Thus the triple is
independent.

Even granting all the size and denominator bounds with the
favorable exponents $\rho_j=2$, substitution of
$n=2$, $d=3$ and $\ell=4$ into
\eqref{eq:cc-rankcriterion} would require
\[
 8>12+6\lambda,
\]
which fails even for $\lambda=0$. The actual geometry is
less favorable: the $\asymp N^4$ points obtained by bounding
the four integer coefficients lie in a ball of radius $O(N)$
in a real plane. Packing separated discs in that plane gives
$N^4\ll N^{2+2\lambda+o(1)}$, hence $\lambda\ge1$.
Removing the rank-two set of poles does not change this count.
Larger growth exponents only worsen the inequality. This
checks the particular published criterion rather than asserting
that every possible multivariable transcendence method fails.
The missing unconditional step in C.7(c) remains the exclusion
of an algebraic value on the diagonal compact leaf.

\subsection{Division values and a variable-degree interpolation test}

\begin{statement}[C.7.12; Proved route test, with an explicit limitation]
Let $E:y^2+y=x^3-38x+90$, $Q=(4,-2)$, and suppose, only in
this statement, that $N(Q)$ is algebraic.  Write
\[
 \delta=19^{1/4},\qquad b=\eta_1/\omega_1,\qquad
 g(z)=\delta\sigma(z)e^{-bz^2/2},\qquad Q=\exp_E(u),
\]
where $u,\omega_1$ are real as in C.7.9.  Thus
$|g(z)|=N(\exp_E z)$ for real $z$.

Every point in the CM division hull of $Q$ has algebraic norm.
There are constants $D,H$, depending on the fixed curve, $Q$ and
the hypothetical algebraic number $N(Q)$, such that, for
$nR=mQ$, $|m|\leq n$, $R\ne O$,
\[
 h(N(R))\leq H,
 \qquad
 [\Q(R,N(R)):\Q]\leq
 \begin{cases}Dn^4,&R\in E(\R),\\2Dn^6,&R\in E(\C).
 \end{cases}
\]
Here $h$ is absolute logarithmic height; the complex line means
the chosen complex embedding of an algebraic point.

Under the same hypothesis, any finite collection of real translates
$g(z+a)$ by division-hull logarithms has function-field
transcendence degree at most three over $\overline\Q$, after
adjoining the elliptic coordinates.  Three translates already
attain this cap.  The analogous holomorphic extensions in one
variable of $N(\exp_E(z+a))^2$ for complex division-hull shifts
have the same cap.  A safe common degree envelope for a fixed
finite family on the real division grids is $O(n^5)$, with
uniformly bounded heights.

These proved degree and height envelopes do not verify the
variable-degree criterion of Gramain--Mignotte--Waldschmidt
\citep[Theorem 2.1]{gramain1986values}.  This is a failure of the
available sufficient certificate, not a proof that the actual
degrees violate every possible application of that theorem.
The arithmetic-germ hypotheses in the criteria of
\citet[Theorem 6.1]{herblot2012meromorphic} and
\citet[Theorem 4.16]{gasbarri2010analytic} also remain unavailable;
the relevant nontorsion germs have a transcendental tangent.
\end{statement}

\begin{proof}
\emph{Propagation and degrees.}
The full division identity in C.4.2 is
\[
 N(R)^{n^2}=N(nR)\delta^{n^2-1}/|\psi_n(R)|.
\]
For nontorsion $R$, its last factor is a nonzero algebraic number.
Consequently algebraicity propagates through multiplication
and division.  The CM multiplication formula in C.7.6 gives
the same conclusion for the CM division hull.  All nonzero
torsion values were already proved algebraic in C.4.
For example, if $T$ has order dividing $n$ and $P$ is algebraic
and nontorsion, division at $P$ and $P+T$ gives directly
\[
 \left(\frac{N(P+T)}{N(P)}\right)^{n^2}
       =\frac{|\psi_n(P)|}{|\psi_n(P+T)|}\in\overline\Q.
\]
No undetermined harmonic normalization is needed.

Put $K_0=\Q(\delta,N(Q))$ and $D=[K_0:\Q]$.
For every nonzero integer $m$, $N(mQ)\in K_0$, since
$\psi_m(Q)$ is rational and its absolute value is rational.
The fibre of $[n]$ over $mQ$ has $n^2$ points, so
$[K_0(R):K_0]\leq n^2$.
If $R$ is real, $|\psi_n(R)|=\pm\psi_n(R)$; adjoining
$N(R)$ therefore costs at most a further factor $n^2$.
For a complex $R$, adjoin both $R$ and $\bar R$, at cost
at most $n^4$, and use
\[
 N(R)^{2n^2}=
 \frac{N(mQ)^2\delta^{2(n^2-1)}}
      {\psi_n(R)\psi_n(\bar R)}.
\]
This proves the stated degree envelopes when $m\ne0$.
For $m=0$, take the limit in the division identity at a nonzero
$n$-torsion point.  The zero of $\psi_n$ is simple in the
uniformizing coordinate, and $N(z)=\delta|z|+o(|z|)$ at zero;
hence
\[
 N(R)^{n^2}=
       \frac{n\delta^{n^2}}{|\partial_z\psi_n(R)|}.
\]
The derivative is algebraic, since the invariant derivation on
the Weierstrass model is algebraic.  The same degree bounds follow.

\emph{Uniform heights.}
We give the coefficient estimate needed here, so the height
assertion does not rely on an unspecified division-polynomial
bound.  On the fixed integral model, $\psi_n$ is a polynomial
in $x$ for odd $n$, and $(2y+1)$ times a polynomial in $x$
for even $n$.  These polynomials have degree $O(n^2)$.
Choose a fixed circle $|x|=A>1$ avoiding the finite branch values.
On its compact inverse image in $E(\C)$, $N$ has a positive
minimum $m_A$, and $|2y+1|$ has a positive minimum.  Let $M$
be the maximum of $N$ on the whole compact elliptic curve.
The rearranged division identity gives
\[
 |\psi_n(R)|\leq M\delta^{n^2-1}m_A^{-n^2}
       \qquad(|x(R)|=A).
\]
Cauchy's coefficient estimate, and division by $2y+1$ for
even $n$, show that the logarithm of the coefficient
$\ell^1$ norm is $O_E(n^2)$.  The same assertion holds for
$\partial_z\psi_n$, because differentiation and the fixed
Weierstrass equation add only polynomial factors in $n$.

The elementary height estimate for a fixed degree-four
duplication map, followed by telescoping, gives
$h(x(R))=2\widehat h(R)+O_E(1)$ uniformly over
$E(\overline\Q)$.  Since
$\widehat h(R)=(m/n)^2\widehat h(Q)$, the coordinate heights
are uniformly bounded when $|m|\leq n$.
The polynomial evaluation height inequality
\citep[Lemma 1.0]{gramain1986values} now gives
$h(\psi_n(R))=O_{E,Q}(n^2)$, and the same estimate for its
invariant derivative at torsion points.  Also
\[
 h(N(mQ))\leq m^2h(N(Q))+O_{E,Q}(m^2).
\]
Divide the height of the division identity by $n^2$.
For a complex value one uses
$h(|\alpha|)\leq h(\alpha)$, which follows by squaring and
using $|\alpha|^2=\alpha\bar\alpha$.
This proves a uniform bound for $h(N(R))$.

\emph{The translate count is finite in function-field dimension.}
For a real logarithm $a$ of a nonzero division point define
\[
 h_a(z)=\frac{\sigma(z+a)}{\sigma(z)\sigma(a)}e^{-baz}.
\]
Then
\[
 g(z+a)=\delta^{-1}g(z)g(a)h_a(z),\qquad
 h_a(z+\omega_1)=h_a(z),\quad
 h_a(z+\omega_2)=e^{-2\pi ia/\omega_1}h_a(z).
\]
The last two identities are exact sigma automorphy and Legendre's
relation.  A real point in the CM division hull is in the ordinary
division hull: if $nR=\alpha Q$ is real, then
$(\alpha-\bar\alpha)Q=O$, so $\alpha=\bar\alpha\in\Z$.
Indeed a nonzero endomorphism has finite kernel, whereas $Q$
is nontorsion.  The real locus of this negative-discriminant
curve is connected, so every real point has a real logarithm.
Consequently its real division-hull logarithms are precisely
of the form $(mu+k\omega_1)/n$, and lie in
$\Q u+\Q\omega_1$.
For finitely many shifts choose a common denominator
$q$ and write $a_j=(m_ju+k_j\omega_1)/q$; put $v=u/q$.
Then $h_{a_j}/h_v^{m_j}$ is elliptic for
$\Lambda_q=\Z\omega_1+\Z q\omega_2$.
The curve $\C/\Lambda_q$ is isogenous to $E$ and has an
algebraic model with this differential normalization.
Its displayed elliptic functions have algebraic divisors.
Each therefore equals a rational function over $\overline\Q$
times a complex scalar.  Evaluate at a real division point
avoiding their finitely many zeros and poles.  The formula
for $h_a$ in terms of $g$ and the assumed algebraicity of all
real division values make that scalar algebraic too.
Consequently all these translates are algebraic over
\[
 \overline\Q(E)\bigl(g,h_v\bigr).
\]
This proves the upper cap three over $\overline\Q$, not merely
over $\C$.

The cap is attained.  The multiplier
$e^{-2\pi iv/\omega_1}$ is not a root of unity, since $Q$
is nontorsion.  Applying the shifts $z\mapsto z+m\omega_2$
to a hypothetical polynomial relation shows that $h_v$ is
transcendental over the elliptic function field: at a fixed
generic $z$, a nonzero polynomial cannot vanish at infinitely
many distinct powers of this multiplier.  The same shifts
give $|g(z+m\omega_2)|=\exp(c m^2+O(m))$ with $c>0$.
After clearing denominators, any nonzero highest coefficient in
a polynomial in $g$, with coefficients in the elliptic field and
$h_v$, stays bounded
away from zero on a subsequence of the irrational rotation.
The highest power of $g$ then dominates, a contradiction.
Thus $g,h_v$ and the elliptic coordinate have transcendence
degree three.  Finally sigma addition gives
\[
 \wp(z)=\wp(v)-
 \frac{\delta^2g(z+v)g(z-v)}{g(z)^2g(v)^2},
\]
so the three translates $g(z),g(z+v),g(z-v)$ already attain it.

For a complex division shift $a$, the holomorphic extension
from real $z$ of the squared norm is
\[
 F_a(z)=\frac{N(a)^2}{\delta^2}
                  g(z)^2h_a(z)h_{\bar a}(z).
\]
The multiplier of $h_ah_{\bar a}$ depends only on
$a+\bar a\in\Q u+\Q\omega_1$ for this real CM curve.
The preceding finite-index-lattice argument applies again;
evaluating $F_a$ at a real division point makes the remaining
scalar algebraic.  Thus complex shifts of the norm do not
provide arbitrarily many independent functions either.

For completeness, the useful joint degree envelope for the
three generators is $O(n^5)$.  The map
$z\mapsto(\exp_Ez,h_v(z))$ describes the complex analytic
subgroup $\exp_{G_V}(z,(\zeta(v)-bv)z)$ of the algebraic
extension $G_V$, $V=\exp_E(v)$, in the meromorphic algebraic
section coordinate used in C.7.1.  Its restriction to real
$z$ is a compact circle; no compactness or closedness of the
whole complex analytic subgroup is asserted.
Its value at $z=u$ is algebraic, because
$h_v(u)=\delta g(u+v)/(g(u)g(v))$.
Its real period is $\omega_1$.  Hence for
$z=(mu+k\omega_1)/n$ its value is an $n$-division point of
a fixed-field algebraic multiple of that point.
The multiplication map on $G_V$ has degree $n^3$: its kernel
is an extension of $E[n]$ by $\mathbb G_m[n]$.
Adjoining $g(z)$ costs at most another $n^2$.
The fixed isogeny and the finitely many algebraic constants
affect only the implied constant.  Uniform height bounds for
the real translates follow from the same norm-bound proof with
denominator $nq$ and a fixed larger bound on the ratio of
numerator to denominator; quotient expressions give the same
bound for $h_v$.  For fixed complex shifts, use instead the
proved expression for $F_a$ as a rational elliptic function
on a fixed isogeny, times $g^2h_v^m$, with an algebraic scalar.
The heights of the isogenous elliptic coordinates, $g$ and
$h_v$ are bounded, so the rational-function evaluation height
inequality bounds $h(F_a)$ uniformly.  Removable zeros or poles
of this particular expression lie in a fixed finite set and
can be handled by their finitely many actual values.  Taking
the positive square root to recover the complex-shifted norm
divides its height by two and costs at most a fixed factor two
in degree for each of the fixed finite number of shifts.

\emph{The criterion has been tested, not assumed.}
For $k$ analytic functions, GMW Theorem 2.1 uses a point count
$\mu$, joint degree envelope $d$, height envelopes $\psi_j$
and growth envelopes $\theta_j$; its decisive sufficient
inequality is
\[
 d(N)\prod_{j=1}^k
  \max\{d(N)\psi_j(N),\theta_j(R(N))\}
 \leq c_0\mu(N)^{k-1}\varphi(N)^k,
\]
where
\[
 \varphi(N)=\log
 \frac{R(N)^2+r(N)r(N+1)}
      {R(N)(r(N)+r(N+1))}.
\]
We checked this formula and the other hypotheses in the
original printed pages 101--102.  On a fixed analytic disc
with fixed smaller sampling disc, $\varphi$ and the growth
envelopes are constant.  The height envelopes just proved
are constant as well.  Thus the required comparison reduces
to $d^{k+1}\ll\mu^{k-1}$.

At denominator $n$ there are $O(n^2)$ real points with
$nR=mQ$, $0\leq m\leq n$; through denominator $N$ there are
$O(N^3)$.  For $(\wp,g)$ on a disc without a pole, our proved
degree envelope $d=O(N^4)$ leads to the attempted comparison
$N^{12}\ll N^3$.  For three independent translates, the
safe joint envelope $d=O(N^5)$ leads to
$N^{20}\ll N^6$.  Neither comparison is available from
these estimates.  Adding further translates does not repair
this certificate: beyond three, the dependence conclusion
would already be true and would not contradict the norm
hypothesis.  We do not infer an actual degree lower bound
from an upper bound, and do not rule out a sharper arithmetic
estimate or a differently organized interpolation argument.

The alternative variable-degree theorems of Herblot and
Gasbarri require a formal germ over the residue number field,
with additional arithmetic estimates.  At a real nontorsion
division point with coordinate $z$ and algebraic $g(z)$,
\[
 \frac{dg}{dx}=\frac{g(z)}{\wp'(z)}
                         \bigl(\zeta(z)-bz\bigr).
\]
The slope in parentheses is transcendental by the period-line
argument of C.7.9 applied to that nontorsion algebraic point.
The other two factors are nonzero algebraic numbers.
Thus the germ is not even defined over $\overline\Q$;
the required arithmetic-germ hypothesis has not been gained
by passing to division values.
\end{proof}

We also checked the Mahler-functional-equation route against
\citet[Theorem 1.1 and system (1.2)]{adamczewski2017mahler}.
That theorem requires algebraic Taylor coefficients, a linear
power-map functional system and an algebraic regular argument.
The compact normalization here has the transcendental coefficient
$b$, its elliptic logarithm is transcendental, and its duplication
in the algebraic coordinate is a Latt\`es map.  We have not obtained
a representation satisfying the stated hypotheses; we do not
assert that every alternative representation is impossible.
The entropy identity of
\citet[Theorem 5.1]{einsiedler2001entropy} uses a sequence of
transformations built from the denominators of $x(2^nQ)$,
with scale $4^n$.  It does not identify the height with the
entropy of a single finite-dimensional algebraic automorphism,
and supplies no algebraic value for its exponential.

The period-refined algebraic-group rank theorem has also been
checked in the original sources
\citep[Theorem 4.1 and Corollary 4.3]{waldschmidt1988periodrank}
\citep[main theorem, p.~363]{roy1993periodrank}.
Its coefficient is
$\alpha(G)=2\dim G-2d_0(G)-d_1(G)$, where $d_1$
counts multiplicative \emph{quotients}, not the dimension of
the torus subgroup.  For the nonsplit $G_Q$ of C.7.9,
$d_0=d_1=0$: a nonzero character on its torus would kill
a nonzero integer multiple of the nontorsion extension class.
Thus $\alpha(G_Q)=4$.  The line $\ell_c$ has no nonzero
algebraic vector, has period rank $\kappa=1$, and the norm
hypothesis supplies one independent algebraic-point logarithm
in addition to that period.  The theorem allows total rank
$3=(4-1)\dim\ell_c$; the available rank $2$ therefore gives
no contradiction.  The period rank is exactly one because
a second period on this line would, by Legendre's relation,
force $u/\omega_1$ to be rational, contrary to the
nontorsion of $Q$.

\bibliographystyle{apalike}
\bibliography{S6_05_ARITHMETIC_references}
\end{document}
